% Options for packages loaded elsewhere
\PassOptionsToPackage{unicode}{hyperref}
\PassOptionsToPackage{hyphens}{url}
\PassOptionsToPackage{dvipsnames,svgnames,x11names}{xcolor}
%
\documentclass[
  letterpaper,
  DIV=11,
  numbers=noendperiod]{scrreprt}

\usepackage{amsmath,amssymb}
\usepackage{iftex}
\ifPDFTeX
  \usepackage[T1]{fontenc}
  \usepackage[utf8]{inputenc}
  \usepackage{textcomp} % provide euro and other symbols
\else % if luatex or xetex
  \usepackage{unicode-math}
  \defaultfontfeatures{Scale=MatchLowercase}
  \defaultfontfeatures[\rmfamily]{Ligatures=TeX,Scale=1}
\fi
\usepackage{lmodern}
\ifPDFTeX\else  
    % xetex/luatex font selection
\fi
% Use upquote if available, for straight quotes in verbatim environments
\IfFileExists{upquote.sty}{\usepackage{upquote}}{}
\IfFileExists{microtype.sty}{% use microtype if available
  \usepackage[]{microtype}
  \UseMicrotypeSet[protrusion]{basicmath} % disable protrusion for tt fonts
}{}
\makeatletter
\@ifundefined{KOMAClassName}{% if non-KOMA class
  \IfFileExists{parskip.sty}{%
    \usepackage{parskip}
  }{% else
    \setlength{\parindent}{0pt}
    \setlength{\parskip}{6pt plus 2pt minus 1pt}}
}{% if KOMA class
  \KOMAoptions{parskip=half}}
\makeatother
\usepackage{xcolor}
\setlength{\emergencystretch}{3em} % prevent overfull lines
\setcounter{secnumdepth}{5}
% Make \paragraph and \subparagraph free-standing
\makeatletter
\ifx\paragraph\undefined\else
  \let\oldparagraph\paragraph
  \renewcommand{\paragraph}{
    \@ifstar
      \xxxParagraphStar
      \xxxParagraphNoStar
  }
  \newcommand{\xxxParagraphStar}[1]{\oldparagraph*{#1}\mbox{}}
  \newcommand{\xxxParagraphNoStar}[1]{\oldparagraph{#1}\mbox{}}
\fi
\ifx\subparagraph\undefined\else
  \let\oldsubparagraph\subparagraph
  \renewcommand{\subparagraph}{
    \@ifstar
      \xxxSubParagraphStar
      \xxxSubParagraphNoStar
  }
  \newcommand{\xxxSubParagraphStar}[1]{\oldsubparagraph*{#1}\mbox{}}
  \newcommand{\xxxSubParagraphNoStar}[1]{\oldsubparagraph{#1}\mbox{}}
\fi
\makeatother

\usepackage{color}
\usepackage{fancyvrb}
\newcommand{\VerbBar}{|}
\newcommand{\VERB}{\Verb[commandchars=\\\{\}]}
\DefineVerbatimEnvironment{Highlighting}{Verbatim}{commandchars=\\\{\}}
% Add ',fontsize=\small' for more characters per line
\usepackage{framed}
\definecolor{shadecolor}{RGB}{241,243,245}
\newenvironment{Shaded}{\begin{snugshade}}{\end{snugshade}}
\newcommand{\AlertTok}[1]{\textcolor[rgb]{0.68,0.00,0.00}{#1}}
\newcommand{\AnnotationTok}[1]{\textcolor[rgb]{0.37,0.37,0.37}{#1}}
\newcommand{\AttributeTok}[1]{\textcolor[rgb]{0.40,0.45,0.13}{#1}}
\newcommand{\BaseNTok}[1]{\textcolor[rgb]{0.68,0.00,0.00}{#1}}
\newcommand{\BuiltInTok}[1]{\textcolor[rgb]{0.00,0.23,0.31}{#1}}
\newcommand{\CharTok}[1]{\textcolor[rgb]{0.13,0.47,0.30}{#1}}
\newcommand{\CommentTok}[1]{\textcolor[rgb]{0.37,0.37,0.37}{#1}}
\newcommand{\CommentVarTok}[1]{\textcolor[rgb]{0.37,0.37,0.37}{\textit{#1}}}
\newcommand{\ConstantTok}[1]{\textcolor[rgb]{0.56,0.35,0.01}{#1}}
\newcommand{\ControlFlowTok}[1]{\textcolor[rgb]{0.00,0.23,0.31}{\textbf{#1}}}
\newcommand{\DataTypeTok}[1]{\textcolor[rgb]{0.68,0.00,0.00}{#1}}
\newcommand{\DecValTok}[1]{\textcolor[rgb]{0.68,0.00,0.00}{#1}}
\newcommand{\DocumentationTok}[1]{\textcolor[rgb]{0.37,0.37,0.37}{\textit{#1}}}
\newcommand{\ErrorTok}[1]{\textcolor[rgb]{0.68,0.00,0.00}{#1}}
\newcommand{\ExtensionTok}[1]{\textcolor[rgb]{0.00,0.23,0.31}{#1}}
\newcommand{\FloatTok}[1]{\textcolor[rgb]{0.68,0.00,0.00}{#1}}
\newcommand{\FunctionTok}[1]{\textcolor[rgb]{0.28,0.35,0.67}{#1}}
\newcommand{\ImportTok}[1]{\textcolor[rgb]{0.00,0.46,0.62}{#1}}
\newcommand{\InformationTok}[1]{\textcolor[rgb]{0.37,0.37,0.37}{#1}}
\newcommand{\KeywordTok}[1]{\textcolor[rgb]{0.00,0.23,0.31}{\textbf{#1}}}
\newcommand{\NormalTok}[1]{\textcolor[rgb]{0.00,0.23,0.31}{#1}}
\newcommand{\OperatorTok}[1]{\textcolor[rgb]{0.37,0.37,0.37}{#1}}
\newcommand{\OtherTok}[1]{\textcolor[rgb]{0.00,0.23,0.31}{#1}}
\newcommand{\PreprocessorTok}[1]{\textcolor[rgb]{0.68,0.00,0.00}{#1}}
\newcommand{\RegionMarkerTok}[1]{\textcolor[rgb]{0.00,0.23,0.31}{#1}}
\newcommand{\SpecialCharTok}[1]{\textcolor[rgb]{0.37,0.37,0.37}{#1}}
\newcommand{\SpecialStringTok}[1]{\textcolor[rgb]{0.13,0.47,0.30}{#1}}
\newcommand{\StringTok}[1]{\textcolor[rgb]{0.13,0.47,0.30}{#1}}
\newcommand{\VariableTok}[1]{\textcolor[rgb]{0.07,0.07,0.07}{#1}}
\newcommand{\VerbatimStringTok}[1]{\textcolor[rgb]{0.13,0.47,0.30}{#1}}
\newcommand{\WarningTok}[1]{\textcolor[rgb]{0.37,0.37,0.37}{\textit{#1}}}

\providecommand{\tightlist}{%
  \setlength{\itemsep}{0pt}\setlength{\parskip}{0pt}}\usepackage{longtable,booktabs,array}
\usepackage{calc} % for calculating minipage widths
% Correct order of tables after \paragraph or \subparagraph
\usepackage{etoolbox}
\makeatletter
\patchcmd\longtable{\par}{\if@noskipsec\mbox{}\fi\par}{}{}
\makeatother
% Allow footnotes in longtable head/foot
\IfFileExists{footnotehyper.sty}{\usepackage{footnotehyper}}{\usepackage{footnote}}
\makesavenoteenv{longtable}
\usepackage{graphicx}
\makeatletter
\def\maxwidth{\ifdim\Gin@nat@width>\linewidth\linewidth\else\Gin@nat@width\fi}
\def\maxheight{\ifdim\Gin@nat@height>\textheight\textheight\else\Gin@nat@height\fi}
\makeatother
% Scale images if necessary, so that they will not overflow the page
% margins by default, and it is still possible to overwrite the defaults
% using explicit options in \includegraphics[width, height, ...]{}
\setkeys{Gin}{width=\maxwidth,height=\maxheight,keepaspectratio}
% Set default figure placement to htbp
\makeatletter
\def\fps@figure{htbp}
\makeatother

\AtBeginDocument{\floatplacement{codelisting}{H}}
\KOMAoption{captions}{tableheading}
\makeatletter
\@ifpackageloaded{tcolorbox}{}{\usepackage[skins,breakable]{tcolorbox}}
\@ifpackageloaded{fontawesome5}{}{\usepackage{fontawesome5}}
\definecolor{quarto-callout-color}{HTML}{909090}
\definecolor{quarto-callout-note-color}{HTML}{0758E5}
\definecolor{quarto-callout-important-color}{HTML}{CC1914}
\definecolor{quarto-callout-warning-color}{HTML}{EB9113}
\definecolor{quarto-callout-tip-color}{HTML}{00A047}
\definecolor{quarto-callout-caution-color}{HTML}{FC5300}
\definecolor{quarto-callout-color-frame}{HTML}{acacac}
\definecolor{quarto-callout-note-color-frame}{HTML}{4582ec}
\definecolor{quarto-callout-important-color-frame}{HTML}{d9534f}
\definecolor{quarto-callout-warning-color-frame}{HTML}{f0ad4e}
\definecolor{quarto-callout-tip-color-frame}{HTML}{02b875}
\definecolor{quarto-callout-caution-color-frame}{HTML}{fd7e14}
\makeatother
\makeatletter
\@ifpackageloaded{bookmark}{}{\usepackage{bookmark}}
\makeatother
\makeatletter
\@ifpackageloaded{caption}{}{\usepackage{caption}}
\AtBeginDocument{%
\ifdefined\contentsname
  \renewcommand*\contentsname{Table of contents}
\else
  \newcommand\contentsname{Table of contents}
\fi
\ifdefined\listfigurename
  \renewcommand*\listfigurename{List of Figures}
\else
  \newcommand\listfigurename{List of Figures}
\fi
\ifdefined\listtablename
  \renewcommand*\listtablename{List of Tables}
\else
  \newcommand\listtablename{List of Tables}
\fi
\ifdefined\figurename
  \renewcommand*\figurename{Figure}
\else
  \newcommand\figurename{Figure}
\fi
\ifdefined\tablename
  \renewcommand*\tablename{Table}
\else
  \newcommand\tablename{Table}
\fi
}
\@ifpackageloaded{float}{}{\usepackage{float}}
\floatstyle{ruled}
\@ifundefined{c@chapter}{\newfloat{codelisting}{h}{lop}}{\newfloat{codelisting}{h}{lop}[chapter]}
\floatname{codelisting}{Listing}
\newcommand*\listoflistings{\listof{codelisting}{List of Listings}}
\makeatother
\makeatletter
\makeatother
\makeatletter
\@ifpackageloaded{caption}{}{\usepackage{caption}}
\@ifpackageloaded{subcaption}{}{\usepackage{subcaption}}
\makeatother

\ifLuaTeX
  \usepackage{selnolig}  % disable illegal ligatures
\fi
\usepackage{bookmark}

\IfFileExists{xurl.sty}{\usepackage{xurl}}{} % add URL line breaks if available
\urlstyle{same} % disable monospaced font for URLs
\hypersetup{
  pdftitle={GeneralEquilibriumModeling.jl and Structural Equilibrium Models},
  pdfauthor={LI Wu (liwu@staff.shu.edu.cn)},
  colorlinks=true,
  linkcolor={blue},
  filecolor={Maroon},
  citecolor={Blue},
  urlcolor={Blue},
  pdfcreator={LaTeX via pandoc}}


\title{GeneralEquilibriumModeling.jl and Structural Equilibrium Models}
\usepackage{etoolbox}
\makeatletter
\providecommand{\subtitle}[1]{% add subtitle to \maketitle
  \apptocmd{\@title}{\par {\large #1 \par}}{}{}
}
\makeatother
\subtitle{A Cookbook}
\author{LI Wu (liwu@staff.shu.edu.cn)}
\date{2026-09-10}

\begin{document}
\maketitle

\renewcommand*\contentsname{Table of contents}
{
\hypersetup{linkcolor=}
\setcounter{tocdepth}{2}
\tableofcontents
}

\bookmarksetup{startatroot}

\chapter*{Preface}\label{preface}
\addcontentsline{toc}{chapter}{Preface}

\markboth{Preface}{Preface}

\begin{tcolorbox}[enhanced jigsaw, opacityback=0, colbacktitle=quarto-callout-warning-color!10!white, arc=.35mm, coltitle=black, leftrule=.75mm, left=2mm, toptitle=1mm, bottomrule=.15mm, colback=white, breakable, rightrule=.15mm, title=\textcolor{quarto-callout-warning-color}{\faExclamationTriangle}\hspace{0.5em}{Warning}, toprule=.15mm, colframe=quarto-callout-warning-color-frame, bottomtitle=1mm, titlerule=0mm, opacitybacktitle=0.6]

🚧 Work in Progress

This manuscript is currently a draft. If you spot any errors or typos,
or have suggestions, please feel free to submit feedback via
\href{https://github.com/LiwuR/gembook/issues}{GitHub Issues}.

\end{tcolorbox}

This book provides a concise introduction to applied general equilibrium
modeling. It aims to explain, as clearly and directly as possible, the
basic structure of general equilibrium models, methods for constructing
such models, and their numerical solution. The book has the following
main features:

\begin{enumerate}
\def\labelenumi{\arabic{enumi}.}
\item
  The mathematical models presented in this book are primarily based on
  the structural equilibrium model (SEM) developed by the author (Li,
  2019). The SEM extends the von Neumann equilibrium model and provides
  a matrix-based framework for describing economic agents and their
  activities, input-output relationships, and market equilibrium
  conditions in a unified manner. A structural equilibrium model can
  also be transformed into a complementarity problem. Compared with its
  complementarity formulation, the structural equilibrium representation
  is generally more intuitive and makes the economic structure of the
  model and its equilibrium conditions easier to understand.
\item
  The numerical models in this book are primarily constructed and solved
  using \texttt{GeneralEquilibriumModeling.jl}, a Julia package
  developed by the author. The package contains two main modules:
  \texttt{GEM}, which provides the core representation and solution
  framework for general equilibrium models, and \texttt{GEMB}, which
  provides more user-friendly, concise, and economically oriented
  modeling tools built on \texttt{GEM}. The \texttt{GEM} module uses
  \texttt{PATHSolver.jl} as its underlying solver for the
  complementarity problems associated with general equilibrium models.
\item
  Compared with \texttt{MPSGE.jl}, which emphasizes concise construction
  and solution of computable general equilibrium (CGE) models, and
  \texttt{JCGE.jl}, which is oriented toward modular CGE modeling,
  \texttt{GeneralEquilibriumModeling.jl} places greater emphasis on the
  direct correspondence between economic theory and numerical models. In
  particular, its \texttt{GEMB} module is well suited for learning and
  teaching general equilibrium modeling. It allows users to begin with
  the behavior of economic agents and equilibrium conditions, understand
  structural equilibrium models and their complementarity
  representations, and gradually construct applied general equilibrium
  models.
\item
  This book is organized largely as a handbook of examples. A sequence
  of models, progressing from simple to more complex settings, is used
  to introduce the construction, representation, numerical solution, and
  economic interpretation of general equilibrium models.
\item
  The book is intended as a concise overview. Its emphasis is on the
  basic ideas, model structures, and practical methods of general
  equilibrium modeling rather than on extensive literature reviews,
  detailed theoretical proofs, or systematic instruction in programming
  syntax. Some Julia examples in this book have been adapted from
  examples in the author's earlier R packages, \texttt{GE} and
  \texttt{CGE}.
\end{enumerate}

\bookmarksetup{startatroot}

\chapter{The von Neumann Equilibrium Model and Structural Equilibrium
Models}\label{the-von-neumann-equilibrium-model-and-structural-equilibrium-models}

\section{The Eigensystem Form of General Equilibrium
Models}\label{the-eigensystem-form-of-general-equilibrium-models}

\subsection{Eigenvalue Equations and the von Neumann Equilibrium
Model}\label{eigenvalue-equations-and-the-von-neumann-equilibrium-model}

Applied general equilibrium models can be represented in a variety of
mathematical forms. An eigensystem is one important representation. It
reveals the dual structure between prices and activity levels; when the
model takes the form of an eigenvalue inequality system, it can also
characterize the associated complementarity relations.

For a square matrix \(\mathbf A\), the corresponding left and right
eigenvalue equations can be written as

\[
\begin{aligned}
\mathbf p^{\top}\mathbf A
&=\rho\mathbf p^{\top}\\
\mathbf A\mathbf z
&=\rho\mathbf z
\end{aligned}
\]

where \(\mathbf p\) and \(\mathbf z\) are the left and right
eigenvectors, respectively, and \(\rho\) is the eigenvalue. Extending
the eigenvalue equations for a square matrix to two \(n\times m\)
matrices, \(\mathbf A\) and \(\mathbf B\), and relaxing the equalities
to inequalities gives the generalized eigenvalue inequality system

\begin{equation}\phantomsection\label{eq-vonNeumannModel}{
\begin{aligned}
\mathbf p^{\top}\mathbf A
&\geq\rho\mathbf p^{\top}\mathbf B\\
\mathbf A\mathbf z
&\leq\rho\mathbf B\mathbf z
\end{aligned}
}\end{equation}

The von Neumann general equilibrium model is precisely such an
eigenvalue inequality system (von Neumann, 1945; Kemeny, Morgenstern,
and Thompson, 1956). Here, \(\mathbf A\) and \(\mathbf B\) are the input
coefficient matrix and output coefficient matrix, respectively,
\(\mathbf p\geq\mathbf 0\) is the price vector, and
\(\mathbf z\geq\mathbf 0\) is the activity-level vector. The first
inequality states that the input value of each production activity is no
less than its discounted output value, while the second states that the
total input of each commodity does not exceed its discounted total
output. The scalar \(\rho\) is the eigenvalue of the system and,
economically, represents the discount factor corresponding to the
equilibrium growth rate. If the equilibrium growth rate is \(\gamma\),
then

\[
\rho=\frac{1}{1+\gamma}
\]

\subsection{Structural Equilibrium
Models}\label{structural-equilibrium-models}

The (strict-sense) structural equilibrium model extends the input and
output coefficient matrices in the von Neumann model to matrix-valued
functions which usually involve consumers' utilities (Li, 2019, chapter
3):

\[
\begin{aligned}
\mathbf p^{\top}\mathbf A(\mathbf p,\mathbf u,\mathbf z)
&\geq
\rho\mathbf p^{\top}\mathbf B(\mathbf p,\mathbf u,\mathbf z)\\
\mathbf A(\mathbf p,\mathbf u,\mathbf z)\mathbf z
&\leq
\rho\mathbf B(\mathbf p,\mathbf u,\mathbf z)\mathbf z
\end{aligned}
\]

where \(\mathbf A(\mathbf p,\mathbf u,\mathbf z)\) and
\(\mathbf B(\mathbf p,\mathbf u,\mathbf z)\) are the unit demand matrix
and unit supply matrix, respectively, and \(\mathbf u\) is the vector of
consumers' utility levels. Each row of the two matrices corresponds to a
commodity, and each column corresponds to an economic agent. The price
vector remains on the left-hand side of the system and is associated
with the unit revenue-expenditure balance conditions, whereas the
activity-level vector remains on the right-hand side and is associated
with the supply-demand balance conditions (i.e., market-clearing
conditions). The structural equilibrium model therefore remains an
eigensystem. The activity-level vector in a strict-sense structural
equilibrium model typically consists of producers' activity levels and
the numbers of consumers of different types.

For notational simplicity, when \(\rho\neq1\) in a structural
equilibrium model, \(\rho\) can be incorporated into
\(\mathbf B(\cdot)\). In this case, \(\mathbf B(\cdot)\) becomes the
discount-adjusted unit supply matrix, and the structural equilibrium
model can be written as

\[
\begin{aligned}
\mathbf p^{\top}\mathbf A(\cdot)
&\geq
\mathbf p^{\top}\mathbf B(\cdot)\\
\mathbf A(\cdot)\mathbf z
&\leq
\mathbf B(\cdot)\mathbf z
\end{aligned}
\]

Let \(\operatorname{Diag}(\mathbf z)\) denote the diagonal matrix with
\(\mathbf z\) on its main diagonal. The demand and supply matrices
generated by the unit demand and unit supply matrices are

\[
\mathbf D(\cdot)=\mathbf A(\cdot)\operatorname{Diag}(\mathbf z),
\qquad
\mathbf S(\cdot)=\mathbf B(\cdot)\operatorname{Diag}(\mathbf z)
\]

respectively. Thus, each column of \(\mathbf D(\cdot)\) or
\(\mathbf S(\cdot)\) gives the demand or supply of the corresponding
economic agent at its activity level. The structural equilibrium
conditions imply
\(\mathbf p^{\top}\mathbf D(\cdot)=\mathbf p^{\top}\mathbf S(\cdot)\)
and \(\mathbf D(\cdot)\mathbf 1\leq\mathbf S(\cdot)\mathbf 1\). This
leads directly to the mixed-form wide-sense structural equilibrium model

\[
\begin{aligned}
\mathbf p^{\top}\mathbf D(\cdot)
&=
\mathbf p^{\top}\mathbf S(\cdot)\\
\mathbf D(\cdot)\mathbf 1
&\leq
\mathbf S(\cdot)\mathbf 1
\end{aligned}
\]

The first condition represents total revenue-expenditure balance for
each economic agent, while the second states that the demand for each
commodity does not exceed its supply.

If there are no commodities in excess supply at equilibrium, so that
supply equals demand for every commodity, the supply-demand inequality
becomes an equality. This gives the equality-form wide-sense structural
equilibrium model

\[
\begin{aligned}
\mathbf p^{\top}\mathbf D(\cdot)
&=
\mathbf p^{\top}\mathbf S(\cdot)\\
\mathbf D(\cdot)\mathbf 1
&=
\mathbf S(\cdot)\mathbf 1
\end{aligned}
\]

The first equation states that every economic agent satisfies total
revenue-expenditure balance in the accounting sense, while the second
states that supply equals demand for every commodity.

Thus, the equality-form wide-sense structural equilibrium model has the
form of an eigenvalue equation system, whereas the mixed-form wide-sense
structural equilibrium model has the form of an eigenvalue inequality
system.

\section{Complementarity Relations in the von Neumann Equilibrium
Model}\label{complementarity-relations-in-the-von-neumann-equilibrium-model}

The von Neumann equilibrium model Equation~\ref{eq-vonNeumannModel}
implies complementary slackness conditions (Kemeny, Morgenstern, and
Thompson, 1956). Right-multiplying the first inequality by \(\mathbf z\)
and left-multiplying the second inequality by \(\mathbf p^{\top}\) gives

\[
\mathbf p^{\top}\mathbf A\mathbf z
=
\rho\mathbf p^{\top}\mathbf B\mathbf z
\]

That is, the total input value equals the discounted total output value.
The von Neumann equilibrium model therefore implies the following two
sets of complementarity conditions:

\[
\begin{aligned}
\mathbf 0
&\leq
\mathbf A^{\top}\mathbf p-\rho\mathbf B^{\top}\mathbf p
\perp
\mathbf z
\geq
\mathbf 0\\
\mathbf 0
&\leq
\rho\mathbf B\mathbf z-\mathbf A\mathbf z
\perp
\mathbf p
\geq
\mathbf 0
\end{aligned}
\]

The first set describes the complementarity relation between production
activity levels and unit-value losses. If the unit cost of an activity
is strictly greater than its discounted unit revenue, that activity must
be inactive. Conversely, if an activity operates at a positive activity
level, its unit cost must equal its discounted unit revenue.

The second set describes the complementarity relation between commodity
prices and excess supply. If a commodity is in excess supply, its
equilibrium price must be zero. Conversely, if the equilibrium price of
a commodity is positive, its supply must equal its demand.

The von Neumann equilibrium model therefore has two sets of
complementarity relations. On the economic-agent side, each activity
level is complementary to the corresponding unit value loss. On the
commodity side, each commodity price is complementary to the
corresponding excess supply.

\section{Mixed Complementarity Problems and Their Solution in
Julia}\label{mixed-complementarity-problems-and-their-solution-in-julia}

\subsection{Mixed Complementarity Formulation and Numerical
Solution}\label{mixed-complementarity-formulation-and-numerical-solution}

A mixed complementarity problem is defined by a mapping \(\mathbf F\)
and lower and upper bounds on the variable vector \(\mathbf x\).

Given a function

\[
\mathbf F:\mathbb R^N\rightarrow\mathbb R^N
\]

and lower and upper bounds \(\mathbf l\) and \(\mathbf u\) for the
variable vector \(\mathbf x\), a mixed complementarity problem seeks a
vector \(\mathbf x\) satisfying \(\mathbf l\leq\mathbf x\leq\mathbf u\)
such that, for each \(i=1,\ldots,N\),

\[
\begin{cases}
F_i(\mathbf x)\geq0, & x_i=l_i\\
F_i(\mathbf x)=0, & l_i<x_i<u_i\\
F_i(\mathbf x)\leq0, & x_i=u_i
\end{cases}
\]

These conditions can be written compactly in vector form as \[
\mathbf F(\mathbf x)\perp\mathbf l\leq\mathbf x\leq\mathbf u
\]

When \(l_i=0\) and \(u_i=+\infty\), this reduces to the usual
nonnegative complementarity condition

\[
0\leq F_i(\mathbf x)\perp x_i\geq0
\]

When \(l_i=-\infty\) and \(u_i=+\infty\), \(x_i\) is a free variable,
and the corresponding complementarity condition reduces to the equation

\[
F_i(\mathbf x)=0
\]

A mixed complementarity problem can therefore represent, within a single
model, nonnegative variables, variables with lower and upper bounds,
free variables, equations, and complementarity conditions.

Complementarity problems have a standardized mathematical form and can
be solved directly using specialized solvers such as PATH. They are
therefore particularly suitable as the computational representation of
general equilibrium models (Ferris and Munson, 2000). From an economic
perspective, however, complementarity formulations are often less
intuitive because economic relationships are decomposed into variables,
residual functions, and their associated complementarity conditions.

By contrast, a structural equilibrium model represents the basic
structure of an economic system directly in terms of demand matrices,
supply matrices, revenue-expenditure balance conditions, and
supply-demand balance conditions. This representation makes the economic
relationships relatively transparent and facilitates theoretical
analysis and the derivation of analytical solutions. A useful modeling
strategy is therefore to construct and analyze an economic model first
as a structural equilibrium model and then transform it into a
complementarity problem for numerical solution.

In Julia, mixed complementarity models can be constructed with
\texttt{JuMP} and solved using the PATH solver through
\texttt{PATHSolver.jl}.

Because a general equilibrium model typically determines only relative
prices, one commodity must normally be selected as the numeraire before
solving the model with PATH, or an explicit price-normalization
condition must be imposed. By Walras' law, one market-clearing condition
is typically redundant. Thus, when the price of a commodity is fixed as
the numeraire, the corresponding redundant market-clearing condition
normally needs to be removed or replaced. In models with multiple
equilibria, the particular equilibrium found by PATH may depend on the
initial values. After solving the model, it is therefore advisable to
check complementarity residuals, market-clearing errors,
revenue-expenditure balance errors, and the solver termination status.

\subsection{General Equilibrium Modeling
Frameworks}\label{general-equilibrium-modeling-frameworks}

\texttt{GeneralEquilibriumModeling.jl} is a Julia package for
constructing and solving general equilibrium models. It emphasizes the
correspondence between economic-theoretical models and their numerical
solution forms. In the GEMB high-level framework, a model is specified
in terms of commodities, economic agents, activity levels,
revenue-expenditure balance conditions, and supply-demand balance
conditions. In the GEMB low-level framework, these conditions are
represented more directly at the computational level. The resulting
equilibrium system is transformed into a mixed complementarity problem,
represented numerically through \texttt{JuMP}, and solved with PATH
through \texttt{PATHSolver.jl}.

\texttt{GeneralEquilibriumModeling.jl} provides three modeling
frameworks.

\begin{enumerate}
\def\labelenumi{\arabic{enumi}.}
\item
  The \texttt{GEM} framework is the lowest-level core representation and
  solution framework for general equilibrium models. Users can directly
  define economic agents' net supply functions, endogenous variables,
  and associated complementarity conditions, and may also introduce
  auxiliary variables and auxiliary equations. \texttt{GEM} constructs
  the resulting mixed complementarity problem through \texttt{JuMP} and
  solves it using \texttt{PATHSolver.jl}.
\item
  The \texttt{GEMB} low-level framework uses interfaces such as
  \texttt{CESSpec}, \texttt{ActivityDemandSpec}, and
  \texttt{build\_agent} to transform specifications of production,
  consumption, and other economic behavior into
  \texttt{GEM.NetSupplyAgent} objects.
\item
  The \texttt{GEMB} high-level framework organizes commodities and
  economic agents through \texttt{GEMBModel}, commodity names or
  \texttt{CommodityRef}, and interfaces such as \texttt{add\_agent!}.
  Users therefore normally do not need to work directly with commodity
  indices, net supply functions, or low-level complementarity
  conditions.
\end{enumerate}

Models constructed at all three levels are ultimately transformed into
\texttt{GEM} equilibrium models and solved through a common numerical
solution interface.

\texttt{GeneralEquilibriumModeling.jl}, \texttt{MPSGE.jl}, and
\texttt{JCGE.jl} can all be used to construct and solve general
equilibrium models, but their primary orientations differ.
\texttt{GeneralEquilibriumModeling.jl} takes economic agents, net supply
relations, and equilibrium conditions as its basic modeling objects. It
emphasizes the direct correspondence between economic theory and
numerical models and is intended primarily for learning and teaching
general equilibrium theory, while also supporting research on structural
equilibrium models and nonstandard equilibrium mechanisms.

\texttt{MPSGE.jl} is oriented primarily toward CGE modeling. It
emphasizes concise representations of calibrated economic structures
through production sectors, commodities, consumers, and nested
production and utility functions. \texttt{JCGE.jl}, by contrast, places
greater emphasis on the complete workflow of applied CGE modeling, using
modular blocks and standardized run specifications to support model
assembly, data calibration, scenario simulation, validation, and
comparison of results.

The \texttt{GE} package for R solves general equilibrium models more
directly on the basis of structural equilibrium models and structural
dynamic models. It simulates the period-by-period adjustment of prices,
activity levels, and resource allocation, allowing the economic system
to approach an equilibrium state gradually. Compared with solving a
mixed complementarity problem directly with PATH, this iterative
approach is generally less computationally efficient, but it makes the
dynamic process through which an economy moves from disequilibrium
toward equilibrium observable. It can therefore be useful for studying
the mechanism through which equilibrium is formed.

\bookmarksetup{startatroot}

\chapter{The GEMB High-Level Framework, the GEMB Low-Level Framework,
and the GEM
Framework}\label{the-gemb-high-level-framework-the-gemb-low-level-framework-and-the-gem-framework}

\texttt{GeneralEquilibriumModeling.jl} provides three interconnected
modeling frameworks: the GEM framework, the GEMB low-level framework,
and the GEMB high-level framework. The GEM framework allows users to
construct economic agents' net supply functions and their associated
complementarity conditions directly. The GEMB low-level framework
generates GEM economic agents through behavioral specifications and
agent builders. The GEMB high-level framework organizes complete
equilibrium models using commodity names, economic agent names, and
high-level model interfaces. In most cases, the GEMB high-level
framework is sufficient, and it is therefore the framework on which this
book primarily focuses.

This chapter introduces the basic modeling methods of
\texttt{GeneralEquilibriumModeling.jl} through some simple equilibrium
examples.

\section{\texorpdfstring{Basic Example: A \(2\times2\) Cobb--Douglas
Equilibrium with
Production}{Basic Example: A 2\textbackslash times2 Cobb--Douglas Equilibrium with Production}}\label{basic-example-a-2times2-cobbdouglas-equilibrium-with-production}

Assume the following:

\begin{enumerate}
\def\labelenumi{\arabic{enumi}.}
\item
  The economy contains two commodities, product and labor, and two
  economic agents, a firm and a consumer. Product is taken as the
  numeraire commodity.
\item
  The firm uses product and labor to produce product. Its production
  function is \(y=2x_1^{1/2}x_2^{1/2}\), where \(x_1\) and \(x_2\)
  denote the inputs of product and labor, respectively, and \(y\)
  denotes total product output.
\item
  The consumer is endowed with \(\omega>0\) units of labor. The utility
  function is \(u=x_{1}^{1/2}x_{2}^{1/2}\), where \(x_{1}\) and
  \(x_{2}\) denote the quantities of product and labor consumed by the
  consumer, respectively.
\end{enumerate}

Let the price vector be \(\mathbf p=(1,p_2)^\top\), the firm's activity
level be \(y\), and the consumer's activity level be the utility level
\(u\). The demand and supply matrices are

\[
\mathbf D(\mathbf p,u,y)
=
\begin{bmatrix}
\frac{y}{2}\sqrt{p_2} & u\sqrt{p_2}\\[8pt]
\frac{y}{2\sqrt{p_2}} & \frac{u}{\sqrt{p_2}}
\end{bmatrix}
\]

and

\[
\mathbf S(y)
=
\begin{bmatrix}
y & 0\\
0 & \omega
\end{bmatrix}
\]

Consider the following equality-form wide-sense structural equilibrium
model:

\[
\begin{aligned}
\mathbf p^\top\mathbf D(\mathbf p,u,y)
&=\mathbf p^\top\mathbf S(y)\\
\mathbf D(\mathbf p,u,y)\mathbf 1
&=\mathbf S(y)\mathbf 1
\end{aligned}
\]

The equilibrium price vector, firm activity level, and consumer utility
level are, respectively,

\[
\mathbf p^*=(1,1)^\top,\qquad
y^*=\omega,\qquad
u^*=\frac{\omega}{2}
\]

The equilibrium demand and supply matrices are

\[
\mathbf D^*
=
\begin{bmatrix}
\frac{\omega}{2} & \frac{\omega}{2}\\[6pt]
\frac{\omega}{2} & \frac{\omega}{2}
\end{bmatrix},
\qquad
\mathbf S^*
=
\begin{bmatrix}
\omega & 0\\
0 & \omega
\end{bmatrix}
\]

In particular, when \(\omega=100\),

\[
\mathbf p^*=(1,1)^\top,\qquad
y^*=100,\qquad
u^*=50
\]

In the next few sections, we will construct the basic example above
using the GEMB high-level framework, the GEMB low-level framework, and
the GEM framework, respectively.

\section{GEMB High-Level Framework}\label{gemb-high-level-framework}

The GEMB high-level framework is built on the GEMB low-level framework
and the GEM framework. Users specify a model using commodity names,
behavioral specifications, and declarations of economic agents. They do
not need to assign commodity indices manually or directly construct
\texttt{GEM.NetSupplyAgent} or \texttt{GEM.NetSupplyEquilibriumModel}
objects. GEMB progressively converts the high-level model into a GEM
equilibrium model and then calls the GEM solver.

The GEMB high-level framework mainly involves the following steps:

\begin{enumerate}
\def\labelenumi{\arabic{enumi}.}
\item
  Use \texttt{GEMBModel} to specify the set of commodities, the
  numeraire commodity, and price bounds.
\item
  Use \texttt{CESSpec}, \texttt{DCESSpec}, or other behavioral
  specifications to describe production and consumption behavior.
\item
  Use \texttt{add\_agent!} to add firms, consumers, and other economic
  agents.
\item
  Use \texttt{solve} to convert the high-level model into a GEM
  equilibrium model and solve it.
\end{enumerate}

In the GEMB low-level framework, commodities are represented by integer
indices. The GEMB high-level framework allows commodities to be referred
to directly by name. GEMB automatically converts commodity names into
the integer indices required by the underlying model.

During the solution process, GEMB sequentially performs commodity-name
resolution, agent construction, aggregation of net supplies, and
assembly of the GEM equilibrium model, and finally calls the solver
provided through GEM. The solution is returned using the common
equilibrium result type, from which information such as prices, activity
levels, agents' net supplies, and market residuals can be obtained.

The main feature of the GEMB high-level framework is that it organizes
models using economically meaningful commodity names and agent
declarations. It hides technical details such as commodity indices,
\texttt{NetSupplyAgent} construction, and low-level equilibrium-model
assembly. For problems that require custom net supply functions or
specialized complementarity conditions, users can instead work with the
GEMB low-level framework or the GEM framework.

The GEMB low-level framework provides reusable construction of economic
agents, whereas the GEMB high-level framework treats GEMBModel as a
one-shot model representation. For each economic agent, the high-level
model retains both the original economic specification and its resolved
commodity mappings, together with the constructed net-supply agent used
for equilibrium computation.

\begin{codelisting}

\caption{\label{lst-gemb_basicExample_CD_nc2_na2}A \(2\times2\)
Cobb--Douglas Equilibrium Model: The GEMB High-Level Framework}

\centering{

\begin{Shaded}
\begin{Highlighting}[]
\CommentTok{\# GEMB high{-}level framework: use commodity names and behavioral specifications.}
\ImportTok{using} \BuiltInTok{GeneralEquilibriumModeling.GEMB}

\NormalTok{alpha }\OperatorTok{=} \FloatTok{2.0}
\NormalTok{omega }\OperatorTok{=} \FloatTok{100.0}
\NormalTok{beta }\OperatorTok{=}\NormalTok{ [}\FloatTok{0.5}\NormalTok{, }\FloatTok{0.5}\NormalTok{]}

\NormalTok{model }\OperatorTok{=} \FunctionTok{GEMBModel}\NormalTok{(}
\NormalTok{    [}\OperatorTok{:}\NormalTok{product, }\OperatorTok{:}\NormalTok{labor];}
\NormalTok{    numeraire}\OperatorTok{=:}\NormalTok{product,}
\NormalTok{)}

\FunctionTok{add\_agent!}\NormalTok{(}
\NormalTok{    model,}
    \FunctionTok{CESSpec}\NormalTok{(beta; es}\OperatorTok{=}\FloatTok{1.0}\NormalTok{, alpha}\OperatorTok{=}\NormalTok{alpha);}
\NormalTok{    outputs}\OperatorTok{=:}\NormalTok{product,}
\NormalTok{    demands}\OperatorTok{=}\NormalTok{[}\OperatorTok{:}\NormalTok{product, }\OperatorTok{:}\NormalTok{labor],}
\NormalTok{    activity\_start}\OperatorTok{=}\FloatTok{100.0}\NormalTok{,}
\NormalTok{    name}\OperatorTok{=:}\NormalTok{firm,}
\NormalTok{)}

\FunctionTok{add\_agent!}\NormalTok{(}
\NormalTok{    model,}
    \FunctionTok{CESSpec}\NormalTok{(beta; es}\OperatorTok{=}\FloatTok{1.0}\NormalTok{, alpha}\OperatorTok{=}\FloatTok{1.0}\NormalTok{);}
\NormalTok{    demands}\OperatorTok{=}\NormalTok{[}\OperatorTok{:}\NormalTok{product, }\OperatorTok{:}\NormalTok{labor],}
\NormalTok{    endowments}\OperatorTok{=:}\NormalTok{labor,}
\NormalTok{    endowment\_quantities}\OperatorTok{=}\NormalTok{omega,}
\NormalTok{    activity\_start}\OperatorTok{=}\FloatTok{100.0}\NormalTok{,}
\NormalTok{    name}\OperatorTok{=:}\NormalTok{consumer,}
\NormalTok{)}

\NormalTok{result }\OperatorTok{=} \FunctionTok{solve}\NormalTok{(}
\NormalTok{    model;}
\NormalTok{    p0}\OperatorTok{=}\NormalTok{[}\FloatTok{1.0}\NormalTok{, }\FloatTok{2.0}\NormalTok{],}
\NormalTok{    residual\_tol}\OperatorTok{=}\FloatTok{1.0e{-}8}\NormalTok{,}
\NormalTok{    silent}\OperatorTok{=}\ConstantTok{true}\NormalTok{,}
\NormalTok{)}

\NormalTok{stats }\OperatorTok{=} \FunctionTok{equilibrium\_statistics}\NormalTok{(model, result)}

\NormalTok{output }\OperatorTok{=}\NormalTok{ stats.activity\_levels[}\FloatTok{1}\NormalTok{]}
\NormalTok{utility }\OperatorTok{=}\NormalTok{ stats.activity\_levels[}\FloatTok{2}\NormalTok{]}

\NormalTok{flows }\OperatorTok{=} \FunctionTok{demand\_supply\_matrices}\NormalTok{(model, result)}

\NormalTok{D }\OperatorTok{=}\NormalTok{ flows.demand}
\NormalTok{S }\OperatorTok{=}\NormalTok{ flows.supply}
\NormalTok{Sbar }\OperatorTok{=}\NormalTok{ flows.net\_supply}

\NormalTok{DV }\OperatorTok{=}\NormalTok{ flows.demand\_value}
\NormalTok{SV }\OperatorTok{=}\NormalTok{ flows.supply\_value}

\PreprocessorTok{@assert}\NormalTok{ S }\OperatorTok{{-}}\NormalTok{ D }\OperatorTok{≈}\NormalTok{ Sbar}
\PreprocessorTok{@assert}\NormalTok{ flows.total\_demand }\OperatorTok{≈}\NormalTok{ flows.total\_supply}
\PreprocessorTok{@assert}\NormalTok{ flows.agent\_expenditure }\OperatorTok{≈}\NormalTok{ flows.agent\_revenue}

\FunctionTok{print\_equilibrium\_statistics}\NormalTok{(model, result)}

\FunctionTok{println}\NormalTok{(}\StringTok{"}\SpecialCharTok{\textbackslash{}n}\StringTok{Demand matrix D:"}\NormalTok{)}
\FunctionTok{display}\NormalTok{(}\FunctionTok{hcat}\NormalTok{(D, flows.total\_demand))}

\FunctionTok{println}\NormalTok{(}\StringTok{"}\SpecialCharTok{\textbackslash{}n}\StringTok{Supply matrix S:"}\NormalTok{)}
\FunctionTok{display}\NormalTok{(}\FunctionTok{hcat}\NormalTok{(S, flows.total\_supply))}

\FunctionTok{println}\NormalTok{(}\StringTok{"}\SpecialCharTok{\textbackslash{}n}\StringTok{Demand value matrix DV:"}\NormalTok{)}
\FunctionTok{display}\NormalTok{([}
    \FunctionTok{hcat}\NormalTok{(DV, flows.total\_demand\_value)}
    \FunctionTok{hcat}\NormalTok{(}\FunctionTok{permutedims}\NormalTok{(flows.agent\_expenditure), }\FunctionTok{sum}\NormalTok{(flows.agent\_expenditure))}
\NormalTok{])}

\FunctionTok{println}\NormalTok{(}\StringTok{"}\SpecialCharTok{\textbackslash{}n}\StringTok{Supply value matrix SV:"}\NormalTok{)}
\FunctionTok{display}\NormalTok{([}
    \FunctionTok{hcat}\NormalTok{(SV, flows.total\_supply\_value)}
    \FunctionTok{hcat}\NormalTok{(}\FunctionTok{permutedims}\NormalTok{(flows.agent\_revenue), }\FunctionTok{sum}\NormalTok{(flows.agent\_revenue))}
\NormalTok{])}
\end{Highlighting}
\end{Shaded}

}

\end{codelisting}%

\section{GEMB Low-Level Framework}\label{gemb-low-level-framework}

The GEMB low-level framework lies between the GEM framework and the GEMB
high-level framework. The GEM framework requires users to define
economic agents' net supply functions and equilibrium conditions
directly. The GEMB low-level framework instead allows users to first
specify production, consumption, and other economic behaviors and then
use \texttt{build\_agent} to convert these behavioral specifications
into \texttt{GEM.NetSupplyAgent} objects.

Compared with the GEMB high-level framework, however, the GEMB low-level
framework still requires users to work relatively close to the
underlying GEM representation. In particular, input commodities, output
commodities, and endowments are generally specified by commodity
indices, and the resulting agents are assembled and solved using the GEM
framework.

The central idea of the GEMB low-level framework is to separate economic
modeling into two steps: \textbf{behavioral specification} and
\textbf{agent construction}.

A behavioral specification describes how an economic agent determines
its demand as prices and its activity level change. For example,
\texttt{ActivityDemandSpec} allows the user to provide a demand function
directly. Common specifications such as \texttt{CESSpec} and
\texttt{DCESSpec} provide more specialized representations for CES,
Cobb--Douglas, and Leontief demand structures.

After the behavioral specification has been defined,
\texttt{build\_agent} converts it into a GEM economic agent. During this
step, the user specifies such information as the relevant commodity
indices, outputs, endowments, activity variable, initial value, and the
condition rule associated with the agent.

For a firm, the activity variable usually represents its production
level. Given prices and an activity level, the behavioral specification
determines the firm's input demand. The firm's net supply is then
obtained by subtracting its input demand from its output supply. Under
constant returns to scale, input demand changes proportionally with the
production activity level.

GEMB can associate the firm's activity level with a unit
revenue-expenditure balance condition. Economically, this condition
compares the expenditure required for one unit of production activity
with the revenue generated by that activity. If unit expenditure exceeds
unit revenue, the firm cannot operate at a positive activity level. If
the firm operates at a positive activity level, unit expenditure and
unit revenue must be equal.

For a consumer, the activity variable usually represents utility. Given
prices and a utility level, the behavioral specification determines the
compensated demand required to attain that utility level. The consumer's
net supply is obtained by subtracting this consumption demand from the
consumer's endowment.

The consumer's utility level can be associated with a total
revenue-expenditure balance condition. This condition compares total
consumption expenditure with the income generated by the consumer's
endowment. When the consumer attains a positive utility level, total
expenditure equals endowment income.

The basic modeling procedure of the GEMB low-level framework is as
follows:

\begin{enumerate}
\def\labelenumi{\arabic{enumi}.}
\item
  Define the economic agent's behavioral specification using a type such
  as \texttt{ActivityDemandSpec}, \texttt{CESSpec}, or
  \texttt{DCESSpec}.
\item
  Use \texttt{build\_agent} to specify the agent's input commodities,
  output commodities, endowments, activity variable, initial value, and
  other relevant settings, and convert the behavioral specification into
  a \texttt{GEM.NetSupplyAgent}.
\item
  Use \texttt{GEM.NetSupplyEquilibriumModel} to assemble the resulting
  economic agents and commodity markets, and then call a GEM solver to
  solve the equilibrium model.
\end{enumerate}

In the GEM framework and the GEMB low-level framework, the user may
explicitly map commodity names to commodity indices before passing them
to \texttt{NetSupplyAgent} or \texttt{build\_agent}, respectively. In
the GEMB high-level framework, \texttt{GEMBModel} performs this mapping
automatically and also assembles the resulting agents into the
equilibrium model.

The GEMB low-level framework therefore retains relatively direct control
over commodity indices, economic agents, activity variables, and
equilibrium conditions, while reducing the need to write net supply
functions manually. It is useful when users want more control than the
GEMB high-level framework provides but do not want to construct every
net supply function and complementarity condition directly in GEM.

\begin{codelisting}

\caption{\label{lst-gemb_basicExample_CD_lowLevel_nc2_na2}A \(2\times2\)
Cobb--Douglas Equilibrium Model: The GEMB Low-Level Framework}

\centering{

\begin{Shaded}
\begin{Highlighting}[]
\CommentTok{\# GEMB low{-}level framework: build NetSupplyAgents from ActivityDemandSpec.}
\ImportTok{using} \BuiltInTok{GeneralEquilibriumModeling}

\NormalTok{alpha }\OperatorTok{=} \FloatTok{2.0}
\NormalTok{omega }\OperatorTok{=} \FloatTok{100.0}
\NormalTok{beta }\OperatorTok{=}\NormalTok{ [}\FloatTok{0.5}\NormalTok{, }\FloatTok{0.5}\NormalTok{]}

\NormalTok{c }\OperatorTok{=}\NormalTok{ (product}\OperatorTok{=}\FloatTok{1}\NormalTok{, labor}\OperatorTok{=}\FloatTok{2}\NormalTok{)}

\NormalTok{firm }\OperatorTok{=}\NormalTok{ GEMB.}\FunctionTok{build\_agent}\NormalTok{(}
\NormalTok{    GEMB.}\FunctionTok{ActivityDemandSpec}\NormalTok{(}
\NormalTok{        (y, p) }\OperatorTok{{-}\textgreater{}}\NormalTok{ GEMB.}\FunctionTok{CES\_input}\NormalTok{(beta, y, p; es}\OperatorTok{=}\FloatTok{1.0}\NormalTok{, alpha}\OperatorTok{=}\NormalTok{alpha),}
\NormalTok{    );}
\NormalTok{    output\_indices}\OperatorTok{=}\NormalTok{[c.product],}
\NormalTok{    output\_coefficients}\OperatorTok{=}\NormalTok{[}\FloatTok{1.0}\NormalTok{],}
\NormalTok{    demand\_indices}\OperatorTok{=}\NormalTok{[c.product, c.labor],}
\NormalTok{    activity\_start}\OperatorTok{=}\FloatTok{100.0}\NormalTok{,}
\NormalTok{    name}\OperatorTok{=:}\NormalTok{firm,}
\NormalTok{)}

\NormalTok{consumer }\OperatorTok{=}\NormalTok{ GEMB.}\FunctionTok{build\_agent}\NormalTok{(}
\NormalTok{    GEMB.}\FunctionTok{ActivityDemandSpec}\NormalTok{(}
\NormalTok{        (u, p) }\OperatorTok{{-}\textgreater{}}\NormalTok{ GEMB.}\FunctionTok{CES\_input}\NormalTok{(beta, u, p; es}\OperatorTok{=}\FloatTok{1.0}\NormalTok{, alpha}\OperatorTok{=}\FloatTok{1.0}\NormalTok{),}
\NormalTok{    );}
\NormalTok{    demand\_indices}\OperatorTok{=}\NormalTok{[c.product, c.labor],}
\NormalTok{    endowment\_indices}\OperatorTok{=}\NormalTok{[c.labor],}
\NormalTok{    endowment\_quantities}\OperatorTok{=}\NormalTok{[omega],}
\NormalTok{    activity\_start}\OperatorTok{=}\FloatTok{100.0}\NormalTok{,}
\NormalTok{    name}\OperatorTok{=:}\NormalTok{consumer,}
\NormalTok{)}

\NormalTok{model }\OperatorTok{=}\NormalTok{ GEM.}\FunctionTok{NetSupplyEquilibriumModel}\NormalTok{(}
\NormalTok{    GEM.AbstractNetSupplyAgent[firm, consumer],}
\NormalTok{    [}\OperatorTok{:}\NormalTok{product, }\OperatorTok{:}\NormalTok{labor];}
\NormalTok{    numeraire\_index}\OperatorTok{=}\NormalTok{c.product,}
\NormalTok{    numeraire\_value}\OperatorTok{=}\FloatTok{1.0}\NormalTok{,}
\NormalTok{    price\_lower\_bounds}\OperatorTok{=}\NormalTok{[}\FloatTok{1.0e{-}10}\NormalTok{, }\FloatTok{1.0e{-}10}\NormalTok{],}
\NormalTok{)}

\NormalTok{result }\OperatorTok{=}\NormalTok{ GEM.}\FunctionTok{solve\_equilibrium\_model\_mcp\_jump}\NormalTok{(}
\NormalTok{    model;}
\NormalTok{    p0}\OperatorTok{=}\NormalTok{[}\FloatTok{1.0}\NormalTok{, }\FloatTok{2.0}\NormalTok{],}
\NormalTok{    residual\_tol}\OperatorTok{=}\FloatTok{1.0e{-}8}\NormalTok{,}
\NormalTok{    silent}\OperatorTok{=}\ConstantTok{true}\NormalTok{,}
\NormalTok{)}

\NormalTok{output }\OperatorTok{=}\NormalTok{ result.agent\_variable\_values[}\FloatTok{1}\NormalTok{][}\FloatTok{1}\NormalTok{]}
\NormalTok{utility }\OperatorTok{=}\NormalTok{ result.agent\_variable\_values[}\FloatTok{2}\NormalTok{][}\FloatTok{1}\NormalTok{]}

\FunctionTok{println}\NormalTok{(}\StringTok{"Prices = "}\NormalTok{, result.prices)}
\FunctionTok{println}\NormalTok{(}\StringTok{"Output = "}\NormalTok{, output)}
\FunctionTok{println}\NormalTok{(}\StringTok{"Utility = "}\NormalTok{, utility)}
\end{Highlighting}
\end{Shaded}

}

\end{codelisting}%

\section{GEM Framework}\label{gem-framework}

The GEM framework is the foundational modeling and solution framework of
\texttt{GeneralEquilibriumModeling.jl}. It takes the net supply
functions of economic agents as its basic modeling objects and
transforms agent behavioral conditions, commodity market-clearing
conditions, and other auxiliary conditions into a unified mixed
complementarity problem. The resulting problem is constructed with
\texttt{JuMP} and solved numerically using \texttt{PATHSolver.jl}.
Models constructed through the GEMB low-level framework and the GEMB
high-level framework are also ultimately converted into GEM equilibrium
models.

\subsection{Net Supply Representation}\label{net-supply-representation}

In a wide-sense structural equilibrium model, each row of the demand
matrix and the supply matrix corresponds to a commodity, and each column
corresponds to an economic agent. Let the demand matrix and the supply
matrix be denoted by \(\mathbf D(\cdot)\) and \(\mathbf S(\cdot)\),
respectively. The \textbf{net supply matrix} is defined as

\[
\bar{\mathbf S}(\cdot)
:=
\mathbf S(\cdot)-\mathbf D(\cdot)
\]

The element \(\bar S_{ij}\) represents the net supply of commodity \(i\)
by economic agent \(j\). A positive value represents net supply, whereas
a negative value represents net demand. The net supply matrix therefore
combines the demand matrix and the supply matrix into a single
representation.

Let \(\mathbf D_{\cdot j}\) and \(\mathbf S_{\cdot j}\) denote the
demand vector and supply vector of economic agent \(j\), respectively.
Column \(j\) of the net supply matrix is the net supply vector of
economic agent \(j\):

\[
\bar{\mathbf s}_j
:=
\mathbf S_{\cdot j}-\mathbf D_{\cdot j}
\]

Thus,

\[
\bar{\mathbf S}
=
\begin{bmatrix}
\bar{\mathbf s}_1 &
\bar{\mathbf s}_2 &
\cdots &
\bar{\mathbf s}_m
\end{bmatrix}
\]

where \(\bar{\mathbf s}_j\) describes the net supplies of all
commodities by economic agent \(j\).

Using the net supply matrix, the mixed-form wide-sense structural
equilibrium model can be written compactly as

\[
\mathbf p^\top\bar{\mathbf S}(\cdot)=\mathbf 0
\]

\[
\bar{\mathbf S}(\cdot)\mathbf 1\geq\mathbf 0
\]

The first condition states that every economic agent satisfies
revenue-expenditure balance. Since column \(j\) of \(\bar{\mathbf S}\)
is \(\bar{\mathbf s}_j\), it is equivalent to

\[
\mathbf p^\top\bar{\mathbf s}_j=0,
\qquad
j=1,\ldots,m
\]

The second condition states that the total net supply of every commodity
is nonnegative, or equivalently, that total supply is no less than total
demand.

The total net supply vector of all economic agents is

\[
\bar{\mathbf s}
:=
\sum_{j=1}^{m}\bar{\mathbf s}_j
=
\bar{\mathbf S}\mathbf 1
\]

so the commodity-side condition can also be written as

\[
\bar{\mathbf s}\geq\mathbf 0
\]

When free commodities and excess supply are allowed, the commodity
market condition can be expressed more generally as

\[
\mathbf 0
\leq
\bar{\mathbf S}(\cdot)\mathbf 1
\perp
\mathbf p
\geq
\mathbf 0
\]

or equivalently,

\[
\mathbf 0
\leq
\bar{\mathbf s}
\perp
\mathbf p
\geq
\mathbf 0
\]

That is,

\[
\bar{\mathbf s}\geq\mathbf 0,\qquad
\mathbf p\geq\mathbf 0,\qquad
\mathbf p^\top\bar{\mathbf s}=0
\]

If the price of commodity \(i\) is positive, its market must clear. If
commodity \(i\) has positive excess supply, its equilibrium price must
be zero. When all equilibrium commodity prices are positive, the market
condition reduces to

\[
\bar{\mathbf s}=\mathbf 0
\]

For an economic agent with a constant-returns-to-scale production
structure, let \(z_j\) denote its activity level and let
\(\bar{\mathbf s}_j^{u}\) denote its net supply vector per unit of
activity. Then

\[
\bar{\mathbf s}_j
=
z_j\bar{\mathbf s}_j^{u}
\]

When \(z_j>0\), the unit revenue-expenditure balance condition is

\[
\mathbf p^\top\bar{\mathbf s}_j^{u}=0
\]

If the economic agent is allowed to be inactive, the more general unit
revenue-expenditure balance condition is

\[
0
\leq
-\mathbf p^\top\bar{\mathbf s}_j^{u}
\perp
z_j
\geq
0
\]

where \(-\mathbf p^\top\bar{\mathbf s}_j^{u}\) is expenditure minus
revenue per unit of activity, that is, the unit loss or unit excess
cost.

The net supply representation therefore does not change the economic
meaning of the wide-sense structural equilibrium model. It simply
combines demand and supply into a single matrix. Because the GEM
framework directly computes the net supply of each economic agent, it
generally does not need to construct the complete demand matrix and
supply matrix explicitly.

\subsection{Economic Agents}\label{economic-agents}

The GEM framework uses \texttt{NetSupplyAgent} to represent an economic
agent. Each economic agent mainly contains the following information:

\begin{enumerate}
\def\labelenumi{\arabic{enumi}.}
\item
  The commodities associated with the agent.
\item
  The agent's endogenous variables, together with their lower bounds,
  upper bounds, and initial values.
\item
  A function that calculates net supply from the agent variables,
  commodity prices, and observed variables.
\item
  Complementarity conditions paired with the agent variables.
\item
  Names used to identify the agent and its variables.
\end{enumerate}

Let \(\mathbf v_j\) denote the local variable vector of economic agent
\(j\), \(\mathbf p_j\) the local price vector for the commodities
associated with that agent, and \(\mathbf o_j\) the vector of observed
variables. The agent's net supply function can then be written generally
as

\[
\bar{\mathbf s}_j
=
\bar{\mathbf s}_j(\mathbf v_j,\mathbf p_j,\mathbf o_j)
\]

Local variables belong only to the economic agent itself, whereas
observed variables may refer to variables of other agents, commodity
prices, or auxiliary variables. This mechanism makes it possible to
represent inter-agent relationships, endogenous tax rates, endogenous
interest rates, and other economic relationships that must be determined
jointly within the equilibrium model.

\texttt{NetSupplyAgent} is the unified agent representation used by the
GEM framework. Firms, consumers, and condition agents may have different
economic interpretations, but numerically they can all be represented as
combinations of local variables, net supply functions, and local
complementarity conditions.

\subsection{Agent Conditions}\label{agent-conditions}

For economic agent \(j\), let the condition function paired with the
local variable vector \(\mathbf v_j\) be

\[
\mathbf c_j
=
\mathbf c_j(
\mathbf v_j,
\mathbf p_j,
\bar{\mathbf s}_j,
\mathbf o_j
)
\]

If the lower and upper bounds of variable \(v_{jk}\) are
\(\underline v_{jk}\) and \(\overline v_{jk}\), respectively, the
corresponding mixed complementarity condition is

\[
c_{jk}
\perp
\underline v_{jk}
\leq
v_{jk}
\leq
\overline v_{jk}
\]

which means

\[
\begin{cases}
c_{jk}\geq0,
&v_{jk}=\underline v_{jk}\\
c_{jk}=0,
&\underline v_{jk}<v_{jk}<\overline v_{jk}\\
c_{jk}\leq0,
&v_{jk}=\overline v_{jk}
\end{cases}
\]

The GEM framework provides several commonly used condition rules:

\begin{enumerate}
\def\labelenumi{\arabic{enumi}.}
\item
  \texttt{UnitRevenueExpenditureBalanceConditions} pairs unit loss with
  the agent's activity variable, which typically represents a production
  activity level. It is mainly applicable to constant-returns-to-scale
  activities for which unit net supply can be defined.
\item
  \texttt{TotalRevenueExpenditureBalanceConditions} pairs total loss
  with the agent's activity variable, which may, for example, represent
  a consumer's utility level.
\item
  \texttt{MarginalUtilityConsumerConditions} constructs the KKT
  conditions for a consumer's utility-maximization problem using
  marginal utilities, commodity prices, and a budget multiplier.
\item
  \texttt{ProductionStationarityConditions} constructs stationary
  production conditions from a production function and marginal product
  functions. A stationary point obtained from these conditions is not
  necessarily a cost-minimizing point.
\item
  \texttt{CostMinimizationKKTConditions} constructs the KKT conditions
  for a firm's cost-minimization problem.
\item
  \texttt{ExplicitAgentConditions} allows the user to define an agent's
  condition function directly when the required relationship cannot be
  represented by one of the existing condition rules.
\end{enumerate}

The GEM framework therefore supports both standard condition rules and
user-defined complementarity conditions.

\subsection{Commodity Market
Conditions}\label{commodity-market-conditions}

The GEM framework maps the local net supply vectors of all economic
agents into a common commodity space and then aggregates them. The total
net supply of commodity \(i\) is

\[
\bar s_i
=
\sum_{j=1}^{m}\bar s_{ij}
\]

When the lower bound of the commodity price is zero and there is no
finite upper bound, the commodity market condition is

\[
0
\leq
\bar s_i
\perp
p_i
\geq
0
\]

If \(p_i>0\), then \(\bar s_i=0\), so the market for commodity \(i\)
clears. If \(\bar s_i>0\), then \(p_i=0\), so a commodity in excess
supply may become a free commodity.

The GEM framework also allows other lower and upper bounds to be imposed
on prices. Once a numeraire commodity is selected, its price is fixed at
a specified value, thereby eliminating the proportional indeterminacy of
the general equilibrium price vector.

\subsection{Auxiliary Variables and Auxiliary
Equations}\label{auxiliary-variables-and-auxiliary-equations}

In addition to economic agent variables and commodity prices, the GEM
framework allows auxiliary variables and auxiliary equations to be
introduced through \texttt{AuxiliaryVariable} and
\texttt{AuxiliaryEquation}.

Let \(q_k\) denote an auxiliary variable. A corresponding auxiliary
equation may be written as

\[
g_k(\mathbf p,\mathbf v,\mathbf q)=0
\]

Auxiliary variables can be used to represent exogenous price conditions,
activity-level normalization conditions, and other model-closure
relationships. Economic agents may observe auxiliary variables through
\texttt{AuxiliaryVariableRef}, variables of other agents through
\texttt{AgentVariableRef}, and commodity prices through
\texttt{PriceVariableRef}.

\subsection{Equilibrium Model and
Solution}\label{equilibrium-model-and-solution}

After the economic agents have been constructed, they can be combined
into a complete equilibrium model using
\texttt{NetSupplyEquilibriumModel}. The model mainly contains the set of
economic agents, the set of commodities, price bounds, the numeraire,
and optional auxiliary variables and auxiliary equations.

The GEM framework combines the following components into a single mixed
complementarity problem:

\begin{enumerate}
\def\labelenumi{\arabic{enumi}.}
\item
  The local complementarity conditions of the economic agents.
\item
  The market-clearing conditions for the commodities.
\item
  The auxiliary equations associated with auxiliary variables.
\item
  The lower and upper bound conditions for prices and other equilibrium
  variables.
\end{enumerate}

The model can be solved using
\texttt{solve\_equilibrium\_model\_mcp\_jump}. This function constructs
the mixed complementarity problem with \texttt{JuMP} and calls
\texttt{PATHSolver.jl} to obtain a numerical solution.

The solution is returned as an \texttt{EquilibriumResult}, which
contains information such as equilibrium prices, agent variable values,
agent net supplies, total net supply, auxiliary variable values, and
complementarity residuals.

\begin{codelisting}

\caption{\label{lst-gem_basicExample_CD_nc2_na2}A \(2\times2\)
Cobb--Douglas Equilibrium Model: The GEM Framework}

\centering{

\begin{Shaded}
\begin{Highlighting}[]
\CommentTok{\# GEM framework: define net supplies and condition rules directly.}
\ImportTok{using} \BuiltInTok{GeneralEquilibriumModeling.GEM}

\NormalTok{alpha }\OperatorTok{=} \FloatTok{2.0}
\NormalTok{omega }\OperatorTok{=} \FloatTok{100.0}

\NormalTok{c }\OperatorTok{=}\NormalTok{ (product}\OperatorTok{=}\FloatTok{1}\NormalTok{, labor}\OperatorTok{=}\FloatTok{2}\NormalTok{)}

\NormalTok{firm }\OperatorTok{=}\NormalTok{ GEM.}\FunctionTok{NetSupplyAgent}\NormalTok{(}
\NormalTok{    [c.product, c.labor],}
\NormalTok{    (v, p) }\OperatorTok{{-}\textgreater{}} \ControlFlowTok{begin}
\NormalTok{        y }\OperatorTok{=}\NormalTok{ v[}\FloatTok{1}\NormalTok{]}
\NormalTok{        p\_product, p\_labor }\OperatorTok{=}\NormalTok{ p}
\NormalTok{        r }\OperatorTok{=} \FunctionTok{sqrt}\NormalTok{(p\_labor }\OperatorTok{/}\NormalTok{ p\_product)}
\NormalTok{        x }\OperatorTok{=}\NormalTok{ (y }\OperatorTok{/}\NormalTok{ alpha) }\OperatorTok{.*}\NormalTok{ [r, }\FloatTok{1} \OperatorTok{/}\NormalTok{ r]}
\NormalTok{        [y }\OperatorTok{{-}}\NormalTok{ x[}\FloatTok{1}\NormalTok{], }\OperatorTok{{-}}\NormalTok{x[}\FloatTok{2}\NormalTok{]]}
    \ControlFlowTok{end}\NormalTok{;}
\NormalTok{    variable\_names}\OperatorTok{=}\NormalTok{[}\OperatorTok{:}\NormalTok{output],}
\NormalTok{    variable\_start}\OperatorTok{=}\NormalTok{[}\FloatTok{100.0}\NormalTok{],}
\NormalTok{    condition\_rule}\OperatorTok{=}\NormalTok{GEM.}\FunctionTok{UnitRevenueExpenditureBalanceConditions}\NormalTok{(),}
\NormalTok{    name}\OperatorTok{=:}\NormalTok{firm,}
\NormalTok{)}

\NormalTok{consumer }\OperatorTok{=}\NormalTok{ GEM.}\FunctionTok{NetSupplyAgent}\NormalTok{(}
\NormalTok{    [c.product, c.labor],}
\NormalTok{    (v, p) }\OperatorTok{{-}\textgreater{}} \ControlFlowTok{begin}
\NormalTok{        u }\OperatorTok{=}\NormalTok{ v[}\FloatTok{1}\NormalTok{]}
\NormalTok{        p\_product, p\_labor }\OperatorTok{=}\NormalTok{ p}
\NormalTok{        r }\OperatorTok{=} \FunctionTok{sqrt}\NormalTok{(p\_labor }\OperatorTok{/}\NormalTok{ p\_product)}
\NormalTok{        x }\OperatorTok{=}\NormalTok{ u }\OperatorTok{.*}\NormalTok{ [r, }\FloatTok{1} \OperatorTok{/}\NormalTok{ r]}
\NormalTok{        [}\OperatorTok{{-}}\NormalTok{x[}\FloatTok{1}\NormalTok{], omega }\OperatorTok{{-}}\NormalTok{ x[}\FloatTok{2}\NormalTok{]]}
    \ControlFlowTok{end}\NormalTok{;}
\NormalTok{    variable\_names}\OperatorTok{=}\NormalTok{[}\OperatorTok{:}\NormalTok{utility],}
\NormalTok{    variable\_start}\OperatorTok{=}\NormalTok{[}\FloatTok{100.0}\NormalTok{],}
\NormalTok{    condition\_rule}\OperatorTok{=}\NormalTok{GEM.}\FunctionTok{TotalRevenueExpenditureBalanceConditions}\NormalTok{(),}
\NormalTok{    name}\OperatorTok{=:}\NormalTok{consumer,}
\NormalTok{)}

\NormalTok{model }\OperatorTok{=}\NormalTok{ GEM.}\FunctionTok{NetSupplyEquilibriumModel}\NormalTok{(}
\NormalTok{    [firm, consumer],}
\NormalTok{    [}\OperatorTok{:}\NormalTok{product, }\OperatorTok{:}\NormalTok{labor];}
\NormalTok{    numeraire\_index}\OperatorTok{=}\NormalTok{c.product,}
\NormalTok{    numeraire\_value}\OperatorTok{=}\FloatTok{1.0}\NormalTok{,}
\NormalTok{    price\_lower\_bounds}\OperatorTok{=}\NormalTok{[}\FloatTok{1.0e{-}10}\NormalTok{, }\FloatTok{1.0e{-}10}\NormalTok{],}
\NormalTok{)}

\NormalTok{result }\OperatorTok{=}\NormalTok{ GEM.}\FunctionTok{solve\_equilibrium\_model\_mcp\_jump}\NormalTok{(}
\NormalTok{    model;}
\NormalTok{    p0}\OperatorTok{=}\NormalTok{[}\FloatTok{1.0}\NormalTok{, }\FloatTok{2.0}\NormalTok{],}
\NormalTok{    residual\_tol}\OperatorTok{=}\FloatTok{1.0e{-}8}\NormalTok{,}
\NormalTok{    silent}\OperatorTok{=}\ConstantTok{true}\NormalTok{,}
\NormalTok{)}

\NormalTok{output }\OperatorTok{=}\NormalTok{ result.agent\_variable\_values[}\FloatTok{1}\NormalTok{][}\FloatTok{1}\NormalTok{]}
\NormalTok{utility }\OperatorTok{=}\NormalTok{ result.agent\_variable\_values[}\FloatTok{2}\NormalTok{][}\FloatTok{1}\NormalTok{]}

\FunctionTok{println}\NormalTok{(}\StringTok{"Prices = "}\NormalTok{, result.prices)}
\FunctionTok{println}\NormalTok{(}\StringTok{"Output = "}\NormalTok{, output, }\StringTok{", utility = "}\NormalTok{, utility)}
\end{Highlighting}
\end{Shaded}

}

\end{codelisting}%

\section{Total Revenue-Expenditure Balance Conditions (TREBC) and
Wide-Sense
Equilibrium}\label{total-revenue-expenditure-balance-conditions-trebc-and-wide-sense-equilibrium}

In a strict-sense structural equilibrium model with constant returns to
scale, a firm's unit loss and activity level satisfy a complementarity
relation. This condition rules out the possibility that a firm remains
inactive when a profitable production opportunity exists.

In GEM, \texttt{TotalRevenueExpenditureBalanceConditions} constructs the
relevant condition using the firm's total loss rather than its unit
loss. As a result, \texttt{TotalRevenueExpenditureBalanceConditions}
permits a broader notion of equilibrium: a firm with constant returns to
scale may remain inactive even when a positive-profit opportunity
exists.

This point can be illustrated by a slight extension of the basic
example. Suppose the economy contains two commodities, product and
labor, with product as the numeraire commodity. There are two firms,
both using Cobb--Douglas production technologies with input shares
\((0.5,0.5)^\top\).

Firm 1 has efficiency parameter \(\alpha_1=2\), while firm 2 has
efficiency parameter \(\alpha_2=1\). The consumer is endowed with
\(\omega=100\) units of labor.

If the usual unit-loss condition is used, firm 1 is more efficient.
Therefore, at equilibrium, firm 1 produces while firm 2 remains
inactive. This equilibrium can be regarded as a strict-sense equilibrium
because no inactive firm has an unexploited positive-profit opportunity.

If \texttt{TotalRevenueExpenditureBalanceConditions} is used instead,
then, in addition to the strict-sense equilibrium above, there is
another equilibrium:

\[
\mathbf p=\left(1,\frac14\right)^\top,\qquad
z_1=0,\qquad
z_2=25, \qquad
u=25
\]

At this equilibrium, firm 1 could earn a positive profit if it entered
production. However, because its activity level is zero, its total loss
is also zero. \texttt{TotalRevenueExpenditureBalanceConditions}
therefore allows firm 1 to remain inactive.

This equilibrium satisfies the total revenue-expenditure balance
conditions and the market-clearing conditions, but it does not satisfy
the unit revenue-expenditure balance conditions of the strict-sense
structural equilibrium model. It can therefore be called a
\textbf{wide-sense equilibrium}, whereas an equilibrium satisfying the
conditions of the strict-sense structural equilibrium model can be
called a \textbf{strict-sense equilibrium}.

Because wide-sense equilibria may involve multiple solutions and
degenerate cases, the solution returned by PATH may be relatively
sensitive to the initial values.

For some production technologies, particularly nonconvex technologies,
the firm's net supply cannot be represented naturally as an activity
level multiplied by an activity-independent unit net-supply vector. In
such cases, unit revenue-expenditure balance conditions may no longer
provide an appropriate representation, whereas total revenue-expenditure
balance conditions can be imposed directly on the firm's total revenue
and total expenditure at its current activity level.

The program for solving the wide-sense equilibrium described above is as
follows.

\begin{codelisting}

\caption{\label{lst-gemb_basicExample_CD_totalRevenueExpenditureBalance_wideSense_nc2_na3}A
\(2\times3\) Cobb--Douglas Equilibrium Model: Total Revenue-Expenditure
Balance Conditions}

\centering{

\begin{Shaded}
\begin{Highlighting}[]
\CommentTok{\# GEMB high{-}level framework: total revenue{-}expenditure balance conditions.}
\ImportTok{using} \BuiltInTok{GeneralEquilibriumModeling}

\NormalTok{alpha }\OperatorTok{=} \FloatTok{2.0}
\NormalTok{omega }\OperatorTok{=} \FloatTok{100.0}
\NormalTok{beta }\OperatorTok{=}\NormalTok{ [}\FloatTok{0.5}\NormalTok{, }\FloatTok{0.5}\NormalTok{]}

\NormalTok{model }\OperatorTok{=}\NormalTok{ GEMB.}\FunctionTok{GEMBModel}\NormalTok{(}
\NormalTok{    [}\OperatorTok{:}\NormalTok{product, }\OperatorTok{:}\NormalTok{labor];}
\NormalTok{    numeraire}\OperatorTok{=:}\NormalTok{product,}
\NormalTok{    numeraire\_value}\OperatorTok{=}\FloatTok{1.0}\NormalTok{,}
\NormalTok{)}

\NormalTok{GEMB.}\FunctionTok{add\_agent!}\NormalTok{(}
\NormalTok{    model,}
\NormalTok{    GEMB.}\FunctionTok{CESSpec}\NormalTok{(beta; es}\OperatorTok{=}\FloatTok{1.0}\NormalTok{, alpha}\OperatorTok{=}\NormalTok{alpha);}
\NormalTok{    outputs}\OperatorTok{=:}\NormalTok{product,}
\NormalTok{    demands}\OperatorTok{=}\NormalTok{[}\OperatorTok{:}\NormalTok{product, }\OperatorTok{:}\NormalTok{labor],}
\NormalTok{    activity\_start}\OperatorTok{=}\FloatTok{0.0}\NormalTok{,}
\NormalTok{    condition\_rule}\OperatorTok{=}\NormalTok{GEM.}\FunctionTok{TotalRevenueExpenditureBalanceConditions}\NormalTok{(),}
\NormalTok{    name}\OperatorTok{=:}\NormalTok{firm1,}
\NormalTok{)}

\NormalTok{GEMB.}\FunctionTok{add\_agent!}\NormalTok{(}
\NormalTok{    model,}
\NormalTok{    GEMB.}\FunctionTok{CESSpec}\NormalTok{(beta; es}\OperatorTok{=}\FloatTok{1.0}\NormalTok{, alpha}\OperatorTok{=}\FloatTok{1.0}\NormalTok{);}
\NormalTok{    demands}\OperatorTok{=}\NormalTok{[}\OperatorTok{:}\NormalTok{product, }\OperatorTok{:}\NormalTok{labor],}
\NormalTok{    endowments}\OperatorTok{=:}\NormalTok{labor,}
\NormalTok{    endowment\_quantities}\OperatorTok{=}\NormalTok{omega,}
\NormalTok{    name}\OperatorTok{=:}\NormalTok{consumer,}
\NormalTok{)}

\NormalTok{GEMB.}\FunctionTok{add\_agent!}\NormalTok{(}
\NormalTok{    model,}
\NormalTok{    GEMB.}\FunctionTok{CESSpec}\NormalTok{(beta; es}\OperatorTok{=}\FloatTok{1.0}\NormalTok{, alpha}\OperatorTok{=}\FloatTok{1.0}\NormalTok{);}
\NormalTok{    outputs}\OperatorTok{=:}\NormalTok{product,}
\NormalTok{    demands}\OperatorTok{=}\NormalTok{[}\OperatorTok{:}\NormalTok{product, }\OperatorTok{:}\NormalTok{labor],}
\NormalTok{    activity\_start}\OperatorTok{=}\FloatTok{1.0}\NormalTok{,}
\NormalTok{    condition\_rule}\OperatorTok{=}\NormalTok{GEM.}\FunctionTok{TotalRevenueExpenditureBalanceConditions}\NormalTok{(),}
\NormalTok{    name}\OperatorTok{=:}\NormalTok{firm2,}
\NormalTok{)}

\NormalTok{result }\OperatorTok{=}\NormalTok{ GEMB.}\FunctionTok{solve}\NormalTok{(}
\NormalTok{    model;}
\NormalTok{    p0}\OperatorTok{=}\NormalTok{[}\FloatTok{1.0}\NormalTok{, }\FloatTok{0.27}\NormalTok{],}
\NormalTok{    residual\_tol}\OperatorTok{=}\FloatTok{1.0e{-}8}\NormalTok{,}
\NormalTok{    silent}\OperatorTok{=}\ConstantTok{true}\NormalTok{,}
\NormalTok{)}

\NormalTok{stats }\OperatorTok{=}\NormalTok{ GEMB.}\FunctionTok{equilibrium\_statistics}\NormalTok{(model, result)}

\NormalTok{output1 }\OperatorTok{=}\NormalTok{ stats.activity\_levels[}\FloatTok{1}\NormalTok{]}
\NormalTok{utility }\OperatorTok{=}\NormalTok{ stats.activity\_levels[}\FloatTok{2}\NormalTok{]}
\NormalTok{output2 }\OperatorTok{=}\NormalTok{ stats.activity\_levels[}\FloatTok{3}\NormalTok{]}

\NormalTok{GEMB.}\FunctionTok{print\_equilibrium\_statistics}\NormalTok{(model, result)}

\FunctionTok{println}\NormalTok{(}\StringTok{"Output1 = "}\NormalTok{, output1,}
        \StringTok{", output2 = "}\NormalTok{, output2,}
        \StringTok{", utility = "}\NormalTok{, utility)}
\end{Highlighting}
\end{Shaded}

}

\end{codelisting}%

\section{Cost-Minimization Conditions and Marginal-Utility
Conditions}\label{cost-minimization-conditions-and-marginal-utility-conditions}

A general equilibrium model does not have to rely only on demand
functions derived in advance. The first-order conditions of firms' and
consumers' optimization problems can also be incorporated directly into
the equilibrium system. GEM provides
\texttt{CostMinimizationKKTConditions} and
\texttt{MarginalUtilityConsumerConditions} for constructing the
complementarity conditions associated with firms' cost-minimization
problems and consumers' utility-maximization problems, respectively.
This approach does not require input demand functions or consumer demand
functions to be derived beforehand, and it can naturally accommodate
corner solutions in which some input or consumption quantities are zero.

In GEM, \texttt{CostMinimizationKKTConditions} pairs the firm's activity
level with its total loss. It therefore incorporates three types of
relations: the firm's revenue-expenditure balance condition, the
cost-minimization conditions for individual inputs, and the production
constraint. The corresponding endogenous variables are, respectively,
the activity level, the quantities of the various inputs, and the
multiplier associated with the production constraint.

It should be noted that KKT conditions are, in general, first-order
necessary conditions for a local optimum. For cost-minimization problems
with the usual appropriate convexity properties, they can characterize a
cost-minimizing solution. For nonconvex production technologies,
however, satisfying the KKT conditions alone does not necessarily imply
that a global cost minimum has been obtained.

\texttt{MarginalUtilityConsumerConditions} constructs
utility-maximization conditions from consumers' marginal utilities. It
does not require the demand functions to be derived in advance. Instead,
the consumption quantity of each commodity is paired directly with its
corresponding marginal-utility condition, while the budget multiplier is
paired with the budget constraint.

The following basic example uses cost-minimization conditions and
marginal-utility conditions.

\begin{codelisting}

\caption{\label{lst-gemb_basicExample_CD_costMinimization_marginalUtility_nc2_na2}A
\(2\times2\) Cobb--Douglas Equilibrium Model: Cost Minimization and
Marginal Utility}

\centering{

\begin{Shaded}
\begin{Highlighting}[]
\CommentTok{\# GEMB high{-}level framework:}
\CommentTok{\# firm uses CostMinimizationKKTConditions;}
\CommentTok{\# consumer uses MarginalUtilityConsumerConditions.}
\ImportTok{using} \BuiltInTok{GeneralEquilibriumModeling}

\NormalTok{alpha }\OperatorTok{=} \FloatTok{2.0}
\NormalTok{omega }\OperatorTok{=} \FloatTok{100.0}
\NormalTok{beta }\OperatorTok{=}\NormalTok{ [}\FloatTok{0.5}\NormalTok{, }\FloatTok{0.5}\NormalTok{]}

\NormalTok{production\_function }\OperatorTok{=}
\NormalTok{    x }\OperatorTok{{-}\textgreater{}}\NormalTok{ alpha }\OperatorTok{*} \FunctionTok{sqrt}\NormalTok{(x[}\FloatTok{1}\NormalTok{] }\OperatorTok{*}\NormalTok{ x[}\FloatTok{2}\NormalTok{])}

\NormalTok{marginal\_product\_function }\OperatorTok{=}
\NormalTok{    x }\OperatorTok{{-}\textgreater{}}\NormalTok{ [}
        \FloatTok{0.5} \OperatorTok{*}\NormalTok{ alpha }\OperatorTok{*} \FunctionTok{sqrt}\NormalTok{(x[}\FloatTok{2}\NormalTok{] }\OperatorTok{/}\NormalTok{ x[}\FloatTok{1}\NormalTok{]),}
        \FloatTok{0.5} \OperatorTok{*}\NormalTok{ alpha }\OperatorTok{*} \FunctionTok{sqrt}\NormalTok{(x[}\FloatTok{1}\NormalTok{] }\OperatorTok{/}\NormalTok{ x[}\FloatTok{2}\NormalTok{]),}
\NormalTok{    ]}

\NormalTok{model }\OperatorTok{=}\NormalTok{ GEMB.}\FunctionTok{GEMBModel}\NormalTok{(}
\NormalTok{    [}\OperatorTok{:}\NormalTok{product, }\OperatorTok{:}\NormalTok{labor];}
\NormalTok{    numeraire}\OperatorTok{=:}\NormalTok{product,}
\NormalTok{)}

\NormalTok{GEMB.}\FunctionTok{add\_agent!}\NormalTok{(}
\NormalTok{    model,}
\NormalTok{    GEMB.}\FunctionTok{ProductionFunctionSpec}\NormalTok{(}
\NormalTok{        production\_function,}
\NormalTok{        marginal\_product\_function;}
\NormalTok{        behavior}\OperatorTok{=:}\NormalTok{cost\_min,}
\NormalTok{    );}
\NormalTok{    outputs}\OperatorTok{=:}\NormalTok{product,}
\NormalTok{    demands}\OperatorTok{=}\NormalTok{[}\OperatorTok{:}\NormalTok{product, }\OperatorTok{:}\NormalTok{labor],}
\NormalTok{    activity\_start}\OperatorTok{=}\FloatTok{100.0}\NormalTok{,}
\NormalTok{    demand\_start}\OperatorTok{=}\NormalTok{[}\FloatTok{100.0}\NormalTok{, }\FloatTok{100.0}\NormalTok{],}
\NormalTok{    name}\OperatorTok{=:}\NormalTok{firm,}
\NormalTok{)}

\NormalTok{GEMB.}\FunctionTok{add\_agent!}\NormalTok{(}
\NormalTok{    model,}
\NormalTok{    GEMB.}\FunctionTok{CESMarginalUtilitySpec}\NormalTok{(beta; es}\OperatorTok{=}\FloatTok{1.0}\NormalTok{);}
\NormalTok{    demands}\OperatorTok{=}\NormalTok{[}\OperatorTok{:}\NormalTok{product, }\OperatorTok{:}\NormalTok{labor],}
\NormalTok{    endowments}\OperatorTok{=:}\NormalTok{labor,}
\NormalTok{    endowment\_quantities}\OperatorTok{=}\NormalTok{omega,}
\NormalTok{    demand\_start}\OperatorTok{=}\NormalTok{[}\FloatTok{100.0}\NormalTok{, }\FloatTok{100.0}\NormalTok{],}
\NormalTok{    demand\_lower\_bounds}\OperatorTok{=}\NormalTok{[}\FloatTok{1.0e{-}10}\NormalTok{, }\FloatTok{1.0e{-}10}\NormalTok{],}
\NormalTok{    multiplier\_start}\OperatorTok{=}\FloatTok{0.01}\NormalTok{,}
\NormalTok{    name}\OperatorTok{=:}\NormalTok{consumer,}
\NormalTok{)}

\NormalTok{result }\OperatorTok{=}\NormalTok{ GEMB.}\FunctionTok{solve}\NormalTok{(}
\NormalTok{    model;}
\NormalTok{    p0}\OperatorTok{=}\NormalTok{[}\FloatTok{1.0}\NormalTok{, }\FloatTok{2.0}\NormalTok{],}
\NormalTok{    residual\_tol}\OperatorTok{=}\FloatTok{1.0e{-}8}\NormalTok{,}
\NormalTok{    silent}\OperatorTok{=}\ConstantTok{true}\NormalTok{,}
\NormalTok{)}

\NormalTok{stats }\OperatorTok{=}\NormalTok{ GEMB.}\FunctionTok{equilibrium\_statistics}\NormalTok{(model, result)}

\NormalTok{output }\OperatorTok{=}\NormalTok{ stats.activity\_levels[}\FloatTok{1}\NormalTok{]}

\NormalTok{firm\_variables }\OperatorTok{=}\NormalTok{ result.agent\_variable\_values[}\FloatTok{1}\NormalTok{]}
\NormalTok{firm\_inputs }\OperatorTok{=}\NormalTok{ firm\_variables[}\FloatTok{2}\OperatorTok{:}\FloatTok{3}\NormalTok{]}
\NormalTok{production\_multiplier }\OperatorTok{=}\NormalTok{ firm\_variables[}\FloatTok{4}\NormalTok{]}

\NormalTok{consumer\_variables }\OperatorTok{=}\NormalTok{ result.agent\_variable\_values[}\FloatTok{2}\NormalTok{]}
\NormalTok{consumption }\OperatorTok{=}\NormalTok{ consumer\_variables[}\FloatTok{1}\OperatorTok{:}\FloatTok{2}\NormalTok{]}
\NormalTok{budget\_multiplier }\OperatorTok{=}\NormalTok{ consumer\_variables[}\FloatTok{3}\NormalTok{]}

\NormalTok{utility }\OperatorTok{=} \FunctionTok{sqrt}\NormalTok{(}\FunctionTok{prod}\NormalTok{(consumption))}

\NormalTok{GEMB.}\FunctionTok{print\_equilibrium\_statistics}\NormalTok{(model, result)}
\FunctionTok{println}\NormalTok{(}\StringTok{"Firm inputs = "}\NormalTok{, firm\_inputs)}
\FunctionTok{println}\NormalTok{(}\StringTok{"Production multiplier = "}\NormalTok{, production\_multiplier)}
\FunctionTok{println}\NormalTok{(}\StringTok{"Consumption = "}\NormalTok{, consumption)}
\FunctionTok{println}\NormalTok{(}\StringTok{"Budget multiplier = "}\NormalTok{, budget\_multiplier)}
\FunctionTok{println}\NormalTok{(}\StringTok{"Utility = "}\NormalTok{, utility)}
\end{Highlighting}
\end{Shaded}

}

\end{codelisting}%

\section{Production Stationarity Conditions and Explicit Agent
Conditions}\label{production-stationarity-conditions-and-explicit-agent-conditions}

\texttt{ProductionStationarityConditions} is used to represent a firm's
production stationarity conditions for a given production technology. It
incorporates the production function, marginal products, and their
relationships with commodity prices directly into the equilibrium
system, so that the firm's output level, input quantities, and
associated production multiplier jointly satisfy the relevant
first-order conditions.

Unlike \texttt{CostMinimizationKKTConditions},
\texttt{ProductionStationarityConditions} does not require the firm to
minimize cost conditional on a given output level. Its economic
interpretation is nevertheless closely related to the underlying
production technology: for each input that is actually used, its
marginal contribution must be consistent with the corresponding
commodity price, while actual output must be consistent with the output
determined by the production function. Thus,
\texttt{ProductionStationarityConditions} provides a way to construct a
firm's equilibrium conditions directly from the first-order stationarity
relations of its production technology. The current formulation assumes
an interior solution, so the marginal-pricing conditions are imposed as
equalities.

\texttt{ExplicitAgentConditions} allows users to define an economic
agent's equilibrium conditions directly. When existing condition rules
such as \texttt{UnitRevenueExpenditureBalanceConditions},
\texttt{CostMinimizationKKTConditions}, or
\texttt{MarginalUtilityConsumerConditions} cannot conveniently represent
a particular economic relationship, users can supply their own condition
functions and pair them with the agent's endogenous variables.
\texttt{ExplicitAgentConditions} therefore provides substantial modeling
flexibility and is useful for representing special behavioral
constraints, institutional rules, and other nonstandard equilibrium
conditions.

The following basic example illustrates the use of production
stationarity conditions together with explicit agent conditions.

\begin{codelisting}

\caption{\label{lst-gemb_basicExample_productionStationarity_explicitConsumer_nc2_na2}A
\(2\times2\) Cobb--Douglas Equilibrium Model: Production Stationarity
Conditions and Explicit Agent Conditions}

\centering{

\begin{Shaded}
\begin{Highlighting}[]
\CommentTok{\# production stationarity conditions and explicit consumer conditions.}
\ImportTok{using} \BuiltInTok{GeneralEquilibriumModeling}

\NormalTok{alpha }\OperatorTok{=} \FloatTok{2.0}
\NormalTok{omega }\OperatorTok{=} \FloatTok{100.0}
\NormalTok{beta }\OperatorTok{=}\NormalTok{ [}\FloatTok{0.5}\NormalTok{, }\FloatTok{0.5}\NormalTok{]}

\NormalTok{production\_function }\OperatorTok{=}\NormalTok{ x }\OperatorTok{{-}\textgreater{}}\NormalTok{ alpha }\OperatorTok{*} \FunctionTok{sqrt}\NormalTok{(x[}\FloatTok{1}\NormalTok{] }\OperatorTok{*}\NormalTok{ x[}\FloatTok{2}\NormalTok{])}

\NormalTok{marginal\_product\_function }\OperatorTok{=}\NormalTok{ x }\OperatorTok{{-}\textgreater{}}\NormalTok{ [}
    \FloatTok{0.5} \OperatorTok{*}\NormalTok{ alpha }\OperatorTok{*} \FunctionTok{sqrt}\NormalTok{(x[}\FloatTok{2}\NormalTok{] }\OperatorTok{/}\NormalTok{ x[}\FloatTok{1}\NormalTok{]),}
    \FloatTok{0.5} \OperatorTok{*}\NormalTok{ alpha }\OperatorTok{*} \FunctionTok{sqrt}\NormalTok{(x[}\FloatTok{1}\NormalTok{] }\OperatorTok{/}\NormalTok{ x[}\FloatTok{2}\NormalTok{]),}
\NormalTok{]}

\NormalTok{model }\OperatorTok{=}\NormalTok{ GEMB.}\FunctionTok{GEMBModel}\NormalTok{(}
\NormalTok{    [}\OperatorTok{:}\NormalTok{product, }\OperatorTok{:}\NormalTok{labor];}
\NormalTok{    numeraire}\OperatorTok{=:}\NormalTok{product,}
\NormalTok{    numeraire\_value}\OperatorTok{=}\FloatTok{1.0}\NormalTok{,}
\NormalTok{)}

\NormalTok{GEMB.}\FunctionTok{add\_agent!}\NormalTok{(}
\NormalTok{    model,}
\NormalTok{    GEMB.}\FunctionTok{ProductionFunctionSpec}\NormalTok{(}
\NormalTok{        production\_function,}
\NormalTok{        marginal\_product\_function;}
\NormalTok{        behavior}\OperatorTok{=:}\NormalTok{stationary,}
\NormalTok{    );}
\NormalTok{    outputs}\OperatorTok{=:}\NormalTok{product,}
\NormalTok{    demands}\OperatorTok{=}\NormalTok{[}\OperatorTok{:}\NormalTok{product, }\OperatorTok{:}\NormalTok{labor],}
\NormalTok{    activity\_start}\OperatorTok{=}\FloatTok{100.0}\NormalTok{,}
\NormalTok{    demand\_start}\OperatorTok{=}\NormalTok{[}\FloatTok{100.0}\NormalTok{, }\FloatTok{100.0}\NormalTok{],}
\NormalTok{    demand\_lower\_bounds}\OperatorTok{=}\NormalTok{[}\FloatTok{1.0e{-}10}\NormalTok{, }\FloatTok{1.0e{-}10}\NormalTok{],}
\NormalTok{    production\_multiplier\_start}\OperatorTok{=}\FloatTok{1.0}\NormalTok{,}
\NormalTok{    name}\OperatorTok{=:}\NormalTok{firm,}
\NormalTok{)}

\NormalTok{consumer\_net\_supply }\OperatorTok{=} \KeywordTok{function}\NormalTok{ (v, p, observed\_values}\OperatorTok{=}\DataTypeTok{Any}\NormalTok{[])}
\NormalTok{    [}\OperatorTok{{-}}\NormalTok{v[}\FloatTok{1}\NormalTok{], omega }\OperatorTok{{-}}\NormalTok{ v[}\FloatTok{2}\NormalTok{]]}
\KeywordTok{end}

\NormalTok{consumer\_conditions }\OperatorTok{=} \KeywordTok{function}\NormalTok{ (v, p, net\_supply)}
\NormalTok{    x }\OperatorTok{=}\NormalTok{ v[}\FloatTok{1}\OperatorTok{:}\FloatTok{2}\NormalTok{]}
\NormalTok{    mu }\OperatorTok{=}\NormalTok{ v[}\FloatTok{3}\NormalTok{]}

    \ControlFlowTok{return}\NormalTok{ [}
\NormalTok{        mu }\OperatorTok{*}\NormalTok{ p[}\FloatTok{1}\NormalTok{] }\OperatorTok{{-}}\NormalTok{ beta[}\FloatTok{1}\NormalTok{] }\OperatorTok{/}\NormalTok{ x[}\FloatTok{1}\NormalTok{],}
\NormalTok{        mu }\OperatorTok{*}\NormalTok{ p[}\FloatTok{2}\NormalTok{] }\OperatorTok{{-}}\NormalTok{ beta[}\FloatTok{2}\NormalTok{] }\OperatorTok{/}\NormalTok{ x[}\FloatTok{2}\NormalTok{],}
        \FunctionTok{sum}\NormalTok{(p }\OperatorTok{.*}\NormalTok{ net\_supply),}
\NormalTok{    ]}
\KeywordTok{end}

\NormalTok{GEMB.}\FunctionTok{add\_net\_supply\_agent!}\NormalTok{(}
\NormalTok{    model,}
\NormalTok{    consumer\_net\_supply;}
\NormalTok{    commodities}\OperatorTok{=}\NormalTok{[}\OperatorTok{:}\NormalTok{product, }\OperatorTok{:}\NormalTok{labor],}
\NormalTok{    variable\_names}\OperatorTok{=}\NormalTok{[}\OperatorTok{:}\NormalTok{demand\_product, }\OperatorTok{:}\NormalTok{demand\_labor, }\OperatorTok{:}\NormalTok{budget\_multiplier],}
\NormalTok{    variable\_lower\_bounds}\OperatorTok{=}\NormalTok{[}\FloatTok{1.0e{-}10}\NormalTok{, }\FloatTok{1.0e{-}10}\NormalTok{, }\FloatTok{0.0}\NormalTok{],}
\NormalTok{    variable\_upper\_bounds}\OperatorTok{=}\NormalTok{[}\ConstantTok{Inf}\NormalTok{, }\ConstantTok{Inf}\NormalTok{, }\ConstantTok{Inf}\NormalTok{],}
\NormalTok{    variable\_start}\OperatorTok{=}\NormalTok{[}\FloatTok{100.0}\NormalTok{, }\FloatTok{100.0}\NormalTok{, }\FloatTok{0.01}\NormalTok{],}
\NormalTok{    condition\_rule}\OperatorTok{=}\NormalTok{GEM.}\FunctionTok{ExplicitAgentConditions}\NormalTok{(consumer\_conditions),}
\NormalTok{    name}\OperatorTok{=:}\NormalTok{consumer,}
\NormalTok{)}

\NormalTok{result }\OperatorTok{=}\NormalTok{ GEMB.}\FunctionTok{solve}\NormalTok{(}
\NormalTok{    model;}
\NormalTok{    p0}\OperatorTok{=}\NormalTok{[}\FloatTok{1.0}\NormalTok{, }\FloatTok{2.0}\NormalTok{],}
\NormalTok{    residual\_tol}\OperatorTok{=}\FloatTok{1.0e{-}8}\NormalTok{,}
\NormalTok{    silent}\OperatorTok{=}\ConstantTok{true}\NormalTok{,}
\NormalTok{)}

\NormalTok{output }\OperatorTok{=}\NormalTok{ result.agent\_variable\_values[}\FloatTok{1}\NormalTok{][}\FloatTok{1}\NormalTok{]}
\NormalTok{consumption }\OperatorTok{=}\NormalTok{ result.agent\_variable\_values[}\FloatTok{2}\NormalTok{][}\FloatTok{1}\OperatorTok{:}\FloatTok{2}\NormalTok{]}
\NormalTok{utility }\OperatorTok{=} \FunctionTok{sqrt}\NormalTok{(}\FunctionTok{prod}\NormalTok{(consumption))}

\FunctionTok{println}\NormalTok{(}\StringTok{"Prices = "}\NormalTok{, result.prices)}
\FunctionTok{println}\NormalTok{(}\StringTok{"Output = "}\NormalTok{, output, }\StringTok{", utility = "}\NormalTok{, utility)}

\PreprocessorTok{@assert}\NormalTok{ result.solved}
\PreprocessorTok{@assert} \FunctionTok{isapprox}\NormalTok{(result.prices, [}\FloatTok{1.0}\NormalTok{, }\FloatTok{1.0}\NormalTok{]; atol}\OperatorTok{=}\FloatTok{1.0e{-}7}\NormalTok{)}
\PreprocessorTok{@assert} \FunctionTok{isapprox}\NormalTok{(output, }\FloatTok{100.0}\NormalTok{; atol}\OperatorTok{=}\FloatTok{1.0e{-}6}\NormalTok{)}
\PreprocessorTok{@assert} \FunctionTok{isapprox}\NormalTok{(utility, }\FloatTok{50.0}\NormalTok{; atol}\OperatorTok{=}\FloatTok{1.0e{-}7}\NormalTok{)}
\PreprocessorTok{@assert} \FunctionTok{maximum}\NormalTok{(abs, result.total\_net\_supply) }\OperatorTok{\textless{}=} \FloatTok{1.0e{-}7}
\end{Highlighting}
\end{Shaded}

}

\end{codelisting}%

\section{Activity-Demand
Specifications}\label{activity-demand-specifications}

An important class of behavioral specifications in the GEMB high-level
framework is the \emph{activity-demand specification}. It applies to
economic agents whose commodity demands can be expressed directly as
functions of an activity level and commodity prices.

For a class of homogeneous consumers represented without aggregation,
the activity level may be interpreted as the number, or population, of
identical consumers. This interpretation is common in strict-sense
structural equilibrium models, where identical consumers need not be
aggregated. In computational models, however, homogeneous consumers are
usually aggregated into a representative consumer. When such a
representative consumer is modeled through a utility-based demand
specification, its activity level typically represents the utility level
rather than the number of consumers.

Let \(z\) denote the activity level of an agent, \(\mathbf p\) the
relevant price vector, and \(\mathbf x\) its commodity demand vector. An
activity-demand specification can be written abstractly as

\[
\mathbf x=\mathbf d(z,\mathbf p)
\]

where \(\mathbf d\) is an activity-demand function. In more general
cases, the demand function may also depend on observed variables.

The interpretation of the activity level depends on the type of economic
agent. For a producer, the activity level usually represents its
production level or production scale. For example, if a producer has
constant returns to scale and requires the unit demand bundle
\(\mathbf a(\mathbf p)\) at prices \(\mathbf p\), its demand can be
written as

\[
\mathbf x=z\mathbf a(\mathbf p)
\]

where \(z\) is the production activity level.

The same representation can also be used for consumers. Suppose a
consumer has a utility function, and let \(\mathbf x^h(\mathbf p,u)\)
denote its compensated demand at prices \(\mathbf p\) and utility level
\(u\). Then

\[
\mathbf x=\mathbf x^h(\mathbf p,u)
\]

can be written in the same mathematical form as an activity-demand
function. The consumer's utility level \(u\) can therefore be regarded
as its \textbf{activity level}. This is the same general concept of an
activity level, although its economic interpretation depends on how the
consumer is represented. For an unaggregated class of homogeneous
consumers, the activity level may instead represent the number or
population of consumers.

This common representation makes it possible to treat producers and
compensated-demand consumers within the same activity-demand framework.
The essential distinction lies in the economic interpretation of the
activity variable and in the conditions determining its equilibrium
value, rather than in the mathematical form of the demand function.

In \texttt{GeneralEquilibriumModeling.jl}, activity-demand
specifications belong to the abstract type
\texttt{AbstractActivityDemandSpec}. The general-purpose specification
is \texttt{ActivityDemandSpec}, whose demand function accepts either

\begin{Shaded}
\begin{Highlighting}[]
\FunctionTok{demand\_function}\NormalTok{(activity, prices)}
\end{Highlighting}
\end{Shaded}

or

\begin{Shaded}
\begin{Highlighting}[]
\FunctionTok{demand\_function}\NormalTok{(activity, prices, observed\_values)}
\end{Highlighting}
\end{Shaded}

and returns the corresponding commodity demand vector. Thus,
\texttt{ActivityDemandSpec} can be used whenever conditional commodity
demand is known explicitly as a function of an activity level and
prices.

GEMB also provides several built-in activity-demand specifications.
\texttt{CESSpec} represents CES activity-demand relationships, while
\texttt{DCESSpec} extends this representation to displaced CES
relationships. These specifications can be used for both producers and
compensated-demand consumers. \texttt{PowerProductionSpec} represents a
one-input power production technology and is intended specifically for
producers.

When an activity-demand specification is used for a producer, the
activity variable must be associated with an appropriate producer
condition. For homogeneous activity-demand relationships without fixed
endowments, GEMB can use unit revenue-expenditure balance conditions.
For specifications involving displacement terms or fixed endowments,
total revenue-expenditure balance conditions are generally more
appropriate. These defaults can be overridden when an alternative
condition rule is required.

For a compensated-demand consumer, the construction is different. The
consumer has no output commodity in the producer sense. Its net supply
is obtained from its endowment minus its compensated demand, and its
equilibrium utility activity is determined together with a total
expenditure-income balance condition. Consequently, the activity-demand
framework provides a unified way to describe two apparently different
economic behaviors:

\begin{itemize}
\tightlist
\item
  a producer chooses a production activity level and demands the
  commodities required to support that production level;
\item
  a consumer attains a utility activity level and demands the
  commodities required to support that utility level.
\end{itemize}

This interpretation is especially useful in structural equilibrium
modeling because both production and consumption can be represented
through the relationship between an activity level, commodity prices,
and the associated demand bundle.

The function \texttt{activity\_supply} is the common output-supply
interface for activity-based producers in GEMB. It maps an activity
level, output prices, and optionally observed equilibrium variables into
the corresponding output quantities. Built-in output specifications such
as \texttt{CETSpec} implement this interface automatically, while users
can define new specifications by creating a subtype of
\texttt{AbstractActivitySupplySpec} and extending
\texttt{activity\_supply}. Such a custom specification can then be
supplied through the \texttt{output\_spec} argument when constructing a
producer.

\section{Custom Net-Supply Agents and Condition
Agents}\label{custom-net-supply-agents-and-condition-agents}

\subsection{Custom Net-Supply Agents and Multi-Activity
Producers}\label{custom-net-supply-agents-and-multi-activity-producers}

The standard behavioral specifications in the GEMB high-level framework,
such as \texttt{CESSpec}, \texttt{DCESSpec}, and
\texttt{MarshallDemandConsumerSpec}, provide convenient representations
for commonly used forms of economic behavior. When an agent cannot be
represented naturally by one of these specifications, GEMB also allows
its net-supply mapping to be defined directly through
\texttt{add\_net\_supply\_agent!}.

A custom net-supply agent may have its own endogenous variables,
variable bounds and starting values, observed equilibrium variables, and
condition rule. Its behavior is specified directly by a net-supply
function that returns positive quantities for supplies and negative
quantities for demands. This provides a flexible high-level interface
for economic agents whose behavior does not fit one of the standard GEMB
specifications.

A custom net-supply producer may also contain multiple endogenous
activity levels. For example, a single firm may control an activity
vector \(\mathbf z=(z_1,\ldots,z_m)^\top\) representing several
production processes. When
\texttt{UnitRevenueExpenditureBalanceConditions} is used, one unit
revenue-expenditure balance condition is constructed for each activity
variable. Thus, the number of production activities need not equal the
number of economic agents. By contrast, the current standard
activity-demand producer specifications, including \texttt{CESSpec} and
\texttt{DCESSpec}, construct one activity level for each agent.

\subsection{Condition Agents}\label{condition-agents}

A second extension mechanism is provided by \texttt{ConditionAgentSpec}.
A condition agent provides the GEMB high-level representation of
endogenous auxiliary variables together with the explicit conditions
that determine them. It may contain one or more endogenous equilibrium
variables, with one condition paired with each variable, but it does not
generate commodity demand or supply. It is therefore useful for
representing variables that would typically be introduced through
auxiliary variable--equation pairs in the GEM framework, including
closure variables, endogenous policy parameters, normalization
variables, and other equilibrium variables determined by explicit
conditions.

\subsection{A Two-Commodity von Neumann Model with Three
Activities}\label{a-two-commodity-von-neumann-model-with-three-activities}

The following von Neumann model illustrates how a custom net-supply
agent and a condition agent can be used together.

Consider a pure-production economy with two commodities, wheat and iron,
and three production processes (Li, 2019, Example 2.3). The input and
output coefficient matrices are

\[
\mathbf A=
\begin{pmatrix}
0.8&0.5&0.06\\
2&2&0.4
\end{pmatrix},
\qquad
\mathbf B=
\begin{pmatrix}
1&1&0\\
0&0&1
\end{pmatrix}
\]

Each column represents a production process. Processes 1 and 2 both
produce wheat, whereas process 3 produces iron. Processes 1 and 2
therefore compete in wheat production. Let
\(\mathbf z=(z_1,z_2,z_3)^\top\) denote the activity-level vector,
\(\mathbf p=(p_1,p_2)^\top\) the commodity-price vector, and \(\rho\)
the equilibrium discount factor.

Processes 1 and 2 produce the same output, one unit of wheat, and use
the same amount of iron, but process 1 requires \(0.8\) units of wheat
as input whereas process 2 requires only \(0.5\). Process 1 therefore
has a higher unit cost for the same unit revenue and is technologically
disadvantaged.

Taking wheat as the numeraire, one equilibrium is

\[
\mathbf p^*=
\begin{pmatrix}
1\\
0.15
\end{pmatrix},
\qquad
\mathbf z^*=
\begin{pmatrix}
0\\
100\\
500
\end{pmatrix},
\qquad
\rho^*=0.8
\]

The corresponding balanced-growth factor is \(1/\rho^*=1.25\), so the
equilibrium growth rate is

\[
\frac{1}{\rho^*}-1=0.25
\]

In the GEMB high-level implementation, the three production processes
are represented as three activity variables of a single custom
net-supply producer. Its aggregate net supply is defined as

\[
(\rho\mathbf B-\mathbf A)\mathbf z
\]

The producer is constructed with \texttt{add\_net\_supply\_agent!} and
uses \texttt{UnitRevenueExpenditureBalanceConditions}. Consequently,
activity \(j\) is paired with the unit excess cost

\[
\mathbf p^\top\mathbf A_{\cdot j}
-
\rho\mathbf p^\top\mathbf B_{\cdot j}
\]

A positive activity level therefore requires the corresponding unit
excess cost to be zero, whereas a production process with a positive
unit excess cost may remain inactive.

The equilibrium discount factor \(\rho\) is represented as the
endogenous variable of a \texttt{ConditionAgentSpec}. Since the von
Neumann activity vector is homogeneous, its scale is fixed by the
normalization

\[
z_1+z_2+z_3=600
\]

which is used as the condition paired with \(\rho\). Together with the
price normalization \(p_1=1\), the model simultaneously determines the
relative commodity price, the three activity levels, and the equilibrium
discount factor.

Thus, this example demonstrates two complementary extensions of the GEMB
high-level framework: \texttt{add\_net\_supply\_agent!} allows the
direct construction of a multi-activity producer, while
\texttt{ConditionAgentSpec} represents an endogenous auxiliary variable
and the explicit condition that determines it.

\begin{codelisting}

\caption{\label{lst-gemb_vonNeumann_multiActivityFirm_nc2_na1}A
Two-Commodity Von Neumann Equilibrium Model with Three Activities}

\centering{

\begin{Shaded}
\begin{Highlighting}[]
\ImportTok{using} \BuiltInTok{GeneralEquilibriumModeling}

\NormalTok{A }\OperatorTok{=}\NormalTok{ [}
    \FloatTok{0.8}  \FloatTok{0.5}  \FloatTok{0.06}
    \FloatTok{2.0}  \FloatTok{2.0}  \FloatTok{0.40}
\NormalTok{]}

\NormalTok{B }\OperatorTok{=}\NormalTok{ [}
    \FloatTok{1.0}  \FloatTok{1.0}  \FloatTok{0.0}
    \FloatTok{0.0}  \FloatTok{0.0}  \FloatTok{1.0}
\NormalTok{]}

\NormalTok{model }\OperatorTok{=}\NormalTok{ GEMB.}\FunctionTok{GEMBModel}\NormalTok{(}
\NormalTok{    [}\OperatorTok{:}\NormalTok{wheat, }\OperatorTok{:}\NormalTok{iron];}
\NormalTok{    numeraire}\OperatorTok{=:}\NormalTok{wheat,}
\NormalTok{    numeraire\_value}\OperatorTok{=}\FloatTok{1.0}\NormalTok{,}
\NormalTok{)}

\NormalTok{GEMB.}\FunctionTok{add\_net\_supply\_agent!}\NormalTok{(}
\NormalTok{    model,}
    \KeywordTok{function}\NormalTok{ (variables, prices, observed\_values)}
\NormalTok{        z }\OperatorTok{=}\NormalTok{ variables}
\NormalTok{        rho }\OperatorTok{=}\NormalTok{ observed\_values[}\FloatTok{1}\NormalTok{]}

        \ControlFlowTok{return}\NormalTok{ (rho }\OperatorTok{.*}\NormalTok{ B }\OperatorTok{{-}}\NormalTok{ A) }\OperatorTok{*}\NormalTok{ z}
    \KeywordTok{end}\NormalTok{;}
\NormalTok{    commodities}\OperatorTok{=}\NormalTok{[}\OperatorTok{:}\NormalTok{wheat, }\OperatorTok{:}\NormalTok{iron],}
\NormalTok{    variable\_names}\OperatorTok{=}\NormalTok{[}\OperatorTok{:}\NormalTok{activity1, }\OperatorTok{:}\NormalTok{activity2, }\OperatorTok{:}\NormalTok{activity3],}
\NormalTok{    variable\_lower\_bounds}\OperatorTok{=}\NormalTok{[}\FloatTok{0.0}\NormalTok{, }\FloatTok{0.0}\NormalTok{, }\FloatTok{0.0}\NormalTok{],}
\NormalTok{    variable\_upper\_bounds}\OperatorTok{=}\NormalTok{[}\ConstantTok{Inf}\NormalTok{, }\ConstantTok{Inf}\NormalTok{, }\ConstantTok{Inf}\NormalTok{],}
\NormalTok{    variable\_start}\OperatorTok{=}\NormalTok{[}\FloatTok{100.0}\NormalTok{, }\FloatTok{100.0}\NormalTok{, }\FloatTok{100.0}\NormalTok{],}
\NormalTok{    observed\_variables}\OperatorTok{=}\NormalTok{[}
\NormalTok{        GEM.}\FunctionTok{AgentVariableRef}\NormalTok{(}\OperatorTok{:}\NormalTok{growth\_condition, }\OperatorTok{:}\NormalTok{rho),}
\NormalTok{    ],}
\NormalTok{    condition\_rule}\OperatorTok{=}\NormalTok{GEM.}\FunctionTok{UnitRevenueExpenditureBalanceConditions}\NormalTok{(),}
\NormalTok{    name}\OperatorTok{=:}\NormalTok{firm,}
\NormalTok{)}

\NormalTok{GEMB.}\FunctionTok{add\_agent!}\NormalTok{(}
\NormalTok{    model,}
\NormalTok{    GEMB.}\FunctionTok{ConditionAgentSpec}\NormalTok{(}
\NormalTok{        (variables, observed\_values) }\OperatorTok{{-}\textgreater{}}\NormalTok{ [}
            \FunctionTok{sum}\NormalTok{(observed\_values) }\OperatorTok{{-}} \FloatTok{600.0}\NormalTok{,}
\NormalTok{        ],}
\NormalTok{    );}
\NormalTok{    variable\_names}\OperatorTok{=}\NormalTok{[}\OperatorTok{:}\NormalTok{rho],}
\NormalTok{    variable\_start}\OperatorTok{=}\NormalTok{[}\FloatTok{0.8}\NormalTok{],}
\NormalTok{    variable\_lower\_bounds}\OperatorTok{=}\NormalTok{[}\FloatTok{0.0}\NormalTok{],}
\NormalTok{    variable\_upper\_bounds}\OperatorTok{=}\NormalTok{[}\ConstantTok{Inf}\NormalTok{],}
\NormalTok{    observed\_variables}\OperatorTok{=}\NormalTok{[}
\NormalTok{        GEM.}\FunctionTok{AgentVariableRef}\NormalTok{(}\OperatorTok{:}\NormalTok{firm, }\OperatorTok{:}\NormalTok{activity1),}
\NormalTok{        GEM.}\FunctionTok{AgentVariableRef}\NormalTok{(}\OperatorTok{:}\NormalTok{firm, }\OperatorTok{:}\NormalTok{activity2),}
\NormalTok{        GEM.}\FunctionTok{AgentVariableRef}\NormalTok{(}\OperatorTok{:}\NormalTok{firm, }\OperatorTok{:}\NormalTok{activity3),}
\NormalTok{    ],}
\NormalTok{    name}\OperatorTok{=:}\NormalTok{growth\_condition,}
\NormalTok{)}

\NormalTok{result }\OperatorTok{=}\NormalTok{ GEMB.}\FunctionTok{solve}\NormalTok{(}
\NormalTok{    model;}
\NormalTok{    p0}\OperatorTok{=}\NormalTok{[}\FloatTok{1.0}\NormalTok{, }\FloatTok{0.15}\NormalTok{],}
\NormalTok{    residual\_tol}\OperatorTok{=}\FloatTok{1.0e{-}8}\NormalTok{,}
\NormalTok{    silent}\OperatorTok{=}\ConstantTok{true}\NormalTok{,}
\NormalTok{)}

\NormalTok{prices }\OperatorTok{=}\NormalTok{ result.prices}
\NormalTok{activities }\OperatorTok{=}\NormalTok{ result.agent\_variable\_values[}\FloatTok{1}\NormalTok{]}
\NormalTok{rho\_star }\OperatorTok{=}\NormalTok{ result.agent\_variable\_values[}\FloatTok{2}\NormalTok{][}\FloatTok{1}\NormalTok{]}
\NormalTok{growth\_rate }\OperatorTok{=} \FunctionTok{inv}\NormalTok{(rho\_star) }\OperatorTok{{-}} \FloatTok{1.0}

\FunctionTok{println}\NormalTok{(}\StringTok{"Prices:      "}\NormalTok{, prices)}
\FunctionTok{println}\NormalTok{(}\StringTok{"Activities:  "}\NormalTok{, activities)}
\FunctionTok{println}\NormalTok{(}\StringTok{"rho:         "}\NormalTok{, rho\_star)}
\FunctionTok{println}\NormalTok{(}\StringTok{"Growth rate: "}\NormalTok{, growth\_rate)}

\ConstantTok{nothing}
\end{Highlighting}
\end{Shaded}

}

\end{codelisting}%

\section{Equilibrium Results and
Statistics}\label{equilibrium-results-and-statistics}

After a GEMB model has been constructed, it can normally be solved
directly by

\begin{Shaded}
\begin{Highlighting}[]
\NormalTok{result }\OperatorTok{=}\NormalTok{ GEMB.}\FunctionTok{solve}\NormalTok{(model)}
\end{Highlighting}
\end{Shaded}

The resulting object contains the equilibrium solution and solver
information. For economic analysis, GEMB also provides post-solve
functions that organize the solution into prices, activity levels,
demand and supply matrices, and related value statistics.

\subsection{Equilibrium Statistics}\label{equilibrium-statistics}

General equilibrium statistics can be obtained by

\begin{Shaded}
\begin{Highlighting}[]
\NormalTok{stats }\OperatorTok{=}\NormalTok{ GEMB.}\FunctionTok{equilibrium\_statistics}\NormalTok{(model, result)}
\end{Highlighting}
\end{Shaded}

The equilibrium price vector is stored in \texttt{stats.prices}.
Explicitly represented activity levels are collected in
\texttt{stats.activity\_levels}, while
\texttt{stats.activity\_level\_refs} identifies the corresponding
agents. The economic interpretation of an activity level depends on the
representation of the agent. For example, it may represent a production
activity level or a consumer's utility level. A multi-activity agent may
therefore contribute several entries to \texttt{activity\_levels},
whereas an agent without an explicitly represented activity level
contributes none.

GEMB represents each economic agent internally through its net supply.
Let there be \(n\) commodities and \(m\) agents, and let the demand and
supply matrices be

\[
\mathbf D=[d_{ij}]
\qquad\text{and}\qquad
\mathbf S=[s_{ij}]
\]

where rows correspond to commodities and columns correspond to agents.
The net-supply matrix is defined as

\[
\bar{\mathbf S}
=
\mathbf S-\mathbf D
\]

Thus, \(\bar s_{ij}>0\) means that agent \(j\) is a net supplier of
commodity \(i\), while \(\bar s_{ij}<0\) means that the agent is a net
demander of that commodity. The total net supply of each commodity is

\[
\bar{\mathbf S}\mathbf 1
=
\mathbf S\mathbf 1-\mathbf D\mathbf 1
\]

The corresponding information is available from
\texttt{stats.net\_supply\_matrix} and
\texttt{stats.total\_net\_supply}.

The net-supply value matrix at the equilibrium price vector
\(\mathbf p\) is

\[
\bar{\mathbf S}_v
=
\operatorname{Diag}(\mathbf p)\bar{\mathbf S}
\]

and is available from \texttt{stats.net\_supply\_value\_matrix}. Summing
this matrix over commodities gives the net-supply value of each agent,

\[
\bar{\mathbf S}^{\mathsf T}\mathbf p
\]

which is stored in \texttt{stats.agent\_net\_supply\_values}. For an
agent satisfying a revenue-expenditure balance condition, this value
should be zero up to numerical error.

\subsection{Demand and Supply
Matrices}\label{demand-and-supply-matrices}

For models consisting of agents with activity-demand specifications, the
demand and supply components can be obtained separately by

\begin{Shaded}
\begin{Highlighting}[]
\NormalTok{flows }\OperatorTok{=}\NormalTok{ GEMB.}\FunctionTok{demand\_supply\_matrices}\NormalTok{(model, result)}
\end{Highlighting}
\end{Shaded}

The main matrices are

\begin{Shaded}
\begin{Highlighting}[]
\NormalTok{D }\OperatorTok{=}\NormalTok{ flows.demand}
\NormalTok{S }\OperatorTok{=}\NormalTok{ flows.supply}
\NormalTok{Sbar }\OperatorTok{=}\NormalTok{ flows.net\_supply}

\NormalTok{DV }\OperatorTok{=}\NormalTok{ flows.demand\_value}
\NormalTok{SV }\OperatorTok{=}\NormalTok{ flows.supply\_value}
\end{Highlighting}
\end{Shaded}

where \texttt{D}, \texttt{S}, and \texttt{Sbar} correspond respectively
to \(\mathbf D\), \(\mathbf S\), and \(\bar{\mathbf S}\). Hence,

\begin{Shaded}
\begin{Highlighting}[]
\NormalTok{S }\OperatorTok{{-}}\NormalTok{ D }\OperatorTok{≈}\NormalTok{ Sbar}
\end{Highlighting}
\end{Shaded}

up to numerical error.

The total demand and supply vectors are

\[
\mathbf d
=
\mathbf D\mathbf 1
\qquad\text{and}\qquad
\mathbf s
=
\mathbf S\mathbf 1
\]

and are available as \texttt{flows.total\_demand} and
\texttt{flows.total\_supply}. Their relationship reflects the
commodity-market conditions imposed in the model. In an equality-form
equilibrium model,

\[
\mathbf D\mathbf 1
=
\mathbf S\mathbf 1
\]

whereas a model allowing free disposal or other inequality-form market
conditions may permit positive total net supply for some commodities.

GEMB also reports demand and supply in value terms. Define

\[
\mathbf{DV}
=
\operatorname{Diag}(\mathbf p)\mathbf D
\]

and

\[
\mathbf{SV}
=
\operatorname{Diag}(\mathbf p)\mathbf S
\]

The element \(DV_{ij}\) is agent \(j\)'s expenditure on commodity \(i\),
while \(SV_{ij}\) is the revenue associated with its supply of that
commodity. Consequently, the expenditure and revenue vectors of the
agents are

\[
\mathbf e
=
\mathbf D^{\mathsf T}\mathbf p
\]

and

\[
\mathbf r
=
\mathbf S^{\mathsf T}\mathbf p
\]

These quantities are available as \texttt{flows.agent\_expenditure} and
\texttt{flows.agent\_revenue}. Their difference satisfies

\[
\mathbf r-\mathbf e
=
\bar{\mathbf S}^{\mathsf T}\mathbf p
\]

Thus, the quantity matrices and value matrices provide two complementary
views of an equilibrium: the former describe the physical allocation of
commodities, while the latter describe the corresponding revenue and
expenditure flows.

\subsection{Printing and Checking an
Equilibrium}\label{printing-and-checking-an-equilibrium}

For routine inspection, GEMB provides

\begin{Shaded}
\begin{Highlighting}[]
\NormalTok{GEMB.}\FunctionTok{print\_equilibrium\_statistics}\NormalTok{(model, result)}
\end{Highlighting}
\end{Shaded}

which prints a compact summary of equilibrium prices, activity levels,
net supplies, net-supply values, and aggregate net supplies. Numerical
values close to zero can be suppressed in the display by specifying
\texttt{display\_tol}, and the number of significant digits can be
controlled by \texttt{sigdigits}.

A typical post-solve workflow is therefore

\begin{Shaded}
\begin{Highlighting}[]
\NormalTok{result }\OperatorTok{=}\NormalTok{ GEMB.}\FunctionTok{solve}\NormalTok{(model)}

\NormalTok{stats }\OperatorTok{=}\NormalTok{ GEMB.}\FunctionTok{equilibrium\_statistics}\NormalTok{(model, result)}
\NormalTok{flows }\OperatorTok{=}\NormalTok{ GEMB.}\FunctionTok{demand\_supply\_matrices}\NormalTok{(model, result)}

\NormalTok{D }\OperatorTok{=}\NormalTok{ flows.demand}
\NormalTok{S }\OperatorTok{=}\NormalTok{ flows.supply}
\NormalTok{Sbar }\OperatorTok{=}\NormalTok{ flows.net\_supply}

\NormalTok{DV }\OperatorTok{=}\NormalTok{ flows.demand\_value}
\NormalTok{SV }\OperatorTok{=}\NormalTok{ flows.supply\_value}

\PreprocessorTok{@assert}\NormalTok{ S }\OperatorTok{{-}}\NormalTok{ D }\OperatorTok{≈}\NormalTok{ Sbar}

\NormalTok{GEMB.}\FunctionTok{print\_equilibrium\_statistics}\NormalTok{(model, result)}
\end{Highlighting}
\end{Shaded}

For an equality-form model in which all agents satisfy
revenue-expenditure balance, two additional numerical checks are

\begin{Shaded}
\begin{Highlighting}[]
\PreprocessorTok{@assert}\NormalTok{ flows.total\_demand }\OperatorTok{≈}\NormalTok{ flows.total\_supply}
\PreprocessorTok{@assert}\NormalTok{ flows.agent\_expenditure }\OperatorTok{≈}\NormalTok{ flows.agent\_revenue}
\end{Highlighting}
\end{Shaded}

These checks help distinguish the economic interpretation of an
equilibrium from the numerical solution process. \texttt{solve}
determines an equilibrium of the underlying mixed complementarity
problem, while the statistics functions reorganize the resulting
variables into economically meaningful quantities and provide convenient
consistency checks.

\bookmarksetup{startatroot}

\chapter{Static Pure-Exchange
Equilibrium}\label{static-pure-exchange-equilibrium}

\section{\texorpdfstring{CES and DCES Consumers: A \(2\times2\)
Pure-Exchange
Model}{CES and DCES Consumers: A 2\textbackslash times2 Pure-Exchange Model}}\label{ces-and-dces-consumers-a-2times2-pure-exchange-model}

\subsection{CES and DCES
Specifications}\label{ces-and-dces-specifications}

CES and displaced CES functions are two built-in function families in
the GEMB high-level framework. They are represented by \texttt{CESSpec}
and \texttt{DCESSpec}, respectively. Both are activity-demand
specifications: given an activity level and the prices of demanded
commodities, GEMB computes the corresponding conditional demand vector.
For a producer, the activity level represents output or production
scale; for a consumer, it can represent utility.

A CES specification is constructed by

\begin{Shaded}
\begin{Highlighting}[]
\NormalTok{GEMB.}\FunctionTok{CESSpec}\NormalTok{(}
\NormalTok{    beta;}
\NormalTok{    es}\OperatorTok{=}\FloatTok{1.0}\NormalTok{,}
\NormalTok{    alpha}\OperatorTok{=}\FloatTok{1.0}\NormalTok{,}
\NormalTok{)}
\end{Highlighting}
\end{Shaded}

where \texttt{beta} is the vector of distribution parameters,
\texttt{es} is the elasticity of substitution, and \texttt{alpha} is the
scale parameter. Let
\(\boldsymbol{\beta}=(\beta_1,\ldots,\beta_n)^\top\), \(e_s\) denote
\texttt{es}, and \(\alpha>0\). For \(e_s\ne0,1\), the corresponding CES
function takes the form

\[
f(\mathbf x)
=
\alpha
\left[
\sum_{i=1}^n
\beta_i x_i^{(e_s-1)/e_s}
\right]^{e_s/(e_s-1)}
\]

The parameter \texttt{es} directly represents the elasticity of
substitution. When \texttt{es=1}, GEMB uses the Cobb--Douglas limit and
requires \(\sum_i\beta_i=1\). When \texttt{es=0}, it uses the Leontief
limit. Hence a single \texttt{CESSpec} interface covers the Leontief,
Cobb--Douglas, and ordinary CES cases.

For example,

\begin{Shaded}
\begin{Highlighting}[]
\NormalTok{spec }\OperatorTok{=}\NormalTok{ GEMB.}\FunctionTok{CESSpec}\NormalTok{(}
\NormalTok{    [}\FloatTok{0.4}\NormalTok{, }\FloatTok{0.6}\NormalTok{];}
\NormalTok{    es}\OperatorTok{=}\FloatTok{1.0}\NormalTok{,}
\NormalTok{    alpha}\OperatorTok{=}\FloatTok{1.0}\NormalTok{,}
\NormalTok{)}
\end{Highlighting}
\end{Shaded}

defines a two-commodity Cobb--Douglas specification. Changing only
\texttt{es} changes the degree of substitution without changing the
general model-building interface.

If the activity level is \(z\) and the relevant price vector is
\(\mathbf p\), GEMB evaluates the conditional demand as

\begin{Shaded}
\begin{Highlighting}[]
\NormalTok{GEMB.}\FunctionTok{activity\_demand}\NormalTok{(spec, z, p)}
\end{Highlighting}
\end{Shaded}

For \(e_s\ne0,1\), its \(i\)th component is

\[
x_i
=
\frac{z}{\alpha}
\beta_i^{e_s} p_i^{-e_s}
\left(
\sum_{j=1}^n
\beta_j^{e_s} p_j^{1-e_s}
\right)^{e_s/(1-e_s)}
\]

When \texttt{CESSpec} is used for a consumer, the same
conditional-demand relation is interpreted as compensated demand, with
the activity level representing utility.

\texttt{DCESSpec} extends this construction by introducing a
displacement vector (Fullerton, 1989). It is constructed by

\begin{Shaded}
\begin{Highlighting}[]
\NormalTok{GEMB.}\FunctionTok{DCESSpec}\NormalTok{(}
\NormalTok{    beta;}
\NormalTok{    es}\OperatorTok{=}\FloatTok{1.0}\NormalTok{,}
\NormalTok{    alpha}\OperatorTok{=}\FloatTok{1.0}\NormalTok{,}
\NormalTok{    xi}\OperatorTok{=}\FunctionTok{zeros}\NormalTok{(}\FunctionTok{length}\NormalTok{(beta)),}
\NormalTok{)}
\end{Highlighting}
\end{Shaded}

where \texttt{xi} is the displacement vector
\(\boldsymbol{\xi}=(\xi_1,\ldots,\xi_n)^\top\). For \(e_s>0\) and
\(e_s\ne1\), the DCES function is parameterized as

\[
f(\mathbf x)
=
\alpha
\left[
\sum_{i=1}^n
\beta_i^{1/e_s}
(x_i-\xi_i)^{(e_s-1)/e_s}
\right]^{e_s/(e_s-1)}
\]

When \(e_s\to0\), the DCES function converges to the displaced Leontief
form

\[
f(\mathbf x)
=
\alpha
\min_{i:\beta_i>0}
\left\{
\frac{x_i-\xi_i}{\beta_i}
\right\}
\]

When \(e_s=1\), the corresponding Cobb--Douglas limit is used.

The displacement parameters shift the quantities entering the CES
aggregator. The conditional demand consequently has the form \[
\mathbf x
=
\boldsymbol{\xi}
+
\frac{z}{\alpha}\mathbf a(\mathbf p)
\]

where \(\mathbf a(\mathbf p)\) is the price-dependent unit demand
vector. In particular, for \(e_s\ne1\),

\[
x_i
=
\xi_i
+
\frac{z}{\alpha}
\beta_i p_i^{-e_s}
\left(
\sum_{j=1}^n
\beta_jp_j^{1-e_s}
\right)^{e_s/(1-e_s)}
\]

\texttt{CESSpec} and \texttt{DCESSpec} use different parameterizations
of the CES coefficients. Consequently, setting \texttt{xi=zeros(...)} in
a \texttt{DCESSpec} produces a homogeneous CES-family specification, but
it should not in general be interpreted as numerically identical to a
\texttt{CESSpec} constructed with the same values of \texttt{beta} and
\texttt{alpha}. The two specifications should therefore be regarded as
related built-in parameterizations rather than as interchangeable
constructors.

When these specifications are passed to \texttt{GEMB.add\_agent!}, GEMB
combines the demand specification with the agent's outputs, endowments,
claims, and equilibrium conditions. For a \texttt{CESSpec} producer
without fixed endowments, the default condition is
\texttt{UnitRevenueExpenditureBalanceConditions}. A producer with fixed
endowments instead uses
\texttt{TotalRevenueExpenditureBalanceConditions}. For
\texttt{DCESSpec}, a producer with a nonzero displacement parameter also
uses \texttt{TotalRevenueExpenditureBalanceConditions}, since its net
supply is no longer proportional to the activity level. These defaults
can be replaced by explicitly supplying another condition rule.

Claims are not part of the CES or DCES function itself. If an ad valorem
claim is required, it is attached separately when the agent is
constructed through arguments such as \texttt{claim} and
\texttt{claim\_rate}. This separation allows the same CES or DCES
behavioral specification to be reused under different institutional
arrangements.

\subsection{\texorpdfstring{A \(2\times2\) Pure-Exchange
Model}{A 2\textbackslash times2 Pure-Exchange Model}}\label{a-2times2-pure-exchange-model}

The following pure-exchange model illustrates how \texttt{CESSpec} and
\texttt{DCESSpec} can be used simultaneously for different consumers.
Assume the following:

\begin{enumerate}
\def\labelenumi{\arabic{enumi}.}
\item
  The economy contains two commodities and two consumers. Commodity 2 is
  taken as the numeraire commodity.
\item
  Consumer 1 is endowed with \(\omega_1>0\) units of commodity 1 and no
  commodity 2. The consumer has a CES specification with
  \texttt{beta={[}beta,\ 1-beta{]}}, \texttt{es=1}, and
  \texttt{alpha=1}, where \(0<\beta<1\). Thus, the CES specification
  reduces to the Cobb--Douglas case.
\item
  Consumer 2 is endowed with \(\omega_2>0\) units of commodity 2 and no
  commodity 1. The consumer has a DCES specification with the same
  \texttt{beta}, \texttt{es}, and \texttt{alpha}, but with displacement
  vector \(\boldsymbol{\xi}=(0,\xi)^\top\), where \(0\leq\xi<\omega_2\).
\end{enumerate}

Let the price vector be \(\mathbf p=(p,1)^\top\). The demand and supply
matrices are

\[
\mathbf D(p)
=
\begin{bmatrix}
\beta\omega_1 &
\frac{\beta(\omega_2-\xi)}{p}
\\
(1-\beta)p\omega_1 &
\xi+(1-\beta)(\omega_2-\xi)
\end{bmatrix}
\]

and

\[
\mathbf S
=
\begin{bmatrix}
\omega_1&0\\
0&\omega_2
\end{bmatrix}
\]

where rows represent commodities and columns represent consumers.

The equality-form wide-sense structural equilibrium model requires the
revenue-expenditure balance condition

\[
\mathbf p^\top\mathbf D(p)
=
\mathbf p^\top\mathbf S
\]

and the supply-demand balance condition

\[
\mathbf D(p)\mathbf 1
=
\mathbf S\mathbf 1
\]

The equilibrium price vector is

\[
\mathbf p^*
=
\left(
\frac{\beta(\omega_2-\xi)}
{(1-\beta)\omega_1},
1
\right)^\top
\]

The equilibrium demand matrix is

\[
\mathbf D^*
=
\begin{bmatrix}
\beta\omega_1&(1-\beta)\omega_1\\
\beta(\omega_2-\xi)&
\xi+(1-\beta)(\omega_2-\xi)
\end{bmatrix}
\]

For example, let \(\beta=0.5\), \(\omega_1=\omega_2=100\), and
\(\xi=20\). Then

\[
\mathbf p^*=(0.8,1)^\top
\]

and

\[
\mathbf D^*
=
\begin{bmatrix}
50&50\\
40&60
\end{bmatrix},
\qquad
\mathbf S
=
\begin{bmatrix}
100&0\\
0&100
\end{bmatrix}
\]

The first consumer therefore consumes \((50,40)^\top\), while the second
consumes \((50,60)^\top\). The 20-unit displacement in commodity 2
directly illustrates how \texttt{DCESSpec} differs from an ordinary
\texttt{CESSpec}.

The corresponding GEMB implementation is given in
Listing~\ref{lst-gemb_pureExchange_CES_DCES}.

\begin{codelisting}

\caption{\label{lst-gemb_pureExchange_CES_DCES}CES and DCES Consumers: A
\(2\times2\) Pure-Exchange Model}

\centering{

\begin{Shaded}
\begin{Highlighting}[]
\ImportTok{using} \BuiltInTok{GeneralEquilibriumModeling}

\NormalTok{beta }\OperatorTok{=}\NormalTok{ [}\FloatTok{0.5}\NormalTok{, }\FloatTok{0.5}\NormalTok{]}

\NormalTok{model }\OperatorTok{=}\NormalTok{ GEMB.}\FunctionTok{GEMBModel}\NormalTok{(}
\NormalTok{    [}\OperatorTok{:}\NormalTok{commodity1, }\OperatorTok{:}\NormalTok{commodity2];}
\NormalTok{    numeraire}\OperatorTok{=:}\NormalTok{commodity2,}
\NormalTok{)}

\NormalTok{GEMB.}\FunctionTok{add\_agent!}\NormalTok{(}
\NormalTok{    model,}
\NormalTok{    GEMB.}\FunctionTok{CESSpec}\NormalTok{(}
\NormalTok{        beta;}
\NormalTok{        es}\OperatorTok{=}\FloatTok{1.0}\NormalTok{,}
\NormalTok{        alpha}\OperatorTok{=}\FloatTok{1.0}\NormalTok{,}
\NormalTok{    );}
\NormalTok{    demands}\OperatorTok{=}\NormalTok{[}\OperatorTok{:}\NormalTok{commodity1, }\OperatorTok{:}\NormalTok{commodity2],}
\NormalTok{    endowments}\OperatorTok{=:}\NormalTok{commodity1,}
\NormalTok{    endowment\_quantities}\OperatorTok{=}\FloatTok{100.0}\NormalTok{,}
\NormalTok{    activity\_start}\OperatorTok{=}\FloatTok{100.0}\NormalTok{,}
\NormalTok{    name}\OperatorTok{=:}\NormalTok{consumer1,}
\NormalTok{)}

\NormalTok{GEMB.}\FunctionTok{add\_agent!}\NormalTok{(}
\NormalTok{    model,}
\NormalTok{    GEMB.}\FunctionTok{DCESSpec}\NormalTok{(}
\NormalTok{        beta;}
\NormalTok{        es}\OperatorTok{=}\FloatTok{1.0}\NormalTok{,}
\NormalTok{        alpha}\OperatorTok{=}\FloatTok{1.0}\NormalTok{,}
\NormalTok{        xi}\OperatorTok{=}\NormalTok{[}\FloatTok{0.0}\NormalTok{, }\FloatTok{20.0}\NormalTok{],}
\NormalTok{    );}
\NormalTok{    demands}\OperatorTok{=}\NormalTok{[}\OperatorTok{:}\NormalTok{commodity1, }\OperatorTok{:}\NormalTok{commodity2],}
\NormalTok{    endowments}\OperatorTok{=:}\NormalTok{commodity2,}
\NormalTok{    endowment\_quantities}\OperatorTok{=}\FloatTok{100.0}\NormalTok{,}
\NormalTok{    activity\_start}\OperatorTok{=}\FloatTok{100.0}\NormalTok{,}
\NormalTok{    name}\OperatorTok{=:}\NormalTok{consumer2,}
\NormalTok{)}

\NormalTok{result }\OperatorTok{=}\NormalTok{ GEMB.}\FunctionTok{solve}\NormalTok{(}
\NormalTok{    model;}
\NormalTok{    p0}\OperatorTok{=}\NormalTok{[}\FloatTok{1.0}\NormalTok{, }\FloatTok{1.0}\NormalTok{],}
\NormalTok{    silent}\OperatorTok{=}\ConstantTok{true}\NormalTok{,}
\NormalTok{)}

\NormalTok{GEMB.}\FunctionTok{print\_equilibrium\_statistics}\NormalTok{(model, result)}

\NormalTok{flows }\OperatorTok{=}\NormalTok{ GEMB.}\FunctionTok{demand\_supply\_matrices}\NormalTok{(model, result)}

\NormalTok{D }\OperatorTok{=}\NormalTok{ flows.demand}
\NormalTok{S }\OperatorTok{=}\NormalTok{ flows.supply}

\FunctionTok{println}\NormalTok{(}\StringTok{"}\SpecialCharTok{\textbackslash{}n}\StringTok{Demand matrix D:"}\NormalTok{)}
\FunctionTok{display}\NormalTok{(D)}

\FunctionTok{println}\NormalTok{(}\StringTok{"}\SpecialCharTok{\textbackslash{}n}\StringTok{Supply matrix S:"}\NormalTok{)}
\FunctionTok{display}\NormalTok{(S)}
\end{Highlighting}
\end{Shaded}

}

\end{codelisting}%

\section{\texorpdfstring{Marshallian-Demand Consumers: A \(2\times2\)
Pure-Exchange
Model}{Marshallian-Demand Consumers: A 2\textbackslash times2 Pure-Exchange Model}}\label{marshallian-demand-consumers-a-2times2-pure-exchange-model}

When a consumer's Marshallian demand is available explicitly as a
function of income and prices, GEMB allows the demand function to be
supplied directly through \texttt{MarshallDemandConsumerSpec}. This
avoids introducing a utility function, marginal-utility conditions, or a
utility activity variable.

A Marshallian-demand specification is constructed as

\begin{Shaded}
\begin{Highlighting}[]
\NormalTok{GEMB.}\FunctionTok{MarshallDemandConsumerSpec}\NormalTok{(demand\_function)}
\end{Highlighting}
\end{Shaded}

where \texttt{demand\_function} takes income and the relevant commodity
prices as arguments and returns a demand vector. A typical function has
the form

\begin{Shaded}
\begin{Highlighting}[]
\NormalTok{demand\_function }\OperatorTok{=}\NormalTok{ (income, prices) }\OperatorTok{{-}\textgreater{}} \ControlFlowTok{begin}
    \CommentTok{\# compute demands}
\NormalTok{    [}\OperatorTok{...}\NormalTok{]}
\ControlFlowTok{end}
\end{Highlighting}
\end{Shaded}

GEMB computes the consumer's income from the value of its endowments and
then evaluates this demand function at the current prices. Thus,
\texttt{MarshallDemandConsumerSpec} is particularly convenient when the
demand system itself, rather than the underlying utility maximization
problem, is the natural starting point of the model.

Unlike an activity-demand consumer, a consumer constructed from
\texttt{MarshallDemandConsumerSpec} does not require an endogenous
utility or activity variable. Its income is determined from its
endowments at the current prices, and its commodity demands are then
obtained directly from the supplied Marshallian-demand function. The
consumer is therefore incorporated into the equilibrium system through
its resulting net supply, while equilibrium prices are determined
jointly by the commodity-market conditions.

This interface is useful not only for demands derived analytically from
standard utility functions, but also whenever a Marshallian demand
system has already been obtained independently and the researcher wishes
to use that demand system directly in a general equilibrium model.

Consider the following pure-exchange economy:

\begin{enumerate}
\def\labelenumi{\arabic{enumi}.}
\item
  The economy contains two commodities and two consumers. Commodity 2 is
  taken as the numeraire commodity.
\item
  Consumer 1 is endowed with \(\omega_1>0\) units of commodity 1 and no
  commodity 2, while consumer 2 is endowed with \(\omega_2>0\) units of
  commodity 2 and no commodity 1.
\item
  Consumer \(i\) allocates a share \(\beta_i\) of income to commodity 1
  and the remaining share \(1-\beta_i\) to commodity 2, where
  \(0<\beta_i<1\). Its Marshallian demand is therefore

  \[
  \mathbf x_i(w_i,\mathbf p)
  =
  \left(
  \frac{\beta_iw_i}{p_1},
  \frac{(1-\beta_i)w_i}{p_2}
  \right)^\top
  \]
\end{enumerate}

Let the price vector be \(\mathbf p=(p,1)^\top\). The incomes of the two
consumers are \(w_1=p\omega_1\) and \(w_2=\omega_2\), respectively.
Hence the demand and supply matrices are

\[
\mathbf D(p)
=
\begin{bmatrix}
\beta_1\omega_1&
\frac{\beta_2\omega_2}{p}\\
(1-\beta_1)p\omega_1&
(1-\beta_2)\omega_2
\end{bmatrix}
\]

and

\[
\mathbf S
=
\begin{bmatrix}
\omega_1&0\\
0&\omega_2
\end{bmatrix}
\]

By solving the equality-form wide-sense structural equilibrium model,
the equilibrium price vector is

\[
\mathbf p^*
=
\left(
\frac{\beta_2\omega_2}
{(1-\beta_1)\omega_1},
1
\right)^\top
\]

For a numerical example, let \(\beta_1=0.5\), \(\beta_2=0.25\), and
\(\omega_1=\omega_2=100\). The equilibrium is then

\[
\mathbf p^*=(0.5,1)^\top
\]

and the equilibrium demand and supply matrices are

\[
\mathbf D^*
=
\begin{bmatrix}
50&50\\
25&75
\end{bmatrix},
\qquad
\mathbf S
=
\begin{bmatrix}
100&0\\
0&100
\end{bmatrix}
\]

The corresponding GEMB implementation is given in
Listing~\ref{lst-gemb_pureExchange_MarshallDemand}.

\begin{codelisting}

\caption{\label{lst-gemb_pureExchange_MarshallDemand}Marshallian-Demand
Consumers: A \(2\times2\) Pure-Exchange Model}

\centering{

\begin{Shaded}
\begin{Highlighting}[]
\ImportTok{using} \BuiltInTok{GeneralEquilibriumModeling}

\NormalTok{beta1 }\OperatorTok{=} \FloatTok{0.5}
\NormalTok{beta2 }\OperatorTok{=} \FloatTok{0.25}
\NormalTok{omega1 }\OperatorTok{=} \FloatTok{100.0}
\NormalTok{omega2 }\OperatorTok{=} \FloatTok{100.0}

\FunctionTok{marshall\_demand}\NormalTok{(beta) }\OperatorTok{=}\NormalTok{ (income, prices) }\OperatorTok{{-}\textgreater{}}\NormalTok{ [}
\NormalTok{    beta }\OperatorTok{*}\NormalTok{ income }\OperatorTok{/}\NormalTok{ prices[}\FloatTok{1}\NormalTok{],}
\NormalTok{    (}\FloatTok{1.0} \OperatorTok{{-}}\NormalTok{ beta) }\OperatorTok{*}\NormalTok{ income }\OperatorTok{/}\NormalTok{ prices[}\FloatTok{2}\NormalTok{],}
\NormalTok{]}

\NormalTok{model }\OperatorTok{=}\NormalTok{ GEMB.}\FunctionTok{GEMBModel}\NormalTok{(}
\NormalTok{    [}
\NormalTok{        GEMB.}\FunctionTok{CommoditySpec}\NormalTok{(}
            \OperatorTok{:}\NormalTok{commodity1;}
\NormalTok{            price\_lower\_bound}\OperatorTok{=}\FloatTok{0.05}\NormalTok{,}
\NormalTok{        ),}
\NormalTok{        GEMB.}\FunctionTok{CommoditySpec}\NormalTok{(}
            \OperatorTok{:}\NormalTok{commodity2;}
\NormalTok{            price\_lower\_bound}\OperatorTok{=}\FloatTok{0.05}\NormalTok{,}
\NormalTok{        ),}
\NormalTok{    ];}
\NormalTok{    numeraire}\OperatorTok{=:}\NormalTok{commodity2,}
\NormalTok{)}

\NormalTok{GEMB.}\FunctionTok{add\_agent!}\NormalTok{(}
\NormalTok{    model,}
\NormalTok{    GEMB.}\FunctionTok{MarshallDemandConsumerSpec}\NormalTok{(}
        \FunctionTok{marshall\_demand}\NormalTok{(beta1),}
\NormalTok{    );}
\NormalTok{    demands}\OperatorTok{=}\NormalTok{[}\OperatorTok{:}\NormalTok{commodity1, }\OperatorTok{:}\NormalTok{commodity2],}
\NormalTok{    endowments}\OperatorTok{=:}\NormalTok{commodity1,}
\NormalTok{    endowment\_quantities}\OperatorTok{=}\NormalTok{omega1,}
\NormalTok{    name}\OperatorTok{=:}\NormalTok{consumer1,}
\NormalTok{)}

\NormalTok{GEMB.}\FunctionTok{add\_agent!}\NormalTok{(}
\NormalTok{    model,}
\NormalTok{    GEMB.}\FunctionTok{MarshallDemandConsumerSpec}\NormalTok{(}
        \FunctionTok{marshall\_demand}\NormalTok{(beta2),}
\NormalTok{    );}
\NormalTok{    demands}\OperatorTok{=}\NormalTok{[}\OperatorTok{:}\NormalTok{commodity1, }\OperatorTok{:}\NormalTok{commodity2],}
\NormalTok{    endowments}\OperatorTok{=:}\NormalTok{commodity2,}
\NormalTok{    endowment\_quantities}\OperatorTok{=}\NormalTok{omega2,}
\NormalTok{    name}\OperatorTok{=:}\NormalTok{consumer2,}
\NormalTok{)}

\NormalTok{result }\OperatorTok{=}\NormalTok{ GEMB.}\FunctionTok{solve}\NormalTok{(}
\NormalTok{    model;}
\NormalTok{    p0}\OperatorTok{=}\NormalTok{[}\FloatTok{1.0}\NormalTok{, }\FloatTok{1.0}\NormalTok{],}
\NormalTok{    silent}\OperatorTok{=}\ConstantTok{true}\NormalTok{,}
\NormalTok{)}

\NormalTok{GEMB.}\FunctionTok{print\_equilibrium\_statistics}\NormalTok{(model, result)}
\end{Highlighting}
\end{Shaded}

}

\end{codelisting}%

\section{\texorpdfstring{Multiple Equilibria: A \(2\times2\)
Model}{Multiple Equilibria: A 2\textbackslash times2 Model}}\label{multiple-equilibria-a-2times2-model}

The example in this section (Mas-Colell, Whinston, Green, 1995, Exercise
15.B.6) illustrates the possibility of multiple equilibria: even with
smooth, strictly monotone, and strictly convex CES preferences, a
general equilibrium need not be unique.

Consider a pure exchange economy. Let \(k=12/37\), and assume the
following:

\begin{enumerate}
\def\labelenumi{\arabic{enumi}.}
\item
  The economy contains two commodities and two consumers.
\item
  Consumer 1 has the utility function
  \(u_1(\mathbf x)=(x_1^{-2}+k^3x_2^{-2})^{-1/2}\) and is endowed with
  one unit of commodity 1.
\item
  Consumer 2 has the utility function
  \(u_2(\mathbf x)=(k^3x_1^{-2}+x_2^{-2})^{-1/2}\) and is endowed with
  one unit of commodity 2.
\end{enumerate}

Let the price vector be \(\mathbf p=(p_1,p_2)^\top\), and define

\[
t=\left(\frac{p_1}{p_2}\right)^{1/3}
\]

The incomes of consumers 1 and 2 are \(w_1=p_1\) and \(w_2=p_2\),
respectively. Utility maximization yields the following demand and
supply matrices:

\[
\mathbf D(t)=
\begin{bmatrix}
\frac{t^2}{t^2+k} & \frac{k}{t(kt^2+1)}\\
\frac{kt^3}{t^2+k} & \frac{1}{kt^2+1}
\end{bmatrix},
\qquad
\mathbf S=
\begin{bmatrix}
1&0\\
0&1
\end{bmatrix}
\]

Using the following equality-form wide-sense structural equilibrium
model,

\[
\mathbf p^\top\mathbf D(t)=\mathbf p^\top\mathbf S
\]

\[
\mathbf D(t)\mathbf 1=\mathbf S\mathbf 1
\]

we obtain

\[
(t-1)[kt^2+(k-1)t+k]=0
\]

Substituting \(k=12/37\) yields

\[
(t-1)(4t-3)(3t-4)=0
\]

and hence

\[
t=\frac{3}{4},\quad 1,\quad \frac{4}{3}
\]

Since

\[
\frac{p_1}{p_2}=t^3
\]

the three equilibrium relative prices are

\[
\frac{27}{64},\quad 1,\quad \frac{64}{27}
\]

Taking commodity 2 as the numeraire, so that \(p_2=1\), the three
equilibrium price vectors are

\[
\mathbf p^{(1)}=\left(\frac{27}{64},1\right)^\top,\qquad
\mathbf p^{(2)}=(1,1)^\top,\qquad
\mathbf p^{(3)}=\left(\frac{64}{27},1\right)^\top
\]

The corresponding equilibrium demand matrices are

\[
\mathbf D^{(1)}
=
\begin{bmatrix}
\frac{111}{175} & \frac{64}{175}\\
\frac{27}{175} & \frac{148}{175}
\end{bmatrix},
\qquad
\mathbf D^{(2)}
=
\begin{bmatrix}
\frac{37}{49} & \frac{12}{49}\\
\frac{12}{49} & \frac{37}{49}
\end{bmatrix},
\qquad
\mathbf D^{(3)}
=
\begin{bmatrix}
\frac{148}{175} & \frac{27}{175}\\
\frac{64}{175} & \frac{111}{175}
\end{bmatrix}
\]

For \(k=12/37\), this multiplicity takes a particularly simple form: the
market-equilibrium equation has exactly three positive solutions for the
relative price.

\begin{codelisting}

\caption{\label{lst-gemb_pureExchange_MWG15B6_nc2_na2}A \(2\times2\)
Equilibrium Model: Multiple Equilibria}

\centering{

\begin{Shaded}
\begin{Highlighting}[]
\ImportTok{using} \BuiltInTok{GeneralEquilibriumModeling.GEM}
\ImportTok{using} \BuiltInTok{GeneralEquilibriumModeling.GEMB}

\NormalTok{k }\OperatorTok{=} \FloatTok{12} \OperatorTok{/} \FloatTok{37}

\NormalTok{model }\OperatorTok{=} \FunctionTok{GEMBModel}\NormalTok{(}
\NormalTok{    [}\OperatorTok{:}\NormalTok{commodity1, }\OperatorTok{:}\NormalTok{commodity2];}
\NormalTok{    numeraire}\OperatorTok{=:}\NormalTok{commodity2,}
\NormalTok{)}

\FunctionTok{add\_agent!}\NormalTok{(}
\NormalTok{    model,}
    \FunctionTok{CESSpec}\NormalTok{([}\FloatTok{1.0}\NormalTok{, k}\OperatorTok{\^{}}\FloatTok{3}\NormalTok{]; es}\OperatorTok{=}\FloatTok{1} \OperatorTok{/} \FloatTok{3}\NormalTok{);}
\NormalTok{    demands}\OperatorTok{=}\NormalTok{[}\OperatorTok{:}\NormalTok{commodity1, }\OperatorTok{:}\NormalTok{commodity2],}
\NormalTok{    endowments}\OperatorTok{=:}\NormalTok{commodity1,}
\NormalTok{    endowment\_quantities}\OperatorTok{=}\FloatTok{1.0}\NormalTok{,}
\NormalTok{    name}\OperatorTok{=:}\NormalTok{consumer1,}
\NormalTok{)}

\FunctionTok{add\_agent!}\NormalTok{(}
\NormalTok{    model,}
    \FunctionTok{CESSpec}\NormalTok{([k}\OperatorTok{\^{}}\FloatTok{3}\NormalTok{, }\FloatTok{1.0}\NormalTok{]; es}\OperatorTok{=}\FloatTok{1} \OperatorTok{/} \FloatTok{3}\NormalTok{);}
\NormalTok{    demands}\OperatorTok{=}\NormalTok{[}\OperatorTok{:}\NormalTok{commodity1, }\OperatorTok{:}\NormalTok{commodity2],}
\NormalTok{    endowments}\OperatorTok{=:}\NormalTok{commodity2,}
\NormalTok{    endowment\_quantities}\OperatorTok{=}\FloatTok{1.0}\NormalTok{,}
\NormalTok{    name}\OperatorTok{=:}\NormalTok{consumer2,}
\NormalTok{)}

\ControlFlowTok{for}\NormalTok{ p0 }\KeywordTok{in}\NormalTok{ (}
\NormalTok{    [}\FloatTok{0.30}\NormalTok{, }\FloatTok{1.0}\NormalTok{],}
\NormalTok{    [}\FloatTok{0.80}\NormalTok{, }\FloatTok{1.0}\NormalTok{],}
\NormalTok{    [}\FloatTok{1.20}\NormalTok{, }\FloatTok{1.0}\NormalTok{],}
\NormalTok{    [}\FloatTok{3.00}\NormalTok{, }\FloatTok{1.0}\NormalTok{],}
\NormalTok{)}
\NormalTok{    result }\OperatorTok{=} \FunctionTok{solve}\NormalTok{(model; p0}\OperatorTok{=}\NormalTok{p0, residual\_tol}\OperatorTok{=}\FloatTok{1e{-}10}\NormalTok{, silent}\OperatorTok{=}\ConstantTok{true}\NormalTok{)}
    \FunctionTok{println}\NormalTok{(}\StringTok{"p0 = "}\NormalTok{, p0, }\StringTok{"  =\textgreater{}  p* = "}\NormalTok{, result.prices)}
\ControlFlowTok{end}
\end{Highlighting}
\end{Shaded}

}

\end{codelisting}%

\section{\texorpdfstring{Money Illusion and Absolute-Price-Dependent
Utility: A \(2\times2\)
Model}{Money Illusion and Absolute-Price-Dependent Utility: A 2\textbackslash times2 Model}}\label{money-illusion-and-absolute-price-dependent-utility-a-2times2-model}

In standard consumer theory, preferences are defined independently of
prices. Accordingly, demand functions are homogeneous of degree zero in
prices and income: if all prices and income are multiplied by the same
positive factor, optimal demand remains unchanged. In this sense, only
relative prices matter, and the absolute price level has no effect on
the equilibrium allocation. A violation of this property may be
interpreted as \textbf{money illusion}.

Consider the following absolute-price-dependent utility function:

\[
u(\mathbf x;p_1)=(x_1-2p_1)x_2
\]

Price-dependent preferences have also been studied in the literature
(see Pollak, 1977). Restrict attention to positive prices and positive
income. If all income is spent, the corresponding demand function is

\[
\mathbf x(\mathbf p,w)=
\begin{cases}
\left(\frac{w}{2p_1}+p_1,\ \frac{w}{2p_2}-\frac{p_1^2}{p_2}\right)^\top,
& w>2p_1^2\\[0.8em]
\left(\frac{w}{p_1},\ 0\right)^\top,
& w\leq2p_1^2
\end{cases}
\]

For example, when \(\lambda>1\) and \(w>2\lambda p_1^2\), the first
commodity demand satisfies
\(x_1(\lambda\mathbf p,\lambda w)=w/(2p_1)+\lambda p_1\neq x_1(\mathbf p,w)\).
Hence the demand function is not homogeneous of degree zero in prices
and income. The absolute price level therefore affects consumer demand.

Now consider a pure-exchange economy with two commodities and two
consumers. The supply matrix is

\[
\mathbf S=
\begin{pmatrix}
s_1&0\\
0&s_2
\end{pmatrix},
\qquad s_1,s_2>0
\]

Consumer 1 owns commodity 1 and has utility function \(u_1=x_{21}\), so
the consumer demands only commodity 2. Consumer 2 owns commodity 2 and
has the absolute-price-dependent utility function introduced above. Let
\(\mathbf p=(p_1,p_2)^\top\), let \(u_1\) denote consumer 1's utility
level, and let \(w_2\) denote consumer 2's income. Restrict attention to
the interior case \(w_2>2p_1^2\). The demand matrix is then

\[
\mathbf D(\mathbf p,u_1,w_2)=
\begin{pmatrix}
0&\frac{w_2}{2p_1}+p_1\\
u_1&\frac{w_2}{2p_2}-\frac{p_1^2}{p_2}
\end{pmatrix}
\]

The equality-form wide-sense structural equilibrium model requires each
consumer's revenue and expenditure to balance and each commodity market
to clear. Because consumer 2's demand depends on the absolute value of
\(p_1\), the equilibrium price vector cannot in general be rescaled
without changing the allocation. Thus, unlike a standard Walrasian
model, fixing one price here is not merely an economically neutral
normalization.

Let \(p_1^*=q\) and regard \(q\) as a given absolute price. The
remaining endogenous variables are \(p_2\), \(u_1\), and \(w_2\).
Consumer 1's revenue-expenditure balance condition is \(p_2u_1=p_1s_1\),
while consumer 2's condition is \(w_2=p_2s_2\). Commodity 2 market
clearing requires

\[
u_1+\frac{w_2}{2p_2}-\frac{p_1^2}{p_2}=s_2
\]

Solving these three conditions for a given \(q\) yields

\[
\mathbf p^*(q)
=
\left(
q,\frac{2q(s_1-q)}{s_2}
\right)^\top,
\qquad
u_1^*(q)=\frac{s_1s_2}{2(s_1-q)},
\qquad
w_2^*(q)=2q(s_1-q)
\]

The interior-solution condition \(w_2^*>2(p_1^*)^2\) implies
\(0<q<s_1/2\). Over this interval, \(0<p_2^*<s_1^2/(2s_2)\). Hence the
model possesses a continuum of equilibria indexed by the absolute price
\(q\). These equilibria are economically distinct rather than
proportional representations of the same relative-price equilibrium.

The model can be calculated directly with the GEMB high-level framework
using \texttt{MarshallDemandConsumerSpec}. For a selected value of
\(q\), the program fixes the absolute price \(p_1=q\) by setting
\texttt{numeraire=:commodity1} and \texttt{numeraire\_value=q}. Because
consumer 2's demand depends on the absolute value of \(p_1\), this
operation should be interpreted as fixing an absolute price rather than
as a simple price normalization.

For consumer 1, GEMB automatically calculates income from the value of
the endowment, \(w_1=p_1s_1\), and the Marshall demand for commodity 2
is \(x_{21}=w_1/p_2\). Hence consumer 1's utility is obtained directly
as \(u_1=x_{21}\). For consumer 2, GEMB similarly calculates income as
\(w_2=p_2s_2\) and passes this income together with the current prices
to the user-defined Marshall demand function

\[
\mathbf x_2(\mathbf p,w_2)
=
\left(
\frac{w_2}{2p_1}+p_1,
\frac{w_2}{2p_2}-\frac{p_1^2}{p_2}
\right)^\top
\]

Thus, neither \(u_1\) nor \(w_2\) needs to be introduced as an
additional endogenous agent variable. Given \(p_1=q\), GEMB constructs
the consumers' net supplies from their endowments and Marshall demands
and determines the remaining price \(p_2\) from market clearing.

The program below takes \(s_1=s_2=4\) and solves the model for two
different absolute prices. When \(p_1=q=1\), the analytical solution is
\(\mathbf p^*=(1,1.5)^\top\), \(u_1^*=8/3\), and \(w_2^*=6\). When
\(p_1=q=1.5\), the analytical solution becomes
\(\mathbf p^*=(1.5,1.875)^\top\), \(u_1^*=3.2\), and \(w_2^*=7.5\). Both
values satisfy the interior-solution condition \(0<q<s_1/2\). The
resulting changes in \(p_2\), \(u_1\), and \(w_2\) demonstrate that
changing the absolute value of \(p_1\) selects economically different
equilibria rather than merely rescaling the same equilibrium price
vector.

\begin{codelisting}

\caption{\label{lst-gemb_moneyIllusion_absolutePriceUtility_nc2_na2}A
\(2\times2\) Pure-Exchange Equilibrium with Absolute-Price-Dependent
Utility}

\centering{

\begin{Shaded}
\begin{Highlighting}[]
\ImportTok{using} \BuiltInTok{GeneralEquilibriumModeling}

\KeywordTok{const}\NormalTok{ S1 }\OperatorTok{=} \FloatTok{4.0}
\KeywordTok{const}\NormalTok{ S2 }\OperatorTok{=} \FloatTok{4.0}
\KeywordTok{const}\NormalTok{ Q }\OperatorTok{=} \FloatTok{1.5}

\NormalTok{model }\OperatorTok{=}\NormalTok{ GEMB.}\FunctionTok{GEMBModel}\NormalTok{(}
\NormalTok{    [}\OperatorTok{:}\NormalTok{commodity1, }\OperatorTok{:}\NormalTok{commodity2];}
\NormalTok{    numeraire}\OperatorTok{=:}\NormalTok{commodity1,}
\NormalTok{    numeraire\_value}\OperatorTok{=}\NormalTok{Q,}
\NormalTok{)}

\NormalTok{consumer1\_demand }\OperatorTok{=}\NormalTok{ (income, prices) }\OperatorTok{{-}\textgreater{}}\NormalTok{ [}
\NormalTok{    income }\OperatorTok{/}\NormalTok{ prices[}\FloatTok{1}\NormalTok{],}
\NormalTok{]}

\NormalTok{GEMB.}\FunctionTok{add\_agent!}\NormalTok{(}
\NormalTok{    model,}
\NormalTok{    GEMB.}\FunctionTok{MarshallDemandConsumerSpec}\NormalTok{(consumer1\_demand);}
\NormalTok{    demands}\OperatorTok{=:}\NormalTok{commodity2,}
\NormalTok{    endowments}\OperatorTok{=:}\NormalTok{commodity1,}
\NormalTok{    endowment\_quantities}\OperatorTok{=}\NormalTok{[S1],}
\NormalTok{    name}\OperatorTok{=:}\NormalTok{consumer1,}
\NormalTok{)}

\NormalTok{consumer2\_demand }\OperatorTok{=} \KeywordTok{function}\NormalTok{ (income, prices)}
\NormalTok{    p1, p2 }\OperatorTok{=}\NormalTok{ prices}

    \ControlFlowTok{return}\NormalTok{ [}
\NormalTok{        income }\OperatorTok{/}\NormalTok{ (}\FloatTok{2}\NormalTok{p1) }\OperatorTok{+}\NormalTok{ p1,}
\NormalTok{        income }\OperatorTok{/}\NormalTok{ (}\FloatTok{2}\NormalTok{p2) }\OperatorTok{{-}}\NormalTok{ p1}\OperatorTok{\^{}}\FloatTok{2} \OperatorTok{/}\NormalTok{ p2,}
\NormalTok{    ]}
\KeywordTok{end}

\NormalTok{GEMB.}\FunctionTok{add\_agent!}\NormalTok{(}
\NormalTok{    model,}
\NormalTok{    GEMB.}\FunctionTok{MarshallDemandConsumerSpec}\NormalTok{(consumer2\_demand);}
\NormalTok{    demands}\OperatorTok{=}\NormalTok{[}\OperatorTok{:}\NormalTok{commodity1, }\OperatorTok{:}\NormalTok{commodity2],}
\NormalTok{    endowments}\OperatorTok{=:}\NormalTok{commodity2,}
\NormalTok{    endowment\_quantities}\OperatorTok{=}\NormalTok{[S2],}
\NormalTok{    name}\OperatorTok{=:}\NormalTok{consumer2,}
\NormalTok{)}

\NormalTok{result }\OperatorTok{=}\NormalTok{ GEMB.}\FunctionTok{solve}\NormalTok{(}
\NormalTok{    model;}
\NormalTok{    p0}\OperatorTok{=}\NormalTok{[Q, }\FloatTok{1.0}\NormalTok{],}
\NormalTok{    residual\_tol}\OperatorTok{=}\FloatTok{1.0e{-}10}\NormalTok{,}
\NormalTok{    silent}\OperatorTok{=}\ConstantTok{true}\NormalTok{,}
\NormalTok{)}

\FunctionTok{println}\NormalTok{(}\StringTok{"Equilibrium prices: "}\NormalTok{, result.prices)}
\end{Highlighting}
\end{Shaded}

}

\end{codelisting}%

\section{\texorpdfstring{Relative-Price-Dependent Utility: A
\(2\times2\)
Model}{Relative-Price-Dependent Utility: A 2\textbackslash times2 Model}}\label{relative-price-dependent-utility-a-2times2-model}

Consider the utility function

\[
u(\mathbf x;\mathbf p)
=
x_1^{\frac{p_1^2}{p_1^2+p_2^2}}
x_2^{\frac{p_2^2}{p_1^2+p_2^2}}
\]

This function has the form of a Cobb--Douglas utility function, but its
expenditure shares depend on the relative price \(p_1/p_2\). It is
therefore a \textbf{relative-price-dependent utility function}. As the
relative price of commodity 1 rises, the consumer assigns a larger
weight to commodity 1.

Relative-price-dependent preferences can be used to represent several
types of price-dependent behavior (see Pollak, 1977). For example,
consumers may use prices as signals of commodity quality, a higher
relative price may increase the attractiveness of a commodity through
conspicuous consumption, or price movements may affect speculative
demand.

Unlike the absolute-price-dependent utility function considered in the
previous section, this utility function is unchanged when all prices are
multiplied by the same positive factor. The preference therefore depends
on relative rather than absolute prices. The associated Marshall demand
is also homogeneous of degree zero in prices and income.

Consider a pure-exchange economy with two commodities and two consumers.
Commodity 2 is the numeraire. The supply matrix is

\[
\mathbf S=
\begin{pmatrix}
s_1&0\\
0&s_2
\end{pmatrix},
\qquad s_1,s_2>0
\]

Consumer 1 owns \(s_1\) units of commodity 1 and has utility function
\(u_1=x_{21}\), so this consumer demands only commodity 2. Consumer 2
owns \(s_2\) units of commodity 2 and has the relative-price-dependent
utility function introduced above.

Let \(\mathbf p=(p_1,1)^\top\). Consumer 1's income is \(p_1s_1\), so
its demand for commodity 2, and hence its utility level, is
\(u_1=p_1s_1\). Consumer 2's income is \(w_2=s_2\), and its Marshall
demands are \(x_{12}=p_1w_2/(p_1^2+1)\) and \(x_{22}=w_2/(p_1^2+1)\).
The demand matrix can therefore be written as

\[
\mathbf D(\mathbf p,u_1,w_2)=
\begin{pmatrix}
0&\frac{p_1w_2}{p_1^2+1}\\
u_1&\frac{w_2}{p_1^2+1}
\end{pmatrix}
\]

For the equality-form wide-sense structural equilibrium model, market
clearing for commodity 1 requires \(p_1s_2/(p_1^2+1)=s_1\).
Equivalently,

\[
s_1p_1^2-s_2p_1+s_1=0
\]

Hence the positive equilibrium prices are

\[
\mathbf p^*
=
\left(
\frac{s_2\pm\sqrt{s_2^2-4s_1^2}}{2s_1},
1
\right)^\top
\]

with \(w_2^*=s_2\) and \(u_1^*=p_1^*s_1\).

The number of positive equilibria depends on the relative magnitudes of
the two supplies. If \(s_2>2s_1\), there are two positive equilibrium
values of \(p_1\), and their product is one, so they are reciprocals. If
\(s_2=2s_1\), there is a unique positive equilibrium with \(p_1^*=1\).
If \(s_2<2s_1\), no positive price can clear both markets.

This result can also be seen directly from consumer 2's demand for
commodity 1, \(x_{12}=p_1s_2/(p_1^2+1)\). As a function of \(p_1\), this
demand first rises and then falls, reaching its maximum value \(s_2/2\)
at \(p_1=1\). Hence commodity 1 can be fully absorbed only if
\(s_1\leq s_2/2\), or equivalently \(s_2\geq2s_1\). When the inequality
is strict, the horizontal supply level \(s_1\) intersects this
nonmonotonic demand function twice, producing two positive equilibrium
prices.

The model can be calculated directly with the GEMB high-level framework
using \texttt{MarshallDemandConsumerSpec}. Since commodity 2 is the
numeraire, the program sets \texttt{numeraire=:commodity2} and
\texttt{numeraire\_value=1.0}. GEMB automatically computes consumer 1's
income from its endowment and price and then evaluates its Marshall
demand for commodity 2. For consumer 2, GEMB similarly calculates
\(w_2=p_2s_2=s_2\) and passes this income together with the current
prices to the user-defined Marshall demand function. No additional
income or utility variables need to be introduced.

An especially useful feature of this example is that the same GEMB model
may have more than one positive equilibrium. For example, let \(s_1=1\)
and \(s_2=3\). The two analytical equilibrium prices of commodity 1 are

\[
p_{1L}^*=\frac{3-\sqrt5}{2}\approx0.381966,
\qquad
p_{1H}^*=\frac{3+\sqrt5}{2}\approx2.618034
\]

The corresponding utility levels of consumer 1 are
\(u_{1L}^*\approx0.381966\) and \(u_{1H}^*\approx2.618034\), while
consumer 2's income is \(w_2^*=3\) in both equilibria. The corresponding
demand matrices are approximately

\[
\mathbf D_L^*=
\begin{pmatrix}
0&1\\
0.381966&2.618034
\end{pmatrix},
\qquad
\mathbf D_H^*=
\begin{pmatrix}
0&1\\
2.618034&0.381966
\end{pmatrix}
\]

Thus, the two equilibrium prices are not merely alternative numerical
representations of the same allocation. They support two different
equilibrium allocations. In numerical computation, the two equilibria
can be obtained by solving the same GEMB model with initial price
vectors chosen near the lower-price and higher-price branches,
respectively. This also provides a simple example of how the initial
value supplied to a nonlinear equilibrium solver can determine which
member of a multiple-equilibrium set is obtained.

If the equality-form wide-sense structural equilibrium model is replaced
by the corresponding mixed-form wide-sense structural equilibrium model,
an additional boundary equilibrium is possible:

\[
\mathbf p^*=(0,1)^\top,\qquad
u_1^*=0,\qquad
w_2^*=s_2
\]

At this equilibrium, both commodities satisfy the supply-demand balance
conditions. For commodity 1, supply exceeds demand, but its equilibrium
price is zero, so the supply-demand balance condition is satisfied
through complementarity between the zero price and excess supply. For
commodity 2, supply equals demand. Thus, the two commodity markets are
both in supply-demand balance, although supply and demand are not equal
for commodity 1. This boundary solution belongs to the mixed-form
wide-sense structural equilibrium model rather than to the
positive-price equality-form equilibrium considered above.

\begin{codelisting}

\caption{\label{lst-gemb_relativePriceDependentUtility_multipleEquilibria_nc2_na2}A
\(2\times2\) Pure-Exchange Model with Relative-Price-Dependent Utility
and Multiple Equilibria}

\centering{

\begin{Shaded}
\begin{Highlighting}[]
\ImportTok{using} \BuiltInTok{GeneralEquilibriumModeling}

\KeywordTok{const}\NormalTok{ S1 }\OperatorTok{=} \FloatTok{1.0}
\KeywordTok{const}\NormalTok{ S2 }\OperatorTok{=} \FloatTok{3.0}

\NormalTok{model }\OperatorTok{=}\NormalTok{ GEMB.}\FunctionTok{GEMBModel}\NormalTok{(}
\NormalTok{    [}\OperatorTok{:}\NormalTok{commodity1, }\OperatorTok{:}\NormalTok{commodity2];}
\NormalTok{    numeraire}\OperatorTok{=:}\NormalTok{commodity2,}
\NormalTok{    numeraire\_value}\OperatorTok{=}\FloatTok{1.0}\NormalTok{,}
\NormalTok{)}

\NormalTok{consumer1\_demand }\OperatorTok{=}\NormalTok{ (income, prices) }\OperatorTok{{-}\textgreater{}}\NormalTok{ [}
\NormalTok{    income }\OperatorTok{/}\NormalTok{ prices[}\FloatTok{1}\NormalTok{],}
\NormalTok{]}

\NormalTok{GEMB.}\FunctionTok{add\_agent!}\NormalTok{(}
\NormalTok{    model,}
\NormalTok{    GEMB.}\FunctionTok{MarshallDemandConsumerSpec}\NormalTok{(consumer1\_demand);}
\NormalTok{    demands}\OperatorTok{=:}\NormalTok{commodity2,}
\NormalTok{    endowments}\OperatorTok{=:}\NormalTok{commodity1,}
\NormalTok{    endowment\_quantities}\OperatorTok{=}\NormalTok{[S1],}
\NormalTok{    name}\OperatorTok{=:}\NormalTok{consumer1,}
\NormalTok{)}

\NormalTok{consumer2\_demand }\OperatorTok{=} \KeywordTok{function}\NormalTok{ (income, prices)}
\NormalTok{    p1, p2 }\OperatorTok{=}\NormalTok{ prices}

    \ControlFlowTok{return}\NormalTok{ [}
\NormalTok{        p1 }\OperatorTok{*}\NormalTok{ income }\OperatorTok{/}\NormalTok{ (p1}\OperatorTok{\^{}}\FloatTok{2} \OperatorTok{+}\NormalTok{ p2}\OperatorTok{\^{}}\FloatTok{2}\NormalTok{),}
\NormalTok{        p2 }\OperatorTok{*}\NormalTok{ income }\OperatorTok{/}\NormalTok{ (p1}\OperatorTok{\^{}}\FloatTok{2} \OperatorTok{+}\NormalTok{ p2}\OperatorTok{\^{}}\FloatTok{2}\NormalTok{),}
\NormalTok{    ]}
\KeywordTok{end}

\NormalTok{GEMB.}\FunctionTok{add\_agent!}\NormalTok{(}
\NormalTok{    model,}
\NormalTok{    GEMB.}\FunctionTok{MarshallDemandConsumerSpec}\NormalTok{(consumer2\_demand);}
\NormalTok{    demands}\OperatorTok{=}\NormalTok{[}\OperatorTok{:}\NormalTok{commodity1, }\OperatorTok{:}\NormalTok{commodity2],}
\NormalTok{    endowments}\OperatorTok{=:}\NormalTok{commodity2,}
\NormalTok{    endowment\_quantities}\OperatorTok{=}\NormalTok{[S2],}
\NormalTok{    name}\OperatorTok{=:}\NormalTok{consumer2,}
\NormalTok{)}

\NormalTok{result\_low }\OperatorTok{=}\NormalTok{ GEMB.}\FunctionTok{solve}\NormalTok{(}
\NormalTok{    model;}
\NormalTok{    p0}\OperatorTok{=}\NormalTok{[}\FloatTok{0.3}\NormalTok{, }\FloatTok{1.0}\NormalTok{],}
\NormalTok{    residual\_tol}\OperatorTok{=}\FloatTok{1.0e{-}10}\NormalTok{,}
\NormalTok{    silent}\OperatorTok{=}\ConstantTok{true}\NormalTok{,}
\NormalTok{)}

\NormalTok{result\_high }\OperatorTok{=}\NormalTok{ GEMB.}\FunctionTok{solve}\NormalTok{(}
\NormalTok{    model;}
\NormalTok{    p0}\OperatorTok{=}\NormalTok{[}\FloatTok{3.0}\NormalTok{, }\FloatTok{1.0}\NormalTok{],}
\NormalTok{    residual\_tol}\OperatorTok{=}\FloatTok{1.0e{-}10}\NormalTok{,}
\NormalTok{    silent}\OperatorTok{=}\ConstantTok{true}\NormalTok{,}
\NormalTok{)}

\NormalTok{result\_boundary }\OperatorTok{=}\NormalTok{ GEMB.}\FunctionTok{solve}\NormalTok{(}
\NormalTok{    model;}
\NormalTok{    p0}\OperatorTok{=}\NormalTok{[}\FloatTok{0.0}\NormalTok{, }\FloatTok{1.0}\NormalTok{],}
\NormalTok{    residual\_tol}\OperatorTok{=}\FloatTok{1.0e{-}10}\NormalTok{,}
\NormalTok{    silent}\OperatorTok{=}\ConstantTok{true}\NormalTok{,}
\NormalTok{)}

\NormalTok{p1\_low\_exact }\OperatorTok{=}\NormalTok{ (S2 }\OperatorTok{{-}} \FunctionTok{sqrt}\NormalTok{(S2}\OperatorTok{\^{}}\FloatTok{2} \OperatorTok{{-}} \FloatTok{4}\NormalTok{S1}\OperatorTok{\^{}}\FloatTok{2}\NormalTok{)) }\OperatorTok{/}\NormalTok{ (}\FloatTok{2}\NormalTok{S1)}
\NormalTok{p1\_high\_exact }\OperatorTok{=}\NormalTok{ (S2 }\OperatorTok{+} \FunctionTok{sqrt}\NormalTok{(S2}\OperatorTok{\^{}}\FloatTok{2} \OperatorTok{{-}} \FloatTok{4}\NormalTok{S1}\OperatorTok{\^{}}\FloatTok{2}\NormalTok{)) }\OperatorTok{/}\NormalTok{ (}\FloatTok{2}\NormalTok{S1)}

\FunctionTok{println}\NormalTok{(}\StringTok{"Lower positive equilibrium prices: "}\NormalTok{, result\_low.prices)}
\FunctionTok{println}\NormalTok{(}\StringTok{"Higher positive equilibrium prices:"}\NormalTok{, result\_high.prices)}
\FunctionTok{println}\NormalTok{(}\StringTok{"Boundary equilibrium prices:       "}\NormalTok{, result\_boundary.prices)}
\FunctionTok{println}\NormalTok{(}
    \StringTok{"Boundary net supply:        "}\NormalTok{,}
\NormalTok{    result\_boundary.total\_net\_supply,}
\NormalTok{)}

\ConstantTok{nothing}
\end{Highlighting}
\end{Shaded}

}

\end{codelisting}%

\section{\texorpdfstring{Asset-Exchange Equilibrium with AMSD Utility: A
\(2\times2\)
Model}{Asset-Exchange Equilibrium with AMSD Utility: A 2\textbackslash times2 Model}}\label{asset-exchange-equilibrium-with-amsd-utility-a-2times2-model}

Suppose there are \(n\) assets and \(m\) investors. Let \(\mathbf x_i\)
denote the portfolio held by investor \(i\), \(\boldsymbol\mu_i\) the
vector of subjective expected asset payoffs, and \(\boldsymbol\Sigma_i\)
the corresponding covariance matrix. Under additive
mean-standard-deviation (AMSD) preferences, investor \(i\) has utility

\[
U_i(\mathbf x_i)
=
\boldsymbol\mu_i^\top\mathbf x_i
-
\gamma_i
\sqrt{\mathbf x_i^\top\boldsymbol\Sigma_i\mathbf x_i}
\]

where \(\gamma_i\geq0\) is the risk-aversion coefficient (Nakamura,
2015). Thus, expected payoff increases utility, whereas portfolio risk
measured by the standard deviation reduces utility.

\texttt{GeneralEquilibriumModeling.jl} provides the high-level function
\texttt{GEMB.solve\_asset\_equilibrium\_amsd} for computing this
equilibrium. Short selling is excluded in the implementation. Suppose
\texttt{Supply}, \texttt{PMP}, and \texttt{PSD} are \(n\times m\)
matrices containing initial asset endowments, subjective expected
payoffs, and subjective payoff standard deviations, respectively;
\texttt{Cor} is the \(n\times n\) payoff correlation matrix; and
\texttt{gamma} contains the \(m\) risk-aversion coefficients. The
equilibrium can then be computed directly by

\begin{Shaded}
\begin{Highlighting}[]
\ImportTok{using} \BuiltInTok{GeneralEquilibriumModeling}

\NormalTok{result }\OperatorTok{=}\NormalTok{ GEMB.}\FunctionTok{solve\_asset\_equilibrium\_amsd}\NormalTok{(}
\NormalTok{    Supply}\OperatorTok{=}\NormalTok{Supply,}
\NormalTok{    PMP}\OperatorTok{=}\NormalTok{PMP,}
\NormalTok{    PSD}\OperatorTok{=}\NormalTok{PSD,}
\NormalTok{    Cor}\OperatorTok{=}\NormalTok{Cor,}
\NormalTok{    gamma}\OperatorTok{=}\NormalTok{gamma,}
\NormalTok{)}
\end{Highlighting}
\end{Shaded}

For each investor, GEMB internally constructs a
\texttt{MeanStandardDeviationMarginalUtilitySpec}, builds the
corresponding pure-exchange consumer, and solves the resulting mixed
complementarity problem through GEM. For numerical stability, the
software uses a small smoothing parameter in the portfolio standard
deviation.

By default, the last asset is used as the numeraire. The main results
include \texttt{result.p} for equilibrium asset prices,
\texttt{result.D} for the equilibrium holding matrix, and
\texttt{result.lambda} for the investors' budget multipliers. Quantities
such as \texttt{result.VMU}, \texttt{result.market\_residual}, and
\texttt{result.kkt\_passed} can be used to examine the marginal-utility
and equilibrium conditions.

\begin{codelisting}

\caption{\label{lst-gemb_assetEquilibriumAMSD_nc2_na2}A \(2\times2\)
Asset-Exchange Equilibrium Model}

\centering{

\begin{Shaded}
\begin{Highlighting}[]
\ImportTok{using} \BuiltInTok{GeneralEquilibriumModeling}

\NormalTok{Supply }\OperatorTok{=}\NormalTok{ [}
    \FloatTok{100.0}  \FloatTok{100.0}
    \FloatTok{100.0}  \FloatTok{100.0}
\NormalTok{]}

\NormalTok{PMP }\OperatorTok{=}\NormalTok{ [}
    \FloatTok{1.20}  \FloatTok{1.05}
    \FloatTok{1.05}  \FloatTok{1.20}
\NormalTok{]}

\NormalTok{PSD }\OperatorTok{=}\NormalTok{ [}
    \FloatTok{0.20}  \FloatTok{0.20}
    \FloatTok{0.30}  \FloatTok{0.30}
\NormalTok{]}

\NormalTok{Cor }\OperatorTok{=}\NormalTok{ [}
    \FloatTok{1.0}  \FloatTok{0.0}
    \FloatTok{0.0}  \FloatTok{1.0}
\NormalTok{]}

\NormalTok{gamma }\OperatorTok{=}\NormalTok{ [}\FloatTok{0.7}\NormalTok{, }\FloatTok{1.0}\NormalTok{]}

\NormalTok{result }\OperatorTok{=}\NormalTok{ GEMB.}\FunctionTok{solve\_asset\_equilibrium\_amsd}\NormalTok{(}
\NormalTok{    Supply}\OperatorTok{=}\NormalTok{Supply,}
\NormalTok{    gamma}\OperatorTok{=}\NormalTok{gamma,}
\NormalTok{    PMP}\OperatorTok{=}\NormalTok{PMP,}
\NormalTok{    PSD}\OperatorTok{=}\NormalTok{PSD,}
\NormalTok{    Cor}\OperatorTok{=}\NormalTok{Cor,}
\NormalTok{    x0}\OperatorTok{=}\FunctionTok{fill}\NormalTok{(}\FloatTok{100.0}\NormalTok{, }\FloatTok{2}\NormalTok{, }\FloatTok{2}\NormalTok{),}
\NormalTok{    silent}\OperatorTok{=}\ConstantTok{true}\NormalTok{,}
\NormalTok{)}

\FunctionTok{println}\NormalTok{(}\StringTok{"Prices: "}\NormalTok{, result.p)}
\FunctionTok{println}\NormalTok{(}\StringTok{"Holdings:"}\NormalTok{)}
\FunctionTok{display}\NormalTok{(result.D)}
\end{Highlighting}
\end{Shaded}

}

\end{codelisting}%

\section{\texorpdfstring{Matthew Effect in Asset-Exchange Equilibrium: A
\(2\times2\)
Model}{Matthew Effect in Asset-Exchange Equilibrium: A 2\textbackslash times2 Model}}\label{matthew-effect-in-asset-exchange-equilibrium-a-2times2-model}

The Matthew effect refers to a process in which an initially small
economic disparity becomes larger through subsequent interactions. In
asset markets, such an effect may arise when investors' attitudes toward
risk depend on their wealth. A relatively poor investor may become more
risk averse and consequently hold more low-risk, low-expected-payoff
assets, whereas a relatively rich investor may tolerate more risk and
hold more high-risk, high-expected-payoff assets (Björk, Murgoci, Zhou,
2014). Asset exchange may therefore widen an initially small wealth
difference.

The following example illustrates this mechanism using the AMSD
preferences introduced in the previous section. There are two assets and
two investors. Their initial asset-endowment matrix is

\[
\mathbf S=
\begin{pmatrix}
0.49&0.51\\
0.49&0.51
\end{pmatrix}
\]

where rows represent assets and columns represent investors. Thus,
investor 1 initially owns \(0.49\) units of each asset, whereas investor
2 owns \(0.51\) units of each asset.

Both investors have the same subjective expected asset payoffs,

\[
\boldsymbol{\mu}_1
=
\boldsymbol{\mu}_2
=
\begin{pmatrix}
120\\
100
\end{pmatrix}
\]

and the same payoff standard deviations,

\[
\boldsymbol{\sigma}_1
=
\boldsymbol{\sigma}_2
=
\begin{pmatrix}
50\\
10
\end{pmatrix}
\]

The two asset payoffs are assumed to be uncorrelated, and the investors'
common correlation matrix is

\[
\mathbf R=
\begin{pmatrix}
1&0\\
0&1
\end{pmatrix}
\]

Asset 1 therefore has both the higher expected payoff and the higher
risk, whereas asset 2 has the lower expected payoff and the lower risk.

The distinguishing feature of this example is that the risk-aversion
coefficient depends on the investor's wealth. In the present example,
expected portfolio payoff \(w_i=\boldsymbol{\mu}_i^\top\mathbf x_i\) is
used as a simple indicator of investor \(i\)'s wealth. Taking \(110\) as
the benchmark, set \(\gamma_i=0.8\) when \(w_i<110/1.02\),
\(\gamma_i=0.4\) when \(110/1.02\leq w_i\leq110\times1.02\), and
\(\gamma_i=0.2\) when \(w_i>110\times1.02\). Thus, lower expected wealth
is associated with greater risk aversion, whereas higher expected wealth
is associated with lower risk aversion.

Under the initial endowments, the expected portfolio payoffs are
\((107.8,112.2)\), so the initial payoff gap is only \(4.4\). The
corresponding initial risk-aversion coefficients are \((0.8,0.4)\).

For a given vector of risk-aversion coefficients,
\texttt{GEMB.solve\_asset\_equilibrium\_amsd} directly computes the
asset-exchange equilibrium. The expected portfolio payoffs implied by
the resulting holdings are then used to update the risk-aversion
coefficients. This process is repeated until the risk-aversion
coefficients are consistent with the equilibrium portfolios that they
generate.

In this example, the procedure reaches a fixed point with
\(\boldsymbol{\gamma}^*=(0.8,0.2)\). Using asset 2 as the numeraire, the
equilibrium asset-price vector is approximately

\[
\mathbf p^*
=
\begin{pmatrix}
1.10036\\
1
\end{pmatrix}
\]

and the equilibrium portfolio matrix is approximately

\[
\mathbf D^*
=
\begin{pmatrix}
0.09245&0.90755\\
0.92745&0.07255
\end{pmatrix}
\]

Investor 1, who starts with the lower expected wealth and has the higher
risk-aversion coefficient, ends up holding mainly the lower-risk asset
2. Investor 2, who starts with the higher expected wealth and eventually
has the lower risk-aversion coefficient, ends up holding mainly the
higher-risk and higher-expected-payoff asset 1.

The equilibrium expected portfolio payoffs are approximately
\((103.839,116.161)\). Hence the payoff gap increases from \(4.4\) to
approximately \(12.322\). Asset exchange therefore amplifies the initial
difference in expected wealth in this example. The mechanism is
self-reinforcing: lower expected wealth is associated with greater risk
aversion and a portfolio tilted toward the lower-risk and
lower-expected-payoff asset, whereas higher expected wealth is
associated with lower risk aversion and greater exposure to the
higher-risk and higher-expected-payoff asset.

Computationally, the model has a simple two-step structure. Conditional
on a given \(\boldsymbol{\gamma}\),
\texttt{GEMB.solve\_asset\_equilibrium\_amsd} solves the asset-exchange
equilibrium directly. The wealth-dependent risk-aversion rule then
updates \(\boldsymbol{\gamma}\), and the calculation continues until
this rule and the asset-market equilibrium are mutually consistent.

\begin{codelisting}

\caption{\label{lst-gemb_assetExchange_MatthewEffect_AMSD_nc2_na2}A
\(2\times2\) Asset-Exchange Equilibrium Model Illustrating the Matthew
Effect}

\centering{

\begin{Shaded}
\begin{Highlighting}[]
\ImportTok{using} \BuiltInTok{GeneralEquilibriumModeling}


\KeywordTok{function} \FunctionTok{gamma\_from\_payoff}\NormalTok{(expected\_payoff, low\_threshold, high\_threshold)}
\NormalTok{    tol }\OperatorTok{=} \FloatTok{1.0e{-}10}

    \ControlFlowTok{if}\NormalTok{ expected\_payoff }\OperatorTok{\textgreater{}}\NormalTok{ high\_threshold }\OperatorTok{+}\NormalTok{ tol}
        \ControlFlowTok{return} \FloatTok{0.2}
    \ControlFlowTok{elseif}\NormalTok{ expected\_payoff }\OperatorTok{\textless{}}\NormalTok{ low\_threshold }\OperatorTok{{-}}\NormalTok{ tol}
        \ControlFlowTok{return} \FloatTok{0.8}
    \ControlFlowTok{else}
        \ControlFlowTok{return} \FloatTok{0.4}
    \ControlFlowTok{end}
\KeywordTok{end}


\FunctionTok{portfolio\_expected\_payoff}\NormalTok{(PMP, D) }\OperatorTok{=}
    \FunctionTok{vec}\NormalTok{(}\FunctionTok{sum}\NormalTok{(PMP }\OperatorTok{.*}\NormalTok{ D; dims}\OperatorTok{=}\FloatTok{1}\NormalTok{))}


\KeywordTok{function} \FunctionTok{main}\NormalTok{()}
\NormalTok{    Supply }\OperatorTok{=}\NormalTok{ [}
        \FloatTok{0.49}  \FloatTok{0.51}
        \FloatTok{0.49}  \FloatTok{0.51}
\NormalTok{    ]}

\NormalTok{    PMP }\OperatorTok{=}\NormalTok{ [}
        \FloatTok{120.0}  \FloatTok{120.0}
        \FloatTok{100.0}  \FloatTok{100.0}
\NormalTok{    ]}

\NormalTok{    PSD }\OperatorTok{=}\NormalTok{ [}
        \FloatTok{50.0}  \FloatTok{50.0}
        \FloatTok{10.0}  \FloatTok{10.0}
\NormalTok{    ]}

\NormalTok{    Cor }\OperatorTok{=}\NormalTok{ [}
        \FloatTok{1.0}  \FloatTok{0.0}
        \FloatTok{0.0}  \FloatTok{1.0}
\NormalTok{    ]}

\NormalTok{    payoff\_benchmark }\OperatorTok{=} \FloatTok{110.0}
\NormalTok{    low\_threshold }\OperatorTok{=}\NormalTok{ payoff\_benchmark }\OperatorTok{/} \FloatTok{1.02}
\NormalTok{    high\_threshold }\OperatorTok{=}\NormalTok{ payoff\_benchmark }\OperatorTok{*} \FloatTok{1.02}

\NormalTok{    initial\_expected\_payoff }\OperatorTok{=}
        \FunctionTok{portfolio\_expected\_payoff}\NormalTok{(PMP, Supply)}

\NormalTok{    gamma }\OperatorTok{=}
        \FunctionTok{gamma\_from\_payoff}\NormalTok{.(}
\NormalTok{            initial\_expected\_payoff,}
\NormalTok{            low\_threshold,}
\NormalTok{            high\_threshold,}
\NormalTok{        )}

\NormalTok{    p0 }\OperatorTok{=}\NormalTok{ [}\FloatTok{1.0}\NormalTok{, }\FloatTok{1.0}\NormalTok{]}
\NormalTok{    x0 }\OperatorTok{=} \FunctionTok{copy}\NormalTok{(Supply)}
\NormalTok{    lambda0 }\OperatorTok{=} \FunctionTok{ones}\NormalTok{(}\FloatTok{2}\NormalTok{)}

\NormalTok{    final\_result }\OperatorTok{=} \ConstantTok{nothing}
\NormalTok{    final\_expected\_payoff }\OperatorTok{=} \ConstantTok{nothing}
\NormalTok{    iterations }\OperatorTok{=} \FloatTok{0}

\NormalTok{    seen }\OperatorTok{=} \FunctionTok{Set}\DataTypeTok{\{Tuple\{Float64,Float64\}\}}\NormalTok{()}

    \ControlFlowTok{for}\NormalTok{ iteration }\KeywordTok{in} \FloatTok{1}\OperatorTok{:}\FloatTok{10}
\NormalTok{        iterations }\OperatorTok{=}\NormalTok{ iteration}
\NormalTok{        state }\OperatorTok{=}\NormalTok{ (gamma[}\FloatTok{1}\NormalTok{], gamma[}\FloatTok{2}\NormalTok{])}

\NormalTok{        state }\KeywordTok{in}\NormalTok{ seen }\OperatorTok{\&\&} \FunctionTok{error}\NormalTok{(}
            \StringTok{"Risk{-}aversion updating entered a cycle before reaching a fixed point."}
\NormalTok{        )}
        \FunctionTok{push!}\NormalTok{(seen, state)}

\NormalTok{        result }\OperatorTok{=}\NormalTok{ GEMB.}\FunctionTok{solve\_asset\_equilibrium\_amsd}\NormalTok{(}
\NormalTok{            Supply}\OperatorTok{=}\NormalTok{Supply,}
\NormalTok{            gamma}\OperatorTok{=}\NormalTok{gamma,}
\NormalTok{            PMP}\OperatorTok{=}\NormalTok{PMP,}
\NormalTok{            PSD}\OperatorTok{=}\NormalTok{PSD,}
\NormalTok{            Cor}\OperatorTok{=}\NormalTok{Cor,}
\NormalTok{            numeraire\_index}\OperatorTok{=}\FloatTok{2}\NormalTok{,}
\NormalTok{            numeraire\_value}\OperatorTok{=}\FloatTok{1.0}\NormalTok{,}
\NormalTok{            p0}\OperatorTok{=}\NormalTok{p0,}
\NormalTok{            x0}\OperatorTok{=}\NormalTok{x0,}
\NormalTok{            lambda0}\OperatorTok{=}\NormalTok{lambda0,}
\NormalTok{            residual\_tol}\OperatorTok{=}\FloatTok{1.0e{-}8}\NormalTok{,}
\NormalTok{            silent}\OperatorTok{=}\ConstantTok{true}\NormalTok{,}
\NormalTok{        )}

\NormalTok{        result.solved }\OperatorTok{||} \FunctionTok{error}\NormalTok{(}
            \StringTok{"The conditional AMSD asset equilibrium was not solved."}
\NormalTok{        )}

\NormalTok{        expected\_payoff }\OperatorTok{=}
            \FunctionTok{portfolio\_expected\_payoff}\NormalTok{(PMP, result.D)}

\NormalTok{        new\_gamma }\OperatorTok{=}
            \FunctionTok{gamma\_from\_payoff}\NormalTok{.(}
\NormalTok{                expected\_payoff,}
\NormalTok{                low\_threshold,}
\NormalTok{                high\_threshold,}
\NormalTok{            )}

\NormalTok{        final\_result }\OperatorTok{=}\NormalTok{ result}
\NormalTok{        final\_expected\_payoff }\OperatorTok{=}\NormalTok{ expected\_payoff}

        \ControlFlowTok{if}\NormalTok{ new\_gamma }\OperatorTok{==}\NormalTok{ gamma}
            \ControlFlowTok{break}
        \ControlFlowTok{end}

\NormalTok{        gamma }\OperatorTok{=}\NormalTok{ new\_gamma}
\NormalTok{        p0 }\OperatorTok{=} \FunctionTok{copy}\NormalTok{(result.p)}
\NormalTok{        x0 }\OperatorTok{=} \FunctionTok{copy}\NormalTok{(result.D)}
\NormalTok{        lambda0 }\OperatorTok{=} \FunctionTok{copy}\NormalTok{(result.lambda)}
    \ControlFlowTok{end}

\NormalTok{    final\_result }\OperatorTok{===} \ConstantTok{nothing} \OperatorTok{\&\&}
        \FunctionTok{error}\NormalTok{(}\StringTok{"No equilibrium was computed."}\NormalTok{)}

    \PreprocessorTok{@assert}\NormalTok{ final\_result.solved}
    \PreprocessorTok{@assert}\NormalTok{ final\_result.kkt\_passed}
    \PreprocessorTok{@assert}\NormalTok{ final\_result.value\_marginal\_utility\_conditions\_passed}

\NormalTok{    initial\_gap }\OperatorTok{=}
        \FunctionTok{maximum}\NormalTok{(initial\_expected\_payoff) }\OperatorTok{{-}}
        \FunctionTok{minimum}\NormalTok{(initial\_expected\_payoff)}

\NormalTok{    final\_gap }\OperatorTok{=}
        \FunctionTok{maximum}\NormalTok{(final\_expected\_payoff) }\OperatorTok{{-}}
        \FunctionTok{minimum}\NormalTok{(final\_expected\_payoff)}

    \FunctionTok{println}\NormalTok{(}\StringTok{"========== AMSD asset equilibrium =========="}\NormalTok{)}
    \FunctionTok{println}\NormalTok{(}\StringTok{"Iterations:        "}\NormalTok{, iterations)}
    \FunctionTok{println}\NormalTok{(}\StringTok{"Prices:            "}\NormalTok{, final\_result.p)}
    \FunctionTok{println}\NormalTok{(}\StringTok{"Holdings:"}\NormalTok{)}
    \FunctionTok{display}\NormalTok{(final\_result.D)}
    \FunctionTok{println}\NormalTok{(}\StringTok{"Expected payoffs:  "}\NormalTok{, final\_expected\_payoff)}
    \FunctionTok{println}\NormalTok{(}\StringTok{"Gamma:             "}\NormalTok{, gamma)}
    \FunctionTok{println}\NormalTok{(}\StringTok{"Payoff gap:        "}\NormalTok{, initial\_gap, }\StringTok{" {-}\textgreater{} "}\NormalTok{, final\_gap)}
    \FunctionTok{println}\NormalTok{(}\StringTok{"Run completed successfully."}\NormalTok{)}

    \ControlFlowTok{return}\NormalTok{ final\_result}
\KeywordTok{end}


\NormalTok{result }\OperatorTok{=} \FunctionTok{main}\NormalTok{()}
\ConstantTok{nothing}
\end{Highlighting}
\end{Shaded}

}

\end{codelisting}%

\bookmarksetup{startatroot}

\chapter{Static Equilibrium with Production: Basic
Models}\label{static-equilibrium-with-production-basic-models}

\section{Production and Transformation
Specifications}\label{production-and-transformation-specifications}

GEMB provides several specifications for describing production
technologies and the transformation of an activity level into outputs.
These specifications can be divided into two broad classes. A production
specification describes the relationship between production activity and
inputs, whereas a transformation specification describes the
relationship between activity level and outputs.

\subsection{Production-Function
Specification}\label{production-function-specification}

\texttt{ProductionFunctionSpec} can be used when the production function
and its marginal products are known explicitly. Let the production
function be

\[
y=f(\mathbf x)
\]

where \(\mathbf x=(x_1,\ldots,x_n)^\top\) is the input vector. The user
supplies both \(f(\mathbf x)\) and a marginal-product vector. The latter
may be either the exact gradient

\[
\nabla f(\mathbf x)
=
\left(
\frac{\partial f}{\partial x_1},
\ldots,
\frac{\partial f}{\partial x_n}
\right)^\top
\]

or this vector multiplied by a common strictly positive factor, since
such a scaling does not change the production stationarity conditions.

For example, consider the Cobb--Douglas production function

\[
f(\mathbf x)=x_1^{0.6}x_2^{0.4}
\]

Its marginal-product vector can be multiplied by the common positive
factor \(x_1^{0.4}x_2^{0.6}\), giving the simpler proportional vector
\((0.6x_2,0.4x_1)^\top\). The specification can therefore be written as

\begin{Shaded}
\begin{Highlighting}[]
\NormalTok{production\_function }\OperatorTok{=}\NormalTok{ x }\OperatorTok{{-}\textgreater{}}
\NormalTok{    x[}\FloatTok{1}\NormalTok{]}\OperatorTok{\^{}}\FloatTok{0.6} \OperatorTok{*}\NormalTok{ x[}\FloatTok{2}\NormalTok{]}\OperatorTok{\^{}}\FloatTok{0.4}

\NormalTok{marginal\_product\_function }\OperatorTok{=}\NormalTok{ x }\OperatorTok{{-}\textgreater{}}\NormalTok{ [}
    \FloatTok{0.6} \OperatorTok{*}\NormalTok{ x[}\FloatTok{2}\NormalTok{],}
    \FloatTok{0.4} \OperatorTok{*}\NormalTok{ x[}\FloatTok{1}\NormalTok{],}
\NormalTok{]}

\NormalTok{spec }\OperatorTok{=}\NormalTok{ GEMB.}\FunctionTok{ProductionFunctionSpec}\NormalTok{(}
\NormalTok{    production\_function,}
\NormalTok{    marginal\_product\_function;}
\NormalTok{    behavior}\OperatorTok{=:}\NormalTok{stationary,}
\NormalTok{)}
\end{Highlighting}
\end{Shaded}

The keyword \texttt{behavior} determines how the production technology
is incorporated into the equilibrium system. With
\texttt{behavior=:stationary}, GEMB constructs production stationarity
conditions. The current stationary formulation assumes an interior
solution, so the marginal-pricing conditions are imposed as equalities.
With \texttt{behavior=:cost\_min}, GEMB instead constructs
cost-minimization KKT conditions, which allow corner solutions. Thus,
the production function describes the technology, while
\texttt{behavior} specifies the conditions imposed on the producer.

Commodity mappings are not contained in \texttt{ProductionFunctionSpec}.
They are supplied when the producer is added to the model. For example,

\begin{Shaded}
\begin{Highlighting}[]
\NormalTok{GEMB.}\FunctionTok{add\_agent!}\NormalTok{(}
\NormalTok{    model,}
\NormalTok{    spec;}
\NormalTok{    outputs}\OperatorTok{=:}\NormalTok{product,}
\NormalTok{    demands}\OperatorTok{=}\NormalTok{[}\OperatorTok{:}\NormalTok{labor, }\OperatorTok{:}\NormalTok{capital],}
\NormalTok{    activity\_start}\OperatorTok{=}\FloatTok{100.0}\NormalTok{,}
\NormalTok{    demand\_start}\OperatorTok{=}\NormalTok{[}\FloatTok{100.0}\NormalTok{, }\FloatTok{100.0}\NormalTok{],}
\NormalTok{    name}\OperatorTok{=:}\NormalTok{firm,}
\NormalTok{)}
\end{Highlighting}
\end{Shaded}

The resulting producer contains an activity level, endogenous input
quantities, and a production multiplier.

\subsection{Power-Production
Specification}\label{power-production-specification}

For the commonly used one-input power production function

\[
y=\alpha x^\theta,
\qquad
\alpha>0,\quad\theta>0
\]

GEMB provides the simpler built-in \texttt{PowerProductionSpec}:

\begin{Shaded}
\begin{Highlighting}[]
\NormalTok{GEMB.}\FunctionTok{PowerProductionSpec}\NormalTok{(}
\NormalTok{    alpha}\OperatorTok{=}\FloatTok{1.0}\NormalTok{,}
\NormalTok{    theta}\OperatorTok{=}\FloatTok{1.0}\NormalTok{,}
\NormalTok{)}
\end{Highlighting}
\end{Shaded}

If the production activity is \(y\), the required input is obtained
directly from the inverse production function,

\[
x
=
\left(
\frac{y}{\alpha}
\right)^{1/\theta}
\]

Thus, unlike \texttt{ProductionFunctionSpec},
\texttt{PowerProductionSpec} does not introduce the input quantity
through a separate production-function constraint and marginal-product
conditions. Instead, the input demand is computed directly as a function
of the activity level. It is therefore an activity-demand specification
and requires exactly one input commodity.

A producer can, for example, be constructed by

\begin{Shaded}
\begin{Highlighting}[]
\NormalTok{GEMB.}\FunctionTok{add\_agent!}\NormalTok{(}
\NormalTok{    model,}
\NormalTok{    GEMB.}\FunctionTok{PowerProductionSpec}\NormalTok{(}
\NormalTok{        alpha}\OperatorTok{=}\FloatTok{2.0}\NormalTok{,}
\NormalTok{        theta}\OperatorTok{=}\FloatTok{0.5}\NormalTok{,}
\NormalTok{    );}
\NormalTok{    outputs}\OperatorTok{=:}\NormalTok{product,}
\NormalTok{    demands}\OperatorTok{=:}\NormalTok{labor,}
\NormalTok{    activity\_start}\OperatorTok{=}\FloatTok{100.0}\NormalTok{,}
\NormalTok{    name}\OperatorTok{=:}\NormalTok{firm,}
\NormalTok{)}
\end{Highlighting}
\end{Shaded}

By default, a producer based on \texttt{PowerProductionSpec} uses
\texttt{TotalRevenueExpenditureBalanceConditions}. Hence the default
specification imposes total revenue-expenditure balance at a positive
activity level. When \(\theta\ne1\), this condition should not be
confused with marginal-cost pricing or the first-order condition of
unconstrained profit maximization. Alternative producer conditions may
be supplied explicitly when required by the economic model.

\texttt{ProductionFunctionSpec} and \texttt{PowerProductionSpec}
therefore provide two different ways to represent production. The former
explicitly introduces input quantities and production first-order
conditions, whereas the latter substitutes the inverse production
function directly into the producer's conditional demand.

\subsection{CET Transformation
Specification}\label{cet-transformation-specification}

When one production activity generates several possible outputs, their
proportions need not be fixed. GEMB provides \texttt{CETSpec} for a
constant-elasticity-of-transformation (CET) output specification:

\begin{Shaded}
\begin{Highlighting}[]
\NormalTok{GEMB.}\FunctionTok{CETSpec}\NormalTok{(}
\NormalTok{    beta;}
\NormalTok{    et}\OperatorTok{=}\FloatTok{1.0}\NormalTok{,}
\NormalTok{    alpha}\OperatorTok{=}\FloatTok{1.0}\NormalTok{,}
\NormalTok{)}
\end{Highlighting}
\end{Shaded}

Let \(\boldsymbol{\beta}=(\beta_1,\ldots,\beta_m)^\top\) denote the
output weights, \(\eta\geq0\) denote the elasticity of transformation
corresponding to \texttt{et}, and \(\alpha>0\) denote the scale
parameter. For an output-price vector
\(\mathbf p=(p_1,\ldots,p_m)^\top\), GEMB uses the unit-revenue function

\[
r(\mathbf p)
=
\frac{1}{\alpha}
\left(
\sum_{i=1}^m
\beta_i p_i^{1+\eta}
\right)^{1/(1+\eta)}
\]

If the activity level is \(z\), the conditional supply of output \(i\)
is

\[
y_i
=
\frac{z}{\alpha}
\beta_i p_i^\eta
\left(
\sum_{j=1}^m
\beta_jp_j^{1+\eta}
\right)^{-\eta/(1+\eta)}
\]

Hence, for two outputs with positive quantities and prices,

\[
\frac{y_i}{y_j}
=
\frac{\beta_i}{\beta_j}
\left(
\frac{p_i}{p_j}
\right)^\eta
\]

so \texttt{et} directly measures the elasticity of the output ratio with
respect to the relative output price. A larger \texttt{et} therefore
makes the output composition more responsive to relative prices. When
\texttt{et=0}, the output proportions are fixed:

\[
y_i
=
\frac{\beta_i}{\alpha}z
\]

\texttt{CETSpec} is an output-side specification and is supplied through
the \texttt{output\_spec} argument of an activity-demand producer. For
example, a producer may use a CES specification for its conditional
input demands and a CET specification for its conditional output
supplies:

\begin{Shaded}
\begin{Highlighting}[]
\NormalTok{GEMB.}\FunctionTok{add\_agent!}\NormalTok{(}
\NormalTok{    model,}
\NormalTok{    GEMB.}\FunctionTok{CESSpec}\NormalTok{(}
\NormalTok{        [}\FloatTok{0.6}\NormalTok{, }\FloatTok{0.4}\NormalTok{];}
\NormalTok{        es}\OperatorTok{=}\FloatTok{1.0}\NormalTok{,}
\NormalTok{        alpha}\OperatorTok{=}\FloatTok{1.0}\NormalTok{,}
\NormalTok{    );}
\NormalTok{    outputs}\OperatorTok{=}\NormalTok{[}\OperatorTok{:}\NormalTok{product1, }\OperatorTok{:}\NormalTok{product2],}
\NormalTok{    output\_spec}\OperatorTok{=}\NormalTok{GEMB.}\FunctionTok{CETSpec}\NormalTok{(}
\NormalTok{        [}\FloatTok{0.5}\NormalTok{, }\FloatTok{0.5}\NormalTok{];}
\NormalTok{        et}\OperatorTok{=}\FloatTok{2.0}\NormalTok{,}
\NormalTok{        alpha}\OperatorTok{=}\FloatTok{1.0}\NormalTok{,}
\NormalTok{    ),}
\NormalTok{    demands}\OperatorTok{=}\NormalTok{[}\OperatorTok{:}\NormalTok{labor, }\OperatorTok{:}\NormalTok{land],}
\NormalTok{    activity\_start}\OperatorTok{=}\FloatTok{100.0}\NormalTok{,}
\NormalTok{    name}\OperatorTok{=:}\NormalTok{firm,}
\NormalTok{)}
\end{Highlighting}
\end{Shaded}

In this construction, the CES specification determines the inputs
required for a given activity level, while the CET specification
determines how that activity is transformed into the two outputs. If no
\texttt{output\_spec} is supplied, output quantities are instead
proportional to the activity level through fixed output coefficients.
Consequently, \texttt{output\_spec} and fixed
\texttt{output\_coefficients} represent alternative ways of describing
the output side and are not used simultaneously.

This separation between input and output specifications is particularly
useful for joint production. The input side determines the resources
required to support an activity level, while the output side determines
the composition of the commodities supplied by that activity. The same
activity variable therefore links production and transformation within a
single economic agent.

\section{Equivalent Representations Using Auxiliary Commodities and
Virtual
Producers}\label{equivalent-representations-using-auxiliary-commodities-and-virtual-producers}

In addition to representing an economic structure directly through the
specifications described above, GEMB models can sometimes be constructed
more conveniently by introducing auxiliary commodities, additional
factors, or virtual producers. Such transformations enlarge the formal
representation of the model while preserving the underlying equilibrium
relationships for the original commodities and economic agents.

\begin{enumerate}
\def\labelenumi{\arabic{enumi}.}
\item
  A production function with nonconstant returns to scale can be
  represented as a constant-returns-to-scale technology by introducing
  an additional commodity with fixed aggregate supply. We refer to this
  commodity as a scale claim. This well-known \textbf{degree-one
  homogenization} method applies to decreasing, increasing, or variable
  returns to scale (see Section~\ref{sec-Homogenization} for details).
  When the original production function exhibits decreasing returns to
  scale, the scale claim may also be interpreted as an entrepreneurial
  factor (McKenzie, 2002, p.~216).
\item
  Nested production or utility functions can be represented using
  virtual producers. Each nest is treated as a separate production
  activity whose output is an intermediate virtual commodity used by the
  next nest. In the case of consumer preferences, a consumer can
  equivalently be decomposed into a virtual utility producer and a
  utility consumer. The utility producer uses the consumer's ordinary
  commodities as inputs and produces a virtual utility commodity
  according to the original utility function, while the utility consumer
  consumes only this utility commodity.
\item
  Virtual producers can also be used to represent taxes imposed on only
  part of a producer's inputs. Taxed inputs can be distinguished from
  their corresponding untaxed commodities by introducing virtual
  taxed-input commodities. Virtual producers convert the original
  commodities into the corresponding taxed-input commodities while
  incorporating the relevant tax payments. The actual producer then
  demands these taxed-input commodities instead of the corresponding
  original commodities, whereas untaxed inputs continue to enter the
  production activity directly.
\end{enumerate}

These constructions share a common modeling principle: auxiliary
commodities, additional factors, or virtual producers are introduced to
transform a nonstandard structure into an equivalent system of standard
activities and market relationships.

\section{\texorpdfstring{CES-CET Production: A \(4\times2\)
Model}{CES-CET Production: A 4\textbackslash times2 Model}}\label{ces-cet-production-a-4times2-model}

Assume the following:

\begin{enumerate}
\def\labelenumi{\arabic{enumi}.}
\item
  The economy contains four commodities and two economic agents, a firm
  and a consumer. Commodities 1 and 2 are products, while commodities 3
  and 4 are labor and land, respectively. Commodity 1 is taken as the
  numeraire commodity.
\item
  The firm uses labor and land to generate the activity level \(z\)
  according to the production function \(z=x_3^{2/3}x_4^{1/3}\), where
  \(x_3\) and \(x_4\) denote labor and land inputs, respectively. The
  activity is transformed into the two products according to the CET
  transformation function \(y_1^2+2y_2^2=3z^2\).
\item
  The consumer is endowed with \(\omega>0\) units of labor and
  \(\omega'>0\) units of land and consumes all four commodities. Its
  utility function is \(u=c_1^{1/6}c_2^{1/3}c_3^{1/3}c_4^{1/6}\), where
  \(c_i\) denotes consumption of commodity \(i\).
\end{enumerate}

Let the price vector be \(\mathbf p=(1,p_2,p_3,p_4)^\top\). The firm's
conditional factor demand and product supply vectors are

\[
\mathbf x(\mathbf p,z)
=
z
\left(
\left(\frac{2p_4}{p_3}\right)^{1/3},
\left(\frac{p_3}{2p_4}\right)^{2/3}
\right)^\top
\]

and

\[
\mathbf y(\mathbf p,z)
=
\frac{3z}{\left(3+\frac32p_2^2\right)^{1/2}}
\left(
1,\frac12p_2
\right)^\top
\]

The consumer's income is \(w=\omega p_3+\omega'p_4\), and its demand
vector is

\[
\mathbf c(\mathbf p)
=
w
\left(
\frac16,
\frac{1}{3p_2},
\frac{1}{3p_3},
\frac{1}{6p_4}
\right)^\top
\]

Hence the demand and supply matrices are

\[
\mathbf D(\mathbf p,z)
=
\begin{bmatrix}
0&c_1\\
0&c_2\\
x_3&c_3\\
x_4&c_4
\end{bmatrix}
\]

and

\[
\mathbf S(\mathbf p,z)
=
\begin{bmatrix}
y_1&0\\
y_2&0\\
0&\omega\\
0&\omega'
\end{bmatrix}
\]

Consider the following equality-form wide-sense structural equilibrium
model:

\[
\begin{aligned}
\mathbf p^\top\mathbf D(\mathbf p,z)
&=\mathbf p^\top\mathbf S(\mathbf p,z)\\
\mathbf D(\mathbf p,z)\mathbf1
&=\mathbf S(\mathbf p,z)\mathbf1
\end{aligned}
\]

The equilibrium price vector and firm activity level are, respectively,

\[
\mathbf p^*
=
\left(
1,\,
2,\,
2\left(\frac{\omega'}{\omega}\right)^{1/3},\,
\left(\frac{\omega}{\omega'}\right)^{2/3}
\right)^\top
\]

and

\[
z^*
=
\frac12\omega^{2/3}{\omega'}^{1/3}
\]

The equilibrium product supply vector, factor demand vector, and
consumer demand vector are

\[
\mathbf y^*
=
z^*(1,1)^\top
\]

\[
\mathbf x^*
=
\left(
\frac{\omega}{2},
\frac{\omega'}{2}
\right)^\top
\]

and

\[
\mathbf c^*
=
\left(
z^*,
z^*,
\frac{\omega}{2},
\frac{\omega'}{2}
\right)^\top
\]

Hence the equilibrium demand and supply matrices are

\[
\mathbf D^*
=
\begin{bmatrix}
0&z^*\\
0&z^*\\
\frac{\omega}{2}&\frac{\omega}{2}\\
\frac{\omega'}{2}&\frac{\omega'}{2}
\end{bmatrix},
\qquad
\mathbf S^*
=
\begin{bmatrix}
z^*&0\\
z^*&0\\
0&\omega\\
0&\omega'
\end{bmatrix}
\]

The consumer's equilibrium utility level is

\[
u^*=z^*
=
\frac12\omega^{2/3}{\omega'}^{1/3}
\]

In particular, when \(\omega=\omega'=100\),

\[
\mathbf p^*=(1,2,2,1)^\top,\qquad
z^*=50,\qquad
u^*=50
\]

and

\[
\mathbf D^*
=
\begin{bmatrix}
0&50\\
0&50\\
50&50\\
50&50
\end{bmatrix},
\qquad
\mathbf S^*
=
\begin{bmatrix}
50&0\\
50&0\\
0&100\\
0&100
\end{bmatrix}
\]

\begin{codelisting}

\caption{\label{lst-gemb_CES_CET_jointProduction_nc4_na2}A \(4\times2\)
Equilibrium Model with CES Inputs and CET Outputs}

\centering{

\begin{Shaded}
\begin{Highlighting}[]
\CommentTok{\# ================================================================}
\CommentTok{\# gemb\_CES\_CET\_jointProduction\_nc4\_na2\_v1.jl}
\CommentTok{\#}
\CommentTok{\# CES{-}CET joint{-}production equilibrium:}
\CommentTok{\# one producer and one consumer who consumes all four commodities.}
\CommentTok{\#}
\CommentTok{\# Commodities:}
\CommentTok{\#   1. product1}
\CommentTok{\#   2. product2}
\CommentTok{\#   3. labor}
\CommentTok{\#   4. land}
\CommentTok{\# ================================================================}

\ImportTok{using} \BuiltInTok{GeneralEquilibriumModeling}


\CommentTok{\# {-}{-}{-}{-}{-}{-}{-}{-}{-}{-}{-}{-}{-}{-}{-}{-}{-}{-}{-}{-}{-}{-}{-}{-}{-}{-}{-}{-}{-}{-}{-}{-}{-}{-}{-}{-}{-}{-}{-}{-}{-}{-}{-}{-}{-}{-}{-}{-}{-}{-}{-}{-}{-}{-}{-}{-}{-}{-}{-}{-}{-}{-}{-}{-}}
\CommentTok{\# 1. Parameters}
\CommentTok{\# {-}{-}{-}{-}{-}{-}{-}{-}{-}{-}{-}{-}{-}{-}{-}{-}{-}{-}{-}{-}{-}{-}{-}{-}{-}{-}{-}{-}{-}{-}{-}{-}{-}{-}{-}{-}{-}{-}{-}{-}{-}{-}{-}{-}{-}{-}{-}{-}{-}{-}{-}{-}{-}{-}{-}{-}{-}{-}{-}{-}{-}{-}{-}{-}}

\NormalTok{omega\_labor }\OperatorTok{=} \FloatTok{100.0}
\NormalTok{omega\_land }\OperatorTok{=} \FloatTok{100.0}

\CommentTok{\# Production:}
\CommentTok{\#}
\CommentTok{\#   z = x\_labor\^{}(2/3) * x\_land\^{}(1/3)}
\CommentTok{\#}
\NormalTok{input\_spec }\OperatorTok{=}\NormalTok{ GEMB.}\FunctionTok{CESSpec}\NormalTok{(}
\NormalTok{    [}\FloatTok{2} \OperatorTok{/} \FloatTok{3}\NormalTok{, }\FloatTok{1} \OperatorTok{/} \FloatTok{3}\NormalTok{];}
\NormalTok{    es}\OperatorTok{=}\FloatTok{1.0}\NormalTok{,}
\NormalTok{    alpha}\OperatorTok{=}\FloatTok{1.0}\NormalTok{,}
\NormalTok{)}

\CommentTok{\# CET transformation:}
\CommentTok{\#}
\CommentTok{\#   y1\^{}2 + 2*y2\^{}2 = 3*z\^{}2}
\CommentTok{\#}
\NormalTok{output\_spec }\OperatorTok{=}\NormalTok{ GEMB.}\FunctionTok{CETSpec}\NormalTok{(}
\NormalTok{    [}\FloatTok{2} \OperatorTok{/} \FloatTok{3}\NormalTok{, }\FloatTok{1} \OperatorTok{/} \FloatTok{3}\NormalTok{];}
\NormalTok{    et}\OperatorTok{=}\FloatTok{1.0}\NormalTok{,}
\NormalTok{    alpha}\OperatorTok{=}\FunctionTok{sqrt}\NormalTok{(}\FloatTok{2}\NormalTok{) }\OperatorTok{/} \FloatTok{3}\NormalTok{,}
\NormalTok{)}

\CommentTok{\# Consumer utility:}
\CommentTok{\#}
\CommentTok{\#   u = c1\^{}(1/6) * c2\^{}(1/3) * c3\^{}(1/3) * c4\^{}(1/6)}
\CommentTok{\#}
\NormalTok{utility\_spec }\OperatorTok{=}\NormalTok{ GEMB.}\FunctionTok{CESSpec}\NormalTok{(}
\NormalTok{    [}\FloatTok{1} \OperatorTok{/} \FloatTok{6}\NormalTok{, }\FloatTok{1} \OperatorTok{/} \FloatTok{3}\NormalTok{, }\FloatTok{1} \OperatorTok{/} \FloatTok{3}\NormalTok{, }\FloatTok{1} \OperatorTok{/} \FloatTok{6}\NormalTok{];}
\NormalTok{    es}\OperatorTok{=}\FloatTok{1.0}\NormalTok{,}
\NormalTok{    alpha}\OperatorTok{=}\FloatTok{1.0}\NormalTok{,}
\NormalTok{)}


\CommentTok{\# {-}{-}{-}{-}{-}{-}{-}{-}{-}{-}{-}{-}{-}{-}{-}{-}{-}{-}{-}{-}{-}{-}{-}{-}{-}{-}{-}{-}{-}{-}{-}{-}{-}{-}{-}{-}{-}{-}{-}{-}{-}{-}{-}{-}{-}{-}{-}{-}{-}{-}{-}{-}{-}{-}{-}{-}{-}{-}{-}{-}{-}{-}{-}{-}}
\CommentTok{\# 2. High{-}level GEMB model}
\CommentTok{\# {-}{-}{-}{-}{-}{-}{-}{-}{-}{-}{-}{-}{-}{-}{-}{-}{-}{-}{-}{-}{-}{-}{-}{-}{-}{-}{-}{-}{-}{-}{-}{-}{-}{-}{-}{-}{-}{-}{-}{-}{-}{-}{-}{-}{-}{-}{-}{-}{-}{-}{-}{-}{-}{-}{-}{-}{-}{-}{-}{-}{-}{-}{-}{-}}

\NormalTok{model }\OperatorTok{=}\NormalTok{ GEMB.}\FunctionTok{GEMBModel}\NormalTok{(}
\NormalTok{    [}\OperatorTok{:}\NormalTok{product1, }\OperatorTok{:}\NormalTok{product2, }\OperatorTok{:}\NormalTok{labor, }\OperatorTok{:}\NormalTok{land];}
\NormalTok{    numeraire}\OperatorTok{=:}\NormalTok{product1,}
\NormalTok{    numeraire\_value}\OperatorTok{=}\FloatTok{1.0}\NormalTok{,}
\NormalTok{)}


\CommentTok{\# {-}{-}{-}{-}{-}{-}{-}{-}{-}{-}{-}{-}{-}{-}{-}{-}{-}{-}{-}{-}{-}{-}{-}{-}{-}{-}{-}{-}{-}{-}{-}{-}{-}{-}{-}{-}{-}{-}{-}{-}{-}{-}{-}{-}{-}{-}{-}{-}{-}{-}{-}{-}{-}{-}{-}{-}{-}{-}{-}{-}{-}{-}{-}{-}}
\CommentTok{\# 3. Producer}
\CommentTok{\# {-}{-}{-}{-}{-}{-}{-}{-}{-}{-}{-}{-}{-}{-}{-}{-}{-}{-}{-}{-}{-}{-}{-}{-}{-}{-}{-}{-}{-}{-}{-}{-}{-}{-}{-}{-}{-}{-}{-}{-}{-}{-}{-}{-}{-}{-}{-}{-}{-}{-}{-}{-}{-}{-}{-}{-}{-}{-}{-}{-}{-}{-}{-}{-}}

\NormalTok{GEMB.}\FunctionTok{add\_agent!}\NormalTok{(}
\NormalTok{    model,}
\NormalTok{    input\_spec;}
\NormalTok{    outputs}\OperatorTok{=}\NormalTok{[}\OperatorTok{:}\NormalTok{product1, }\OperatorTok{:}\NormalTok{product2],}
\NormalTok{    output\_spec}\OperatorTok{=}\NormalTok{output\_spec,}
\NormalTok{    demands}\OperatorTok{=}\NormalTok{[}\OperatorTok{:}\NormalTok{labor, }\OperatorTok{:}\NormalTok{land],}
\NormalTok{    activity\_start}\OperatorTok{=}\FloatTok{100.0}\NormalTok{,}
\NormalTok{    name}\OperatorTok{=:}\NormalTok{firm,}
\NormalTok{)}


\CommentTok{\# {-}{-}{-}{-}{-}{-}{-}{-}{-}{-}{-}{-}{-}{-}{-}{-}{-}{-}{-}{-}{-}{-}{-}{-}{-}{-}{-}{-}{-}{-}{-}{-}{-}{-}{-}{-}{-}{-}{-}{-}{-}{-}{-}{-}{-}{-}{-}{-}{-}{-}{-}{-}{-}{-}{-}{-}{-}{-}{-}{-}{-}{-}{-}{-}}
\CommentTok{\# 4. Consumer}
\CommentTok{\# {-}{-}{-}{-}{-}{-}{-}{-}{-}{-}{-}{-}{-}{-}{-}{-}{-}{-}{-}{-}{-}{-}{-}{-}{-}{-}{-}{-}{-}{-}{-}{-}{-}{-}{-}{-}{-}{-}{-}{-}{-}{-}{-}{-}{-}{-}{-}{-}{-}{-}{-}{-}{-}{-}{-}{-}{-}{-}{-}{-}{-}{-}{-}{-}}

\NormalTok{GEMB.}\FunctionTok{add\_agent!}\NormalTok{(}
\NormalTok{    model,}
\NormalTok{    utility\_spec;}
\NormalTok{    demands}\OperatorTok{=}\NormalTok{[}\OperatorTok{:}\NormalTok{product1, }\OperatorTok{:}\NormalTok{product2, }\OperatorTok{:}\NormalTok{labor, }\OperatorTok{:}\NormalTok{land],}
\NormalTok{    endowments}\OperatorTok{=}\NormalTok{[}\OperatorTok{:}\NormalTok{labor, }\OperatorTok{:}\NormalTok{land],}
\NormalTok{    endowment\_quantities}\OperatorTok{=}\NormalTok{[omega\_labor, omega\_land],}
\NormalTok{    activity\_start}\OperatorTok{=}\FloatTok{100.0}\NormalTok{,}
\NormalTok{    name}\OperatorTok{=:}\NormalTok{consumer,}
\NormalTok{)}


\CommentTok{\# {-}{-}{-}{-}{-}{-}{-}{-}{-}{-}{-}{-}{-}{-}{-}{-}{-}{-}{-}{-}{-}{-}{-}{-}{-}{-}{-}{-}{-}{-}{-}{-}{-}{-}{-}{-}{-}{-}{-}{-}{-}{-}{-}{-}{-}{-}{-}{-}{-}{-}{-}{-}{-}{-}{-}{-}{-}{-}{-}{-}{-}{-}{-}{-}}
\CommentTok{\# 5. Solve}
\CommentTok{\# {-}{-}{-}{-}{-}{-}{-}{-}{-}{-}{-}{-}{-}{-}{-}{-}{-}{-}{-}{-}{-}{-}{-}{-}{-}{-}{-}{-}{-}{-}{-}{-}{-}{-}{-}{-}{-}{-}{-}{-}{-}{-}{-}{-}{-}{-}{-}{-}{-}{-}{-}{-}{-}{-}{-}{-}{-}{-}{-}{-}{-}{-}{-}{-}}

\NormalTok{result }\OperatorTok{=}\NormalTok{ GEMB.}\FunctionTok{solve}\NormalTok{(}
\NormalTok{    model;}
\NormalTok{    p0}\OperatorTok{=}\FunctionTok{ones}\NormalTok{(}\FloatTok{4}\NormalTok{),}
\NormalTok{    residual\_tol}\OperatorTok{=}\FloatTok{1.0e{-}10}\NormalTok{,}
\NormalTok{    silent}\OperatorTok{=}\ConstantTok{true}\NormalTok{,}
\NormalTok{)}


\CommentTok{\# {-}{-}{-}{-}{-}{-}{-}{-}{-}{-}{-}{-}{-}{-}{-}{-}{-}{-}{-}{-}{-}{-}{-}{-}{-}{-}{-}{-}{-}{-}{-}{-}{-}{-}{-}{-}{-}{-}{-}{-}{-}{-}{-}{-}{-}{-}{-}{-}{-}{-}{-}{-}{-}{-}{-}{-}{-}{-}{-}{-}{-}{-}{-}{-}}
\CommentTok{\# 6. Results}
\CommentTok{\# {-}{-}{-}{-}{-}{-}{-}{-}{-}{-}{-}{-}{-}{-}{-}{-}{-}{-}{-}{-}{-}{-}{-}{-}{-}{-}{-}{-}{-}{-}{-}{-}{-}{-}{-}{-}{-}{-}{-}{-}{-}{-}{-}{-}{-}{-}{-}{-}{-}{-}{-}{-}{-}{-}{-}{-}{-}{-}{-}{-}{-}{-}{-}{-}}

\NormalTok{prices }\OperatorTok{=}\NormalTok{ result.prices}
\NormalTok{firm\_activity }\OperatorTok{=}\NormalTok{ result.agent\_variable\_values[}\FloatTok{1}\NormalTok{][}\FloatTok{1}\NormalTok{]}
\NormalTok{consumer\_utility }\OperatorTok{=}\NormalTok{ result.agent\_variable\_values[}\FloatTok{2}\NormalTok{][}\FloatTok{1}\NormalTok{]}

\FunctionTok{println}\NormalTok{(}\StringTok{"Solved:            "}\NormalTok{, result.solved)}
\FunctionTok{println}\NormalTok{(}\StringTok{"Prices:            "}\NormalTok{, prices)}
\FunctionTok{println}\NormalTok{(}\StringTok{"Firm activity:     "}\NormalTok{, firm\_activity)}
\FunctionTok{println}\NormalTok{(}\StringTok{"Consumer utility:  "}\NormalTok{, consumer\_utility)}
\FunctionTok{println}\NormalTok{(}\StringTok{"Firm net supply:   "}\NormalTok{, result.agent\_net\_supplies[}\FloatTok{1}\NormalTok{])}
\FunctionTok{println}\NormalTok{(}\StringTok{"Consumer net supply:"}\NormalTok{, result.agent\_net\_supplies[}\FloatTok{2}\NormalTok{])}
\FunctionTok{println}\NormalTok{(}\StringTok{"Total net supply:  "}\NormalTok{, result.total\_net\_supply)}

\FunctionTok{println}\NormalTok{()}
\NormalTok{GEMB.}\FunctionTok{print\_equilibrium\_statistics}\NormalTok{(}
\NormalTok{    model,}
\NormalTok{    result;}
\NormalTok{    display\_tol}\OperatorTok{=}\FloatTok{1.0e{-}10}\NormalTok{,}
\NormalTok{    sigdigits}\OperatorTok{=}\FloatTok{8}\NormalTok{,}
\NormalTok{)}
\end{Highlighting}
\end{Shaded}

}

\end{codelisting}%

\section{\texorpdfstring{Ad Valorem Tax and Ad Valorem Claim: A
\(3\times2\)
Model}{Ad Valorem Tax and Ad Valorem Claim: A 3\textbackslash times2 Model}}\label{ad-valorem-tax-and-ad-valorem-claim-a-3times2-model}

\subsection{Ad Valorem Claim}\label{ad-valorem-claim}

An ad valorem tax is levied as a fixed proportion of the value of a
transaction, output, input, or other economic base. More generally, this
concept can be extended to an ad valorem claim, which represents a right
to claim a fixed proportion of a specified economic value. This concept
provides a unified way to represent a variety of value-based claims in
equilibrium models.

\begin{enumerate}
\def\labelenumi{\arabic{enumi}.}
\item
  Let \(V\) denote the value base and \(\tau\) the ad valorem claim
  rate. The value of the claim is \(\tau V\). If the claim commodity has
  price \(p_c\) and quantity \(c\), then \(p_c c=\tau V\).
\item
  A positive claim rate, \(\tau>0\), can represent an ad valorem tax,
  fee, or other positive value claim, while a negative claim rate,
  \(\tau<0\), can represent an ad valorem subsidy or other reverse
  claim. Money and equity in equilibrium models can also be interpreted
  as forms of ad valorem claims.
\item
  By treating an ad valorem claim as a commodity and incorporating it
  into demand, supply, and revenue-expenditure balance, ad valorem
  taxes, subsidies, value-sharing arrangements, rent claims, money,
  equity, and related institutions can be represented within a unified
  wide-sense structural equilibrium framework.
\item
  Under an appropriate claim-based representation, profit rates, markup
  rates, saving rates, interest rates, rates of return, and ad valorem
  tax or subsidy rates can all be interpreted as ad valorem claim rates.
  In particular, money may serve as a claim-bearing commodity, while
  saving may be represented as the purchase of claims proportional to
  consumption expenditure.
\end{enumerate}

\subsection{\texorpdfstring{A \(3\times2\) Model with Ad Valorem
Tax}{A 3\textbackslash times2 Model with Ad Valorem Tax}}\label{a-3times2-model-with-ad-valorem-tax}

Assume the following:

\begin{enumerate}
\def\labelenumi{\arabic{enumi}.}
\item
  The economy contains three commodities: a product, labor, and a tax
  certificate. There are two economic agents: a firm and a consumer.
  Labor is taken as the numeraire commodity.
\item
  The firm produces and supplies the product according to the production
  function \(\alpha x_2\), where \(x_2\) denotes labor input and
  \(\alpha>0\).
\item
  The consumer has a labor endowment of \(\omega\) and the utility
  function \(x_1^\beta x_2^{1-\beta}\), where \(x_1\) denotes the
  quantity of the product consumed and \(x_2\) denotes the quantity of
  labor enjoyed by the consumer, or equivalently leisure, with
  \(0<\beta<1\).
\item
  An ad valorem tax is imposed on the product, and the tax revenue is
  transferred to the consumer. The tax rate satisfies \(0<\tau<1\), so
  the firm pays a tax equal to the fraction \(\tau\) of its product
  sales revenue \(p_1y\). The consumer is endowed with one unit of the
  tax certificate.
\end{enumerate}

Let \(\mathbf p:=(p_1,1,p_3)^\top\), let \(y\) denote output, and let
\(w\) denote consumer income. The demand and supply matrices are

\[
\mathbf D(\mathbf p,w,y):=
\begin{bmatrix}
0 & \frac{\beta w}{p_1}\\
\frac{y}{\alpha} & (1-\beta)w\\
\frac{\tau p_1y}{p_3} & 0
\end{bmatrix}
\]

and

\[
\mathbf S(y):=
\begin{bmatrix}
y & 0\\
0 & \omega\\
0 & 1
\end{bmatrix}
\]

Solving the equality-form wide-sense structural equilibrium model gives

\[
\mathbf p^*
=
\left(
\frac{1}{\alpha(1-\tau)},
\,1,\,
\frac{\beta\omega\tau}{1-\beta\tau}
\right)^\top
\]

\[
w^*
=
\frac{\omega}{1-\beta\tau}
\]

and

\[
y^*
=
\frac{\alpha\beta(1-\tau)\omega}{1-\beta\tau}
\]

The consumer's equilibrium utility level is

\[
u^*
=
\frac{
\omega(\alpha\beta)^\beta(1-\beta)^{1-\beta}(1-\tau)^\beta
}{
1-\beta\tau
}
\]

For \(0<\tau<1\),

\[
\frac{\partial\ln u^*}{\partial\tau}
=
-\frac{\beta(1-\beta)\tau}
{(1-\tau)(1-\beta\tau)}
<0
\]

Hence, the consumer's equilibrium utility decreases as the ad valorem
tax rate increases. The tax therefore generates a deadweight loss, or
excess burden, in this example.

For example, let \(\alpha=2\), \(\beta=1/2\), \(\omega=90\), and
\(\tau=1/2\). The equilibrium is then particularly simple:

\[
\mathbf p^*=(1,1,30)^\top,\qquad
w^*=120,\qquad
y^*=60
\]

The consumer demands 60 units of the product and enjoys 60 units of
labor, so the equilibrium utility level is \(u^*=60\). The corresponding
demand and supply matrices are

\[
\mathbf D^*
=
\begin{bmatrix}
0&60\\
30&60\\
1&0
\end{bmatrix},
\qquad
\mathbf S^*
=
\begin{bmatrix}
60&0\\
0&90\\
0&1
\end{bmatrix}
\]

Thus, all commodity markets clear exactly, and the revenue-expenditure
balance condition is satisfied for both the firm and the consumer.

\begin{codelisting}

\caption{\label{lst-gemb_CD_adValoremTax_nc3_na2}A \(3\times2\)
Equilibrium Model with an Ad Valorem Product Tax}

\centering{

\begin{Shaded}
\begin{Highlighting}[]
\ImportTok{using} \BuiltInTok{GeneralEquilibriumModeling}

\NormalTok{alpha }\OperatorTok{=} \FloatTok{2.0}
\NormalTok{beta }\OperatorTok{=} \FloatTok{0.5}
\NormalTok{omega }\OperatorTok{=} \FloatTok{90.0}
\NormalTok{tau }\OperatorTok{=} \FloatTok{0.5}

\NormalTok{claim\_rate }\OperatorTok{=}\NormalTok{ tau }\OperatorTok{/}\NormalTok{ (}\FloatTok{1.0} \OperatorTok{{-}}\NormalTok{ tau)}

\NormalTok{model }\OperatorTok{=}\NormalTok{ GEMB.}\FunctionTok{GEMBModel}\NormalTok{(}
\NormalTok{    [}\OperatorTok{:}\NormalTok{product, }\OperatorTok{:}\NormalTok{labor, }\OperatorTok{:}\NormalTok{tax\_certificate];}
\NormalTok{    numeraire}\OperatorTok{=:}\NormalTok{labor,}
\NormalTok{    numeraire\_value}\OperatorTok{=}\FloatTok{1.0}\NormalTok{,}
\NormalTok{)}

\NormalTok{GEMB.}\FunctionTok{add\_agent!}\NormalTok{(}
\NormalTok{    model,}
\NormalTok{    GEMB.}\FunctionTok{CESSpec}\NormalTok{(}
\NormalTok{        [}\FloatTok{1.0}\NormalTok{];}
\NormalTok{        alpha}\OperatorTok{=}\NormalTok{alpha,}
\NormalTok{    );}
\NormalTok{    outputs}\OperatorTok{=:}\NormalTok{product,}
\NormalTok{    demands}\OperatorTok{=:}\NormalTok{labor,}
\NormalTok{    claim\_rate}\OperatorTok{=}\NormalTok{claim\_rate,}
\NormalTok{    claim}\OperatorTok{=}\NormalTok{GEMB.}\FunctionTok{CommodityRef}\NormalTok{(}\OperatorTok{:}\NormalTok{tax\_certificate),}
\NormalTok{    activity\_start}\OperatorTok{=}\FloatTok{100.0}\NormalTok{,}
\NormalTok{    name}\OperatorTok{=:}\NormalTok{firm,}
\NormalTok{)}

\NormalTok{GEMB.}\FunctionTok{add\_agent!}\NormalTok{(}
\NormalTok{    model,}
\NormalTok{    GEMB.}\FunctionTok{CESSpec}\NormalTok{(}
\NormalTok{        [beta, }\FloatTok{1.0} \OperatorTok{{-}}\NormalTok{ beta];}
\NormalTok{        es}\OperatorTok{=}\FloatTok{1.0}\NormalTok{,}
\NormalTok{        alpha}\OperatorTok{=}\FloatTok{1.0}\NormalTok{,}
\NormalTok{    );}
\NormalTok{    demands}\OperatorTok{=}\NormalTok{[}\OperatorTok{:}\NormalTok{product, }\OperatorTok{:}\NormalTok{labor],}
\NormalTok{    endowments}\OperatorTok{=}\NormalTok{[}\OperatorTok{:}\NormalTok{labor, }\OperatorTok{:}\NormalTok{tax\_certificate],}
\NormalTok{    endowment\_quantities}\OperatorTok{=}\NormalTok{[omega, }\FloatTok{1.0}\NormalTok{],}
\NormalTok{    activity\_start}\OperatorTok{=}\FloatTok{100.0}\NormalTok{,}
\NormalTok{    name}\OperatorTok{=:}\NormalTok{consumer,}
\NormalTok{)}

\NormalTok{result }\OperatorTok{=}\NormalTok{ GEMB.}\FunctionTok{solve}\NormalTok{(}
\NormalTok{    model;}
\NormalTok{    p0}\OperatorTok{=}\NormalTok{[}\FloatTok{1.0}\NormalTok{, }\FloatTok{1.0}\NormalTok{, }\FloatTok{30.0}\NormalTok{],}
\NormalTok{    residual\_tol}\OperatorTok{=}\FloatTok{1.0e{-}8}\NormalTok{,}
\NormalTok{    silent}\OperatorTok{=}\ConstantTok{true}\NormalTok{,}
\NormalTok{)}

\NormalTok{stats }\OperatorTok{=}\NormalTok{ GEMB.}\FunctionTok{equilibrium\_statistics}\NormalTok{(model, result)}

\NormalTok{output }\OperatorTok{=}\NormalTok{ stats.activity\_levels[}\FloatTok{1}\NormalTok{]}
\NormalTok{utility }\OperatorTok{=}\NormalTok{ stats.activity\_levels[}\FloatTok{2}\NormalTok{]}

\NormalTok{GEMB.}\FunctionTok{print\_equilibrium\_statistics}\NormalTok{(model, result)}
\end{Highlighting}
\end{Shaded}

}

\end{codelisting}%

\section{\texorpdfstring{Specific Tax: A \(3\times2\)
Model}{Specific Tax: A 3\textbackslash times2 Model}}\label{specific-tax-a-3times2-model}

Consider the same economy as in the previous section, but replace the ad
valorem tax with a specific tax on the product. Let \(\tau>0\) denote
the tax per unit of output. The firm therefore pays \(\tau y\) in taxes
when producing \(y\) units of the product, and the tax revenue is
transferred to the consumer through the tax certificate.

Let \(\mathbf p:=(p_1,1,p_3)^\top\), where labor is the numeraire
commodity. The demand matrix becomes

\[
\mathbf D(\mathbf p,w,y)
=
\begin{bmatrix}
0 & \frac{\beta w}{p_1}\\
\frac{y}{\alpha} & (1-\beta)w\\
\frac{\tau y}{p_3} & 0
\end{bmatrix}
\]

while the supply matrix remains

\[
\mathbf S(y)
=
\begin{bmatrix}
y & 0\\
0 & \omega\\
0 & 1
\end{bmatrix}
\]

Solving the equality-form wide-sense structural equilibrium model gives

\[
\mathbf p^*
=
\left(
\frac{1}{\alpha}+\tau,
\,1,\,
\frac{\alpha\beta\omega\tau}
{1+\alpha\tau(1-\beta)}
\right)^\top
\]

\[
w^*
=
\frac{(1+\alpha\tau)\omega}
{1+\alpha\tau(1-\beta)}
\]

and

\[
y^*
=
\frac{\alpha\beta\omega}
{1+\alpha\tau(1-\beta)}
\]

The consumer's equilibrium utility level is

\[
u^*
=
\frac{
\omega(\alpha\beta)^\beta
\big((1-\beta)(1+\alpha\tau)\big)^{1-\beta}
}{
1+\alpha(1-\beta)\tau
}
\]

For example, let \(\alpha=2\), \(\beta=1/2\), \(\omega=90\), and
\(\tau=1/2\). The equilibrium is

\[
\mathbf p^*=(1,1,30)^\top,\qquad
w^*=120,\qquad
y^*=60,\qquad
u^*=60
\]

For these parameter values, the numerical equilibrium happens to
coincide with that of the ad valorem tax example in the previous
section, although the definitions of the two tax rates are different.

\begin{codelisting}

\caption{\label{lst-gemb_CD_specificTax_nc3_na2}A \(3\times2\)
Equilibrium Model with a Specific Tax}

\centering{

\begin{Shaded}
\begin{Highlighting}[]
\ImportTok{using} \BuiltInTok{GeneralEquilibriumModeling}

\NormalTok{alpha }\OperatorTok{=} \FloatTok{2.0}
\NormalTok{beta }\OperatorTok{=} \FloatTok{0.5}
\NormalTok{omega }\OperatorTok{=} \FloatTok{90.0}
\NormalTok{tau }\OperatorTok{=} \FloatTok{0.5}

\NormalTok{model }\OperatorTok{=}\NormalTok{ GEMB.}\FunctionTok{GEMBModel}\NormalTok{(}
\NormalTok{    [}\OperatorTok{:}\NormalTok{product, }\OperatorTok{:}\NormalTok{labor, }\OperatorTok{:}\NormalTok{tax\_certificate];}
\NormalTok{    numeraire}\OperatorTok{=:}\NormalTok{labor,}
\NormalTok{    numeraire\_value}\OperatorTok{=}\FloatTok{1.0}\NormalTok{,}
\NormalTok{)}

\NormalTok{firm\_spec }\OperatorTok{=}\NormalTok{ GEMB.}\FunctionTok{ActivityDemandSpec}\NormalTok{(}
    \KeywordTok{function}\NormalTok{ (output, prices)}
\NormalTok{        p\_labor, p\_tax }\OperatorTok{=}\NormalTok{ prices}

        \ControlFlowTok{return}\NormalTok{ [}
\NormalTok{            output }\OperatorTok{/}\NormalTok{ alpha,}
\NormalTok{            tau }\OperatorTok{*}\NormalTok{ output }\OperatorTok{/}\NormalTok{ p\_tax,}
\NormalTok{        ]}
    \KeywordTok{end}\NormalTok{,}
\NormalTok{)}

\NormalTok{GEMB.}\FunctionTok{add\_agent!}\NormalTok{(}
\NormalTok{    model,}
\NormalTok{    firm\_spec;}
\NormalTok{    outputs}\OperatorTok{=:}\NormalTok{product,}
\NormalTok{    demands}\OperatorTok{=}\NormalTok{[}\OperatorTok{:}\NormalTok{labor, }\OperatorTok{:}\NormalTok{tax\_certificate],}
\NormalTok{    activity\_start}\OperatorTok{=}\FloatTok{100.0}\NormalTok{,}
\NormalTok{    name}\OperatorTok{=:}\NormalTok{firm,}
\NormalTok{)}

\NormalTok{GEMB.}\FunctionTok{add\_agent!}\NormalTok{(}
\NormalTok{    model,}
\NormalTok{    GEMB.}\FunctionTok{CESSpec}\NormalTok{(}
\NormalTok{        [beta, }\FloatTok{1.0} \OperatorTok{{-}}\NormalTok{ beta];}
\NormalTok{        es}\OperatorTok{=}\FloatTok{1.0}\NormalTok{,}
\NormalTok{        alpha}\OperatorTok{=}\FloatTok{1.0}\NormalTok{,}
\NormalTok{    );}
\NormalTok{    demands}\OperatorTok{=}\NormalTok{[}\OperatorTok{:}\NormalTok{product, }\OperatorTok{:}\NormalTok{labor],}
\NormalTok{    endowments}\OperatorTok{=}\NormalTok{[}\OperatorTok{:}\NormalTok{labor, }\OperatorTok{:}\NormalTok{tax\_certificate],}
\NormalTok{    endowment\_quantities}\OperatorTok{=}\NormalTok{[omega, }\FloatTok{1.0}\NormalTok{],}
\NormalTok{    activity\_start}\OperatorTok{=}\FloatTok{100.0}\NormalTok{,}
\NormalTok{    name}\OperatorTok{=:}\NormalTok{consumer,}
\NormalTok{)}

\NormalTok{result }\OperatorTok{=}\NormalTok{ GEMB.}\FunctionTok{solve}\NormalTok{(}
\NormalTok{    model;}
\NormalTok{    p0}\OperatorTok{=}\NormalTok{[}\FloatTok{1.0}\NormalTok{, }\FloatTok{1.0}\NormalTok{, }\FloatTok{30.0}\NormalTok{],}
\NormalTok{    residual\_tol}\OperatorTok{=}\FloatTok{1.0e{-}8}\NormalTok{,}
\NormalTok{    silent}\OperatorTok{=}\ConstantTok{true}\NormalTok{,}
\NormalTok{)}

\NormalTok{stats }\OperatorTok{=}\NormalTok{ GEMB.}\FunctionTok{equilibrium\_statistics}\NormalTok{(model, result)}

\NormalTok{output }\OperatorTok{=}\NormalTok{ stats.activity\_levels[}\FloatTok{1}\NormalTok{]}
\NormalTok{utility }\OperatorTok{=}\NormalTok{ stats.activity\_levels[}\FloatTok{2}\NormalTok{]}

\NormalTok{GEMB.}\FunctionTok{print\_equilibrium\_statistics}\NormalTok{(model, result)}
\end{Highlighting}
\end{Shaded}

}

\end{codelisting}%

\section{\texorpdfstring{Non-rival Service and Lindahl Prices: A
\(3\times3\)
Model}{Non-rival Service and Lindahl Prices: A 3\textbackslash times3 Model}}\label{non-rival-service-and-lindahl-prices-a-3times3-model}

Assume the following:

\begin{enumerate}
\def\labelenumi{\arabic{enumi}.}
\item
  The economy contains three commodities: service 1, service 2, and
  labor. There are three economic agents: one firm and two consumers.
  Labor is taken as the numeraire commodity.
\item
  The firm provides a non-rival service, such as residential security
  service. Each unit of labor input provides one unit of personalized
  service to each consumer. These personalized services are referred to
  as service 1 and service 2, respectively.
\item
  Both consumers demand their personalized service and labor and have
  Cobb--Douglas utility functions. For consumer 1, the expenditure
  shares on service 1 and labor are \(\beta_1\) and \(1-\beta_1\),
  respectively. For consumer 2, the corresponding shares on service 2
  and labor are \(\beta_2\) and \(1-\beta_2\).
\item
  Consumers 1 and 2 are endowed with \(\omega_1>0\) and \(\omega_2>0\)
  units of labor, respectively.
\end{enumerate}

Let \(\mathbf p:=(p_1,p_2,1)^\top\), let \(y\) denote the service level,
and let the consumer-income vector be \(\mathbf w=(w_1,w_2)^\top\). The
demand and supply matrices are

\[
\mathbf D(\mathbf p,\mathbf w,y)
:=
\begin{bmatrix}
0 & \frac{\beta_1w_1}{p_1} & 0\\
0 & 0 & \frac{\beta_2w_2}{p_2}\\
y & (1-\beta_1)w_1 & (1-\beta_2)w_2
\end{bmatrix}
\]

and

\[
\mathbf S(y)
:=
\begin{bmatrix}
y & 0 & 0\\
y & 0 & 0\\
0 & \omega_1 & \omega_2
\end{bmatrix}
\]

Solving the equality-form wide-sense structural equilibrium model gives

\[
\mathbf p^*
=
\left(
\frac{\beta_1\omega_1}
{\beta_1\omega_1+\beta_2\omega_2},
\frac{\beta_2\omega_2}
{\beta_1\omega_1+\beta_2\omega_2},
1
\right)^\top
\]

\[
\mathbf w^*
=
(\omega_1,\omega_2)^\top
\]

and

\[
y^*
=
\beta_1\omega_1+\beta_2\omega_2
\]

The personalized equilibrium prices satisfy

\[
p_1^*+p_2^*=1
\]

so the sum of the two consumers' personalized prices equals the marginal
cost of providing the non-rival service. The prices \(p_1^*\) and
\(p_2^*\) can therefore be interpreted as Lindahl prices: each consumer
faces a personalized price for the same underlying non-rival service,
while the sum of these prices covers its unit production cost.

For example, let \(\beta_1=1/4\), \(\beta_2=2/3\), \(\omega_1=120\), and
\(\omega_2=90\). The equilibrium is

\[
\mathbf p^*
=
\left(
\frac13,\frac23,1
\right)^\top,
\qquad
\mathbf w^*
=
(120,90)^\top,
\qquad
y^*=90
\]

Both consumers receive 90 units of the non-rival service. Consumer 1
enjoys 90 units of labor, whereas consumer 2 enjoys 30 units. The
corresponding demand and supply matrices are

\[
\mathbf D^*
=
\begin{bmatrix}
0&90&0\\
0&0&90\\
90&90&30
\end{bmatrix},
\qquad
\mathbf S^*
=
\begin{bmatrix}
90&0&0\\
90&0&0\\
0&120&90
\end{bmatrix}
\]

The two personalized Lindahl prices are therefore \(1/3\) and \(2/3\),
reflecting the consumers' different valuations of the non-rival service,
while their sum remains equal to the unit cost of providing it.

\begin{codelisting}

\caption{\label{lst-gemb_CD_nonRivalService_Lindahl_nc3_na3}A
\(3\times3\) Equilibrium Model with a Non-rival Service and Lindahl
Prices}

\centering{

\begin{Shaded}
\begin{Highlighting}[]
\ImportTok{using} \BuiltInTok{GeneralEquilibriumModeling}

\NormalTok{beta1 }\OperatorTok{=} \FloatTok{1} \OperatorTok{/} \FloatTok{4}
\NormalTok{beta2 }\OperatorTok{=} \FloatTok{2} \OperatorTok{/} \FloatTok{3}
\NormalTok{omega1 }\OperatorTok{=} \FloatTok{120.0}
\NormalTok{omega2 }\OperatorTok{=} \FloatTok{90.0}

\NormalTok{model }\OperatorTok{=}\NormalTok{ GEMB.}\FunctionTok{GEMBModel}\NormalTok{(}
\NormalTok{    [}\OperatorTok{:}\NormalTok{service1, }\OperatorTok{:}\NormalTok{service2, }\OperatorTok{:}\NormalTok{labor];}
\NormalTok{    numeraire}\OperatorTok{=:}\NormalTok{labor,}
\NormalTok{    numeraire\_value}\OperatorTok{=}\FloatTok{1.0}\NormalTok{,}
\NormalTok{)}

\NormalTok{GEMB.}\FunctionTok{add\_agent!}\NormalTok{(}
\NormalTok{    model,}
\NormalTok{    GEMB.}\FunctionTok{CESSpec}\NormalTok{([}\FloatTok{1.0}\NormalTok{]);}
\NormalTok{    outputs}\OperatorTok{=}\NormalTok{[}\OperatorTok{:}\NormalTok{service1, }\OperatorTok{:}\NormalTok{service2],}
\NormalTok{    demands}\OperatorTok{=:}\NormalTok{labor,}
\NormalTok{    activity\_start}\OperatorTok{=}\FloatTok{100.0}\NormalTok{,}
\NormalTok{    name}\OperatorTok{=:}\NormalTok{firm,}
\NormalTok{)}

\NormalTok{GEMB.}\FunctionTok{add\_agent!}\NormalTok{(}
\NormalTok{    model,}
\NormalTok{    GEMB.}\FunctionTok{CESSpec}\NormalTok{([beta1, }\FloatTok{1} \OperatorTok{{-}}\NormalTok{ beta1]);}
\NormalTok{    demands}\OperatorTok{=}\NormalTok{[}\OperatorTok{:}\NormalTok{service1, }\OperatorTok{:}\NormalTok{labor],}
\NormalTok{    endowments}\OperatorTok{=:}\NormalTok{labor,}
\NormalTok{    endowment\_quantities}\OperatorTok{=}\NormalTok{omega1,}
\NormalTok{    activity\_start}\OperatorTok{=}\FloatTok{100.0}\NormalTok{,}
\NormalTok{    name}\OperatorTok{=:}\NormalTok{consumer1,}
\NormalTok{)}

\NormalTok{GEMB.}\FunctionTok{add\_agent!}\NormalTok{(}
\NormalTok{    model,}
\NormalTok{    GEMB.}\FunctionTok{CESSpec}\NormalTok{([beta2, }\FloatTok{1} \OperatorTok{{-}}\NormalTok{ beta2]);}
\NormalTok{    demands}\OperatorTok{=}\NormalTok{[}\OperatorTok{:}\NormalTok{service2, }\OperatorTok{:}\NormalTok{labor],}
\NormalTok{    endowments}\OperatorTok{=:}\NormalTok{labor,}
\NormalTok{    endowment\_quantities}\OperatorTok{=}\NormalTok{omega2,}
\NormalTok{    activity\_start}\OperatorTok{=}\FloatTok{100.0}\NormalTok{,}
\NormalTok{    name}\OperatorTok{=:}\NormalTok{consumer2,}
\NormalTok{)}

\NormalTok{result }\OperatorTok{=}\NormalTok{ GEMB.}\FunctionTok{solve}\NormalTok{(}
\NormalTok{    model;}
\NormalTok{    p0}\OperatorTok{=}\NormalTok{[}\FloatTok{0.5}\NormalTok{, }\FloatTok{0.5}\NormalTok{, }\FloatTok{1.0}\NormalTok{],}
\NormalTok{    residual\_tol}\OperatorTok{=}\FloatTok{1.0e{-}8}\NormalTok{,}
\NormalTok{    silent}\OperatorTok{=}\ConstantTok{true}\NormalTok{,}
\NormalTok{)}

\NormalTok{GEMB.}\FunctionTok{print\_equilibrium\_statistics}\NormalTok{(model, result)}
\end{Highlighting}
\end{Shaded}

}

\end{codelisting}%

\section{\texorpdfstring{Capital Rental and Equity: A \(4\times3\)
Model}{Capital Rental and Equity: A 4\textbackslash times3 Model}}\label{capital-rental-and-equity-a-4times3-model}

Assume the following (see Torres, 2016, Tables 2.1 and 2.2):

\begin{enumerate}
\def\labelenumi{\arabic{enumi}.}
\item
  The economy contains four commodities: a product, labor, capital
  services, and equity. There are three economic agents: a production
  firm, a consumer, and a capital-rental firm. Equity, money, and ad
  valorem tax certificates can all be regarded as ad valorem claims.
\item
  The production firm produces the product according to
  \(y=x_2^{1-\beta_1}x_3^{\beta_1}\), where \(x_2\) and \(x_3\) denote
  labor and capital-service inputs, respectively, with \(0<\beta_1<1\).
\item
  The consumer has utility function \(u=x_1^{\beta_2}x_2^{1-\beta_2}\),
  where \(0<\beta_2<1\), and is endowed with one unit of labor and one
  unit of equity.
\item
  The capital-rental firm uses one unit of the product to provide one
  unit of capital services and recover \(1-\delta\) units of the product
  after depreciation. Let \(r>0\) denote the required rate of return.
  For each unit of activity, the firm demands equity whose value is
  \(rp_1\).
\end{enumerate}

Let \(\mathbf p=(p_1,p_2,p_3,p_4)^\top\), and let \(y\), \(u\), and
\(k\) denote the activity levels of the production firm, the consumer,
and the capital-rental firm, respectively. The demand matrix is

\[
\mathbf D(\mathbf p,y,u,k)
=
\begin{bmatrix}
0
&
u\left(\frac{\beta_2}{1-\beta_2}\frac{p_2}{p_1}\right)^{1-\beta_2}
&
k
\\
y\left(\frac{1-\beta_1}{\beta_1}\frac{p_3}{p_2}\right)^{\beta_1}
&
u\left(\frac{1-\beta_2}{\beta_2}\frac{p_1}{p_2}\right)^{\beta_2}
&
0
\\
y\left(\frac{\beta_1}{1-\beta_1}\frac{p_2}{p_3}\right)^{1-\beta_1}
&
0
&
0
\\
0
&
0
&
\frac{rp_1}{p_4}k
\end{bmatrix}
\]

and the supply matrix is

\[
\mathbf S(y,k)
=
\begin{bmatrix}
y&0&(1-\delta)k\\
0&1&0\\
0&0&k\\
0&1&0
\end{bmatrix}
\]

By solving the wide-sense structural equilibrium model, the analytical
solution for the equilibrium prices can be obtained. Taking the product
as the numeraire commodity, so that \(p_1^*=1\), the analytical
equilibrium prices are

\[
p_3^*=r+\delta,
\qquad
p_2^*
=
(1-\beta_1)
\left(
\frac{\beta_1}{r+\delta}
\right)^{\frac{\beta_1}{1-\beta_1}}
\]

and

\[
p_4^*
=
\frac{
r\beta_1\beta_2
}{
r(1-\beta_1\beta_2)+\delta(1-\beta_1)
}
p_2^*
\]

For \(r=1/0.97-1\approx0.030928\), \(\delta=0.06\), \(\beta_1=0.35\),
and \(\beta_2=0.4\), the equilibrium price vector is

\[
\mathbf p^*
\approx
(1,\ 1.343115,\ 0.090928,\ 0.088654)^\top
\]

\begin{codelisting}

\caption{\label{lst-gemb_CD_capitalRentalEquity_nc4_na3}A \(4\times3\)
Equilibrium Model with Capital Rental and Equity Claims}

\centering{

\begin{Shaded}
\begin{Highlighting}[]
\ImportTok{using} \BuiltInTok{GeneralEquilibriumModeling}

\NormalTok{discount\_factor }\OperatorTok{=} \FloatTok{0.97}
\NormalTok{return\_rate }\OperatorTok{=} \FloatTok{1} \OperatorTok{/}\NormalTok{ discount\_factor }\OperatorTok{{-}} \FloatTok{1}
\NormalTok{delta }\OperatorTok{=} \FloatTok{0.06}
\NormalTok{beta\_firm }\OperatorTok{=} \FloatTok{0.35}
\NormalTok{beta\_consumer }\OperatorTok{=} \FloatTok{0.4}

\NormalTok{model }\OperatorTok{=}\NormalTok{ GEMB.}\FunctionTok{GEMBModel}\NormalTok{(}
\NormalTok{    [}
        \OperatorTok{:}\NormalTok{product,}
        \OperatorTok{:}\NormalTok{labor,}
        \OperatorTok{:}\NormalTok{capital\_service,}
\NormalTok{        GEMB.}\FunctionTok{CommoditySpec}\NormalTok{(}\OperatorTok{:}\NormalTok{equity; price\_lower\_bound}\OperatorTok{=}\FloatTok{1.0e{-}10}\NormalTok{),}
\NormalTok{    ];}
\NormalTok{    numeraire}\OperatorTok{=:}\NormalTok{product,}
\NormalTok{    numeraire\_value}\OperatorTok{=}\FloatTok{1.0}\NormalTok{,}
\NormalTok{)}

\NormalTok{GEMB.}\FunctionTok{add\_agent!}\NormalTok{(}
\NormalTok{    model,}
\NormalTok{    GEMB.}\FunctionTok{CESSpec}\NormalTok{([}\FloatTok{1} \OperatorTok{{-}}\NormalTok{ beta\_firm, beta\_firm]);}
\NormalTok{    outputs}\OperatorTok{=:}\NormalTok{product,}
\NormalTok{    demands}\OperatorTok{=}\NormalTok{[}\OperatorTok{:}\NormalTok{labor, }\OperatorTok{:}\NormalTok{capital\_service],}
\NormalTok{    activity\_start}\OperatorTok{=}\FloatTok{100.0}\NormalTok{,}
\NormalTok{    name}\OperatorTok{=:}\NormalTok{production\_firm,}
\NormalTok{)}

\NormalTok{GEMB.}\FunctionTok{add\_agent!}\NormalTok{(}
\NormalTok{    model,}
\NormalTok{    GEMB.}\FunctionTok{CESSpec}\NormalTok{([beta\_consumer, }\FloatTok{1} \OperatorTok{{-}}\NormalTok{ beta\_consumer]);}
\NormalTok{    demands}\OperatorTok{=}\NormalTok{[}\OperatorTok{:}\NormalTok{product, }\OperatorTok{:}\NormalTok{labor],}
\NormalTok{    endowments}\OperatorTok{=}\NormalTok{[}\OperatorTok{:}\NormalTok{labor, }\OperatorTok{:}\NormalTok{equity],}
\NormalTok{    endowment\_quantities}\OperatorTok{=}\NormalTok{[}\FloatTok{1.0}\NormalTok{, }\FloatTok{1.0}\NormalTok{],}
\NormalTok{    activity\_start}\OperatorTok{=}\FloatTok{100.0}\NormalTok{,}
\NormalTok{    name}\OperatorTok{=:}\NormalTok{consumer,}
\NormalTok{)}

\NormalTok{GEMB.}\FunctionTok{add\_agent!}\NormalTok{(}
\NormalTok{    model,}
\NormalTok{    GEMB.}\FunctionTok{CESSpec}\NormalTok{([}\FloatTok{1.0}\NormalTok{]);}
\NormalTok{    outputs}\OperatorTok{=}\NormalTok{[}\OperatorTok{:}\NormalTok{product, }\OperatorTok{:}\NormalTok{capital\_service],}
\NormalTok{    output\_coefficients}\OperatorTok{=}\NormalTok{[}\FloatTok{1} \OperatorTok{{-}}\NormalTok{ delta, }\FloatTok{1.0}\NormalTok{],}
\NormalTok{    demands}\OperatorTok{=:}\NormalTok{product,}
\NormalTok{    claim\_rate}\OperatorTok{=}\NormalTok{return\_rate,}
\NormalTok{    claim}\OperatorTok{=}\NormalTok{GEMB.}\FunctionTok{CommodityRef}\NormalTok{(}\OperatorTok{:}\NormalTok{equity),}
\NormalTok{    activity\_start}\OperatorTok{=}\FloatTok{100.0}\NormalTok{,}
\NormalTok{    name}\OperatorTok{=:}\NormalTok{capital\_rental\_firm,}
\NormalTok{)}

\NormalTok{result }\OperatorTok{=}\NormalTok{ GEMB.}\FunctionTok{solve}\NormalTok{(}
\NormalTok{    model;}
\NormalTok{    p0}\OperatorTok{=}\NormalTok{[}\FloatTok{1.0}\NormalTok{, }\FloatTok{1.343115}\NormalTok{, }\FloatTok{0.090928}\NormalTok{, }\FloatTok{0.088654}\NormalTok{],}
\NormalTok{    residual\_tol}\OperatorTok{=}\FloatTok{1.0e{-}8}\NormalTok{,}
\NormalTok{    silent}\OperatorTok{=}\ConstantTok{true}\NormalTok{,}
\NormalTok{)}

\NormalTok{stats }\OperatorTok{=}\NormalTok{ GEMB.}\FunctionTok{equilibrium\_statistics}\NormalTok{(model, result)}

\NormalTok{flows }\OperatorTok{=}\NormalTok{ GEMB.}\FunctionTok{demand\_supply\_matrices}\NormalTok{(model, result)}

\NormalTok{D }\OperatorTok{=}\NormalTok{ flows.demand}
\NormalTok{S }\OperatorTok{=}\NormalTok{ flows.supply}
\NormalTok{Sbar }\OperatorTok{=}\NormalTok{ flows.net\_supply}

\NormalTok{GEMB.}\FunctionTok{print\_equilibrium\_statistics}\NormalTok{(model, result)}

\FunctionTok{println}\NormalTok{(}\StringTok{"}\SpecialCharTok{\textbackslash{}n}\StringTok{Demand matrix D:"}\NormalTok{)}
\FunctionTok{display}\NormalTok{(}\FunctionTok{hcat}\NormalTok{(D, flows.total\_demand))}

\FunctionTok{println}\NormalTok{(}\StringTok{"}\SpecialCharTok{\textbackslash{}n}\StringTok{Supply matrix S:"}\NormalTok{)}
\FunctionTok{display}\NormalTok{(}\FunctionTok{hcat}\NormalTok{(S, flows.total\_supply))}
\end{Highlighting}
\end{Shaded}

}

\end{codelisting}%

\bookmarksetup{startatroot}

\chapter{Static Equilibrium with Production: Extended
Models}\label{static-equilibrium-with-production-extended-models}

\section{\texorpdfstring{Pollution and Negative Prices: A \(3\times2\)
Model}{Pollution and Negative Prices: A 3\textbackslash times2 Model}}\label{pollution-and-negative-prices-a-3times2-model}

Assume the following:

\begin{enumerate}
\def\labelenumi{\arabic{enumi}.}
\item
  The economy contains three commodities and two economic agents, a firm
  and a consumer. Commodities 1, 2, and 3 are a product, labor, and
  pollution, respectively. Commodity 2 is taken as the numeraire
  commodity, so \(p_2=1\).
\item
  The firm operates at activity level \(z\). Each unit of activity uses
  \(a>0\) units of labor, produces one unit of the product, and jointly
  generates \(c>0\) units of pollution. The product output per unit of
  activity is normalized to one.
\item
  The consumer is endowed with \(\omega>0\) units of labor and consumes
  \(x\) units of the product, retains \(l\) units of labor, and accepts
  \(q\) units of pollution. Its utility function is
  \(u=\beta\ln x+(1-\beta)\ln l-\delta q\), where \(0<\beta<1\) and
  \(\delta>0\).
\end{enumerate}

For a given price vector \(\mathbf p=(p_1,1,p_3)^\top\), the consumer's
optimal demands are

\[
x(\mathbf p)=-\frac{\beta p_3}{\delta p_1},\qquad
l(\mathbf p)=-\frac{(1-\beta)p_3}{\delta},\qquad
q(\mathbf p)=\frac{1}{\delta}+\frac{\omega}{p_3}
\]

for an interior solution with \(p_3<-\delta\omega<0\). Hence the demand
and supply matrices are

\[
\mathbf D(\mathbf p,z)=
\begin{bmatrix}
0&-\frac{\beta p_3}{\delta p_1}\\
az&-\frac{(1-\beta)p_3}{\delta}\\
0&\frac{1}{\delta}+\frac{\omega}{p_3}
\end{bmatrix},
\qquad
\mathbf S(z)=
\begin{bmatrix}
z&0\\
0&\omega\\
cz&0
\end{bmatrix}
\]

The equilibrium can then be obtained from the equality-form wide-sense
structural equilibrium model.

Define \(\Delta=a+c\delta\omega\). The unique economically relevant
equilibrium activity level is

\[
z^*=
\frac{\Delta-\sqrt{\Delta^2-4ac\delta\beta\omega}}
{2ac\delta}
\]

with \(0<z^*<\min{\omega/a,\beta/(c\delta)}\). The equilibrium
quantities are \(x^*=z^*\), \(l^*=\omega-az^*\), and \(q^*=cz^*\), while

\[
p_1^*=
\frac{\beta(\omega-az^*)}{(1-\beta)z^*},
\qquad
p_2^*=1,
\qquad
p_3^*=
-\frac{\delta(\omega-az^*)}{1-\beta},
\qquad
\lambda^*=
\frac{1-\beta}{\omega-az^*}
\]

Thus, a commodity need not have a nonnegative equilibrium price in a
wide-sense structural equilibrium model. When a commodity is an
undesirable joint product such as pollution, its negative price
represents the compensation required for another economic agent to
accept it.

As a simple numerical case, let \(a=c=1\), \(\beta=1/2\),
\(\delta=1/8\), and \(\omega=6\). Then \(z^*=2\),
\((x^*,l^*,q^*)=(2,4,2)\), and

\[
\mathbf p^*=[2,1,-1]^\top
\]

so the equilibrium price of pollution is \(p_3^*=-1\).

\begin{codelisting}

\caption{\label{lst-gemb_pollution_negativePrice_nc3_na2}A \(3\times2\)
Equilibrium Model with Pollution and a Negative Price}

\centering{

\begin{Shaded}
\begin{Highlighting}[]
\CommentTok{\# GEMB high{-}level framework: pollution with a negative equilibrium price.}
\ImportTok{using} \BuiltInTok{GeneralEquilibriumModeling}

\NormalTok{a }\OperatorTok{=} \FloatTok{1.0}
\NormalTok{c }\OperatorTok{=} \FloatTok{1.0}
\NormalTok{beta }\OperatorTok{=} \FloatTok{0.5}
\NormalTok{delta }\OperatorTok{=} \FloatTok{1.0} \OperatorTok{/} \FloatTok{8.0}
\NormalTok{omega }\OperatorTok{=} \FloatTok{6.0}

\NormalTok{model }\OperatorTok{=}\NormalTok{ GEMB.}\FunctionTok{GEMBModel}\NormalTok{(}
\NormalTok{    [}
        \OperatorTok{:}\NormalTok{product,}
        \OperatorTok{:}\NormalTok{labor,}
\NormalTok{        GEMB.}\FunctionTok{CommoditySpec}\NormalTok{(}
            \OperatorTok{:}\NormalTok{pollution;}
\NormalTok{            price\_lower\_bound}\OperatorTok{={-}}\ConstantTok{Inf}\NormalTok{,}
\NormalTok{            price\_upper\_bound}\OperatorTok{=}\ConstantTok{Inf}\NormalTok{,}
\NormalTok{        ),}
\NormalTok{    ];}
\NormalTok{    numeraire}\OperatorTok{=:}\NormalTok{labor,}
\NormalTok{    numeraire\_value}\OperatorTok{=}\FloatTok{1.0}\NormalTok{,}
\NormalTok{)}

\NormalTok{firm\_spec }\OperatorTok{=}\NormalTok{ GEMB.}\FunctionTok{ActivityDemandSpec}\NormalTok{(}
\NormalTok{    (z, prices) }\OperatorTok{{-}\textgreater{}}\NormalTok{ [a }\OperatorTok{*}\NormalTok{ z],}
\NormalTok{)}

\NormalTok{GEMB.}\FunctionTok{add\_agent!}\NormalTok{(}
\NormalTok{    model,}
\NormalTok{    firm\_spec;}
\NormalTok{    outputs}\OperatorTok{=}\NormalTok{[}\OperatorTok{:}\NormalTok{product, }\OperatorTok{:}\NormalTok{pollution],}
\NormalTok{    output\_coefficients}\OperatorTok{=}\NormalTok{[}\FloatTok{1.0}\NormalTok{, c],}
\NormalTok{    demands}\OperatorTok{=:}\NormalTok{labor,}
\NormalTok{    activity\_start}\OperatorTok{=}\FloatTok{100.0}\NormalTok{,}
\NormalTok{    name}\OperatorTok{=:}\NormalTok{firm,}
\NormalTok{)}

\NormalTok{consumer\_spec }\OperatorTok{=}\NormalTok{ GEMB.}\FunctionTok{MarginalUtilityConsumerSpec}\NormalTok{(}
\NormalTok{    demand }\OperatorTok{{-}\textgreater{}} \ControlFlowTok{begin}
\NormalTok{        x, l, q }\OperatorTok{=}\NormalTok{ demand}
\NormalTok{        [}
\NormalTok{            beta }\OperatorTok{/}\NormalTok{ x,}
\NormalTok{            (}\FloatTok{1.0} \OperatorTok{{-}}\NormalTok{ beta) }\OperatorTok{/}\NormalTok{ l,}
            \OperatorTok{{-}}\NormalTok{delta,}
\NormalTok{        ]}
    \ControlFlowTok{end}\NormalTok{,}
\NormalTok{)}

\NormalTok{GEMB.}\FunctionTok{add\_agent!}\NormalTok{(}
\NormalTok{    model,}
\NormalTok{    consumer\_spec;}
\NormalTok{    demands}\OperatorTok{=}\NormalTok{[}\OperatorTok{:}\NormalTok{product, }\OperatorTok{:}\NormalTok{labor, }\OperatorTok{:}\NormalTok{pollution],}
\NormalTok{    endowments}\OperatorTok{=:}\NormalTok{labor,}
\NormalTok{    endowment\_quantities}\OperatorTok{=}\NormalTok{omega,}
\NormalTok{    demand\_start}\OperatorTok{=}\NormalTok{[}\FloatTok{100.0}\NormalTok{, }\FloatTok{100.0}\NormalTok{, }\FloatTok{100.0}\NormalTok{],}
\NormalTok{    demand\_lower\_bounds}\OperatorTok{=}\NormalTok{[}\FloatTok{1.0e{-}10}\NormalTok{, }\FloatTok{1.0e{-}10}\NormalTok{, }\FloatTok{0.0}\NormalTok{],}
\NormalTok{    multiplier\_start}\OperatorTok{=}\FloatTok{1.0}\NormalTok{,}
\NormalTok{    multiplier\_lower\_bound}\OperatorTok{=}\FloatTok{0.0}\NormalTok{,}
\NormalTok{    name}\OperatorTok{=:}\NormalTok{consumer,}
\NormalTok{)}

\NormalTok{result }\OperatorTok{=}\NormalTok{ GEMB.}\FunctionTok{solve}\NormalTok{(}
\NormalTok{    model;}
\NormalTok{    p0}\OperatorTok{=}\NormalTok{[}\FloatTok{1.0}\NormalTok{, }\FloatTok{1.0}\NormalTok{, }\OperatorTok{{-}}\FloatTok{0.5}\NormalTok{],}
\NormalTok{    residual\_tol}\OperatorTok{=}\FloatTok{1.0e{-}10}\NormalTok{,}
\NormalTok{    silent}\OperatorTok{=}\ConstantTok{true}\NormalTok{,}
\NormalTok{)}

\NormalTok{GEMB.}\FunctionTok{print\_equilibrium\_statistics}\NormalTok{(}
\NormalTok{    model,}
\NormalTok{    result;}
\NormalTok{    display\_tol}\OperatorTok{=}\FloatTok{1.0e{-}10}\NormalTok{,}
\NormalTok{    sigdigits}\OperatorTok{=}\FloatTok{8}\NormalTok{,}
\NormalTok{)}
\end{Highlighting}
\end{Shaded}

}

\end{codelisting}%

\section{\texorpdfstring{Producer-to-Consumer Externality: A
\(2\times2\)
Model}{Producer-to-Consumer Externality: A 2\textbackslash times2 Model}}\label{producer-to-consumer-externality-a-2times2-model}

Assume the following:

\begin{enumerate}
\def\labelenumi{\arabic{enumi}.}
\item
  The economy contains two commodities and two economic agents, a firm
  and a consumer. Commodities 1 and 2 are a product and labor,
  respectively. Commodity 2 is taken as the numeraire commodity, so
  \(p_2=1\).
\item
  The firm uses one unit of labor to produce one unit of the product.
  Let \(y\) denote its output.
\item
  The consumer is endowed with \(\omega>0\) units of labor and has
  utility function \(u=(1-\kappa y)\ln x_1+\ln x_2\), where \(\kappa>0\)
  measures the negative external effect of the firm's production and
  \(1-\kappa y>0\). The consumer takes \(y\) as given when choosing
  consumption.
\end{enumerate}

For a given \(\mathbf p=(p_1,p_2)^\top\) and \(y\), the consumer's
demands are

\[
x_1(\mathbf p,y)=\frac{1-\kappa y}{2-\kappa y}\frac{\omega p_2}{p_1},
\qquad
x_2(\mathbf p,y)=\frac{\omega}{2-\kappa y}
\]

Hence the demand and supply matrices are

\[
\mathbf D(\mathbf p,y)=
\begin{bmatrix}
0&\frac{1-\kappa y}{2-\kappa y}\frac{\omega p_2}{p_1}\\
y&\frac{\omega}{2-\kappa y}
\end{bmatrix},
\qquad
\mathbf S(y)=
\begin{bmatrix}
y&0\\
0&\omega
\end{bmatrix}
\]

Applying the equality-form wide-sense structural equilibrium model
yields \(p_1^*=1\) and

\[
y^*=
\frac{2+\omega\kappa-\sqrt{4+\omega^2\kappa^2}}
{2\kappa}
\]

The equilibrium consumption vector is
\([x_1^*,x_2^*]^\top=[y^*,\omega-y^*]^\top\).

For \(\omega=12\) and \(\kappa=1/8\),

\[
y^*=4,
\qquad
\mathbf p^*=[1,1]^\top
\]

\begin{codelisting}

\caption{\label{lst-gemb_CD_productionExternality_nc2_na2}A \(2\times2\)
Equilibrium Model with Producer-to-Consumer Externality}

\centering{

\begin{Shaded}
\begin{Highlighting}[]
\CommentTok{\# GEMB high{-}level framework: production externality.}
\ImportTok{using} \BuiltInTok{GeneralEquilibriumModeling.GEMB}

\NormalTok{omega }\OperatorTok{=} \FloatTok{12.0}
\NormalTok{kappa }\OperatorTok{=} \FloatTok{1.0} \OperatorTok{/} \FloatTok{8.0}

\NormalTok{model }\OperatorTok{=} \FunctionTok{GEMBModel}\NormalTok{(}
\NormalTok{    [}\OperatorTok{:}\NormalTok{product, }\OperatorTok{:}\NormalTok{labor];}
\NormalTok{    numeraire}\OperatorTok{=:}\NormalTok{labor,}
\NormalTok{)}

\FunctionTok{add\_agent!}\NormalTok{(}
\NormalTok{    model,}
    \FunctionTok{CESSpec}\NormalTok{([}\FloatTok{1.0}\NormalTok{]);}
\NormalTok{    outputs}\OperatorTok{=:}\NormalTok{product,}
\NormalTok{    demands}\OperatorTok{=:}\NormalTok{labor,}
\NormalTok{    activity\_start}\OperatorTok{=}\FloatTok{4.0}\NormalTok{,}
\NormalTok{    name}\OperatorTok{=:}\NormalTok{firm,}
\NormalTok{)}

\FunctionTok{add\_agent!}\NormalTok{(}
\NormalTok{    model,}
    \FunctionTok{MarshallDemandConsumerSpec}\NormalTok{(}
\NormalTok{        (income, prices, observed) }\OperatorTok{{-}\textgreater{}} \ControlFlowTok{begin}
\NormalTok{            activity }\OperatorTok{=}\NormalTok{ observed[}\FloatTok{1}\NormalTok{]}
\NormalTok{            [}
\NormalTok{                (}\FloatTok{1.0} \OperatorTok{{-}}\NormalTok{ kappa }\OperatorTok{*}\NormalTok{ activity) }\OperatorTok{/}\NormalTok{ (}\FloatTok{2.0} \OperatorTok{{-}}\NormalTok{ kappa }\OperatorTok{*}\NormalTok{ activity) }\OperatorTok{*}
\NormalTok{                    income }\OperatorTok{/}\NormalTok{ prices[}\FloatTok{1}\NormalTok{],}
\NormalTok{                income }\OperatorTok{/}\NormalTok{ ((}\FloatTok{2.0} \OperatorTok{{-}}\NormalTok{ kappa }\OperatorTok{*}\NormalTok{ activity) }\OperatorTok{*}\NormalTok{ prices[}\FloatTok{2}\NormalTok{]),}
\NormalTok{            ]}
        \ControlFlowTok{end}\NormalTok{,}
\NormalTok{    );}
\NormalTok{    demands}\OperatorTok{=}\NormalTok{[}\OperatorTok{:}\NormalTok{product, }\OperatorTok{:}\NormalTok{labor],}
\NormalTok{    endowments}\OperatorTok{=:}\NormalTok{labor,}
\NormalTok{    endowment\_quantities}\OperatorTok{=}\NormalTok{omega,}
\NormalTok{    observed\_variables}\OperatorTok{=}\NormalTok{[}
        \FunctionTok{agent\_variable\_ref}\NormalTok{(}\FunctionTok{AgentRef}\NormalTok{(}\OperatorTok{:}\NormalTok{firm), }\OperatorTok{:}\NormalTok{activity),}
\NormalTok{    ],}
\NormalTok{    name}\OperatorTok{=:}\NormalTok{consumer,}
\NormalTok{)}

\NormalTok{result }\OperatorTok{=} \FunctionTok{solve}\NormalTok{(}
\NormalTok{    model;}
\NormalTok{    p0}\OperatorTok{=}\NormalTok{[}\FloatTok{2.0}\NormalTok{, }\FloatTok{1.0}\NormalTok{],}
\NormalTok{    silent}\OperatorTok{=}\ConstantTok{true}\NormalTok{,}
\NormalTok{)}

\FunctionTok{print\_equilibrium\_statistics}\NormalTok{(model, result)}
\end{Highlighting}
\end{Shaded}

}

\end{codelisting}%

\section{Marginal Pricing Principle and Scale Claim
Rates}\label{marginal-pricing-principle-and-scale-claim-rates}

\subsection{Scale Claim Rates}\label{scale-claim-rates}

As noted earlier, profit rates, ad valorem subsidy rates, and similar
value ratios can generally be interpreted as ad valorem claim rates.

A marginal pricing equilibrium (also called marginal cost pricing
equilibrium) can be viewed as a natural extension of competitive
equilibrium to nonconvex economies (Bonnisseau and Cornet, 1990;
Quinzii, 1992). Under appropriate conditions, a Pareto-optimal
allocation can often be represented as a marginal-pricing equilibrium
allocation.

In a marginal-pricing equilibrium, an appropriate claim rate may be
required to reconcile marginal pricing with the firm's
revenue-expenditure balance. When this claim rate is determined by the
firm's scale elasticity, it is referred to as the \textbf{scale claim
rate}. The scale claim rate denotes the value of the scale claim
relative to the value of ordinary inputs. When positive, it can be
interpreted as a cost-based profit rate, or scale profit rate; when
negative, its absolute value can be interpreted as a subsidy rate.

\begin{enumerate}
\def\labelenumi{\arabic{enumi}.}
\tightlist
\item
  Let a firm's production function be \(y=f(\mathbf x)\). Its scale
  elasticity is
\end{enumerate}

\[
\varepsilon
=
\frac{\mathbf x^\top \nabla f(\mathbf x)}{f(\mathbf x)}
\]

and the corresponding scale claim rate is

\[
\tau_s
=
\frac{f(\mathbf x)}{\mathbf x^\top \nabla f(\mathbf x)}-1
=
\frac{1}{\varepsilon}-1
\]

\begin{enumerate}
\def\labelenumi{\arabic{enumi}.}
\setcounter{enumi}{1}
\item
  Under marginal pricing, inputs are valued according to their marginal
  products. The scale claim rate measures the proportional claim
  required to reconcile the value of output with the value of inputs
  evaluated at marginal products. Thus, incorporating a claim at rate
  \(\tau_s\) allows marginal pricing to be represented together with the
  firm's revenue-expenditure balance. The marginal-pricing principle can
  be applied not only to firms but also to consumer decision-making.
\item
  If \(f(\mathbf x)\) is homogeneous of degree \(\theta>0\), Euler's
  theorem gives
  \(\mathbf x^\top \nabla f(\mathbf x)=\theta f(\mathbf x)\), and hence
\end{enumerate}

\[
\tau_s=\frac{1}{\theta}-1
\]

Thus, \(\tau_s=0\) under constant returns to scale, \(\tau_s>0\) under
decreasing returns to scale, and \(-1<\tau_s<0\) under increasing
returns to scale. A positive scale claim rate implies that the firm
distributes profits to its shareholders, whereas a negative scale claim
rate implies that a subsidy must be provided to the firm to support
marginal pricing.

\subsection{\texorpdfstring{Marginal Pricing with a Power Production
Function: A \(3\times2\)
Model}{Marginal Pricing with a Power Production Function: A 3\textbackslash times2 Model}}\label{marginal-pricing-with-a-power-production-function-a-3times2-model}

Consider a simple equilibrium model with marginal pricing and an ad
valorem claim. Assume the following:

\begin{enumerate}
\def\labelenumi{\arabic{enumi}.}
\item
  The economy contains three commodities: a product, labor, and an ad
  valorem claim commodity. There is one firm and one consumer. Labor is
  the numeraire commodity.
\item
  The firm has the production function \(y=\alpha x_2^\theta\), where
  \(x_2\) is labor input and \(\alpha,\theta>0\).
\item
  The consumer has utility function \(u=x_1^\beta x_2^{1-\beta}\), where
  \(0<\beta<1\), and is endowed with \(\omega>0\) units of labor.
\item
  The consumer owns one unit of the claim commodity. A claim at rate
  \(\tau>-1\) is imposed on the firm's input expenditure. A positive
  \(\tau\) can be interpreted as profit distributed by the firm to the
  claim holder, whereas a negative \(\tau\) represents a subsidy
  provided to the firm. The scale-claim representation below assumes
  \(\tau\ne0\). When \(\tau=0\), the scale claim has zero value and can
  simply be omitted from the model.
\end{enumerate}

Let \(\mathbf p=(p_1,1,p_3)^\top\) and let \(w\) denote consumer income.
The demand and supply matrices are

\[
\mathbf D(\mathbf p,w,y)=
\begin{bmatrix}
0&\beta w/p_1\\
(y/\alpha)^{1/\theta}&(1-\beta)w\\
(y/\alpha)^{1/\theta}\tau/p_3&0
\end{bmatrix},
\qquad
\mathbf S(y)=
\begin{bmatrix}
y&0\\
0&\omega\\
0&1
\end{bmatrix}
\]

Applying the equality-form wide-sense structural equilibrium model gives

\[
p_1^*=
\frac{1+\tau}{\alpha}
\left(
\frac{\beta\omega}{1+(1-\beta)\tau}
\right)^{1-\theta},
\qquad
p_3^*=
\frac{\tau\beta\omega}{1+(1-\beta)\tau}
\]

and

\[
w^*=
\frac{\omega(1+\tau)}{1+(1-\beta)\tau},
\qquad
y^*=
\alpha
\left(
\frac{\beta\omega}{1+(1-\beta)\tau}
\right)^\theta
\]

Under marginal pricing, the claim rate should satisfy the
scale-claim-rate consistency condition

\[
\tau=\frac{1}{\theta}-1
\]

The equilibrium then becomes

\[
p_1^*=
\frac{1}{\alpha\theta}
\left(
\frac{\beta\theta\omega}{1-\beta+\beta\theta}
\right)^{1-\theta},
\qquad
p_3^*=
\frac{\beta\omega(1-\theta)}
{1-\beta+\beta\theta}
\]

and

\[
w^*=
\frac{\omega}{1-\beta+\beta\theta},
\qquad
y^*=
\alpha
\left(
\frac{\beta\theta\omega}
{1-\beta+\beta\theta}
\right)^\theta
\]

Thus, the scale claim reconciles marginal pricing with the firm's
revenue-expenditure balance.

For a simple comparison, let \(\beta=1/2\) and \(\omega=3\) in both
cases. Under decreasing returns to scale, take \(\theta=1/2\) and
\(\alpha=2\). The scale claim rate is then \(\tau=1\), and the
equilibrium is \(\mathbf p^*=[1,1,1]^\top\), \(w^*=4\), and \(y^*=2\).
Under increasing returns to scale, take \(\theta=2\) and \(\alpha=1/4\).
The corresponding scale claim rate is \(\tau=-1/2\), and the equilibrium
is \(\mathbf p^*=[1,1,-1]^\top\), \(w^*=2\), and \(y^*=1\).

\begin{codelisting}

\caption{\label{lst-gemb_powerProduction_scaleClaimRate_nc3_na2}A
\(3\times2\) Equilibrium Model with a Power Production Function and a
Scale Claim}

\centering{

\begin{Shaded}
\begin{Highlighting}[]
\ImportTok{using} \BuiltInTok{GeneralEquilibriumModeling.GEMB}

\NormalTok{beta }\OperatorTok{=} \FloatTok{0.5}
\NormalTok{omega }\OperatorTok{=} \FloatTok{3.0}

\NormalTok{increasing\_returns }\OperatorTok{=} \ConstantTok{true} \CommentTok{\# Set to false for the alternative numerical example.}

\NormalTok{alpha }\OperatorTok{=}\NormalTok{ increasing\_returns ? }\FloatTok{0.25} \OperatorTok{:} \FloatTok{2.0}
\NormalTok{theta }\OperatorTok{=}\NormalTok{ increasing\_returns ? }\FloatTok{2.0} \OperatorTok{:} \FloatTok{0.5}
\NormalTok{tau }\OperatorTok{=} \FloatTok{1.0} \OperatorTok{/}\NormalTok{ theta }\OperatorTok{{-}} \FloatTok{1.0}

\NormalTok{model }\OperatorTok{=} \FunctionTok{GEMBModel}\NormalTok{(}
\NormalTok{    [}
        \FunctionTok{CommoditySpec}\NormalTok{(}
            \OperatorTok{:}\NormalTok{product;}
\NormalTok{            price\_lower\_bound}\OperatorTok{=}\FloatTok{1e{-}3}\NormalTok{,}
\NormalTok{        ),}
        \OperatorTok{:}\NormalTok{labor,}
        \FunctionTok{CommoditySpec}\NormalTok{(}
            \OperatorTok{:}\NormalTok{claim;}
\NormalTok{            price\_lower\_bound}\OperatorTok{=}\NormalTok{increasing\_returns ? }\OperatorTok{{-}}\ConstantTok{Inf} \OperatorTok{:} \FloatTok{1e{-}3}\NormalTok{,}
\NormalTok{            price\_upper\_bound}\OperatorTok{=}\NormalTok{increasing\_returns ? }\FloatTok{1e{-}3} \OperatorTok{:} \ConstantTok{Inf}\NormalTok{,}
\NormalTok{        ),}
\NormalTok{    ];}
\NormalTok{    numeraire}\OperatorTok{=:}\NormalTok{labor,}
\NormalTok{)}

\FunctionTok{add\_agent!}\NormalTok{(}
\NormalTok{    model,}
    \FunctionTok{PowerProductionSpec}\NormalTok{(}
\NormalTok{        alpha}\OperatorTok{=}\NormalTok{alpha,}
\NormalTok{        theta}\OperatorTok{=}\NormalTok{theta,}
\NormalTok{    );}
\NormalTok{    outputs}\OperatorTok{=:}\NormalTok{product,}
\NormalTok{    demands}\OperatorTok{=:}\NormalTok{labor,}
\NormalTok{    claim}\OperatorTok{=:}\NormalTok{claim,}
\NormalTok{    claim\_rate}\OperatorTok{=}\NormalTok{tau,}
\NormalTok{    activity\_start}\OperatorTok{=}\NormalTok{increasing\_returns ? }\FloatTok{0.8} \OperatorTok{:} \FloatTok{1.5}\NormalTok{,}
\NormalTok{    activity\_lower\_bound}\OperatorTok{=}\FloatTok{1.0e{-}6}\NormalTok{,}
\NormalTok{    name}\OperatorTok{=:}\NormalTok{firm,}
\NormalTok{)}

\FunctionTok{add\_agent!}\NormalTok{(}
\NormalTok{    model,}
    \FunctionTok{CESSpec}\NormalTok{([beta, }\FloatTok{1.0} \OperatorTok{{-}}\NormalTok{ beta]);}
\NormalTok{    demands}\OperatorTok{=}\NormalTok{[}\OperatorTok{:}\NormalTok{product, }\OperatorTok{:}\NormalTok{labor],}
\NormalTok{    endowments}\OperatorTok{=}\NormalTok{[}\OperatorTok{:}\NormalTok{labor, }\OperatorTok{:}\NormalTok{claim],}
\NormalTok{    endowment\_quantities}\OperatorTok{=}\NormalTok{[omega, }\FloatTok{1.0}\NormalTok{],}
\NormalTok{    activity\_start}\OperatorTok{=}\FloatTok{1.5}\NormalTok{,}
\NormalTok{    name}\OperatorTok{=:}\NormalTok{consumer,}
\NormalTok{)}

\NormalTok{result }\OperatorTok{=} \FunctionTok{solve}\NormalTok{(}
\NormalTok{    model;}
\NormalTok{    p0}\OperatorTok{=}\NormalTok{increasing\_returns ?}
\NormalTok{        [}\FloatTok{1.4}\NormalTok{, }\FloatTok{1.0}\NormalTok{, }\OperatorTok{{-}}\FloatTok{0.8}\NormalTok{] }\OperatorTok{:}
\NormalTok{        [}\FloatTok{1.4}\NormalTok{, }\FloatTok{1.0}\NormalTok{, }\FloatTok{0.7}\NormalTok{],}
\NormalTok{    silent}\OperatorTok{=}\ConstantTok{true}\NormalTok{,}
\NormalTok{)}

\FunctionTok{print\_equilibrium\_statistics}\NormalTok{(model, result)}
\FunctionTok{print}\NormalTok{(result.solved)}
\end{Highlighting}
\end{Shaded}

}

\end{codelisting}%

\subsection{\texorpdfstring{Marginal Pricing with an Exponential
Production Function: A \(3\times2\)
Model}{Marginal Pricing with an Exponential Production Function: A 3\textbackslash times2 Model}}\label{marginal-pricing-with-an-exponential-production-function-a-3times2-model}

Assume the following:

\begin{enumerate}
\def\labelenumi{\arabic{enumi}.}
\item
  The economy contains three commodities and two economic agents, a firm
  and a consumer. Commodities 1, 2, and 3 are a product, labor, and a
  scale claim, respectively. Commodity 2 is taken as the numeraire
  commodity, so \(p_2=1\).
\item
  The firm minimizes cost for a given output \(y\) and has the piecewise
  production function
\end{enumerate}

\[
f(x)=
\begin{cases}
\frac{\alpha\theta^\xi}{\xi}x,&0\leq x<\xi\\[4pt]
\alpha\theta^x,&x\geq\xi
\end{cases}
\]

where \(\alpha>0\), \(\theta>1\), and \(\xi>0\). Only equilibria with
\(x^*>\xi\) are considered. In this region, \(y=\alpha\theta^x\) and
\(x=\ln(y/\alpha)/\ln\theta\).

\begin{enumerate}
\def\labelenumi{\arabic{enumi}.}
\setcounter{enumi}{2}
\item
  The consumer has utility function \(u=x_1^\beta x_2^{1-\beta}\), where
  \(0<\beta<1\), and is endowed with \(\omega>0\) units of labor and one
  unit of the scale claim.
\item
  Under marginal pricing, the scale claim rate depends on the firm's
  input and is given by
\end{enumerate}

\[
\tau(x)=\frac{1}{x\ln\theta}-1
\] The scale-claim representation below assumes \(\tau(x^*)\ne0\). If
\(\tau(x^*)=0\), the scale claim has zero value and can simply be
omitted from the model.

Thus, unlike the power-production example, the scale claim rate is
endogenous and varies with the equilibrium scale of production.

Let \(\mathbf p=(p_1,1,p_3)^\top\) and let \(w\) denote consumer income.
The demand and supply matrices are

\[
\mathbf D(\mathbf p,w,y)=
\begin{bmatrix}
0&\beta w/p_1\\
x&(1-\beta)w\\
x\tau(x)/p_3&0
\end{bmatrix},
\qquad
\mathbf S(y)=
\begin{bmatrix}
y&0\\
0&\omega\\
0&1
\end{bmatrix},
\qquad
x=\frac{\ln(y/\alpha)}{\ln\theta}
\]

Applying the equality-form wide-sense structural equilibrium model
together with the scale-claim-rate consistency condition yields

\[
x^*=
\omega-\frac{1-\beta}{\beta\ln\theta}
\]

which requires

\[
\omega>
\xi+\frac{1-\beta}{\beta\ln\theta}
\]

The corresponding scale claim rate is

\[
\tau^*=
\frac{1-\beta\omega\ln\theta}
{\beta\omega\ln\theta-(1-\beta)}
\]

and the remaining equilibrium values are

\[
p_1^*=
\frac{e^{(1-\beta)/\beta}}
{\alpha(\ln\theta)\theta^\omega},
\qquad
p_3^*=
\frac{1}{\beta\ln\theta}-\omega
\]

\[
w^*=
\frac{1}{\beta\ln\theta},
\qquad
y^*=
\alpha\theta^\omega e^{-(1-\beta)/\beta}
\]

Thus, the scale claim rate is determined jointly with the firm's
production scale.

For a simple comparison, let \(\beta=1/2\), \(\theta=e\), and
\(\xi=1/4\) in both cases. When \(\omega=3/2\) and \(\alpha=e^{-1/2}\),
the equilibrium input is \(x^*=1/2\), the endogenous scale claim rate is
\(\tau^*=1\), and the equilibrium is \(\mathbf p^*=[1,1,1/2]^\top\),
\(w^*=2\), and \(y^*=1\). When \(\omega=3\) and \(\alpha=e^{-2}\), the
equilibrium input is \(x^*=2\), the endogenous scale claim rate is
\(\tau^*=-1/2\), and the equilibrium is \(\mathbf p^*=[1,1,-1]^\top\),
\(w^*=2\), and \(y^*=1\).

In the GEMB implementation, an additional condition agent is introduced
to impose the endogenous scale-claim-rate consistency condition, so the
computational model contains three agents although the economy itself
contains only two economic agents.

\begin{codelisting}

\caption{\label{lst-gemb_exponentialProduction_endogenousScaleClaimRate_nc3_na2}An
Equilibrium Model with Exponential Production and an Endogenous Scale
Claim Rate}

\centering{

\begin{Shaded}
\begin{Highlighting}[]
\ImportTok{using} \BuiltInTok{GeneralEquilibriumModeling.GEMB}

\NormalTok{beta }\OperatorTok{=} \FloatTok{0.5}
\NormalTok{theta }\OperatorTok{=} \FunctionTok{exp}\NormalTok{(}\FloatTok{1.0}\NormalTok{)}
\NormalTok{xi }\OperatorTok{=} \FloatTok{0.25}

\NormalTok{positive\_claim }\OperatorTok{=} \ConstantTok{false} \CommentTok{\# Set to true for the alternative numerical example.}

\NormalTok{omega }\OperatorTok{=}\NormalTok{ positive\_claim ? }\FloatTok{1.5} \OperatorTok{:} \FloatTok{3.0}
\NormalTok{alpha }\OperatorTok{=}\NormalTok{ positive\_claim ? }\FunctionTok{exp}\NormalTok{(}\OperatorTok{{-}}\FloatTok{0.5}\NormalTok{) }\OperatorTok{:} \FunctionTok{exp}\NormalTok{(}\OperatorTok{{-}}\FloatTok{2.0}\NormalTok{)}

\NormalTok{model }\OperatorTok{=} \FunctionTok{GEMBModel}\NormalTok{(}
\NormalTok{    [}
        \FunctionTok{CommoditySpec}\NormalTok{(}
            \OperatorTok{:}\NormalTok{product;}
\NormalTok{            price\_lower\_bound}\OperatorTok{=}\FloatTok{1e{-}3}\NormalTok{,}
\NormalTok{        ),}
        \OperatorTok{:}\NormalTok{labor,}
        \FunctionTok{CommoditySpec}\NormalTok{(}
            \OperatorTok{:}\NormalTok{scale\_claim;}
\NormalTok{            price\_lower\_bound}\OperatorTok{=}\NormalTok{positive\_claim ? }\FloatTok{1e{-}3} \OperatorTok{:} \OperatorTok{{-}}\ConstantTok{Inf}\NormalTok{,}
\NormalTok{            price\_upper\_bound}\OperatorTok{=}\NormalTok{positive\_claim ? }\ConstantTok{Inf} \OperatorTok{:} \OperatorTok{{-}}\FloatTok{1e{-}3}\NormalTok{,}
\NormalTok{        ),}
\NormalTok{    ];}
\NormalTok{    numeraire}\OperatorTok{=:}\NormalTok{labor,}
\NormalTok{)}

\FunctionTok{add\_agent!}\NormalTok{(}
\NormalTok{    model,}
    \FunctionTok{ConditionAgentSpec}\NormalTok{(}
\NormalTok{        (variables, observed) }\OperatorTok{{-}\textgreater{}} \ControlFlowTok{begin}
\NormalTok{            tau }\OperatorTok{=}\NormalTok{ variables[}\FloatTok{1}\NormalTok{]}
\NormalTok{            activity }\OperatorTok{=}\NormalTok{ observed[}\FloatTok{1}\NormalTok{]}

\NormalTok{            [}
\NormalTok{                (}\FloatTok{1.0} \OperatorTok{+}\NormalTok{ tau) }\OperatorTok{*} \FunctionTok{log}\NormalTok{(activity }\OperatorTok{/}\NormalTok{ alpha) }\OperatorTok{{-}} \FloatTok{1.0}\NormalTok{,}
\NormalTok{            ]}
        \ControlFlowTok{end}\NormalTok{,}
\NormalTok{    );}
\NormalTok{    variable\_names}\OperatorTok{=:}\NormalTok{claim\_rate,}
\NormalTok{    variable\_start}\OperatorTok{=}\NormalTok{positive\_claim ? }\FloatTok{0.6} \OperatorTok{:} \OperatorTok{{-}}\FloatTok{0.3}\NormalTok{,}
\NormalTok{    variable\_lower\_bounds}\OperatorTok{={-}}\FloatTok{0.99}\NormalTok{,}
\NormalTok{    observed\_variables}\OperatorTok{=}\NormalTok{[}
        \FunctionTok{agent\_variable\_ref}\NormalTok{(}\OperatorTok{:}\NormalTok{firm, }\OperatorTok{:}\NormalTok{activity),}
\NormalTok{    ],}
\NormalTok{    name}\OperatorTok{=:}\NormalTok{scaleClaimSetter,}
\NormalTok{)}

\FunctionTok{add\_agent!}\NormalTok{(}
\NormalTok{    model,}
    \FunctionTok{ActivityDemandSpec}\NormalTok{(}
\NormalTok{        (activity, prices, observed) }\OperatorTok{{-}\textgreater{}} \ControlFlowTok{begin}
\NormalTok{            p\_labor, p\_claim }\OperatorTok{=}\NormalTok{ prices}
\NormalTok{            tau }\OperatorTok{=}\NormalTok{ observed[}\FloatTok{1}\NormalTok{]}

\NormalTok{            x }\OperatorTok{=} \FunctionTok{log}\NormalTok{(activity }\OperatorTok{/}\NormalTok{ alpha) }\OperatorTok{/} \FunctionTok{log}\NormalTok{(theta)}

\NormalTok{            [}
\NormalTok{                x,}
\NormalTok{                tau }\OperatorTok{*}\NormalTok{ p\_labor }\OperatorTok{*}\NormalTok{ x }\OperatorTok{/}\NormalTok{ p\_claim,}
\NormalTok{            ]}
        \ControlFlowTok{end}\NormalTok{;}
\NormalTok{        producer\_condition\_function}\OperatorTok{=}
\NormalTok{            (variables, prices, net\_supply) }\OperatorTok{{-}\textgreater{}}\NormalTok{ [}
                \FunctionTok{{-}sum}\NormalTok{(prices }\OperatorTok{.*}\NormalTok{ net\_supply),}
\NormalTok{            ],}
\NormalTok{    );}
\NormalTok{    outputs}\OperatorTok{=:}\NormalTok{product,}
\NormalTok{    demands}\OperatorTok{=}\NormalTok{[}\OperatorTok{:}\NormalTok{labor, }\OperatorTok{:}\NormalTok{scale\_claim],}
\NormalTok{    observed\_variables}\OperatorTok{=}\NormalTok{[}
        \FunctionTok{agent\_variable\_ref}\NormalTok{(}
            \OperatorTok{:}\NormalTok{scaleClaimSetter,}
            \OperatorTok{:}\NormalTok{claim\_rate,}
\NormalTok{        ),}
\NormalTok{    ],}
\NormalTok{    activity\_start}\OperatorTok{=}\NormalTok{positive\_claim ? }\FloatTok{0.9} \OperatorTok{:} \FloatTok{0.8}\NormalTok{,}
\NormalTok{    activity\_lower\_bound}\OperatorTok{=}\NormalTok{alpha }\OperatorTok{*}\NormalTok{ theta}\OperatorTok{\^{}}\NormalTok{xi,}
\NormalTok{    name}\OperatorTok{=:}\NormalTok{firm,}
\NormalTok{)}

\FunctionTok{add\_agent!}\NormalTok{(}
\NormalTok{    model,}
    \FunctionTok{CESSpec}\NormalTok{([beta, }\FloatTok{1.0} \OperatorTok{{-}}\NormalTok{ beta]);}
\NormalTok{    demands}\OperatorTok{=}\NormalTok{[}\OperatorTok{:}\NormalTok{product, }\OperatorTok{:}\NormalTok{labor],}
\NormalTok{    endowments}\OperatorTok{=}\NormalTok{[}\OperatorTok{:}\NormalTok{labor, }\OperatorTok{:}\NormalTok{scale\_claim],}
\NormalTok{    endowment\_quantities}\OperatorTok{=}\NormalTok{[omega, }\FloatTok{1.0}\NormalTok{],}
\NormalTok{    activity\_start}\OperatorTok{=}\FloatTok{0.8}\NormalTok{,}
\NormalTok{    name}\OperatorTok{=:}\NormalTok{consumer,}
\NormalTok{)}

\NormalTok{result }\OperatorTok{=} \FunctionTok{solve}\NormalTok{(}
\NormalTok{    model;}
\NormalTok{    p0}\OperatorTok{=}\NormalTok{positive\_claim ?}
\NormalTok{        [}\FloatTok{1.3}\NormalTok{, }\FloatTok{1.0}\NormalTok{, }\FloatTok{0.8}\NormalTok{] }\OperatorTok{:}
\NormalTok{        [}\FloatTok{1.3}\NormalTok{, }\FloatTok{1.0}\NormalTok{, }\OperatorTok{{-}}\FloatTok{0.7}\NormalTok{],}
\NormalTok{    silent}\OperatorTok{=}\ConstantTok{true}\NormalTok{,}
\NormalTok{)}

\NormalTok{claim\_rate }\OperatorTok{=}\NormalTok{ result.agent\_variable\_values[}\FloatTok{1}\NormalTok{][}\FloatTok{1}\NormalTok{]}
\NormalTok{activity }\OperatorTok{=}\NormalTok{ result.agent\_variable\_values[}\FloatTok{2}\NormalTok{][}\FloatTok{1}\NormalTok{]}

\FunctionTok{print\_equilibrium\_statistics}\NormalTok{(model, result)}

\FunctionTok{println}\NormalTok{(}\StringTok{"}\SpecialCharTok{\textbackslash{}n}\StringTok{Scale claim rate: "}\NormalTok{, claim\_rate)}
\FunctionTok{println}\NormalTok{(}\StringTok{"Firm activity: "}\NormalTok{, activity)}
\end{Highlighting}
\end{Shaded}

}

\end{codelisting}%

\section{Degree-One Homogenization of Production
Functions}\label{sec-Homogenization}

\subsection{Degree-One Homogenization and the Scale Claim
Rate}\label{degree-one-homogenization-and-the-scale-claim-rate}

Degree-one homogenization is commonly used to transform a production
function with decreasing returns to scale into one with constant returns
to scale (McKenzie, 2002, p.~216). More generally, the method can also
be applied to production functions with increasing or variable returns
to scale, thereby providing a unified constant-returns-to-scale
representation for different types of returns to scale.

Consider a firm with production function \(y=f(\mathbf x)\), where
\(\mathbf x=(x_1,\ldots,x_n)^\top\) denotes ordinary production inputs.
The firm is assumed to follow the marginal-pricing principle,
represented in GEMB by the first-order stationarity conditions with
\texttt{behavior=:stationary}.

Introduce an additional commodity, called a scale claim, with quantity
\(c>0\), and define the degree-one homogenized production function

\[
g(\mathbf x,c)
=
c f\left(\frac{\mathbf x}{c}\right)
\]

where \(\mathbf x/c=(x_1/c,\ldots,x_n/c)^\top\). The function
\(g(\mathbf x,c)\) is homogeneous of degree one in \((\mathbf x,c)\).
Suppose the aggregate supply of the scale claim is normalized to one.
Market clearing then implies \(c=1\).

Let \(p_y\) denote the output price,
\(\mathbf p_x=(p_{x1},\ldots,p_{xn})^\top\) the prices of ordinary
inputs, and \(p_c\) the price of the scale claim. Under marginal
pricing, the stationarity conditions at \(c=1\) imply

\[
\mathbf p_x
=
p_y\nabla f(\mathbf x)
\]

and

\[
p_c
=
p_y
\left(
f(\mathbf x)
-
\mathbf x^\top\nabla f(\mathbf x)
\right)
\]

Since \(g\) is homogeneous of degree one, the firm's revenue-expenditure
balance is

\[
p_yy
=
\mathbf p_x^\top\mathbf x+p_cc
\]

Define the scale claim rate as the ratio of the value of the scale claim
to the value of ordinary inputs. Then

\[
\tau(\mathbf x)
=
\frac{f(\mathbf x)}
{\mathbf x^\top\nabla f(\mathbf x)}
-1
\]

Let the local scale elasticity of production be

\[
\varepsilon_f(\mathbf x)
=
\frac{\mathbf x^\top\nabla f(\mathbf x)}
{f(\mathbf x)}
\]

The scale claim rate can therefore be written as

\[
\tau(\mathbf x)
=
\frac{1}{\varepsilon_f(\mathbf x)}-1
\]

Thus, locally decreasing returns to scale imply a positive scale claim
rate, constant returns to scale imply a zero rate, and locally
increasing returns to scale imply a negative rate.

\subsection{Equivalence under Degree-One
Homogenization}\label{equivalence-under-degree-one-homogenization}

Degree-one homogenization expands the commodity space by introducing the
scale claim while preserving the original production technology when
\(c=1\). In particular,

\[
g(\mathbf x,1)=f(\mathbf x)
\]

and

\[
\nabla_{\mathbf x}g(\mathbf x,1)
=
\nabla f(\mathbf x)
\]

Therefore, when the scale-claim market clears at \(c=1\), the
technological relation between the original inputs and output is
unchanged, and the marginal-pricing conditions for the ordinary inputs
are identical to those in the original model. The additional
stationarity condition associated with the scale claim implies

\[
p_c
=
p_yf(\mathbf x)
-
\mathbf p_x^\top\mathbf x
\]

Thus, the difference between output value and the value of ordinary
inputs under marginal pricing is represented as the value of the scale
claim. The scale-claim price \(p_c\) is an endogenous equilibrium price
determined jointly with the other prices and quantities.

Hence, on the production side, degree-one homogenization provides an
equivalent representation of the original technology and its
marginal-pricing conditions in an expanded commodity space. Equivalence
of the complete equilibrium system additionally requires consistency in
the treatment of the scale claim. Its aggregate supply must be
normalized so that market clearing implies \(c=1\), and its endowment
must be allocated so that the resulting scale-claim value corresponds to
the distribution of the residual between output value and ordinary-input
value in the original model. All other endowments, technologies,
preferences, and market-clearing conditions remain unchanged.

\subsection{Power and Exponential Production
Functions}\label{power-and-exponential-production-functions}

For the power production function

\[
f(x)=\alpha x^\theta
\]

degree-one homogenization gives

\[
g(x,c)
=
c f\left(\frac{x}{c}\right)
=
\alpha x^\theta c^{1-\theta}
\]

For the exponential branch of the production function,

\[
f(x)=\alpha\theta^x,\qquad x\geq\xi>0
\]

degree-one homogenization gives

\[
g(x,c)
=
c f\left(\frac{x}{c}\right)
=
\alpha c\theta^{x/c},
\qquad
\frac{x}{c}\geq\xi>0,\quad c>0
\]

In GEMB, the scale claim commodity can be included as an additional
production input in a \texttt{ProductionFunctionSpec}, with its
aggregate supply normalized to one and the firm specified with
\texttt{behavior=:stationary}. When the marginal products contain a
common strictly positive factor, this factor may be omitted to obtain a
scaled marginal-product function. Such scaling changes only the
production multiplier and leaves equilibrium quantities and prices
unchanged.

The power production case is implemented in
Listing~\ref{lst-gemb_powerProduction_degreeOneHomogenization}.

\begin{codelisting}

\caption{\label{lst-gemb_powerProduction_degreeOneHomogenization}Degree-One
Homogenization of a Power Production Function}

\centering{

\begin{Shaded}
\begin{Highlighting}[]
\ImportTok{using} \BuiltInTok{GeneralEquilibriumModeling}

\NormalTok{decreasing\_returns }\OperatorTok{=} \ConstantTok{false}  \CommentTok{\# Set to true for the alternative numerical example.}

\NormalTok{beta }\OperatorTok{=} \FloatTok{0.5}
\NormalTok{omega }\OperatorTok{=} \FloatTok{3.0}

\NormalTok{alpha }\OperatorTok{=}\NormalTok{ decreasing\_returns ? }\FloatTok{2.0} \OperatorTok{:} \FloatTok{0.25}
\NormalTok{theta }\OperatorTok{=}\NormalTok{ decreasing\_returns ? }\FloatTok{0.5} \OperatorTok{:} \FloatTok{2.0}
\NormalTok{tau }\OperatorTok{=} \FloatTok{1.0} \OperatorTok{/}\NormalTok{ theta }\OperatorTok{{-}} \FloatTok{1.0}

\NormalTok{production\_function }\OperatorTok{=}\NormalTok{ inputs }\OperatorTok{{-}\textgreater{}} \ControlFlowTok{begin}
\NormalTok{    x, c }\OperatorTok{=}\NormalTok{ inputs}
\NormalTok{    alpha }\OperatorTok{*}\NormalTok{ x}\OperatorTok{\^{}}\NormalTok{theta }\OperatorTok{*}\NormalTok{ c}\OperatorTok{\^{}}\NormalTok{(}\FloatTok{1.0} \OperatorTok{{-}}\NormalTok{ theta)}
\ControlFlowTok{end}

\NormalTok{marginal\_product\_function }\OperatorTok{=}\NormalTok{ inputs }\OperatorTok{{-}\textgreater{}} \ControlFlowTok{begin}
\NormalTok{    x, c }\OperatorTok{=}\NormalTok{ inputs}

\NormalTok{    [}
\NormalTok{        alpha }\OperatorTok{*}\NormalTok{ theta }\OperatorTok{*}\NormalTok{ x}\OperatorTok{\^{}}\NormalTok{(theta }\OperatorTok{{-}} \FloatTok{1.0}\NormalTok{) }\OperatorTok{*}\NormalTok{ c}\OperatorTok{\^{}}\NormalTok{(}\FloatTok{1.0} \OperatorTok{{-}}\NormalTok{ theta),}
\NormalTok{        alpha }\OperatorTok{*}\NormalTok{ (}\FloatTok{1.0} \OperatorTok{{-}}\NormalTok{ theta) }\OperatorTok{*}\NormalTok{ x}\OperatorTok{\^{}}\NormalTok{theta }\OperatorTok{*}\NormalTok{ c}\OperatorTok{\^{}}\NormalTok{(}\OperatorTok{{-}}\NormalTok{theta),}
\NormalTok{    ]}
\ControlFlowTok{end}

\NormalTok{model }\OperatorTok{=}\NormalTok{ GEMB.}\FunctionTok{GEMBModel}\NormalTok{(}
\NormalTok{    [}
        \OperatorTok{:}\NormalTok{product,}
\NormalTok{        GEMB.}\FunctionTok{CommoditySpec}\NormalTok{(}\OperatorTok{:}\NormalTok{labor; price\_lower\_bound}\OperatorTok{=}\FloatTok{1.0e{-}10}\NormalTok{),}
\NormalTok{        GEMB.}\FunctionTok{CommoditySpec}\NormalTok{(}
            \OperatorTok{:}\NormalTok{scale\_claim;}
\NormalTok{            price\_lower\_bound}\OperatorTok{={-}}\ConstantTok{Inf}\NormalTok{,}
\NormalTok{            price\_upper\_bound}\OperatorTok{=}\ConstantTok{Inf}\NormalTok{,}
\NormalTok{        ),}
\NormalTok{    ];}
\NormalTok{    numeraire}\OperatorTok{=:}\NormalTok{product,}
\NormalTok{    numeraire\_value}\OperatorTok{=}\FloatTok{1.0}\NormalTok{,}
\NormalTok{)}

\NormalTok{GEMB.}\FunctionTok{add\_agent!}\NormalTok{(}
\NormalTok{    model,}
\NormalTok{    GEMB.}\FunctionTok{ProductionFunctionSpec}\NormalTok{(}
\NormalTok{        production\_function,}
\NormalTok{        marginal\_product\_function;}
\NormalTok{        behavior}\OperatorTok{=:}\NormalTok{stationary,}
\NormalTok{    );}
\NormalTok{    outputs}\OperatorTok{=:}\NormalTok{product,}
\NormalTok{    demands}\OperatorTok{=}\NormalTok{[}\OperatorTok{:}\NormalTok{labor, }\OperatorTok{:}\NormalTok{scale\_claim],}
\NormalTok{    activity\_start}\OperatorTok{=}\NormalTok{decreasing\_returns ? }\FloatTok{2.0} \OperatorTok{:} \FloatTok{1.0}\NormalTok{,}
\NormalTok{    demand\_start}\OperatorTok{=}\NormalTok{[decreasing\_returns ? }\FloatTok{1.0} \OperatorTok{:} \FloatTok{2.0}\NormalTok{, }\FloatTok{1.0}\NormalTok{],}
\NormalTok{    demand\_lower\_bounds}\OperatorTok{=}\NormalTok{[}\FloatTok{1.0e{-}10}\NormalTok{, }\FloatTok{1.0e{-}10}\NormalTok{],}
\NormalTok{    production\_multiplier\_start}\OperatorTok{=}\FloatTok{1.0}\NormalTok{,}
\NormalTok{    name}\OperatorTok{=:}\NormalTok{firm,}
\NormalTok{)}

\NormalTok{GEMB.}\FunctionTok{add\_agent!}\NormalTok{(}
\NormalTok{    model,}
\NormalTok{    GEMB.}\FunctionTok{CESSpec}\NormalTok{([beta, }\FloatTok{1.0} \OperatorTok{{-}}\NormalTok{ beta]; es}\OperatorTok{=}\FloatTok{1.0}\NormalTok{, alpha}\OperatorTok{=}\FloatTok{1.0}\NormalTok{);}
\NormalTok{    demands}\OperatorTok{=}\NormalTok{[}\OperatorTok{:}\NormalTok{product, }\OperatorTok{:}\NormalTok{labor],}
\NormalTok{    endowments}\OperatorTok{=}\NormalTok{[}\OperatorTok{:}\NormalTok{labor, }\OperatorTok{:}\NormalTok{scale\_claim],}
\NormalTok{    endowment\_quantities}\OperatorTok{=}\NormalTok{[omega, }\FloatTok{1.0}\NormalTok{],}
\NormalTok{    name}\OperatorTok{=:}\NormalTok{consumer,}
\NormalTok{)}

\NormalTok{p0 }\OperatorTok{=}\NormalTok{ decreasing\_returns ?}
\NormalTok{    [}\FloatTok{1.0}\NormalTok{, }\FloatTok{1.0}\NormalTok{, }\FloatTok{1.0}\NormalTok{] }\OperatorTok{:}
\NormalTok{    [}\FloatTok{1.0}\NormalTok{, }\FloatTok{1.0}\NormalTok{, }\OperatorTok{{-}}\FloatTok{1.0}\NormalTok{]}

\NormalTok{result }\OperatorTok{=}\NormalTok{ GEMB.}\FunctionTok{solve}\NormalTok{(}
\NormalTok{    model;}
\NormalTok{    p0}\OperatorTok{=}\NormalTok{p0,}
\NormalTok{    residual\_tol}\OperatorTok{=}\FloatTok{1.0e{-}8}\NormalTok{,}
\NormalTok{    silent}\OperatorTok{=}\ConstantTok{true}\NormalTok{,}
\NormalTok{)}

\NormalTok{GEMB.}\FunctionTok{print\_equilibrium\_statistics}\NormalTok{(model, result)}

\NormalTok{\_, labor\_input, scale\_claim\_input, \_ }\OperatorTok{=}
\NormalTok{    result.agent\_variable\_values[}\FloatTok{1}\NormalTok{]}

\NormalTok{implied\_claim\_rate }\OperatorTok{=}
\NormalTok{    result.prices[}\FloatTok{3}\NormalTok{] }\OperatorTok{*}\NormalTok{ scale\_claim\_input }\OperatorTok{/}
\NormalTok{    (result.prices[}\FloatTok{2}\NormalTok{] }\OperatorTok{*}\NormalTok{ labor\_input)}

\FunctionTok{println}\NormalTok{(}\StringTok{"}\SpecialCharTok{\textbackslash{}n}\StringTok{Scale{-}claim rate:"}\NormalTok{)}
\FunctionTok{println}\NormalTok{(}\StringTok{"theoretical = "}\NormalTok{, tau)}
\FunctionTok{println}\NormalTok{(}\StringTok{"implied = "}\NormalTok{, implied\_claim\_rate)}
\end{Highlighting}
\end{Shaded}

}

\end{codelisting}%

The exponential production case illustrates a production function whose
scale elasticity, and hence scale claim rate, varies with the input
level. The corresponding implementation is given in
Listing~\ref{lst-gemb_exponentialProduction_degreeOneHomogenization}.

\begin{codelisting}

\caption{\label{lst-gemb_exponentialProduction_degreeOneHomogenization}Degree-One
Homogenization of an Exponential Production Function}

\centering{

\begin{Shaded}
\begin{Highlighting}[]
\ImportTok{using} \BuiltInTok{GeneralEquilibriumModeling}

\NormalTok{decreasing\_returns }\OperatorTok{=} \ConstantTok{false}  \CommentTok{\# Set to true for the alternative numerical example.}

\NormalTok{beta }\OperatorTok{=} \FloatTok{0.5}
\NormalTok{omega }\OperatorTok{=} \FloatTok{3.0}

\NormalTok{theta }\OperatorTok{=}\NormalTok{ decreasing\_returns ? }\FunctionTok{exp}\NormalTok{(}\FloatTok{0.5}\NormalTok{) }\OperatorTok{:} \FunctionTok{exp}\NormalTok{(}\FloatTok{1.0}\NormalTok{)}
\NormalTok{alpha }\OperatorTok{=}\NormalTok{ decreasing\_returns ? }\FloatTok{2.0} \OperatorTok{/}\NormalTok{ theta }\OperatorTok{:} \FloatTok{1.0} \OperatorTok{/}\NormalTok{ theta}\OperatorTok{\^{}}\FloatTok{2}

\NormalTok{output\_start }\OperatorTok{=} \FloatTok{1.5}
\NormalTok{labor\_start }\OperatorTok{=} \FloatTok{1.5}
\NormalTok{scale\_claim\_start }\OperatorTok{=} \FloatTok{1.2}

\NormalTok{p0 }\OperatorTok{=}\NormalTok{ decreasing\_returns ?}
\NormalTok{    [}\FloatTok{1.0}\NormalTok{, }\FloatTok{0.8}\NormalTok{, }\FloatTok{0.5}\NormalTok{] }\OperatorTok{:}
\NormalTok{    [}\FloatTok{1.0}\NormalTok{, }\FloatTok{0.8}\NormalTok{, }\OperatorTok{{-}}\FloatTok{0.5}\NormalTok{]}

\CommentTok{\# Degree{-}one homogenization:}
\CommentTok{\# g(x, c) = c * f(x / c) = alpha * c * theta\^{}(x / c)}
\NormalTok{production\_function }\OperatorTok{=}\NormalTok{ inputs }\OperatorTok{{-}\textgreater{}} \ControlFlowTok{begin}
\NormalTok{    x, c }\OperatorTok{=}\NormalTok{ inputs}
\NormalTok{    alpha }\OperatorTok{*}\NormalTok{ c }\OperatorTok{*}\NormalTok{ theta}\OperatorTok{\^{}}\NormalTok{(x }\OperatorTok{/}\NormalTok{ c)}
\ControlFlowTok{end}

\CommentTok{\# Scaled marginal product function:}
\CommentTok{\# the common strictly positive factor alpha * theta\^{}(x / c) is omitted.}
\NormalTok{marginal\_product\_function }\OperatorTok{=}\NormalTok{ inputs }\OperatorTok{{-}\textgreater{}} \ControlFlowTok{begin}
\NormalTok{    x, c }\OperatorTok{=}\NormalTok{ inputs}

\NormalTok{    [}
        \FunctionTok{log}\NormalTok{(theta),}
        \FloatTok{1.0} \OperatorTok{{-}}\NormalTok{ (x }\OperatorTok{/}\NormalTok{ c) }\OperatorTok{*} \FunctionTok{log}\NormalTok{(theta),}
\NormalTok{    ]}
\ControlFlowTok{end}

\NormalTok{model }\OperatorTok{=}\NormalTok{ GEMB.}\FunctionTok{GEMBModel}\NormalTok{(}
\NormalTok{    [}
        \OperatorTok{:}\NormalTok{product,}
\NormalTok{        GEMB.}\FunctionTok{CommoditySpec}\NormalTok{(}
            \OperatorTok{:}\NormalTok{labor;}
\NormalTok{            price\_lower\_bound}\OperatorTok{=}\FloatTok{1.0e{-}10}\NormalTok{,}
\NormalTok{        ),}
\NormalTok{        GEMB.}\FunctionTok{CommoditySpec}\NormalTok{(}
            \OperatorTok{:}\NormalTok{scale\_claim;}
\NormalTok{            price\_lower\_bound}\OperatorTok{={-}}\ConstantTok{Inf}\NormalTok{,}
\NormalTok{            price\_upper\_bound}\OperatorTok{=}\ConstantTok{Inf}\NormalTok{,}
\NormalTok{        ),}
\NormalTok{    ];}
\NormalTok{    numeraire}\OperatorTok{=:}\NormalTok{product,}
\NormalTok{    numeraire\_value}\OperatorTok{=}\FloatTok{1.0}\NormalTok{,}
\NormalTok{)}

\NormalTok{GEMB.}\FunctionTok{add\_agent!}\NormalTok{(}
\NormalTok{    model,}
\NormalTok{    GEMB.}\FunctionTok{ProductionFunctionSpec}\NormalTok{(}
\NormalTok{        production\_function,}
\NormalTok{        marginal\_product\_function;}
\NormalTok{        behavior}\OperatorTok{=:}\NormalTok{stationary,}
\NormalTok{    );}
\NormalTok{    outputs}\OperatorTok{=:}\NormalTok{product,}
\NormalTok{    demands}\OperatorTok{=}\NormalTok{[}\OperatorTok{:}\NormalTok{labor, }\OperatorTok{:}\NormalTok{scale\_claim],}
\NormalTok{    activity\_start}\OperatorTok{=}\NormalTok{output\_start,}
\NormalTok{    demand\_start}\OperatorTok{=}\NormalTok{[labor\_start, scale\_claim\_start],}
\NormalTok{    demand\_lower\_bounds}\OperatorTok{=}\NormalTok{[}\FloatTok{1.0e{-}10}\NormalTok{, }\FloatTok{1.0e{-}10}\NormalTok{],}
\NormalTok{    production\_multiplier\_start}\OperatorTok{=}\FloatTok{1.0}\NormalTok{,}
\NormalTok{    name}\OperatorTok{=:}\NormalTok{firm,}
\NormalTok{)}

\NormalTok{GEMB.}\FunctionTok{add\_agent!}\NormalTok{(}
\NormalTok{    model,}
\NormalTok{    GEMB.}\FunctionTok{CESSpec}\NormalTok{(}
\NormalTok{        [beta, }\FloatTok{1.0} \OperatorTok{{-}}\NormalTok{ beta];}
\NormalTok{        es}\OperatorTok{=}\FloatTok{1.0}\NormalTok{,}
\NormalTok{        alpha}\OperatorTok{=}\FloatTok{1.0}\NormalTok{,}
\NormalTok{    );}
\NormalTok{    demands}\OperatorTok{=}\NormalTok{[}\OperatorTok{:}\NormalTok{product, }\OperatorTok{:}\NormalTok{labor],}
\NormalTok{    endowments}\OperatorTok{=}\NormalTok{[}\OperatorTok{:}\NormalTok{labor, }\OperatorTok{:}\NormalTok{scale\_claim],}
\NormalTok{    endowment\_quantities}\OperatorTok{=}\NormalTok{[omega, }\FloatTok{1.0}\NormalTok{],}
\NormalTok{    name}\OperatorTok{=:}\NormalTok{consumer,}
\NormalTok{)}

\NormalTok{result }\OperatorTok{=}\NormalTok{ GEMB.}\FunctionTok{solve}\NormalTok{(}
\NormalTok{    model;}
\NormalTok{    p0}\OperatorTok{=}\NormalTok{p0,}
\NormalTok{    residual\_tol}\OperatorTok{=}\FloatTok{1.0e{-}8}\NormalTok{,}
\NormalTok{    silent}\OperatorTok{=}\ConstantTok{true}\NormalTok{,}
\NormalTok{)}

\NormalTok{GEMB.}\FunctionTok{print\_equilibrium\_statistics}\NormalTok{(model, result)}

\NormalTok{\_, labor\_input, scale\_claim\_input, \_ }\OperatorTok{=}
\NormalTok{    result.agent\_variable\_values[}\FloatTok{1}\NormalTok{]}

\NormalTok{implied\_claim\_rate }\OperatorTok{=}
\NormalTok{    result.prices[}\FloatTok{3}\NormalTok{] }\OperatorTok{*}\NormalTok{ scale\_claim\_input }\OperatorTok{/}
\NormalTok{    (result.prices[}\FloatTok{2}\NormalTok{] }\OperatorTok{*}\NormalTok{ labor\_input)}

\NormalTok{theoretical\_claim\_rate }\OperatorTok{=}
    \FloatTok{1.0} \OperatorTok{/}\NormalTok{ (labor\_input }\OperatorTok{*} \FunctionTok{log}\NormalTok{(theta)) }\OperatorTok{{-}} \FloatTok{1.0}

\FunctionTok{println}\NormalTok{(}\StringTok{"}\SpecialCharTok{\textbackslash{}n}\StringTok{Scale{-}claim check:"}\NormalTok{)}
\FunctionTok{println}\NormalTok{(}\StringTok{"production elasticity = "}\NormalTok{, labor\_input }\OperatorTok{*} \FunctionTok{log}\NormalTok{(theta))}
\FunctionTok{println}\NormalTok{(}\StringTok{"theoretical claim rate = "}\NormalTok{, theoretical\_claim\_rate)}
\FunctionTok{println}\NormalTok{(}\StringTok{"implied claim rate = "}\NormalTok{, implied\_claim\_rate)}
\end{Highlighting}
\end{Shaded}

}

\end{codelisting}%

\subsection{\texorpdfstring{Square-Sum Production Function and Marginal
Pricing Equilibria: A \(4\times2\)
Model}{Square-Sum Production Function and Marginal Pricing Equilibria: A 4\textbackslash times2 Model}}\label{square-sum-production-function-and-marginal-pricing-equilibria-a-4times2-model}

Assume the following:

\begin{enumerate}
\def\labelenumi{\arabic{enumi}.}
\item
  The economy contains four commodities: a product, labor, land, and a
  scale claim. There are two economic agents, a firm and a consumer. The
  product is taken as the numeraire commodity.
\item
  The firm's original production function is \(f(x_2,x_3)=x_2^2+x_3^2\),
  where \(x_2\) and \(x_3\) denote labor and land inputs, respectively.
  Introduce the scale claim input \(x_4>0\) and degree-one homogenize
  the production function as
\end{enumerate}

\[
g(x_2,x_3,x_4)
=
x_4 f\left(\frac{x_2}{x_4},\frac{x_3}{x_4}\right)
=
\frac{x_2^2+x_3^2}{x_4}
\]

The aggregate supply of the scale claim is normalized to one. The firm
follows the marginal-pricing principle, represented in GEMB by the
stationarity conditions with \texttt{behavior=:stationary}.

\begin{enumerate}
\def\labelenumi{\arabic{enumi}.}
\setcounter{enumi}{2}
\tightlist
\item
  The consumer has utility function \(u=c_1^{0.5}c_2^{0.25}c_3^{0.25}\),
  where \(c_1\), \(c_2\), and \(c_3\) denote consumption of the product,
  labor, and land, respectively. The consumer is endowed with 100 units
  of labor, 100 units of land, and one unit of the scale claim.
\end{enumerate}

The Pareto-optimal allocation of the original commodities corresponds to
a maximizer of

\[
F(x_2,x_3)
=
(x_2^2+x_3^2)^{0.5}
(100-x_2)^{0.25}
(100-x_3)^{0.25},
\qquad
0\leq x_2,x_3\leq100
\]

This function has three stationary points. Checking the boundary shows
that the global maximum is attained at

\[
x_2=x_3=\frac{200}{3}
\]

while the other two stationary points are saddle points.

Let the price vector be \(\mathbf p=(1,p_2,p_3,p_4)^\top\). The marginal
products of the degree-one homogenized production function are

\[
g_{x_2}
=
\frac{2x_2}{x_4},
\qquad
g_{x_3}
=
\frac{2x_3}{x_4},
\qquad
g_{x_4}
=
-\frac{x_2^2+x_3^2}{x_4^2}
\]

Since the aggregate supply of the scale claim equals one, market
clearing implies \(x_4=1\). The firm's stationarity conditions therefore
give

\[
p_2=2x_2,
\qquad
p_3=2x_3,
\qquad
p_4=-(x_2^2+x_3^2)
\]

and hence

\[
x_2=\frac{p_2}{2},
\qquad
x_3=\frac{p_3}{2},
\qquad
y=\frac{p_2^2+p_3^2}{4},
\qquad
p_4=-\frac{p_2^2+p_3^2}{4}
\]

Let the consumer's income be

\[
w=100p_2+100p_3+p_4
\]

The demand and supply matrices are

\[
\mathbf D(\mathbf p,w)
=
\begin{bmatrix}
0 & \frac{w}{2}\\
\frac{p_2}{2} & \frac{w}{4p_2}\\
\frac{p_3}{2} & \frac{w}{4p_3}\\
1 & 0
\end{bmatrix}
\]

and

\[
\mathbf S(\mathbf p)
=
\begin{bmatrix}
\frac{p_2^2+p_3^2}{4} & 0\\
0 & 100\\
0 & 100\\
0 & 1
\end{bmatrix}
\]

Solving the corresponding equality-form wide-sense structural
equilibrium model yields three equilibria, corresponding to the three
stationary points of \(F(x_2,x_3)\).

\begin{enumerate}
\def\labelenumi{\arabic{enumi}.}
\tightlist
\item
  The first equilibrium, corresponding to the Pareto-optimal allocation,
  is
\end{enumerate}

\[
\mathbf p^*
=
\left(
1,
\frac{400}{3},
\frac{400}{3},
-\frac{80000}{9}
\right)^\top,
\qquad
w^*
=
\frac{160000}{9}
\]

with

\[
x_2^*=x_3^*=\frac{200}{3}
\]

\begin{enumerate}
\def\labelenumi{\arabic{enumi}.}
\setcounter{enumi}{1}
\tightlist
\item
  The second equilibrium is
\end{enumerate}

\[
\mathbf p^*
=
\left(
1,
100-\frac{100\sqrt3}{3},
100+\frac{100\sqrt3}{3},
-\frac{20000}{3}
\right)^\top,
\qquad
w^*
=
\frac{40000}{3}
\]

with

\[
(x_2^*,x_3^*)
=
\left(
50-\frac{50\sqrt3}{3},
50+\frac{50\sqrt3}{3}
\right)
\]

\begin{enumerate}
\def\labelenumi{\arabic{enumi}.}
\setcounter{enumi}{2}
\tightlist
\item
  The third equilibrium is
\end{enumerate}

\[
\mathbf p^*
=
\left(
1,
100+\frac{100\sqrt3}{3},
100-\frac{100\sqrt3}{3},
-\frac{20000}{3}
\right)^\top,
\qquad
w^*
=
\frac{40000}{3}
\]

with

\[
(x_2^*,x_3^*)
=
\left(
50+\frac{50\sqrt3}{3},
50-\frac{50\sqrt3}{3}
\right)
\]

Hence degree-one homogenization yields three marginal-pricing equilibria
corresponding to the three stationary points of the original allocation
problem. Only the first corresponds to the Pareto-optimal allocation,
while the other two correspond to saddle points.

The three equilibria can be obtained numerically from different starting
values. The corresponding GEMB implementation is given in
Listing~\ref{lst-gemb_squareSumProduction_degreeOneHomogenization}.

\begin{codelisting}

\caption{\label{lst-gemb_squareSumProduction_degreeOneHomogenization}Degree-One
Homogenization of a Square-Sum Production Function with Multiple
Equilibria}

\centering{

\begin{Shaded}
\begin{Highlighting}[]
\ImportTok{using} \BuiltInTok{GeneralEquilibriumModeling}

\NormalTok{beta }\OperatorTok{=}\NormalTok{ [}\FloatTok{0.5}\NormalTok{, }\FloatTok{0.25}\NormalTok{, }\FloatTok{0.25}\NormalTok{]}
\NormalTok{omega }\OperatorTok{=} \FloatTok{100.0}

\CommentTok{\# Degree{-}one homogenization:}
\CommentTok{\# g(x1, x2, x\_4) = (x1\^{}2 + x2\^{}2) / x\_4.}
\NormalTok{production\_function }\OperatorTok{=}\NormalTok{ inputs }\OperatorTok{{-}\textgreater{}} \ControlFlowTok{begin}
\NormalTok{    x1, x2, x\_4 }\OperatorTok{=}\NormalTok{ inputs}
\NormalTok{    (x1}\OperatorTok{\^{}}\FloatTok{2} \OperatorTok{+}\NormalTok{ x2}\OperatorTok{\^{}}\FloatTok{2}\NormalTok{) }\OperatorTok{/}\NormalTok{ x\_4}
\ControlFlowTok{end}

\CommentTok{\# Scaled marginal product function:}
\CommentTok{\# the common strictly positive factor 1 / x\_4 is omitted.}
\NormalTok{marginal\_product\_function }\OperatorTok{=}\NormalTok{ inputs }\OperatorTok{{-}\textgreater{}} \ControlFlowTok{begin}
\NormalTok{    x1, x2, x\_4 }\OperatorTok{=}\NormalTok{ inputs}
\NormalTok{    [}\FloatTok{2.0} \OperatorTok{*}\NormalTok{ x1, }\FloatTok{2.0} \OperatorTok{*}\NormalTok{ x2, }\FunctionTok{{-}}\NormalTok{(x1}\OperatorTok{\^{}}\FloatTok{2} \OperatorTok{+}\NormalTok{ x2}\OperatorTok{\^{}}\FloatTok{2}\NormalTok{) }\OperatorTok{/}\NormalTok{ x\_4]}
\ControlFlowTok{end}

\NormalTok{cases }\OperatorTok{=}\NormalTok{ [}
\NormalTok{    ([}\FloatTok{60.0}\NormalTok{, }\FloatTok{60.0}\NormalTok{], [}\FloatTok{1.0}\NormalTok{, }\FloatTok{100.0}\NormalTok{, }\FloatTok{100.0}\NormalTok{, }\OperatorTok{{-}}\FloatTok{5000.0}\NormalTok{]),}
\NormalTok{    ([}\FloatTok{30.0}\NormalTok{, }\FloatTok{70.0}\NormalTok{], [}\FloatTok{1.0}\NormalTok{, }\FloatTok{60.0}\NormalTok{, }\FloatTok{140.0}\NormalTok{, }\OperatorTok{{-}}\FloatTok{5000.0}\NormalTok{]),}
\NormalTok{    ([}\FloatTok{70.0}\NormalTok{, }\FloatTok{30.0}\NormalTok{], [}\FloatTok{1.0}\NormalTok{, }\FloatTok{140.0}\NormalTok{, }\FloatTok{60.0}\NormalTok{, }\OperatorTok{{-}}\FloatTok{5000.0}\NormalTok{]),}
\NormalTok{]}

\KeywordTok{function} \FunctionTok{solve\_case}\NormalTok{(case\_id, input\_start, p0)}
\NormalTok{    x1\_start, x2\_start }\OperatorTok{=}\NormalTok{ input\_start}

\NormalTok{    model }\OperatorTok{=}\NormalTok{ GEMB.}\FunctionTok{GEMBModel}\NormalTok{(}
\NormalTok{        [}
            \OperatorTok{:}\NormalTok{product,}
\NormalTok{            GEMB.}\FunctionTok{CommoditySpec}\NormalTok{(}\OperatorTok{:}\NormalTok{labor; price\_lower\_bound}\OperatorTok{=}\FloatTok{1.0e{-}10}\NormalTok{),}
\NormalTok{            GEMB.}\FunctionTok{CommoditySpec}\NormalTok{(}\OperatorTok{:}\NormalTok{land; price\_lower\_bound}\OperatorTok{=}\FloatTok{1.0e{-}10}\NormalTok{),}
\NormalTok{            GEMB.}\FunctionTok{CommoditySpec}\NormalTok{(}
                \OperatorTok{:}\NormalTok{scale\_claim;}
\NormalTok{                price\_lower\_bound}\OperatorTok{={-}}\ConstantTok{Inf}\NormalTok{,}
\NormalTok{                price\_upper\_bound}\OperatorTok{=}\ConstantTok{Inf}\NormalTok{,}
\NormalTok{            ),}
\NormalTok{        ];}
\NormalTok{        numeraire}\OperatorTok{=:}\NormalTok{product,}
\NormalTok{        numeraire\_value}\OperatorTok{=}\FloatTok{1.0}\NormalTok{,}
\NormalTok{    )}

\NormalTok{    GEMB.}\FunctionTok{add\_agent!}\NormalTok{(}
\NormalTok{        model,}
\NormalTok{        GEMB.}\FunctionTok{ProductionFunctionSpec}\NormalTok{(}
\NormalTok{            production\_function,}
\NormalTok{            marginal\_product\_function;}
\NormalTok{            behavior}\OperatorTok{=:}\NormalTok{stationary,}
\NormalTok{        );}
\NormalTok{        outputs}\OperatorTok{=:}\NormalTok{product,}
\NormalTok{        demands}\OperatorTok{=}\NormalTok{[}\OperatorTok{:}\NormalTok{labor, }\OperatorTok{:}\NormalTok{land, }\OperatorTok{:}\NormalTok{scale\_claim],}
\NormalTok{        activity\_start}\OperatorTok{=}\NormalTok{x1\_start}\OperatorTok{\^{}}\FloatTok{2} \OperatorTok{+}\NormalTok{ x2\_start}\OperatorTok{\^{}}\FloatTok{2}\NormalTok{,}
\NormalTok{        demand\_start}\OperatorTok{=}\NormalTok{[x1\_start, x2\_start, }\FloatTok{1.0}\NormalTok{],}
\NormalTok{        demand\_lower\_bounds}\OperatorTok{=}\NormalTok{[}\FloatTok{1.0e{-}10}\NormalTok{, }\FloatTok{1.0e{-}10}\NormalTok{, }\FloatTok{1.0e{-}10}\NormalTok{],}
\NormalTok{        production\_multiplier\_start}\OperatorTok{=}\FloatTok{1.0}\NormalTok{,}
\NormalTok{        name}\OperatorTok{=:}\NormalTok{firm,}
\NormalTok{    )}

\NormalTok{    GEMB.}\FunctionTok{add\_agent!}\NormalTok{(}
\NormalTok{        model,}
\NormalTok{        GEMB.}\FunctionTok{CESSpec}\NormalTok{(beta; es}\OperatorTok{=}\FloatTok{1.0}\NormalTok{, alpha}\OperatorTok{=}\FloatTok{1.0}\NormalTok{);}
\NormalTok{        demands}\OperatorTok{=}\NormalTok{[}\OperatorTok{:}\NormalTok{product, }\OperatorTok{:}\NormalTok{labor, }\OperatorTok{:}\NormalTok{land],}
\NormalTok{        endowments}\OperatorTok{=}\NormalTok{[}\OperatorTok{:}\NormalTok{labor, }\OperatorTok{:}\NormalTok{land, }\OperatorTok{:}\NormalTok{scale\_claim],}
\NormalTok{        endowment\_quantities}\OperatorTok{=}\NormalTok{[omega, omega, }\FloatTok{1.0}\NormalTok{],}
\NormalTok{        name}\OperatorTok{=:}\NormalTok{consumer,}
\NormalTok{    )}

\NormalTok{    result }\OperatorTok{=}\NormalTok{ GEMB.}\FunctionTok{solve}\NormalTok{(}
\NormalTok{        model;}
\NormalTok{        p0}\OperatorTok{=}\NormalTok{p0,}
\NormalTok{        residual\_tol}\OperatorTok{=}\FloatTok{1.0e{-}8}\NormalTok{,}
\NormalTok{        silent}\OperatorTok{=}\ConstantTok{true}\NormalTok{,}
\NormalTok{    )}

    \CommentTok{\# Firm variables: output, labor, land, scale claim, multiplier.}
\NormalTok{    \_, x1, x2, x\_4, \_ }\OperatorTok{=}\NormalTok{ result.agent\_variable\_values[}\FloatTok{1}\NormalTok{] }
\NormalTok{    p }\OperatorTok{=}\NormalTok{ result.prices}

\NormalTok{    claim\_rate }\OperatorTok{=}
\NormalTok{        p[}\FloatTok{4}\NormalTok{] }\OperatorTok{*}\NormalTok{ x\_4 }\OperatorTok{/}
\NormalTok{        (p[}\FloatTok{2}\NormalTok{] }\OperatorTok{*}\NormalTok{ x1 }\OperatorTok{+}\NormalTok{ p[}\FloatTok{3}\NormalTok{] }\OperatorTok{*}\NormalTok{ x2)}

    \PreprocessorTok{@assert}\NormalTok{ result.solved}
    \PreprocessorTok{@assert}\NormalTok{ result.mcp\_solved}
    \PreprocessorTok{@assert}\NormalTok{ result.all\_markets\_clear}
    \PreprocessorTok{@assert} \FunctionTok{isapprox}\NormalTok{(x\_4, }\FloatTok{1.0}\NormalTok{; atol}\OperatorTok{=}\FloatTok{1.0e{-}6}\NormalTok{, rtol}\OperatorTok{=}\FloatTok{1.0e{-}6}\NormalTok{)}
    \PreprocessorTok{@assert} \FunctionTok{isapprox}\NormalTok{(claim\_rate, }\OperatorTok{{-}}\FloatTok{0.5}\NormalTok{; atol}\OperatorTok{=}\FloatTok{1.0e{-}6}\NormalTok{, rtol}\OperatorTok{=}\FloatTok{1.0e{-}6}\NormalTok{)}

    \FunctionTok{println}\NormalTok{(}\StringTok{"}\SpecialCharTok{\textbackslash{}n}\StringTok{================ Case }\SpecialCharTok{$}\NormalTok{case\_id}\StringTok{ ================"}\NormalTok{)}
    \FunctionTok{println}\NormalTok{(}\StringTok{"Starting inputs = "}\NormalTok{, input\_start)}
    \FunctionTok{println}\NormalTok{(}\StringTok{"Starting prices = "}\NormalTok{, p0)}

\NormalTok{    GEMB.}\FunctionTok{print\_equilibrium\_statistics}\NormalTok{(model, result)}

    \FunctionTok{println}\NormalTok{(}\StringTok{"}\SpecialCharTok{\textbackslash{}n}\StringTok{Production inputs:"}\NormalTok{)}
    \FunctionTok{println}\NormalTok{(}\StringTok{"labor input = "}\NormalTok{, x1)}
    \FunctionTok{println}\NormalTok{(}\StringTok{"land input = "}\NormalTok{, x2)}
    \FunctionTok{println}\NormalTok{(}\StringTok{"scale{-}claim input x\_4 = "}\NormalTok{, x\_4)}

    \FunctionTok{println}\NormalTok{(}\StringTok{"}\SpecialCharTok{\textbackslash{}n}\StringTok{Implied scale{-}claim rate = "}\NormalTok{, claim\_rate)}

    \ControlFlowTok{return}\NormalTok{ result}
\KeywordTok{end}

\NormalTok{solutions }\OperatorTok{=}\NormalTok{ [}
    \FunctionTok{solve\_case}\NormalTok{(i, case}\OperatorTok{...}\NormalTok{)}
\NormalTok{    for (i, case) }\KeywordTok{in} \FunctionTok{enumerate}\NormalTok{(cases)}
\NormalTok{]}

\ConstantTok{nothing}
\end{Highlighting}
\end{Shaded}

}

\end{codelisting}%

\bookmarksetup{startatroot}

\chapter{Intertemporal Equilibrium}\label{intertemporal-equilibrium}

\section{Intertemporal Model
Construction}\label{intertemporal-model-construction}

The GEMB high-level framework does not use a separate model type for
intertemporal equilibrium. A finite-horizon intertemporal economy is
constructed as an ordinary \texttt{GEMBModel} in which time is
represented by the \texttt{period} axis of dated commodities. Thus, the
same behavioral specifications and equilibrium solver used for static
models can also be used for intertemporal models.

Dated commodities are defined with \texttt{CommoditySpec}. For example,

\begin{Shaded}
\begin{Highlighting}[]
\NormalTok{product }\OperatorTok{=}\NormalTok{ GEMB.}\FunctionTok{CommoditySpec}\NormalTok{(}
    \OperatorTok{:}\NormalTok{product;}
\NormalTok{    axes}\OperatorTok{=}\NormalTok{(period}\OperatorTok{=}\FloatTok{1}\OperatorTok{:}\FloatTok{4}\NormalTok{,),}
\NormalTok{)}
\end{Highlighting}
\end{Shaded}

defines four dated product commodities. \texttt{CommodityRef} selects
particular dated commodities or an entire intertemporal commodity path.
For example, \texttt{CommodityRef(:product;\ period=2)} selects the
period-2 product, whereas \texttt{CommodityRef(:product)} selects the
product in all periods. A dated commodity can therefore be treated in
exactly the same way as an ordinary commodity once its period has been
specified.

A concrete economic agent is added with the ordinary
\texttt{GEMB.add\_agent!} interface. Its identity may contain a period
through \texttt{AgentRef}, while \texttt{RelativePeriod} can be used to
express commodity references relative to that period. In particular,

\begin{Shaded}
\begin{Highlighting}[]
\NormalTok{GEMB.}\FunctionTok{RelativePeriod}\NormalTok{(}\FloatTok{0}\NormalTok{)}
\NormalTok{GEMB.}\FunctionTok{RelativePeriod}\NormalTok{(}\FloatTok{1}\NormalTok{)}
\NormalTok{GEMB.}\FunctionTok{RelativePeriod}\NormalTok{(}\OperatorTok{{-}}\FloatTok{1}\NormalTok{)}
\end{Highlighting}
\end{Shaded}

denote the current, next, and previous periods, respectively. Hence a
producer operating in period \(t\) may demand
\texttt{CommodityRef(:product;\ period=RelativePeriod(0))} and supply
\texttt{CommodityRef(:product;\ period=RelativePeriod(1))}, without
explicitly writing the absolute period in each instance.

When the same economic structure is repeated over several periods,
\texttt{AgentTemplate} provides a more convenient representation. For
example,

\begin{Shaded}
\begin{Highlighting}[]
\NormalTok{producer }\OperatorTok{=}\NormalTok{ GEMB.}\FunctionTok{AgentTemplate}\NormalTok{(}
    \OperatorTok{:}\NormalTok{producer,}
\NormalTok{    spec;}
\NormalTok{    periods}\OperatorTok{=}\FloatTok{1}\OperatorTok{:}\FloatTok{3}\NormalTok{,}
\NormalTok{    demands}\OperatorTok{=}\NormalTok{[}
\NormalTok{        GEMB.}\FunctionTok{CommodityRef}\NormalTok{(}
            \OperatorTok{:}\NormalTok{labor;}
\NormalTok{            period}\OperatorTok{=}\NormalTok{GEMB.}\FunctionTok{RelativePeriod}\NormalTok{(}\FloatTok{0}\NormalTok{),}
\NormalTok{        ),}
\NormalTok{    ],}
\NormalTok{    outputs}\OperatorTok{=}\NormalTok{[}
\NormalTok{        GEMB.}\FunctionTok{CommodityRef}\NormalTok{(}
            \OperatorTok{:}\NormalTok{product;}
\NormalTok{            period}\OperatorTok{=}\NormalTok{GEMB.}\FunctionTok{RelativePeriod}\NormalTok{(}\FloatTok{1}\NormalTok{),}
\NormalTok{        ),}
\NormalTok{    ],}
\NormalTok{)}

\NormalTok{GEMB.}\FunctionTok{add\_agents!}\NormalTok{(model, producer)}
\end{Highlighting}
\end{Shaded}

creates the corresponding producers in periods 1, 2, and 3.
\texttt{AgentTemplate} therefore describes a repeated structural pattern
rather than a new type of economic behavior: \texttt{spec} remains an
ordinary GEMB behavioral specification such as \texttt{CESSpec},
\texttt{CESMarginalUtilitySpec}, \texttt{MarshallDemandConsumerSpec}, or
\texttt{ProductionFunctionSpec}.

Parameters that differ across the repeated periods can be specified with
\texttt{ByPeriod}. For example,

\begin{Shaded}
\begin{Highlighting}[]
\NormalTok{claim\_rate}\OperatorTok{=}\NormalTok{GEMB.}\FunctionTok{ByPeriod}\NormalTok{([}\FloatTok{0.10}\NormalTok{, }\FloatTok{0.15}\NormalTok{, }\FloatTok{0.20}\NormalTok{])}
\end{Highlighting}
\end{Shaded}

assigns the three values to the corresponding template periods. A scalar
value, by contrast, is used unchanged in every period.
\texttt{RelativePeriod} and \texttt{ByPeriod} thus serve different
purposes: the former changes the dated object referred to by an agent,
whereas the latter changes a parameter value across repeated agents.

Intertemporal claims, money, bonds, capital, and other dated objects do
not require special commodity types. They can be represented as ordinary
dated commodities and connected to economic agents through the usual
GEMB demand, supply, endowment, claim, and auxiliary-condition
mechanisms. Cross-period equilibrium relationships that cannot be
represented directly by standard behavioral specifications can also be
introduced through condition agents and observed equilibrium variables.

The basic construction procedure is therefore the same as for a static
GEMB model: define dated commodities with \texttt{CommoditySpec}, create
a \texttt{GEMBModel}, add concrete agents with \texttt{add\_agent!},
expand repeated structures with \texttt{AgentTemplate} and
\texttt{add\_agents!} when necessary, and finally solve the resulting
model with the ordinary GEMB equilibrium solver. The intertemporal tools
mainly provide a convenient way to organize dated commodities, repeated
agents, relative-period references, and period-varying parameters, while
leaving the underlying equilibrium framework unchanged.

\section{A Finite-Horizon Intertemporal Equilibrium
Model}\label{a-finite-horizon-intertemporal-equilibrium-model}

Consider a finite-horizon intertemporal equilibrium model in a closed
economy. Assume the following:

\begin{enumerate}
\def\labelenumi{\arabic{enumi}.}
\item
  The economy contains two types of commodities, a product and labor,
  and two types of economic agents, a firm and a consumer.
\item
  The consumer lives for \(n_p\) periods, and the economy also lasts for
  \(n_p\) periods.
\item
  The economy is closed and has no external borrowing. The initial
  supply of the product is exogenously given.
\item
  In period 1, the consumer is endowed with 150 units of the product.
  From period 1 through period \(n_p-1\), the consumer is endowed with
  100 units of labor in each period.
\item
  The consumer consumes only the product and has demand for the product
  in periods \(1,\ldots,n_p\). Intertemporal preferences are represented
  by a single-level Cobb--Douglas utility function with equal share
  coefficients across periods.
\item
  There is one firm in each period from period 1 through period
  \(n_p-1\). In period \(t\), the firm uses the product and labor as
  inputs to produce the product available in period \(t+1\). Its
  production function is
\end{enumerate}

\[
f(x,l)=2\sqrt{xl}
\]

The product used as an input is completely consumed in production.

The corresponding GEMB implementation is given in
Listing~\ref{lst-gemb_intertemporalEquilibrium_nct2_nat2}.

\begin{codelisting}

\caption{\label{lst-gemb_intertemporalEquilibrium_nct2_nat2}A
Finite-Horizon Intertemporal Equilibrium Model}

\centering{

\begin{Shaded}
\begin{Highlighting}[]
\CommentTok{\# GEMB high{-}level framework: intertemporal equilibrium.}
\ImportTok{using} \BuiltInTok{GeneralEquilibriumModeling}

\NormalTok{np }\OperatorTok{=} \FloatTok{5}

\NormalTok{initial\_product\_endowment }\OperatorTok{=} \FloatTok{150.0}
\NormalTok{labor\_endowment }\OperatorTok{=} \FloatTok{100.0}
\NormalTok{positive\_price\_floor }\OperatorTok{=} \FloatTok{1.0e{-}10}

\NormalTok{product }\OperatorTok{=}\NormalTok{ GEMB.}\FunctionTok{CommoditySpec}\NormalTok{(}
    \OperatorTok{:}\NormalTok{product;}
\NormalTok{    axes}\OperatorTok{=}\NormalTok{(period}\OperatorTok{=}\FloatTok{1}\OperatorTok{:}\NormalTok{np,),}
\NormalTok{    price\_lower\_bound}\OperatorTok{=}\NormalTok{positive\_price\_floor,}
\NormalTok{)}

\NormalTok{labor }\OperatorTok{=}\NormalTok{ GEMB.}\FunctionTok{CommoditySpec}\NormalTok{(}
    \OperatorTok{:}\NormalTok{labor;}
\NormalTok{    axes}\OperatorTok{=}\NormalTok{(period}\OperatorTok{=}\FloatTok{1}\OperatorTok{:}\NormalTok{(np }\OperatorTok{{-}} \FloatTok{1}\NormalTok{),),}
\NormalTok{    price\_lower\_bound}\OperatorTok{=}\NormalTok{positive\_price\_floor,}
\NormalTok{)}

\NormalTok{model }\OperatorTok{=}\NormalTok{ GEMB.}\FunctionTok{GEMBModel}\NormalTok{(}
\NormalTok{    [}
\NormalTok{        product,}
\NormalTok{        labor,}
\NormalTok{    ];}
\NormalTok{    numeraire}\OperatorTok{=}\NormalTok{GEMB.}\FunctionTok{CommodityRef}\NormalTok{(}
        \OperatorTok{:}\NormalTok{product;}
\NormalTok{        period}\OperatorTok{=}\FloatTok{1}\NormalTok{,}
\NormalTok{    ),}
\NormalTok{    numeraire\_value}\OperatorTok{=}\FloatTok{1.0}\NormalTok{,}
\NormalTok{)}

\NormalTok{producer }\OperatorTok{=}\NormalTok{ GEMB.}\FunctionTok{AgentTemplate}\NormalTok{(}
    \OperatorTok{:}\NormalTok{producer,}
\NormalTok{    GEMB.}\FunctionTok{CESSpec}\NormalTok{(}
\NormalTok{        [}\FloatTok{0.5}\NormalTok{, }\FloatTok{0.5}\NormalTok{];}
\NormalTok{        es}\OperatorTok{=}\FloatTok{1.0}\NormalTok{,}
\NormalTok{        alpha}\OperatorTok{=}\FloatTok{2.0}\NormalTok{,}
\NormalTok{    );}
\NormalTok{    periods}\OperatorTok{=}\FloatTok{1}\OperatorTok{:}\NormalTok{(np }\OperatorTok{{-}} \FloatTok{1}\NormalTok{),}
\NormalTok{    demands}\OperatorTok{=}\NormalTok{[}
\NormalTok{        GEMB.}\FunctionTok{CommodityRef}\NormalTok{(}
            \OperatorTok{:}\NormalTok{product;}
\NormalTok{            period}\OperatorTok{=}\NormalTok{GEMB.}\FunctionTok{RelativePeriod}\NormalTok{(}\FloatTok{0}\NormalTok{),}
\NormalTok{        ),}
\NormalTok{        GEMB.}\FunctionTok{CommodityRef}\NormalTok{(}
            \OperatorTok{:}\NormalTok{labor;}
\NormalTok{            period}\OperatorTok{=}\NormalTok{GEMB.}\FunctionTok{RelativePeriod}\NormalTok{(}\FloatTok{0}\NormalTok{),}
\NormalTok{        ),}
\NormalTok{    ],}
\NormalTok{    outputs}\OperatorTok{=}\NormalTok{[}
\NormalTok{        GEMB.}\FunctionTok{CommodityRef}\NormalTok{(}
            \OperatorTok{:}\NormalTok{product;}
\NormalTok{            period}\OperatorTok{=}\NormalTok{GEMB.}\FunctionTok{RelativePeriod}\NormalTok{(}\FloatTok{1}\NormalTok{),}
\NormalTok{        ),}
\NormalTok{    ],}
\NormalTok{    activity\_start}\OperatorTok{=}\FloatTok{100.0}\NormalTok{,}
\NormalTok{)}

\NormalTok{GEMB.}\FunctionTok{add\_agents!}\NormalTok{(model, producer)}

\NormalTok{GEMB.}\FunctionTok{add\_agent!}\NormalTok{(}
\NormalTok{    model,}
\NormalTok{    GEMB.}\FunctionTok{CESSpec}\NormalTok{(}
        \FunctionTok{fill}\NormalTok{(}\FloatTok{1.0} \OperatorTok{/}\NormalTok{ np, np);}
\NormalTok{        es}\OperatorTok{=}\FloatTok{1.0}\NormalTok{,}
\NormalTok{        alpha}\OperatorTok{=}\FloatTok{1.0}\NormalTok{,}
\NormalTok{    );}
\NormalTok{    demands}\OperatorTok{=}\NormalTok{[}
\NormalTok{        GEMB.}\FunctionTok{CommodityRef}\NormalTok{(}\OperatorTok{:}\NormalTok{product),}
\NormalTok{    ],}
\NormalTok{    endowments}\OperatorTok{=}\NormalTok{[}
\NormalTok{        GEMB.}\FunctionTok{CommodityRef}\NormalTok{(}
            \OperatorTok{:}\NormalTok{product;}
\NormalTok{            period}\OperatorTok{=}\FloatTok{1}\NormalTok{,}
\NormalTok{        ),}
\NormalTok{        GEMB.}\FunctionTok{CommodityRef}\NormalTok{(}\OperatorTok{:}\NormalTok{labor),}
\NormalTok{    ],}
\NormalTok{    endowment\_quantities}\OperatorTok{=}\FunctionTok{vcat}\NormalTok{(}
\NormalTok{        initial\_product\_endowment,}
        \FunctionTok{fill}\NormalTok{(labor\_endowment, np }\OperatorTok{{-}} \FloatTok{1}\NormalTok{),}
\NormalTok{    ),}
\NormalTok{    activity\_start}\OperatorTok{=}\FloatTok{100.0}\NormalTok{,}
\NormalTok{    name}\OperatorTok{=:}\NormalTok{consumer,}
\NormalTok{)}

\NormalTok{result }\OperatorTok{=}\NormalTok{ GEMB.}\FunctionTok{solve}\NormalTok{(}
\NormalTok{    model;}
\NormalTok{    p0}\OperatorTok{=}\FunctionTok{ones}\NormalTok{(}\FloatTok{2} \OperatorTok{*}\NormalTok{ np }\OperatorTok{{-}} \FloatTok{1}\NormalTok{),}
\NormalTok{    residual\_tol}\OperatorTok{=}\FloatTok{1.0e{-}8}\NormalTok{,}
\NormalTok{    silent}\OperatorTok{=}\ConstantTok{true}\NormalTok{,}
\NormalTok{)}

\NormalTok{GEMB.}\FunctionTok{print\_equilibrium\_statistics}\NormalTok{(model, result)}

\PreprocessorTok{@assert}\NormalTok{ result.solved}
\PreprocessorTok{@assert}\NormalTok{ result.mcp\_solved}
\PreprocessorTok{@assert}\NormalTok{ result.all\_markets\_clear}
\PreprocessorTok{@assert}\NormalTok{ result.max\_natural\_residual }\OperatorTok{\textless{}=} \FloatTok{1.0e{-}7}
\PreprocessorTok{@assert} \FunctionTok{maximum}\NormalTok{(abs, result.total\_net\_supply) }\OperatorTok{\textless{}=} \FloatTok{1.0e{-}7}
\end{Highlighting}
\end{Shaded}

}

\end{codelisting}%

\section{\texorpdfstring{Endogenous Interest Rates: A Four-Period
\(8\times2\) Pure-Exchange
Model}{Endogenous Interest Rates: A Four-Period 8\textbackslash times2 Pure-Exchange Model}}\label{endogenous-interest-rates-a-four-period-8times2-pure-exchange-model}

\subsection{Endogenous Interest Rates}\label{endogenous-interest-rates}

In a structural equilibrium model, money can be treated as an \textbf{ad
valorem claim}, with the interest rate interpreted as the corresponding
\textbf{ad valorem claim rate} (Li, 2019, Chapter 7). When an economic
agent borrows money to finance an expenditure of value \(V\), an
interest rate \(r\) generates an interest payment proportional to that
value. In this sense, the monetary claim introduces a value wedge in the
same way as other ad valorem claims.

In structural equilibrium models, the price of money in one period
refers to the interest generated by one physical unit of money over that
period. Since the choice of the physical unit of money is arbitrary,
this price depends on the chosen unit of measurement. The capital value
of one physical unit of money is the present value of the stream of
future interest payments generated by that unit of money.

In structural equilibrium models, interest rates are often treated as
exogenous parameters, for example, as rates determined by the central
bank. An economy may involve multiple interest rates: rates may differ
across periods, countries, financial markets, or economic agents. Given
these rates, the equilibrium conditions determine commodity prices,
activity levels, utility levels, and other endogenous variables. If one
or more interest rates are instead treated as endogenous variables, the
equilibrium system generally requires one additional independent
auxiliary condition for each endogenous interest rate. Such auxiliary
conditions may take the form of monetary policy rules, asset-pricing or
no-arbitrage conditions, or intertemporal relationships between the
value of money and the interest it generates.

An auxiliary condition used here is the \textbf{interest rate--interest
relationship}. Let \(v_{\mathrm{int}}^{(t)}\) denote the present value
of the interest generated in period \(t\). Under a competitive money
market and a constant money stock, the value of money at the beginning
of period \(t\) equals the present value of the stream of interest
generated from period \(t\) onward. Hence, the value of money remaining
after the interest of period \(t\) has been received is
\(\sum_{i>t}v_{\mathrm{int}}^{(i)}\). The endogenous interest rate
therefore satisfies

\[
r^{(t)}
=
\frac{v_{\mathrm{int}}^{(t)}}
{\displaystyle\sum_{i>t}v_{\mathrm{int}}^{(i)}}
\]

This interest rate--interest relationship can be added to the structural
equilibrium conditions as an auxiliary condition, thereby allowing the
equilibrium interest rate to be determined endogenously rather than
specified in advance.

\subsection{\texorpdfstring{A Four-Period \(8\times2\) Pure-Exchange
Model}{A Four-Period 8\textbackslash times2 Pure-Exchange Model}}\label{a-four-period-8times2-pure-exchange-model}

Consider a four-period pure-exchange economy. Interest rates are first
treated as exogenous parameters. The interest rate--interest
relationship is then introduced as an auxiliary condition to determine
the interest rates endogenously. The velocity of money is assumed to be
one in every period.

Assume the following:

\begin{enumerate}
\def\labelenumi{\arabic{enumi}.}
\item
  The economy contains eight commodities: labor \(i\) and money \(i\),
  \(i=1,2,3,4\). There are two consumers, a laborer and a money owner
  (rentier). Labor 1 is taken as the numeraire commodity.
\item
  The laborer is endowed with \(\omega_i>0\) units of labor \(i\),
  \(i=1,\ldots,4\), and the money owner is endowed with one physical
  unit of money in each period. Both consumers have the utility function
  \(\prod_{i=1}^4x_i^{\beta_i}\), where \(x_i\) denotes consumption of
  labor \(i\), \(\sum_{i=1}^4\beta_i=1\), and
  \(\beta_1>\beta_2>\beta_3>\beta_4>0\).
\item
  Let \(r_i>0\) denote the interest rate in period \(i\). A purchase of
  labor \(i\) with value \(v\) requires an interest payment of
  \(r_i v\).
\end{enumerate}

Let \(\mathbf p=(1,p_2,\ldots,p_8)^\top\) and
\(\mathbf w=(w_1,w_2)^\top\), where \(w_1\) and \(w_2\) are the incomes
of the two consumers. The demand and supply matrices are

\[
\mathbf D(\mathbf p,\mathbf w)
=
\begin{bmatrix}
\frac{\beta_1w_1}{1+r_1}
&
\frac{\beta_1w_2}{1+r_1}
\\
\frac{\beta_2w_1}{p_2(1+r_2)}
&
\frac{\beta_2w_2}{p_2(1+r_2)}
\\
\frac{\beta_3w_1}{p_3(1+r_3)}
&
\frac{\beta_3w_2}{p_3(1+r_3)}
\\
\frac{\beta_4w_1}{p_4(1+r_4)}
&
\frac{\beta_4w_2}{p_4(1+r_4)}
\\
\frac{r_1\beta_1w_1}{(1+r_1)p_5}
&
\frac{r_1\beta_1w_2}{(1+r_1)p_5}
\\
\frac{r_2\beta_2w_1}{(1+r_2)p_6}
&
\frac{r_2\beta_2w_2}{(1+r_2)p_6}
\\
\frac{r_3\beta_3w_1}{(1+r_3)p_7}
&
\frac{r_3\beta_3w_2}{(1+r_3)p_7}
\\
\frac{r_4\beta_4w_1}{(1+r_4)p_8}
&
\frac{r_4\beta_4w_2}{(1+r_4)p_8}
\end{bmatrix}
\]

and

\[
\mathbf S
=
\begin{bmatrix}
\omega_1&0\\
\omega_2&0\\
\omega_3&0\\
\omega_4&0\\
0&1\\
0&1\\
0&1\\
0&1
\end{bmatrix}
\]

Solving the equality-form wide-sense structural equilibrium model for
given interest rates gives

\[
p_j^*
=
\frac{\beta_j(1+r_1)\omega_1}
{\beta_1(1+r_j)\omega_j},
\qquad j=2,3,4
\]

and

\[
p_{4+i}^*
=
\frac{r_i\beta_i(1+r_1)\omega_1}
{\beta_1(1+r_i)},
\qquad i=1,2,3,4
\]

where \(p_5^*=r_1\omega_1\). The prices \(p_5^*,\ldots,p_8^*\) are the
present values of the interest payments in the four periods. The
equilibrium incomes are

\[
w_1^*
=
\omega_1
+
\frac{(1+r_1)\omega_1}{\beta_1}
\sum_{j=2}^4\frac{\beta_j}{1+r_j}
\]

and

\[
w_2^*
=
r_1\omega_1
+
\frac{(1+r_1)\omega_1}{\beta_1}
\sum_{i=2}^4\frac{r_i\beta_i}{1+r_i}
\]

The interest rates can now be endogenized by adding the interest
rate--interest relationship. Under the terminal condition
\(r_4=+\infty\), the auxiliary conditions are

\[
r_1
=
\frac{p_5^*}{p_6^*+p_7^*+p_8^*},
\qquad
r_2
=
\frac{p_6^*}{p_7^*+p_8^*},
\qquad
r_3
=
\frac{p_7^*}{p_8^*}
\]

Under the terminal condition \(r_4=+\infty\), the ordinary price of
labor 4 converges to zero, while its interest-inclusive price may remain
finite. Define the terminal interest-inclusive price as

\[
\widetilde p_4
:=
\lim_{r_4\to+\infty}(1+r_4)p_4
\]

Since the entire endowment \(\omega_4\) of labor 4 is traded in
equilibrium, the value of the corresponding interest payment satisfies

\[
p_8
=
r_4p_4\omega_4
\]

Therefore,

\[
\widetilde p_4
=
\lim_{r_4\to+\infty}r_4p_4
=
\frac{p_8}{\omega_4}
\]

Thus, although \(r_4\) is infinite and \(p_4=0\) under the terminal
condition, the effective unit cost of purchasing labor 4 remains finite
and is represented by the terminal interest-inclusive price
\(\widetilde p_4\).

Substituting the conditional-equilibrium expressions for
\(p_5^*,\ldots,p_8^*\) into these auxiliary conditions gives

\[
r_1^*
=
\frac{\beta_1}{\beta_2}-1,
\qquad
r_2^*
=
\frac{\beta_2}{\beta_3}-1,
\qquad
r_3^*
=
\frac{\beta_3}{\beta_4}-1
\] Thus, the model illustrates a general way to endogenize interest
rates in structural equilibrium models: first solve the equilibrium
conditional on exogenous interest rates, and then add one independent
auxiliary condition for each interest rate to be determined
endogenously.

For the numerical example, let

\[
(\beta_1,\beta_2,\beta_3,\beta_4)=(0.4,0.3,0.2,0.1)
\]

and suppose that the laborer is endowed with 100 units of each labor
commodity.

The numerical equilibrium gives

\[
p_1=1,\qquad p_2=\frac{2}{3},\qquad p_3=\frac{1}{3},\qquad p_4=0,
\]

\[
p_5=p_6=p_7=p_8=\frac{100}{3},
\]

and

\[
r_1=\frac{1}{3},\qquad r_2=\frac{1}{2},\qquad r_3=1
\]

The corresponding GEMB implementation is given in
Listing~\ref{lst-gemb_intertemporalEquilibrium_endogenousInterestRate_nct2_nat2}.
The two consumers are represented by Marshall-demand specifications,
while a virtual condition agent is used to determine the first three
interest rates and impose the terminal condition.

\begin{codelisting}

\caption{\label{lst-gemb_intertemporalEquilibrium_endogenousInterestRate_nct2_nat2}A
Four-Period \(8\times2\) Pure-Exchange Model with Endogenous Interest
Rates}

\centering{

\begin{Shaded}
\begin{Highlighting}[]
\CommentTok{\# GEMB high{-}level framework: four{-}period pure{-}exchange equilibrium}
\CommentTok{\# with endogenous interest rates.}
\ImportTok{using} \BuiltInTok{GeneralEquilibriumModeling}

\NormalTok{beta }\OperatorTok{=}\NormalTok{ [}\FloatTok{0.4}\NormalTok{, }\FloatTok{0.3}\NormalTok{, }\FloatTok{0.2}\NormalTok{, }\FloatTok{0.1}\NormalTok{]}
\NormalTok{labor\_endowment }\OperatorTok{=} \FunctionTok{fill}\NormalTok{(}\FloatTok{100.0}\NormalTok{, }\FloatTok{4}\NormalTok{)}
\NormalTok{money\_endowment }\OperatorTok{=} \FunctionTok{ones}\NormalTok{(}\FloatTok{4}\NormalTok{)}

\NormalTok{positive\_price\_floor }\OperatorTok{=} \FloatTok{1.0e{-}10}

\NormalTok{labor }\OperatorTok{=}\NormalTok{ GEMB.}\FunctionTok{CommoditySpec}\NormalTok{(}
    \OperatorTok{:}\NormalTok{labor;}
\NormalTok{    axes}\OperatorTok{=}\NormalTok{(period}\OperatorTok{=}\FloatTok{1}\OperatorTok{:}\FloatTok{3}\NormalTok{,),}
\NormalTok{    price\_lower\_bound}\OperatorTok{=}\NormalTok{positive\_price\_floor,}
\NormalTok{)}

\NormalTok{money }\OperatorTok{=}\NormalTok{ GEMB.}\FunctionTok{CommoditySpec}\NormalTok{(}
    \OperatorTok{:}\NormalTok{money;}
\NormalTok{    axes}\OperatorTok{=}\NormalTok{(period}\OperatorTok{=}\FloatTok{1}\OperatorTok{:}\FloatTok{4}\NormalTok{,),}
\NormalTok{    price\_lower\_bound}\OperatorTok{=}\NormalTok{positive\_price\_floor,}
\NormalTok{)}

\NormalTok{model }\OperatorTok{=}\NormalTok{ GEMB.}\FunctionTok{GEMBModel}\NormalTok{(}
\NormalTok{    [}
\NormalTok{        labor,}
        \OperatorTok{:}\NormalTok{labor\_4,}
\NormalTok{        money,}
\NormalTok{    ];}
\NormalTok{    numeraire}\OperatorTok{=}\NormalTok{GEMB.}\FunctionTok{CommodityRef}\NormalTok{(}
        \OperatorTok{:}\NormalTok{labor;}
\NormalTok{        period}\OperatorTok{=}\FloatTok{1}\NormalTok{,}
\NormalTok{    ),}
\NormalTok{    numeraire\_value}\OperatorTok{=}\FloatTok{1.0}\NormalTok{,}
\NormalTok{)}

\NormalTok{interest\_rate\_spec }\OperatorTok{=}\NormalTok{ GEMB.}\FunctionTok{ConditionAgentSpec}\NormalTok{(}
\NormalTok{    (variables, observed\_values) }\OperatorTok{{-}\textgreater{}} \ControlFlowTok{begin}
\NormalTok{        r1, r2, r3, terminal\_interest\_inclusive\_price }\OperatorTok{=}\NormalTok{ variables}
\NormalTok{        p\_money1, p\_money2, p\_money3, p\_money4 }\OperatorTok{=}\NormalTok{ observed\_values}

        \ControlFlowTok{return}\NormalTok{ [}
\NormalTok{            r1 }\OperatorTok{{-}}\NormalTok{ p\_money1 }\OperatorTok{/}\NormalTok{ (p\_money2 }\OperatorTok{+}\NormalTok{ p\_money3 }\OperatorTok{+}\NormalTok{ p\_money4),}
\NormalTok{            r2 }\OperatorTok{{-}}\NormalTok{ p\_money2 }\OperatorTok{/}\NormalTok{ (p\_money3 }\OperatorTok{+}\NormalTok{ p\_money4),}
\NormalTok{            r3 }\OperatorTok{{-}}\NormalTok{ p\_money3 }\OperatorTok{/}\NormalTok{ p\_money4,}
\NormalTok{            terminal\_interest\_inclusive\_price }\OperatorTok{{-}}\NormalTok{ p\_money4 }\OperatorTok{/}\NormalTok{ labor\_endowment[}\FloatTok{4}\NormalTok{],}
\NormalTok{        ]}
    \ControlFlowTok{end}\NormalTok{,}
\NormalTok{)}

\NormalTok{GEMB.}\FunctionTok{add\_agent!}\NormalTok{(}
\NormalTok{    model,}
\NormalTok{    interest\_rate\_spec;}
\NormalTok{    variable\_names}\OperatorTok{=}\NormalTok{[}
        \OperatorTok{:}\NormalTok{r1,}
        \OperatorTok{:}\NormalTok{r2,}
        \OperatorTok{:}\NormalTok{r3,}
        \OperatorTok{:}\NormalTok{terminal\_interest\_inclusive\_price,}
\NormalTok{    ],}
\NormalTok{    variable\_start}\OperatorTok{=}\FloatTok{1.0}\NormalTok{,}
\NormalTok{    variable\_lower\_bounds}\OperatorTok{=}\NormalTok{[}
        \FloatTok{0.0}\NormalTok{,}
        \FloatTok{0.0}\NormalTok{,}
        \FloatTok{0.0}\NormalTok{,}
\NormalTok{        positive\_price\_floor,}
\NormalTok{    ],}
\NormalTok{    variable\_upper\_bounds}\OperatorTok{=}\ConstantTok{Inf}\NormalTok{,}
\NormalTok{    observed\_variables}\OperatorTok{=}\NormalTok{[}
\NormalTok{        GEM.}\FunctionTok{PriceVariableRef}\NormalTok{(}\OperatorTok{:}\NormalTok{money\_1),}
\NormalTok{        GEM.}\FunctionTok{PriceVariableRef}\NormalTok{(}\OperatorTok{:}\NormalTok{money\_2),}
\NormalTok{        GEM.}\FunctionTok{PriceVariableRef}\NormalTok{(}\OperatorTok{:}\NormalTok{money\_3),}
\NormalTok{        GEM.}\FunctionTok{PriceVariableRef}\NormalTok{(}\OperatorTok{:}\NormalTok{money\_4),}
\NormalTok{    ],}
\NormalTok{    name}\OperatorTok{=:}\NormalTok{interestRateDeterminer,}
\NormalTok{)}

\NormalTok{consumer\_spec }\OperatorTok{=}\NormalTok{ GEMB.}\FunctionTok{MarshallDemandConsumerSpec}\NormalTok{(}
    \KeywordTok{function}\NormalTok{ (income, prices, observed\_values)}
\NormalTok{        rates }\OperatorTok{=}\NormalTok{ observed\_values[}\FloatTok{1}\OperatorTok{:}\FloatTok{3}\NormalTok{]}
\NormalTok{        terminal\_interest\_inclusive\_price }\OperatorTok{=}\NormalTok{ observed\_values[}\FloatTok{4}\NormalTok{]}

\NormalTok{        p\_labor }\OperatorTok{=}\NormalTok{ prices[}\FloatTok{1}\OperatorTok{:}\FloatTok{4}\NormalTok{]}
\NormalTok{        p\_money }\OperatorTok{=}\NormalTok{ prices[}\FloatTok{5}\OperatorTok{:}\FloatTok{8}\NormalTok{]}

\NormalTok{        labor\_demand }\OperatorTok{=}
\NormalTok{            beta[}\FloatTok{1}\OperatorTok{:}\FloatTok{3}\NormalTok{] }\OperatorTok{.*}\NormalTok{ income }\OperatorTok{./}
\NormalTok{            (p\_labor[}\FloatTok{1}\OperatorTok{:}\FloatTok{3}\NormalTok{] }\OperatorTok{.*}\NormalTok{ (}\FloatTok{1} \OperatorTok{.+}\NormalTok{ rates))}

\NormalTok{        money\_demand }\OperatorTok{=}
\NormalTok{            rates }\OperatorTok{.*}\NormalTok{ beta[}\FloatTok{1}\OperatorTok{:}\FloatTok{3}\NormalTok{] }\OperatorTok{.*}\NormalTok{ income }\OperatorTok{./}
\NormalTok{            ((}\FloatTok{1} \OperatorTok{.+}\NormalTok{ rates) }\OperatorTok{.*}\NormalTok{ p\_money[}\FloatTok{1}\OperatorTok{:}\FloatTok{3}\NormalTok{])}

\NormalTok{        labor4\_demand }\OperatorTok{=}
\NormalTok{            beta[}\FloatTok{4}\NormalTok{] }\OperatorTok{*}\NormalTok{ income }\OperatorTok{/}\NormalTok{ terminal\_interest\_inclusive\_price}

\NormalTok{        money4\_demand }\OperatorTok{=}
\NormalTok{            beta[}\FloatTok{4}\NormalTok{] }\OperatorTok{*}\NormalTok{ income }\OperatorTok{*}
\NormalTok{            (terminal\_interest\_inclusive\_price }\OperatorTok{{-}}\NormalTok{ p\_labor[}\FloatTok{4}\NormalTok{]) }\OperatorTok{/}
\NormalTok{            (terminal\_interest\_inclusive\_price }\OperatorTok{*}\NormalTok{ p\_money[}\FloatTok{4}\NormalTok{])}

        \ControlFlowTok{return} \FunctionTok{vcat}\NormalTok{(}
\NormalTok{            labor\_demand,}
\NormalTok{            labor4\_demand,}
\NormalTok{            money\_demand,}
\NormalTok{            money4\_demand,}
\NormalTok{        )}
    \KeywordTok{end}\NormalTok{,}
\NormalTok{)}

\NormalTok{observed\_interest\_variables }\OperatorTok{=}\NormalTok{ [}
\NormalTok{    GEMB.}\FunctionTok{agent\_variable\_ref}\NormalTok{(}
        \OperatorTok{:}\NormalTok{interestRateDeterminer,}
        \OperatorTok{:}\NormalTok{r1,}
\NormalTok{    ),}
\NormalTok{    GEMB.}\FunctionTok{agent\_variable\_ref}\NormalTok{(}
        \OperatorTok{:}\NormalTok{interestRateDeterminer,}
        \OperatorTok{:}\NormalTok{r2,}
\NormalTok{    ),}
\NormalTok{    GEMB.}\FunctionTok{agent\_variable\_ref}\NormalTok{(}
        \OperatorTok{:}\NormalTok{interestRateDeterminer,}
        \OperatorTok{:}\NormalTok{r3,}
\NormalTok{    ),}
\NormalTok{    GEMB.}\FunctionTok{agent\_variable\_ref}\NormalTok{(}
        \OperatorTok{:}\NormalTok{interestRateDeterminer,}
        \OperatorTok{:}\NormalTok{terminal\_interest\_inclusive\_price,}
\NormalTok{    ),}
\NormalTok{]}

\NormalTok{all\_commodities }\OperatorTok{=}\NormalTok{ [}
\NormalTok{    GEMB.}\FunctionTok{CommodityRef}\NormalTok{(}\OperatorTok{:}\NormalTok{labor),}
    \OperatorTok{:}\NormalTok{labor\_4,}
\NormalTok{    GEMB.}\FunctionTok{CommodityRef}\NormalTok{(}\OperatorTok{:}\NormalTok{money),}
\NormalTok{]}

\NormalTok{GEMB.}\FunctionTok{add\_agent!}\NormalTok{(}
\NormalTok{    model,}
\NormalTok{    consumer\_spec;}
\NormalTok{    demands}\OperatorTok{=}\NormalTok{all\_commodities,}
\NormalTok{    endowments}\OperatorTok{=}\NormalTok{[}
\NormalTok{        GEMB.}\FunctionTok{CommodityRef}\NormalTok{(}\OperatorTok{:}\NormalTok{labor),}
        \OperatorTok{:}\NormalTok{labor\_4,}
\NormalTok{    ],}
\NormalTok{    endowment\_quantities}\OperatorTok{=}\NormalTok{labor\_endowment,}
\NormalTok{    observed\_variables}\OperatorTok{=}\NormalTok{observed\_interest\_variables,}
\NormalTok{    name}\OperatorTok{=:}\NormalTok{laborer,}
\NormalTok{)}

\NormalTok{GEMB.}\FunctionTok{add\_agent!}\NormalTok{(}
\NormalTok{    model,}
\NormalTok{    consumer\_spec;}
\NormalTok{    demands}\OperatorTok{=}\NormalTok{all\_commodities,}
\NormalTok{    endowments}\OperatorTok{=}\NormalTok{GEMB.}\FunctionTok{CommodityRef}\NormalTok{(}\OperatorTok{:}\NormalTok{money),}
\NormalTok{    endowment\_quantities}\OperatorTok{=}\NormalTok{money\_endowment,}
\NormalTok{    observed\_variables}\OperatorTok{=}\NormalTok{observed\_interest\_variables,}
\NormalTok{    name}\OperatorTok{=:}\NormalTok{moneyOwner,}
\NormalTok{)}

\NormalTok{result }\OperatorTok{=}\NormalTok{ GEMB.}\FunctionTok{solve}\NormalTok{(}
\NormalTok{    model;}
\NormalTok{    p0}\OperatorTok{=}\FunctionTok{ones}\NormalTok{(}\FloatTok{8}\NormalTok{),}
\NormalTok{    residual\_tol}\OperatorTok{=}\FloatTok{1.0e{-}8}\NormalTok{,}
\NormalTok{    silent}\OperatorTok{=}\ConstantTok{true}\NormalTok{,}
\NormalTok{)}

\NormalTok{interest\_variables }\OperatorTok{=}
\NormalTok{    result.agent\_variable\_values[model.agent\_index[}\OperatorTok{:}\NormalTok{interestRateDeterminer]]}

\NormalTok{GEMB.}\FunctionTok{print\_equilibrium\_statistics}\NormalTok{(model, result)}

\FunctionTok{println}\NormalTok{()}
\FunctionTok{println}\NormalTok{(}\StringTok{"Primitive interest rates:          "}\NormalTok{, interest\_variables[}\FloatTok{1}\OperatorTok{:}\FloatTok{3}\NormalTok{])}
\FunctionTok{println}\NormalTok{(}\StringTok{"Terminal interest{-}inclusive price: "}\NormalTok{, interest\_variables[}\FloatTok{4}\NormalTok{])}
\end{Highlighting}
\end{Shaded}

}

\end{codelisting}%

\section{A Basic Overlapping-Generations Trading
Model}\label{a-basic-overlapping-generations-trading-model}

\subsection{The Model}\label{the-model}

Consider a modified version of the basic overlapping-generations (OLG)
trading model of Samuelson (1958). Assume the following:

\begin{enumerate}
\def\labelenumi{\arabic{enumi}.}
\item
  The economy lasts for \(n_p\) periods and contains \(n_p\) dated
  commodities and \(n_p-1\) consumers.
\item
  Consumer 1 has utility function
  \(u_1=x_1^{\beta_1}x_2^{\beta_2}x_3^{1-\beta_1-\beta_2}\), where
  \(\beta_1,\beta_2\geq0\) and \(\beta_1+\beta_2\leq1\).
\item
  Consumer \(i\), \(i=2,\ldots,n_p-2\), has utility function
  \(u_i=(x_i x_{i+1}x_{i+2})^{1/3}\).
\item
  Consumer \(n_p-1\) has utility function
  \(u_{n_p-1}=x_{n_p-1}^{\bar\beta_1}x_{n_p}^{\bar\beta_2}x_1^{1-\bar\beta_1-\bar\beta_2}\),
  where \(\bar\beta_1,\bar\beta_2\geq0\) and
  \(\bar\beta_1+\bar\beta_2\leq1\). When \(\bar\beta_1+\bar\beta_2=1\),
  the model has a time-line structure; when
  \(\bar\beta_1+\bar\beta_2<1\), consumer \(n_p-1\) also demands the
  first-period commodity and the model has a time-circle structure.
\item
  Consumer \(i\), \(i=1,\ldots,n_p-1\), is endowed with 100 units of
  commodity \(i\) and 100 units of commodity \(i+1\).
\end{enumerate}

Consider the case \(n_p=6\) and restrict equilibrium prices to a
geometric path,

\[
\mathbf p=(1,q,q^2,q^3,q^4,q^5)^\top,\qquad q>0
\]

Let \(\mathbf w=(w_1,\ldots,w_5)^\top\) denote consumer incomes. Since
consumer \(i\) owns commodities \(i\) and \(i+1\),

\[
w_i=100(p_i+p_{i+1}),\qquad i=1,\ldots,5
\]

The demand and supply matrices are

\[
\mathbf D(\mathbf p,\mathbf w)=
\begin{bmatrix}
\beta_1 w_1
&0&0&0&(1-\bar\beta_1-\bar\beta_2)w_5\\
\frac{\beta_2w_1}{p_2}
&\frac{w_2}{3p_2}
&0&0&0\\
\frac{(1-\beta_1-\beta_2)w_1}{p_3}
&\frac{w_2}{3p_3}
&\frac{w_3}{3p_3}
&0&0\\
0
&\frac{w_2}{3p_4}
&\frac{w_3}{3p_4}
&\frac{w_4}{3p_4}
&0\\
0&0
&\frac{w_3}{3p_5}
&\frac{w_4}{3p_5}
&\frac{\bar\beta_1w_5}{p_5}\\
0&0&0
&\frac{w_4}{3p_6}
&\frac{\bar\beta_2w_5}{p_6}
\end{bmatrix}
\]

and

\[
\mathbf S=
\begin{bmatrix}
100&0&0&0&0\\
100&100&0&0&0\\
0&100&100&0&0\\
0&0&100&100&0\\
0&0&0&100&100\\
0&0&0&0&100
\end{bmatrix}
\]

The four endpoint preference parameters \(\beta_1\), \(\beta_2\),
\(\bar\beta_1\), and \(\bar\beta_2\) are treated as endogenous. The
equality-form wide-sense structural equilibrium model consists of the
revenue-expenditure balance conditions and commodity-market clearing
conditions, together with the geometric price-path restriction. Two
equilibria are obtained.

The first is a \textbf{time-circle equilibrium}:

\[
q^*=1,\qquad
\mathbf p^*=(1,1,1,1,1,1)^\top,\qquad
\mathbf w^*=(200,200,200,200,200)^\top
\]

with

\[
\beta_1^*=0,\qquad
\beta_2^*=\frac23,\qquad
1-\beta_1^*-\beta_2^*=\frac13
\]

and

\[
\bar\beta_1^*=\frac13,\qquad
\bar\beta_2^*=\frac16,\qquad
1-\bar\beta_1^*-\bar\beta_2^*=\frac12
\]

The corresponding demand matrix is

\[
\mathbf D^*=
\begin{bmatrix}
0&0&0&0&100\\
\frac{400}{3}&\frac{200}{3}&0&0&0\\
\frac{200}{3}&\frac{200}{3}&\frac{200}{3}&0&0\\
0&\frac{200}{3}&\frac{200}{3}&\frac{200}{3}&0\\
0&0&\frac{200}{3}&\frac{200}{3}&\frac{200}{3}\\
0&0&0&\frac{200}{3}&\frac{100}{3}
\end{bmatrix}
\]

Since consumer 5 spends a positive share of income on commodity 1, the
last generation is linked back to the first period, producing a
time-circle structure.

The second is a \textbf{time-line equilibrium}. Let

\[
q^*=\frac{3+\sqrt{13}}{2}
\]

Then

\[
\mathbf p^*=(1,q^*,q^{*2},q^{*3},q^{*4},q^{*5})^\top
\]

and

\[
\mathbf w^*
=
100(1+q^*)
(1,q^*,q^{*2},q^{*3},q^{*4})^\top
\]

The endpoint preference parameters are

\[
\beta_1^*=\frac{5-\sqrt{13}}6,\qquad
\beta_2^*=\frac{-1+\sqrt{13}}6,\qquad
1-\beta_1^*-\beta_2^*=\frac13
\]

and

\[
\bar\beta_1^*=\frac13,\qquad
\bar\beta_2^*=\frac23,\qquad
1-\bar\beta_1^*-\bar\beta_2^*=0
\]

The corresponding demand matrix is

\[
\mathbf D^*
=
\frac{100(1+q^*)}{3}
\begin{bmatrix}
\frac{5-\sqrt{13}}2&0&0&0&0\\
\frac{-1+\sqrt{13}}{2q^*}&1&0&0&0\\
\frac{1}{q^{*2}}&\frac{1}{q^*}&1&0&0\\
0&\frac{1}{q^{*2}}&\frac{1}{q^*}&1&0\\
0&0&\frac{1}{q^{*2}}&\frac{1}{q^*}&1\\
0&0&0&\frac{1}{q^{*2}}&\frac{2}{q^*}
\end{bmatrix}
\]

In this equilibrium, \(1-\bar\beta_1^*-\bar\beta_2^*=0\), so the last
consumer has no demand for the first-period commodity and the
intergenerational trading structure forms a time line rather than a
circle.

Both solutions represent steady-state equilibrium paths. Although the
finite-horizon model contains special endpoint consumers to close the
six-period economy, its interior periods reproduce the same
intergenerational structure under a one-period shift of calendar time
and cohort index. In particular, the endowment and age-specific demand
patterns are unchanged, while prices and incomes are multiplied by the
constant factor \(q^*\). Hence, abstracting from the endpoint closure,
the interior equilibrium structure can be extended indefinitely forward
and backward in time, yielding the corresponding steady-state OLG
equilibrium path.

\subsection{Computational
Implementation}\label{computational-implementation}

The model is implemented by treating products delivered in different
periods as different commodities. For \(n_p=6\), the economy therefore
contains six dated commodities with price vector
\(\mathbf p=(p_1,\ldots,p_6)^\top\), and the first-period product is
chosen as the numeraire, so that \(p_1=1\).

Consumer \(i\) is endowed with 100 units of commodities \(i\) and
\(i+1\). Hence its income is determined directly by prices as
\(w_i=100(p_i+p_{i+1})\) and does not need to be introduced as a
separate equilibrium variable. Consumers 2--4 have equal-share
Cobb--Douglas preferences and are represented directly by
\texttt{CESSpec} with three shares equal to \(1/3\).

The preference shares of consumers 1 and 5 are endogenous. Their demands
are therefore specified by \texttt{MarshallDemandConsumerSpec}. Consumer
1 observes the endogenous variables \(\beta_1\) and \(\beta_2\) and has
expenditure shares \(\beta_1\), \(\beta_2\), and \(1-\beta_1-\beta_2\).
Consumer 5 similarly observes \(\bar\beta_1\) and \(\bar\beta_2\) and
has expenditure shares \(\bar\beta_1\), \(\bar\beta_2\), and
\(1-\bar\beta_1-\bar\beta_2\).

A \texttt{ConditionAgentSpec} is introduced to determine the
steady-state price path and the four endpoint preference parameters. Its
endogenous variables are

\[
(\log q,\beta_1,\beta_2,\bar\beta_1,\bar\beta_2)
\]

where \(q=\exp(\log q)>0\). The five associated conditions impose the
geometric price path

\[
p_{t+1}=qp_t,\qquad t=1,\ldots,5
\]

Together with \(p_1=1\), these conditions imply

\[
\mathbf p=(1,q,q^2,q^3,q^4,q^5)^\top
\]

Although the four preference parameters do not appear explicitly in
these five price-path conditions, they enter the demands of consumers 1
and 5 and therefore affect commodity-market clearing. PATH consequently
solves prices, \(q\), and the four preference parameters jointly so that
both the geometric price conditions and all market-clearing conditions
are satisfied.

After fixing the numeraire, there are five ordinary price variables
\(p_2,\ldots,p_6\) and five additional endogenous variables in the
condition agent. Correspondingly, Walras' law leaves five independent
commodity-market clearing conditions, while the condition agent supplies
five price-path equations. The resulting system is therefore closed.

The model has multiple equilibria. To obtain the two equilibria reported
above, the same initial values \((0.1,0.1,0.1,0.1)\) are used for the
four preference parameters, while only the initial value of \(q\) is
changed. The time-circle equilibrium is obtained from \(q_0=0.5\),
whereas the time-line equilibrium is obtained from \(q_0=2\).

The corresponding GEMB implementation is given in
Listing~\ref{lst-gemb_OLG_trading}.

\begin{codelisting}

\caption{\label{lst-gemb_OLG_trading}A Basic Overlapping-Generations
Trading Model}

\centering{

\begin{Shaded}
\begin{Highlighting}[]
\ImportTok{using} \BuiltInTok{GeneralEquilibriumModeling}

\KeywordTok{const}\NormalTok{ GEM }\OperatorTok{=}\NormalTok{ GeneralEquilibriumModeling.GEM}
\KeywordTok{const}\NormalTok{ GEMB }\OperatorTok{=}\NormalTok{ GeneralEquilibriumModeling.GEMB}

\NormalTok{np }\OperatorTok{=} \FloatTok{6}
\NormalTok{omega }\OperatorTok{=} \FloatTok{100.0}
\FunctionTok{product}\NormalTok{(t) }\OperatorTok{=}\NormalTok{ GEMB.}\FunctionTok{CommodityRef}\NormalTok{(}\OperatorTok{:}\NormalTok{product; period}\OperatorTok{=}\NormalTok{t)}

\KeywordTok{function} \FunctionTok{build\_model}\NormalTok{(q0, beta0)}
\NormalTok{    model }\OperatorTok{=}\NormalTok{ GEMB.}\FunctionTok{GEMBModel}\NormalTok{(}
\NormalTok{        [}
\NormalTok{            GEMB.}\FunctionTok{CommoditySpec}\NormalTok{(}
                \OperatorTok{:}\NormalTok{product;}
\NormalTok{                axes}\OperatorTok{=}\NormalTok{(period}\OperatorTok{=}\FloatTok{1}\OperatorTok{:}\NormalTok{np,),}
\NormalTok{                price\_lower\_bound}\OperatorTok{=}\FloatTok{1.0e{-}10}\NormalTok{,}
\NormalTok{            ),}
\NormalTok{        ];}
\NormalTok{        numeraire}\OperatorTok{=}\FunctionTok{product}\NormalTok{(}\FloatTok{1}\NormalTok{),}
\NormalTok{    )}

    \CommentTok{\# q and endpoint preference shares}
\NormalTok{    GEMB.}\FunctionTok{add\_agent!}\NormalTok{(}
\NormalTok{        model,}
\NormalTok{        GEMB.}\FunctionTok{ConditionAgentSpec}\NormalTok{(}
\NormalTok{            (v, p) }\OperatorTok{{-}\textgreater{}} \ControlFlowTok{begin}
\NormalTok{                q }\OperatorTok{=} \FunctionTok{exp}\NormalTok{(v[}\FloatTok{1}\NormalTok{])}
\NormalTok{                [p[t }\OperatorTok{+} \FloatTok{1}\NormalTok{] }\OperatorTok{{-}}\NormalTok{ q }\OperatorTok{*}\NormalTok{ p[t] for t }\KeywordTok{in} \FloatTok{1}\OperatorTok{:}\NormalTok{(np }\OperatorTok{{-}} \FloatTok{1}\NormalTok{)]}
            \ControlFlowTok{end}\NormalTok{,}
\NormalTok{        );}
\NormalTok{        variable\_names}\OperatorTok{=}\NormalTok{[}\OperatorTok{:}\NormalTok{log\_q, }\OperatorTok{:}\NormalTok{beta1, }\OperatorTok{:}\NormalTok{beta2, }\OperatorTok{:}\NormalTok{beta\_bar1, }\OperatorTok{:}\NormalTok{beta\_bar2],}
\NormalTok{        variable\_start}\OperatorTok{=}\NormalTok{[}\FunctionTok{log}\NormalTok{(q0); beta0],}
\NormalTok{        variable\_lower\_bounds}\OperatorTok{={-}}\ConstantTok{Inf}\NormalTok{,}
\NormalTok{        variable\_upper\_bounds}\OperatorTok{=}\ConstantTok{Inf}\NormalTok{,}
\NormalTok{        observed\_variables}\OperatorTok{=}\NormalTok{[}
\NormalTok{            GEM.}\FunctionTok{PriceVariableRef}\NormalTok{(}\FunctionTok{Symbol}\NormalTok{(}\OperatorTok{:}\NormalTok{product\_, t))}
\NormalTok{            for t }\KeywordTok{in} \FloatTok{1}\OperatorTok{:}\NormalTok{np}
\NormalTok{        ],}
\NormalTok{        name}\OperatorTok{=:}\NormalTok{steady\_state,}
\NormalTok{    )}

    \CommentTok{\# Consumer 1}
\NormalTok{    GEMB.}\FunctionTok{add\_agent!}\NormalTok{(}
\NormalTok{        model,}
\NormalTok{        GEMB.}\FunctionTok{MarshallDemandConsumerSpec}\NormalTok{(}
\NormalTok{            (w, p, b) }\OperatorTok{{-}\textgreater{}}\NormalTok{ [}
\NormalTok{                b[}\FloatTok{1}\NormalTok{] }\OperatorTok{*}\NormalTok{ w }\OperatorTok{/}\NormalTok{ p[}\FloatTok{1}\NormalTok{],}
\NormalTok{                b[}\FloatTok{2}\NormalTok{] }\OperatorTok{*}\NormalTok{ w }\OperatorTok{/}\NormalTok{ p[}\FloatTok{2}\NormalTok{],}
\NormalTok{                (}\FloatTok{1} \OperatorTok{{-}}\NormalTok{ b[}\FloatTok{1}\NormalTok{] }\OperatorTok{{-}}\NormalTok{ b[}\FloatTok{2}\NormalTok{]) }\OperatorTok{*}\NormalTok{ w }\OperatorTok{/}\NormalTok{ p[}\FloatTok{3}\NormalTok{],}
\NormalTok{            ],}
\NormalTok{        );}
\NormalTok{        demands}\OperatorTok{=}\FunctionTok{product}\NormalTok{.(}\FloatTok{1}\OperatorTok{:}\FloatTok{3}\NormalTok{),}
\NormalTok{        endowments}\OperatorTok{=}\FunctionTok{product}\NormalTok{.(}\FloatTok{1}\OperatorTok{:}\FloatTok{2}\NormalTok{),}
\NormalTok{        endowment\_quantities}\OperatorTok{=}\FunctionTok{fill}\NormalTok{(omega, }\FloatTok{2}\NormalTok{),}
\NormalTok{        observed\_variables}\OperatorTok{=}\NormalTok{[}
\NormalTok{            GEM.}\FunctionTok{AgentVariableRef}\NormalTok{(}\OperatorTok{:}\NormalTok{steady\_state, }\OperatorTok{:}\NormalTok{beta1),}
\NormalTok{            GEM.}\FunctionTok{AgentVariableRef}\NormalTok{(}\OperatorTok{:}\NormalTok{steady\_state, }\OperatorTok{:}\NormalTok{beta2),}
\NormalTok{        ],}
\NormalTok{        name}\OperatorTok{=:}\NormalTok{consumer1,}
\NormalTok{    )}

    \CommentTok{\# Consumers 2{-}{-}4}
\CommentTok{\# Consumers 2{-}{-}4}
\ControlFlowTok{for}\NormalTok{ i }\KeywordTok{in} \FloatTok{2}\OperatorTok{:}\FloatTok{4}
\NormalTok{    GEMB.}\FunctionTok{add\_agent!}\NormalTok{(}
\NormalTok{        model,}
\NormalTok{        GEMB.}\FunctionTok{CESSpec}\NormalTok{(}\FunctionTok{fill}\NormalTok{(}\FloatTok{1} \OperatorTok{/} \FloatTok{3}\NormalTok{, }\FloatTok{3}\NormalTok{));}
\NormalTok{        demands}\OperatorTok{=}\FunctionTok{product}\NormalTok{.(i}\OperatorTok{:}\NormalTok{(i }\OperatorTok{+} \FloatTok{2}\NormalTok{)),}
\NormalTok{        endowments}\OperatorTok{=}\FunctionTok{product}\NormalTok{.(i}\OperatorTok{:}\NormalTok{(i }\OperatorTok{+} \FloatTok{1}\NormalTok{)),}
\NormalTok{        endowment\_quantities}\OperatorTok{=}\FunctionTok{fill}\NormalTok{(omega, }\FloatTok{2}\NormalTok{),}
\NormalTok{        name}\OperatorTok{=}\FunctionTok{Symbol}\NormalTok{(}\OperatorTok{:}\NormalTok{consumer, i),}
\NormalTok{    )}
\ControlFlowTok{end}

    \CommentTok{\# Consumer 5}
\NormalTok{    GEMB.}\FunctionTok{add\_agent!}\NormalTok{(}
\NormalTok{        model,}
\NormalTok{        GEMB.}\FunctionTok{MarshallDemandConsumerSpec}\NormalTok{(}
\NormalTok{            (w, p, b) }\OperatorTok{{-}\textgreater{}}\NormalTok{ [}
\NormalTok{                b[}\FloatTok{1}\NormalTok{] }\OperatorTok{*}\NormalTok{ w }\OperatorTok{/}\NormalTok{ p[}\FloatTok{1}\NormalTok{],}
\NormalTok{                b[}\FloatTok{2}\NormalTok{] }\OperatorTok{*}\NormalTok{ w }\OperatorTok{/}\NormalTok{ p[}\FloatTok{2}\NormalTok{],}
\NormalTok{                (}\FloatTok{1} \OperatorTok{{-}}\NormalTok{ b[}\FloatTok{1}\NormalTok{] }\OperatorTok{{-}}\NormalTok{ b[}\FloatTok{2}\NormalTok{]) }\OperatorTok{*}\NormalTok{ w }\OperatorTok{/}\NormalTok{ p[}\FloatTok{3}\NormalTok{],}
\NormalTok{            ],}
\NormalTok{        );}
\NormalTok{        demands}\OperatorTok{=}\NormalTok{[}\FunctionTok{product}\NormalTok{(}\FloatTok{5}\NormalTok{), }\FunctionTok{product}\NormalTok{(}\FloatTok{6}\NormalTok{), }\FunctionTok{product}\NormalTok{(}\FloatTok{1}\NormalTok{)],}
\NormalTok{        endowments}\OperatorTok{=}\FunctionTok{product}\NormalTok{.(}\FloatTok{5}\OperatorTok{:}\FloatTok{6}\NormalTok{),}
\NormalTok{        endowment\_quantities}\OperatorTok{=}\FunctionTok{fill}\NormalTok{(omega, }\FloatTok{2}\NormalTok{),}
\NormalTok{        observed\_variables}\OperatorTok{=}\NormalTok{[}
\NormalTok{            GEM.}\FunctionTok{AgentVariableRef}\NormalTok{(}\OperatorTok{:}\NormalTok{steady\_state, }\OperatorTok{:}\NormalTok{beta\_bar1),}
\NormalTok{            GEM.}\FunctionTok{AgentVariableRef}\NormalTok{(}\OperatorTok{:}\NormalTok{steady\_state, }\OperatorTok{:}\NormalTok{beta\_bar2),}
\NormalTok{        ],}
\NormalTok{        name}\OperatorTok{=:}\NormalTok{consumer5,}
\NormalTok{    )}

    \ControlFlowTok{return}\NormalTok{ model}
\KeywordTok{end}

\KeywordTok{function} \FunctionTok{solve\_case}\NormalTok{(q0, beta0)}
\NormalTok{    model }\OperatorTok{=} \FunctionTok{build\_model}\NormalTok{(q0, beta0)}

\NormalTok{    result }\OperatorTok{=}\NormalTok{ GEMB.}\FunctionTok{solve}\NormalTok{(}
\NormalTok{        model;}
\NormalTok{        p0}\OperatorTok{=}\NormalTok{q0 }\OperatorTok{.\^{}}\NormalTok{ (}\FloatTok{0}\OperatorTok{:}\NormalTok{(np }\OperatorTok{{-}} \FloatTok{1}\NormalTok{)),}
\NormalTok{        silent}\OperatorTok{=}\ConstantTok{true}\NormalTok{,}
\NormalTok{    )}

\NormalTok{    i }\OperatorTok{=}\NormalTok{ model.agent\_index[}\OperatorTok{:}\NormalTok{steady\_state]}
\NormalTok{    v }\OperatorTok{=}\NormalTok{ result.agent\_variable\_values[i]}

\NormalTok{    q }\OperatorTok{=} \FunctionTok{exp}\NormalTok{(v[}\FloatTok{1}\NormalTok{])}
\NormalTok{    beta1, beta2, beta\_bar1, beta\_bar2 }\OperatorTok{=}\NormalTok{ v[}\FloatTok{2}\OperatorTok{:}\FloatTok{5}\NormalTok{]}

    \FunctionTok{println}\NormalTok{(}\StringTok{"}\SpecialCharTok{\textbackslash{}n}\StringTok{q = "}\NormalTok{, q)}
    \FunctionTok{println}\NormalTok{(}\StringTok{"p = "}\NormalTok{, result.prices)}
    \FunctionTok{println}\NormalTok{(}\StringTok{"beta = "}\NormalTok{, [beta1, beta2, beta\_bar1, beta\_bar2])}

\NormalTok{    GEMB.}\FunctionTok{print\_equilibrium\_statistics}\NormalTok{(model, result)}

    \ControlFlowTok{return}\NormalTok{ result}
\KeywordTok{end}

\CommentTok{\# Time{-}circle equilibrium}
\NormalTok{circle }\OperatorTok{=} \FunctionTok{solve\_case}\NormalTok{(}
    \FloatTok{0.5}\NormalTok{,}
    \FunctionTok{fill}\NormalTok{(}\FloatTok{0.1}\NormalTok{, }\FloatTok{4}\NormalTok{),}
\NormalTok{)}

\CommentTok{\# Time{-}line equilibrium}
\NormalTok{line }\OperatorTok{=} \FunctionTok{solve\_case}\NormalTok{(}
    \FloatTok{2.0}\NormalTok{,}
    \FunctionTok{fill}\NormalTok{(}\FloatTok{0.1}\NormalTok{, }\FloatTok{4}\NormalTok{),}
\NormalTok{)}

\ConstantTok{nothing}
\end{Highlighting}
\end{Shaded}

}

\end{codelisting}%

\bookmarksetup{startatroot}

\chapter*{References}\label{references}
\addcontentsline{toc}{chapter}{References}

\markboth{References}{References}

Björk, T., A. Murgoci and X. Y. Zhou (2014) Mean--Variance Portfolio
Optimization with State-Dependent Risk Aversion. Mathematical Finance,
24(1): 1-24.

Bonnisseau, J. M., and Cornet, B. (1990) Existence of Marginal Cost
Pricing Equilibria in Economies with Several Nonconvex Firms.
Econometrica, 58(3): 661-682.

Ferris, Michael C., and Todd S. Munson (2000) Complementarity Problems
in GAMS and the PATH Solver. Journal of Economic Dynamics and Control,
24(2): 165-188.

Fullerton, D. (1989) Notes on Displaced CES Functional Forms.
Unpublished manuscript.

Kemeny, J. G., O. Morgenstern and G. L. Thompson (1956) A Generalization
of the von Neumann Model of an Expanding Economy. Econometrica, 24(2):
115-135.

LI Wu (2019) General Equilibrium and Structural Dynamics: Perspectives
of New Structural Economics. Beijing: Economic Science Press. (In
Chinese)

Mas-Colell, A., Whinston, M. D., Green, J. R. (1995) Microeconomic
Theory. New York: Oxford University Press.

McKenzie, L. W. (2002) Classical General Equilibrium Theory. MIT Press.

Nakamura, Yutaka (2015) Mean-Variance Utility. Journal of Economic
Theory, 160: 536-556.

Pollak, Robert A. (1977) Price Dependent Preferences. American Economic
Review, 67(2): 64--75.

Quinzii, M. (1992) Increasing Returns and Efficiency. Oxford University
Press.

Samuelson, P. A. (1958) An Exact Consumption-Loan Model of Interest with
or without the Social Contrivance of Money. Journal of Political
Economy, 66(6): 467-482.

Torres, Jose L. (2016) Introduction to Dynamic Macroeconomic General
Equilibrium Models (Second Edition). Vernon Press.

von Neumann, J. (1945) A Model of General Economic Equilibrium. The
Review of Economic Studies, 13(1): 1-9. English translation of von
Neumann (1937).




\end{document}
